\documentclass[10pt]{article}
\usepackage[
  paperwidth=5in,
  paperheight=7in,
  top=10mm,
  bottom=9mm,
  left=9.5mm,
  right=9.5mm,
  headheight=18pt,
  headsep=3.5mm
]{geometry}
\usepackage{fontspec}
\usepackage{amsmath,amssymb}
\usepackage{graphicx}
\usepackage{eso-pic}
\usepackage{booktabs}
\usepackage{tabularx}
\usepackage{adjustbox}
\usepackage{fancyhdr}
\usepackage{tikz}
\usepackage{xcolor}
\usepackage{colortbl}
\usepackage[skins]{tcolorbox}

% --- Design palette: monochrome interior (the red cover keeps its own colours) ---
\definecolor{kbvaccent}{HTML}{1A1A1A}
\definecolor{kbvaccentdark}{HTML}{000000}
\definecolor{kbvtint}{HTML}{FFFFFF}
\definecolor{kbvblue}{HTML}{333333}
\definecolor{kbvbluetint}{HTML}{FFFFFF}
\definecolor{kbvgrey}{HTML}{6E6E6E}
\definecolor{kbvrow}{HTML}{FFFFFF}
\definecolor{kbvline}{HTML}{BDBDBD}
\newlength{\kbvcolw}

\setmainfont[
  Path=fonts/,
  Extension=.ttf,
  UprightFont=*-Regular,
  ItalicFont=*-Italic,
  Script=Kannada,
  HyphenChar=None,
  AutoFakeBold=1.4
]{TiroKannada}
\setsansfont[
  Script=Kannada,
  HyphenChar=None,
  BoldFont={Noto Sans Kannada Bold}
]{Noto Sans Kannada}
\newfontfamily\kbvdisplayfont[
  Path=fonts/,
  Extension=.ttf,
  UprightFont=*-ExtraBold,
  Script=Kannada,
  HyphenChar=None
]{BalooTamma2}
\newfontfamily\kbvcoversans[
  Path=fonts/,
  Extension=.ttf,
  UprightFont=*-SemiBold,
  Script=Kannada,
  HyphenChar=None
]{BalooTamma2}
\newfontfamily\kbvchapfont[
  Path=fonts/,
  Extension=.ttf,
  UprightFont=*-Bold,
  Script=Kannada,
  HyphenChar=None
]{BalooTamma2}
\newfontfamily\kbvsecfont[
  Script=Kannada,
  HyphenChar=None
]{Noto Sans Kannada Bold}
% Latin fallback: Noto Sans Kannada carries no Latin letters (used for (i), (ii), U metre marks, names)
\newfontfamily\kbvlatin[Path=fonts/,Extension=.ttf,UprightFont=*-Regular,ItalicFont=*-Italic]{TiroKannada}

\setlength{\parindent}{1.6em}
\setlength{\parskip}{4.5pt plus 6pt minus 1.5pt}
\linespread{1.20}
\sloppy

\newsavebox{\kbvpagebox}
\newcommand{\kbvtypesetpage}[1]{%
  \begin{lrbox}{\kbvpagebox}%
    \begin{minipage}{\textwidth}%
      \fontsize{10.0pt}{13.0pt}\selectfont
      \setlength{\parindent}{1.6em}%
      \setlength{\parskip}{4.5pt}%
      #1%
    \end{minipage}%
  \end{lrbox}%
  \ifdim\dimexpr\ht\kbvpagebox+\dp\kbvpagebox\relax>0.965\textheight
    \message{^^JKBVSCALED page \thepage\space ht=\the\dimexpr\ht\kbvpagebox+\dp\kbvpagebox\relax\space th=\the\textheight^^J}%
    \noindent\begin{adjustbox}{max totalheight=0.965\textheight, max width=\textwidth, keepaspectratio}%
      \usebox{\kbvpagebox}%
    \end{adjustbox}%
  \else
    \noindent\vbox to 0.96\textheight{%
      \fontsize{10.0pt}{13.0pt}\selectfont
      \setlength{\parindent}{1.6em}%
      \setlength{\parskip}{3.5pt plus 1.5pt minus 1pt}%
      #1%
      \par\vfill
    }%
  \fi
}

% --- Running header: chapter name left, page number in a red tab right ---
\newcommand{\kbvchaptername}{}
\fancypagestyle{kbvpage}{%
  \fancyhf{}%
  \fancyhead[L]{{\kbvsecfont\fontsize{8pt}{10pt}\selectfont\color{kbvgrey}\kbvchaptername}}%
  \fancyhead[R]{{\kbvsecfont\fontsize{9.5pt}{11pt}\selectfont\thepage}}%
  \renewcommand{\headrule}{{\color{kbvline}\hrule width\headwidth height 0.4pt}}%
  \renewcommand{\footrulewidth}{0pt}%
}

% --- Design macros (clean, monochrome) ---
% Chapter opener: small grey lesson label, large title, hairline.
\newcommand{\kbvchapter}[2]{%
  \vspace{2pt plus 2pt}%
  \noindent{\kbvsecfont\fontsize{9.5pt}{12pt}\selectfont\color{kbvgrey}#1}\par
  \vspace{2pt}%
  \noindent{\kbvchapfont\fontsize{16pt}{20pt}\selectfont #2}\par
  \vspace{2pt}\noindent{\color{kbvline}\rule{\textwidth}{0.5pt}}\par
  \vspace{5pt plus 4pt minus 1pt}%
}
\newcommand{\kbvchapterplain}[1]{%
  \vspace{2pt plus 2pt}%
  \noindent{\kbvchapfont\fontsize{16pt}{20pt}\selectfont #1}\par
  \vspace{2pt}\noindent{\color{kbvline}\rule{\textwidth}{0.5pt}}\par
  \vspace{5pt plus 4pt minus 1pt}%
}
% Exercise heading: bold sans with a hairline beneath.
\newcommand{\kbvexercise}[1]{%
  \vspace{7pt plus 6pt minus 2pt}%
  \noindent{\kbvsecfont\fontsize{12pt}{15pt}\selectfont #1}\par
  \vspace{1.5pt}\noindent{\color{kbvline}\rule{\textwidth}{0.4pt}}\par
  \vspace{3pt plus 3pt minus 1pt}%
}
% Centered subtitle between two hairlines:  ——— ಲಿಂಗ ———
\newcommand{\kbvsubtitle}[1]{%
  \vspace{4pt plus 3pt minus 1pt}%
  \noindent{\color{kbvline}\hrulefill}\hspace{0.8em}{\kbvsecfont\fontsize{11.5pt}{14.5pt}\selectfont #1}\hspace{0.8em}{\color{kbvline}\hrulefill}\par
  \vspace{3pt plus 3pt minus 1pt}%
}
% Section head: bold sans, nothing else.
\newcommand{\kbvsechead}[1]{%
  \vspace{5pt plus 4pt minus 1pt}%
  \noindent{\kbvsecfont\fontsize{11.5pt}{15pt}\selectfont #1}\par
  \vspace{2pt plus 2pt}%
}
% Prosody marks: laghu = deep U (cup), guru = macron (used from the ಛಂದಸ್ಸು chapter on).
\newcommand{\laghu}{\scalebox{0.85}{\ensuremath{\cup}}}
\newcommand{\guru}{\raisebox{0.62ex}{\rule{0.78em}{0.07em}}}
% Example label, always "ಉದಾ :" in bold sans.
\newcommand{\kbvlabel}[1]{{\kbvsecfont #1\,:}}
% Highlighted result word of a derivation.
\newcommand{\kbvresult}[1]{{\kbvsecfont #1}}
% Item marker (1), (೧), ೧.
\newcommand{\kbvnum}[1]{{\kbvsecfont #1}}
% Proverb / list bullet.
\newcommand{\kbvbullet}{{\color{kbvgrey}\raisebox{0.25ex}{\scriptsize\ensuremath{\blacksquare}}}}
% Grammar rule block: white, no frame (the square bullet marks it).
\newtcolorbox{kbvrule}{enhanced,colback=kbvtint,colframe=kbvtint,boxrule=0pt,arc=0pt,
  left=7pt,right=7pt,top=5pt,bottom=5pt,boxsep=0pt,before skip=5pt plus 3pt,after skip=5pt plus 3pt}
% Note / explanation box: hairline frame on white.
\newtcolorbox{kbvnote}{enhanced,colback=white,colframe=kbvline,boxrule=0.4pt,arc=0pt,
  left=7pt,right=7pt,top=5pt,bottom=5pt,boxsep=0pt,before skip=5pt plus 3pt,after skip=5pt plus 3pt}
% Derivation / word-formation block: white, no frame.
\newtcolorbox{kbvderiv}{enhanced,colback=kbvrow,colframe=kbvrow,boxrule=0pt,arc=0pt,
  left=5pt,right=5pt,top=3pt,bottom=3pt,boxsep=0pt,before skip=5pt plus 3pt,after skip=5pt plus 3pt}
% Verse block: indented, no decoration.
\newtcolorbox{kbvverse}{enhanced,colback=white,colframe=white,boxrule=0pt,arc=0pt,
  left=12pt,right=4pt,top=2pt,bottom=2pt,boxsep=0pt,before skip=4pt plus 3pt,after skip=4pt plus 3pt}

% --- Edition-specific macros (page-faithful) ---
\newcommand{\kbvcontinues}{\par}      % paragraph cut by the original page break: end it here
\newcommand{\kbvcontinued}{\indent}   % ... and start the remainder as a new paragraph
\newcommand{\kbvbooktitle}{%          % book title repeated on the first text page
  \vspace{2pt plus 2pt}\begin{center}{\kbvchapfont\fontsize{20pt}{24pt}\selectfont ಕನ್ನಡ ಬಾಲ ವ್ಯಾಕರಣ}\end{center}\vspace{4pt plus 4pt minus 1pt}}
\begin{document}

% === tex/0001.tex ===
% === File 0001.md: Native TikZ Cover Page ===
\thispagestyle{empty}
\AddToShipoutPictureBG*{%
  \AtPageLowerLeft{%
    \begin{tikzpicture}
      \useasboundingbox (0,0) rectangle (\paperwidth,\paperheight);
      % Define exact cover palette from coverpage.png
      \definecolor{kbvred}{HTML}{DF212B}
      \definecolor{kbvdarkred}{HTML}{9E1319}
      \definecolor{kbvcloudedge}{HTML}{B51921}
      \definecolor{kbvtextred}{HTML}{98181D}
      \definecolor{kbvoutline}{HTML}{2A0507}
      \definecolor{kbvwavepink}{HTML}{EB605A}
      \definecolor{kbvplumegrey}{HTML}{B8C2C4}
      \definecolor{kbvplumedark}{HTML}{3D4345}

      % 1. Full-bleed vermilion red background
      \fill[kbvred] (0,0) rectangle (\paperwidth,\paperheight);

      % 2. Darker crimson band & 3 sinusoidal wavy ribbons at bottom right (6.5 crests)
      \begin{scope}
        \clip (0.430\paperwidth, 0.0700\paperheight) rectangle (\paperwidth, 0.1524\paperheight);
        \fill[kbvdarkred] (0.430\paperwidth, 0.0860\paperheight) rectangle (\paperwidth, 0.1365\paperheight);
        % Top white wavy ribbon
        \draw[white, line width=3.8pt, line cap=butt]
          plot[domain=0.425:1.01, samples=220, smooth]
          (\x*\paperwidth, {0.1348*\paperheight + 0.0062*\paperheight*cos((\x - 0.430)*11.4*360)});
        % Middle coral-pink wavy ribbon
        \draw[kbvwavepink, line width=3.8pt, line cap=butt]
          plot[domain=0.425:1.01, samples=220, smooth]
          (\x*\paperwidth, {0.1112*\paperheight + 0.0062*\paperheight*cos((\x - 0.430)*11.4*360)});
        % Bottom white wavy ribbon
        \draw[white, line width=3.8pt, line cap=butt]
          plot[domain=0.425:1.01, samples=220, smooth]
          (\x*\paperwidth, {0.0876*\paperheight + 0.0062*\paperheight*cos((\x - 0.430)*11.4*360)});
      \end{scope}

      % 3. Six horizontal white rules across the page width
      \draw[white, line width=1.1pt] (0, 0.9440\paperheight) -- (\paperwidth, 0.9440\paperheight);
      \draw[white, line width=1.0pt] (0, 0.7574\paperheight) -- (\paperwidth, 0.7574\paperheight);
      \draw[white, line width=1.0pt] (0, 0.6151\paperheight) -- (\paperwidth, 0.6151\paperheight);
      \draw[white, line width=1.0pt] (0, 0.4191\paperheight) -- (\paperwidth, 0.4191\paperheight);
      \draw[white, line width=1.0pt] (0, 0.1524\paperheight) -- (\paperwidth, 0.1524\paperheight);
      \draw[white, line width=1.0pt] (0, 0.0700\paperheight) -- (\paperwidth, 0.0700\paperheight);

      % 4. Scalloped 11-lobed cloud badge on the right (matching coverpage.png lobe geometry)
      \begin{scope}[shift={(0.750\paperwidth, 0.778\paperheight)}]
        \def\cloudlobes{
          -5.8/14.2/9.2, 5.8/14.2/9.2,
          -14.2/9.2/9.2, 14.2/9.2/9.2,
          -17.4/0.8/9.0, 17.4/0.8/9.0,
          -15.2/-7.5/9.0, 15.2/-7.5/9.0,
          -8.6/-13.2/8.6, 8.6/-13.2/8.6,
          0.0/-15.2/8.4
        }
        % Layer 1: Outer dark-red border
        \foreach \rx/\ry/\rad in \cloudlobes {
          \fill[kbvcloudedge] (\rx mm, \ry mm) circle ({\rad mm + 1.35mm});
        }
        \fill[kbvcloudedge] (0,0) ellipse (19.5mm and 16.5mm);

        % Layer 2: Outer white ring
        \foreach \rx/\ry/\rad in \cloudlobes {
          \fill[white] (\rx mm, \ry mm) circle ({\rad mm + 0.85mm});
        }
        \fill[white] (0,0) ellipse (19.0mm and 16.0mm);

        % Layer 3: Inner red contour line
        \foreach \rx/\ry/\rad in \cloudlobes {
          \fill[kbvred] (\rx mm, \ry mm) circle ({\rad mm - 0.15mm});
        }
        \fill[kbvred] (0,0) ellipse (18.0mm and 15.0mm);

        % Layer 4: Pure white cloud interior
        \foreach \rx/\ry/\rad in \cloudlobes {
          \fill[white] (\rx mm, \ry mm) circle ({\rad mm - 0.65mm});
        }
        \fill[white] (0,0) ellipse (17.5mm and 14.5mm);

        % 6-line centered Kannada subtitle inside the cloud
        \node[align=center, text=kbvtextred, font=\kbvcoversans\fontsize{12.2pt}{14.2pt}\selectfont] at (0, -0.2mm) {%
          ಪ್ರಾಥಮಿಕ\\
          ಶಾಲಾ ಮತ್ತು\\
          ಪ್ರೌಢ ಶಾಲಾ\\
          ತರಗತಿಗಳಿಗೆ\\
          ಉಪಯುಕ್ತವಾದ\\
          ಪಠ್ಯ ಪುಸ್ತಕ%
        };
      \end{scope}

      % 5. Outlined giant 3-line display title: ಕನ್ನಡ / ಬಾಲ / ವ್ಯಾಕರಣ
      \foreach \ang in {0,15,...,345} {
        \node[anchor=west, text=kbvoutline] at ([shift={(\ang:1.6pt)}]0.052\paperwidth, 0.840\paperheight)
          {\kbvdisplayfont\fontsize{64pt}{64pt}\selectfont ಕನ್ನಡ};
      }
      \node[anchor=west, text=white] at (0.052\paperwidth, 0.840\paperheight)
        {\kbvdisplayfont\fontsize{64pt}{64pt}\selectfont ಕನ್ನಡ};

      \foreach \ang in {0,15,...,345} {
        \node[anchor=west, text=kbvoutline] at ([shift={(\ang:1.6pt)}]0.052\paperwidth, 0.684\paperheight)
          {\kbvdisplayfont\fontsize{65pt}{65pt}\selectfont ಬಾಲ};
      }
      \node[anchor=west, text=white] at (0.052\paperwidth, 0.684\paperheight)
        {\kbvdisplayfont\fontsize{65pt}{65pt}\selectfont ಬಾಲ};

      \foreach \ang in {0,15,...,345} {
        \node[anchor=west, text=kbvoutline] at ([shift={(\ang:1.6pt)}]0.052\paperwidth, 0.512\paperheight)
          {\kbvdisplayfont\fontsize{63pt}{63pt}\selectfont ವ್ಯಾಕರಣ};
      }
      \node[anchor=west, text=white] at (0.052\paperwidth, 0.512\paperheight)
        {\kbvdisplayfont\fontsize{63pt}{63pt}\selectfont ವ್ಯಾಕರಣ};

      % 6. Author name on the right
      \node[anchor=east, text=white, font=\kbvcoversans\fontsize{15pt}{17pt}\selectfont]
        at (0.945\paperwidth, 0.286\paperheight) {ಕೆ. ಆರ್. ಕೆ. ಮೂರ್ತಿ};

      % 7. Publisher Emblem (Vahini Prakashana stylized peacock in framed box)
      \begin{scope}[shift={(0.106\paperwidth, 0.201\paperheight)}]
        \fill[kbvdarkred, draw=white, line width=0.65pt] (-4.3mm, -6.6mm) rectangle (4.3mm, 6.6mm);
        \fill[kbvplumegrey, draw=kbvplumedark, line width=0.35pt] (0.55mm, 0.35mm) circle (2.65mm);
        \foreach \a in {0,18,...,342} {
          \draw[kbvplumedark, line width=0.22pt] (0.55mm, 0.35mm) ++(\a:1.15mm) -- ++(\a:1.5mm);
        }
        \fill[kbvdarkred, draw=kbvplumedark, line width=0.35pt] (0.55mm, 0.35mm) circle (1.15mm);
        \fill[kbvplumegrey, draw=kbvplumedark, line width=0.3pt]
          (-0.5mm, 2.7mm) .. controls (1.2mm, 3.0mm) and (2.2mm, 2.5mm) .. (2.6mm, 1.8mm)
          .. controls (1.2mm, 1.2mm) and (-0.8mm, 0.6mm) .. (-1.2mm, -1.2mm)
          .. controls (-1.6mm, 0.5mm) and (-1.4mm, 2.0mm) .. cycle;
        \fill[white, draw=kbvplumedark, line width=0.35pt]
          (0.1mm, 3.4mm)
          .. controls (-2.4mm, 3.3mm) and (-3.5mm, 1.2mm) .. (-3.3mm, -1.0mm)
          .. controls (-3.1mm, -2.8mm) and (-1.2mm, -3.8mm) .. (0.3mm, -4.8mm)
          .. controls (0.5mm, -5.4mm) and (0.4mm, -6.0mm) .. (0.3mm, -6.1mm)
          .. controls (1.2mm, -5.0mm) and (1.4mm, -3.4mm) .. (0.7mm, -2.2mm)
          .. controls (-0.8mm, -1.5mm) and (-1.6mm, -0.2mm) .. (-1.5mm, 1.2mm)
          .. controls (-1.4mm, 2.1mm) and (-0.2mm, 2.3mm) .. (0.6mm, 2.6mm)
          .. controls (0.8mm, 3.0mm) and (0.6mm, 3.4mm) .. cycle;
        \fill[kbvplumedark] (-0.25mm, 2.85mm) ellipse (0.45mm and 0.28mm);
        \fill[white] (-0.25mm, 2.85mm) circle (0.14mm);
        \foreach \ca in {55, 75, 95, 115, 135} {
          \draw[white, line width=0.35pt] (-0.1mm, 3.5mm) -- ++(\ca:1.3mm);
          \fill[white] (-0.1mm, 3.5mm) ++(\ca:1.3mm) circle (0.22mm);
        }
      \end{scope}

      % 8. Publisher Name at bottom left (between Rule 5 and Rule 6)
      \node[anchor=west, text=white, font=\kbvcoversans\fontsize{14pt}{16pt}\selectfont]
        at (0.066\paperwidth, 0.111\paperheight) {ವಾಹಿನಿ ಪ್ರಕಾಶನ};
    \end{tikzpicture}%
  }%
}
\mbox{}
\clearpage


% === tex/0002.tex (Blank Page) ===
\thispagestyle{empty}
\mbox{}
\clearpage

% === tex/0003.tex ===
% === File 0003.md: Inner Title Page (matching full/0003.webp) ===
\thispagestyle{empty}
\vbox to 0.96\textheight{%
  \vspace*{12mm}
  \begin{center}
    {\kbvchapfont\fontsize{25pt}{30pt}\selectfont ಕನ್ನಡ ಬಾಲ ವ್ಯಾಕರಣ}
  \end{center}
  \vfill
  \begin{center}
    \begin{tikzpicture}
      \node[
        inner xsep=7mm,
        inner ysep=5.5mm,
        text width=82mm,
        align=center
      ] (tbox) {%
        \linespread{1.32}\fontsize{11.5pt}{15.5pt}\selectfont
        ಪ್ರಾಥಮಿಕ ಶಾಲೆಯ ಮತ್ತು ಪ್ರೌಢ ಶಾಲೆಯ\\
        ವಿದ್ಯಾರ್ಥಿಗಳಿಗೆ ಹೆಚ್ಚು ಉಪಯುಕ್ತವಾದ\\
        ಸರಳ ವ್ಯಾಕರಣ ಪಠ್ಯ ಪುಸ್ತಕ%
      };
      % Letterpress-style border rules with slight corner gaps matching full/0003.webp
      \draw[line width=0.9pt] ([xshift=1.2mm]tbox.north west) -- ([xshift=-1.2mm]tbox.north east);
      \draw[line width=0.9pt] ([xshift=1.2mm]tbox.south west) -- ([xshift=-1.2mm]tbox.south east);
      \draw[line width=0.9pt] ([yshift=-1.2mm]tbox.north west) -- ([yshift=1.2mm]tbox.south west);
      \draw[line width=0.9pt] ([yshift=-1.2mm]tbox.north east) -- ([yshift=1.2mm]tbox.south east);
    \end{tikzpicture}
  \end{center}
  \vspace{6mm}
  \begin{center}
    {\kbvsecfont\fontsize{13pt}{17pt}\selectfont ಕೆ. ಆರ್. ಕೆ. ಮೂರ್ತಿ}\\[3.5pt]
    {\fontsize{11pt}{15pt}\selectfont {[}ನಿವೃತ್ತ ಪ್ರೌಢಶಾಲಾ ಶಿಕ್ಷಕ{]}}
  \end{center}
  \vfill
  \noindent\rule{\textwidth}{0.95pt}\par
  \vspace{2.5mm}
  \begin{center}
    {\kbvchapfont\fontsize{14.5pt}{18pt}\selectfont ವಾಹಿನಿ ಪ್ರಕಾಶನ : : ಬೆಂಗಳೂರು-22}
  \end{center}
  \vspace{1.5mm}
  \noindent\rule{\textwidth}{0.95pt}\par
  \vspace*{4mm}
}
\clearpage


% === tex/0004.tex ===
% === File 0004.md: Copyright & Imprint Page (matching full/0004.webp) ===
\thispagestyle{empty}
\vbox to 0.96\textheight{%
  \vspace*{16mm}
  \noindent\hspace*{5mm}{\fontsize{11.5pt}{15.5pt}\selectfont ಈ ಗ್ರಂಥದ ಸ್ವಾಮ್ಯ ಲೇಖಕರಿಗೆ ಸೇರಿದ್ದು}\par
  \vspace{28mm}
  \noindent\hspace*{5mm}{\fontsize{11pt}{15pt}\selectfont ಪ್ರತಿಗಳ ಅಧಿಕೃತ ಮಾರಾಟಗಾರರು :}\\[4pt]
  \hspace*{13mm}{\kbvsecfont\fontsize{11.5pt}{15.5pt}\selectfont ಚಾಮುಂಡಿ ಸ್ಟೋರ್ಸ್}\\[2.5pt]
  \hspace*{13mm}{\fontsize{11pt}{15pt}\selectfont ಪುಸ್ತಕ ವ್ಯಾಪಾರಿಗಳು,}\\[2.5pt]
  \hspace*{13mm}{\fontsize{11pt}{15pt}\selectfont ಯಶವಂತಪುರ, ಬೆಂಗಳೂರು-560 022.}\par
  \vfill
  \noindent\hspace*{5mm}{\fontsize{11pt}{15pt}\selectfont ಮುದ್ರಣ :}\\[4pt]
  \hspace*{13mm}{\kbvsecfont\fontsize{11.5pt}{15.5pt}\selectfont ಪದ್ಮ ಪ್ರಿಂಟಿಂಗ್ ಪ್ರೆಸ್}\\[2.5pt]
  \hspace*{13mm}{\fontsize{11pt}{15pt}\selectfont ಬಳೇಪೇಟೆ, ಬೆಂಗಳೂರು-53.}\par
  \vspace*{6mm}
}
\clearpage


% === tex/0005.tex ===
% === File 0005.md: Table of Contents (matching full/0005.webp) ===
\thispagestyle{empty}
\vbox to 0.96\textheight{%
  \vspace*{2mm}
  \begin{center}
    {\kbvchapfont\fontsize{18pt}{22pt}\selectfont ವಿಷಯಾನುಕ್ರಮಣಿಕೆ}\\[3pt]
    {\color{kbvline}\rule{48mm}{0.5pt}}
  \end{center}
  \vspace{3mm}
  \noindent\hspace*{6mm}{\kbvsecfont\fontsize{11pt}{14pt}\selectfont ಪ್ರಕರಣ}\hfill{\kbvsecfont\fontsize{11pt}{14pt}\selectfont ಪುಟ ಸಂಖ್ಯೆ}\par
  \vspace{1.5mm}
  \noindent{\renewcommand{\arraystretch}{1.20}\setlength{\tabcolsep}{4pt}%
  \begin{tabularx}{\textwidth}{r X c r}
    {\kbvsecfont\color{kbvaccent}೧.} & ವಿಷಯ ಪ್ರವೇಶ & {\color{kbvgrey}\ldots\ldots} & {\kbvsecfont 5} \\
    {\kbvsecfont\color{kbvaccent}೨.} & ಅಕ್ಷರ ಮಾಲೆ & {\color{kbvgrey}\ldots\ldots} & {\kbvsecfont 7} \\
    {\kbvsecfont\color{kbvaccent}೩.} & ನಾಮಪದ & {\color{kbvgrey}\ldots\ldots} & {\kbvsecfont 12} \\
    {\kbvsecfont\color{kbvaccent}೪.} & ವಾಕ್ಯರಚನೆ & {\color{kbvgrey}\ldots\ldots} & {\kbvsecfont 14} \\
    {\kbvsecfont\color{kbvaccent}೫.} & ಕ್ರಿಯಾಪದ & {\color{kbvgrey}\ldots\ldots} & {\kbvsecfont 18} \\
    {\kbvsecfont\color{kbvaccent}೬.} & ಕನ್ನಡ ಸಂಧಿಗಳು & {\color{kbvgrey}\ldots\ldots} & {\kbvsecfont 21} \\
    {\kbvsecfont\color{kbvaccent}೭.} & ಸಂಸ್ಕೃತ ಸಂಧಿಗಳು & {\color{kbvgrey}\ldots\ldots} & {\kbvsecfont 25} \\
    {\kbvsecfont\color{kbvaccent}೮.} & ನಾಮಪದದ ವಿಭೇದಗಳು & {\color{kbvgrey}\ldots\ldots} & {\kbvsecfont 33} \\
    {\kbvsecfont\color{kbvaccent}೯.} & ನಾಮಪದಗಳ ರೂಪ ಬದಲಾವಣೆ ({\kbvlatin i}) & {\color{kbvgrey}\ldots\ldots} & {\kbvsecfont 39} \\
    {\kbvsecfont\color{kbvaccent}೧೦.} & ನಾಮಪದಗಳ ರೂಪ ಬದಲಾವಣೆ ({\kbvlatin ii}) & {\color{kbvgrey}\ldots\ldots} & {\kbvsecfont 43} \\
    {\kbvsecfont\color{kbvaccent}೧೧.} & ನಾಮಪದಗಳ ರೂಪ ಬದಲಾವಣೆ ({\kbvlatin iii}) & {\color{kbvgrey}\ldots\ldots} & {\kbvsecfont 46} \\
    {\kbvsecfont\color{kbvaccent}೧೨.} & ಕ್ರಿಯಾಪದ : ಲಿಂಗ, ವಚನ, ಕಾಲ & {\color{kbvgrey}\ldots\ldots} & {\kbvsecfont 52} \\
    {\kbvsecfont\color{kbvaccent}೧೩.} & ಸಮಾಸಗಳು & {\color{kbvgrey}\ldots\ldots} & {\kbvsecfont 62} \\
    {\kbvsecfont\color{kbvaccent}೧೪.} & ಕೃದಂತಗಳು & {\color{kbvgrey}\ldots\ldots} & {\kbvsecfont 70} \\
    {\kbvsecfont\color{kbvaccent}೧೫.} & ತದ್ಧಿತಾಂತಗಳು & {\color{kbvgrey}\ldots\ldots} & {\kbvsecfont 75} \\
    {\kbvsecfont\color{kbvaccent}೧೬.} & ಅವ್ಯಯಗಳು & {\color{kbvgrey}\ldots\ldots} & {\kbvsecfont 79} \\
    {\kbvsecfont\color{kbvaccent}೧೭.} & ದ್ವಿರುಕ್ತಿ & {\color{kbvgrey}\ldots\ldots} & {\kbvsecfont 83} \\
    {\kbvsecfont\color{kbvaccent}೧೮.} & ತತ್ಸಮ-ತದ್ಭವ & {\color{kbvgrey}\ldots\ldots} & {\kbvsecfont 84} \\
    {\kbvsecfont\color{kbvaccent}೧೯.} & ಲೇಖನ ಚಿಹ್ನೆಗಳು & {\color{kbvgrey}\ldots\ldots} & {\kbvsecfont 85} \\
    {\kbvsecfont\color{kbvaccent}೨೦.} & ಅಲಂಕಾರಗಳು & {\color{kbvgrey}\ldots\ldots} & {\kbvsecfont 90} \\
    {\kbvsecfont\color{kbvaccent}೨೧.} & ಛಂದಸ್ಸು & {\color{kbvgrey}\ldots\ldots} & {\kbvsecfont 97} \\
  \end{tabularx}}\par
  \vfill
  \begin{center}
    {\color{kbvline}\rule{28mm}{0.5pt}}
  \end{center}
  \vspace*{2mm}
}
\clearpage


% === tex/0006.tex (Blank Page) ===
\thispagestyle{empty}
\mbox{}
\clearpage

% === tex/0007.tex (Printed Page 5) ===
\renewcommand{\kbvchaptername}{ವಿಷಯ ಪ್ರವೇಶ}
\setcounter{page}{5}
\thispagestyle{kbvpage}
\kbvtypesetpage{\kbvbooktitle

\kbvchapter{ಪಾಠ ೧}{ವಿಷಯ ಪ್ರವೇಶ}

\indent ನಮ್ಮ ಮನಸ್ಸಿನಲ್ಲಿರುವ ಅಭಿಪ್ರಾಯವನ್ನು ಮತ್ತೊಬ್ಬರಿಗೆ ತಿಳಿಸಲು ನೆರವಾಗುವ ಒಂದು ಸಾಧನವೇ ಭಾಷೆ. ನಾವು ಆಡು ಭಾಷೆಯಲ್ಲಿ ಅಭಿಪ್ರಾಯವ್ಯಕ್ತಪಡಿಸುವ ಶಬ್ದರೂಪ ಮಾಧ್ಯಮವೇ ಮಾತು. ಉದಾಹರಣೆಗೆ, “ಕಮಲಳ ಮಗ ಬಸವ”-ಈ ಮಾತಿನ ಅರ್ಥವನ್ನು ನಾವು ಸುಲಭವಾಗಿ ಗ್ರಹಿಸಬಲ್ಲೆವು. ಕಮಲ ಯಾರು? ಬಸವ ಯಾರು? ಅವರಿಬ್ಬರಿಗಿರುವ ಸಂಬಂಧವೇನು? ಎಂಬುದು ನಮಗೆ ತಿಳಿದು ಬರುವುದು; ಇದು ಅರ್ಥಪೂರ್ಣವಾದ ಮಾತು. ಹೀಗೆ ಪೂರ್ಣವಾದ ಅರ್ಥವನ್ನು ಕೊಡುವ ಮಾತಿಗೆ ವಾಕ್ಯ ಎನ್ನುತ್ತೇವೆ.\par

\indent ಮೇಲೆ ಉದಾಹರಿಸಿದ ವಾಕ್ಯದಲ್ಲಿ ಮೂರು ಭಾಗಗಳಿವೆ ; ಕಮಲಳ ಮಗ ಬಸವ-ಇವೇ ಆ ಮೂರು ಭಾಗಗಳು. ಇವುಗಳಲ್ಲಿ ಯಾವುದೆ ಒಂದನ್ನು ಹೇಳಿದರೂ ಮಾತಿನ ಅರ್ಥ ಬರುವುದಿಲ್ಲ; ಇಂತಹ ಬಿಡಿ ಭಾಗಗಳಿಗೆ “ಪದ”ಗಳೆಂದು ಹೆಸರು.\par

\indent ಮೇಲಿನ ಮೂರು ಪದಗಳಲ್ಲಿ ಒಂದನ್ನು ಗಮನಿಸೋಣ ; ಉದಾಹರಣೆಗೆ, “ಬಸವನು” ಈ ಪದದಲ್ಲಿ ಬಿಡಿಯಾದ ಮೂರು ಅಕ್ಷರಗಳಿವೆ. ಅಕ್ಷರಮಾಲೆಯಲ್ಲಿ ಬ, ಸ, ವ ಎಂಬೀ ಅಕ್ಷರಗಳನ್ನು ಸುಲಭವಾಗಿ ಗುರ್ತಿಸಬಹುದು. ಹೀಗೆ ಅಕ್ಷರಗಳ ಸಮಂಜಸವಾದ ಜೋಡಣೆಯಿಂದಾಗಿರುವುದೇ “ಪದ”.\par

\indent ಮನಸ್ಸಿಗೆ ಬಂದಂತೆ ಪದಗಳನ್ನು ಕೂಡಿಸಿದರೆ ವಾಕ್ಯವಾಗುವುದಿಲ್ಲ. ಅಂತೆಯೇ ಮನಸ್ಸಿಗೆ ಬಂದಂತೆ ಅಕ್ಷರಗಳನ್ನು ಜೋಡಿಸಿದರೆ ಪದವಾಗುವುದಿಲ್ಲ.\par

\kbvexercise{ಅಭ್ಯಾಸ ಪತ್ರಿಕೆ ೧}

\vspace{2.5pt plus 3pt minus 0.5pt}\noindent\hspace*{1.2em}\kbvnum{(1)}\hspace{0.55em}ಎರಡು ಅಕ್ಷರಗಳ ಹತ್ತು ಪದಗಳನ್ನು ಬರೆಯಿರಿ.\par

\noindent\hspace*{1.2em}\kbvlabel{ಉದಾ}\hspace{0.6em}ಏತ, ಆಳ\par

\vspace{2.5pt plus 3pt minus 0.5pt}\noindent\hspace*{1.2em}\kbvnum{(2)}\hspace{0.55em}ಮೂರು ಅಕ್ಷರಗಳ ಹತ್ತು ಪದಗಳನ್ನು ಬರೆಯಿರಿ.\par

\noindent\hspace*{1.2em}\kbvlabel{ಉದಾ}\hspace{0.6em}ಅಗಸ, ಸರಳ\par
}
\clearpage

% === tex/0008.tex (Printed Page 6) ===
\renewcommand{\kbvchaptername}{ವಿಷಯ ಪ್ರವೇಶ}
\setcounter{page}{6}
\thispagestyle{kbvpage}
\kbvtypesetpage{\vspace{2.5pt plus 3pt minus 0.5pt}\noindent\hspace*{1.2em}\kbvnum{(3)}\hspace{0.55em}ಎರಡು ಪದಗಳುಳ್ಳ  ಐದು ವಾಕ್ಯಗಳನ್ನು ಜ್ಞಾಪಿಸಿರಿ.\par

\noindent\hspace*{1.2em}\kbvlabel{ಉದಾ}\hspace{0.6em}ಹೂವು ಅರಳಿತು...\par

\vspace{2.5pt plus 3pt minus 0.5pt}\noindent\hspace*{1.2em}\kbvnum{(4)}\hspace{0.55em}ಮೂರು ಪದಗಳುಳ್ಳ ಐದು ವಾಕ್ಯಗಳನ್ನು ಬರೆಯಿರಿ.\par

\noindent\hspace*{1.2em}\kbvlabel{ಉದಾ}\hspace{0.6em}ಬೆಕ್ಕು ಹಾಲು ಕುಡಿಯಿತು...\par

\vspace{2.5pt plus 3pt minus 0.5pt}\noindent\hspace*{1.2em}\kbvnum{(5)}\hspace{0.55em}ನಾಲ್ಕು ಪದಗಳುಳ್ಳ ಮೂರು ವಾಕ್ಯಗಳನ್ನು ಬರೆಯಿರಿ.\par

\vspace{2.5pt plus 3pt minus 0.5pt}\noindent\hspace*{1.2em}\kbvnum{(6)}\hspace{0.55em}ಈ ಪದಗಳನ್ನು ಸ್ಪಷ್ಟವಾಗಿ ಉಚ್ಚರಿಸಿರಿ : ಹಕ್ಕಿ, ಅಕ್ಕಿ, ಅನ್ನ, ಅಣ್ಣ, ಉಷ್ಣ, ತೀಕ್ಷ್ಣ, ಗಾನ, ಜ್ಞಾನ, ರಾಷ್ಟ್ರ, ಆಜ್ಞೆ, ನಕ್ಷತ್ರ, ಉತ್ಥಾನ, ಉತ್ಸವ, ಸಂಸ್ಕೃತ, ಅಷ್ಟಾವಕ್ರ.\par

\vspace{2.5pt plus 3pt minus 0.5pt}\noindent\hspace*{1.2em}\kbvnum{(7)}\hspace{0.55em}ದಿಹಳ, ರಸಳ, ನಿಜರ, ಜನನ, ವತಾರೆ, ಮರಿಪಳ, ನಿವಾಶರ, ಪಾಜಾತರಿ, ಯಲವಾಶಿ, ಮಣರಕವೆಂಟ  - ಈ ಪದಗಳಲ್ಲಿ ಅಕ್ಷರಗಳು ಕಲಸಿಹೋಗಿವೆ ; ಅರ್ಥವತ್ತಾದ ಪದಗಳಾಗುವಂತೆ ಅಕ್ಷರಗಳನ್ನು ಕ್ರಮವಾಗಿ ಜೋಡಿಸಿ ಬರೆಯಿರಿ.\par

\vspace{2.5pt plus 3pt minus 0.5pt}\noindent\hspace*{1.2em}\kbvnum{(8)}\hspace{0.55em}ಅ-ಪಟ್ಟಿಯಲ್ಲಿರುವ ಪದಗಳೊಂದಿಗೆ ಬ-ಪಟ್ಟಿಯಲ್ಲಿರುವ ಪದಗಳನ್ನು ಕೂಡಿಸಿ ಅರ್ಥಪೂರ್ಣವಾದ ವಾಕ್ಯಗಳನ್ನು ಬರೆಯಿರಿ.\par

\vspace{5pt plus 4pt minus 1pt}
\begin{center}
{\renewcommand{\arraystretch}{1.24}\setlength{\tabcolsep}{10pt}
\begin{tabular}{l c l}
\textbf{ಅ-ಪಟ್ಟಿ} & \hspace{12mm} & \textbf{ಬ-ಪಟ್ಟಿ} \\[3pt]
ಕೋಳಿ & \ldots\ldots & ಬೊಗಳಿತು \\
ಕತ್ತೆ & \ldots\ldots & ಘರ್ಜಿಸಿತು \\
ನಾಯಿ & \ldots\ldots & ಅರಚಿತು \\
ಕುದುರೆ & \ldots\ldots & ಕೂಗಿತು \\
ಸಿಂಹ & \ldots\ldots & ಕೆನೆಯಿತು \\
\end{tabular}}
\end{center}
\vspace{5pt plus 4pt minus 1pt}

\vspace{2.5pt plus 3pt minus 0.5pt}\noindent\hspace*{1.2em}\kbvnum{(9)}\hspace{0.55em}ಕೆಳಗೆ ಕೊಟ್ಟಿರುವ ವಾಕ್ಯಗಳಲ್ಲಿ ಪದಗಳು ಸ್ಥಾನಪಲ್ಲಟಗೊಂಡಿವೆ. ಅವುಗಳನ್ನು ಕ್ರಮವಾಗಿ ಜೋಡಿಸಿ ಬರೆಯಿರಿ.\par

\vspace{2.5pt plus 3pt minus 0.5pt}\noindent\hspace*{1.2em}\kbvnum{(೧)}\hspace{0.55em}ಜೋಳದ ಆಹಾರ ರುಚಿಯಾದ ರೊಟ್ಟಿ.\par

\vspace{2.5pt plus 3pt minus 0.5pt}\noindent\hspace*{1.2em}\kbvnum{(೨)}\hspace{0.55em}ಅತ್ಯುತ್ತಮವಾದ ಮಗುವಿಗೆ ತಾಯಿಯ ಹಾಲು ಆಹಾರ.\par

\vspace{2.5pt plus 3pt minus 0.5pt}\noindent\hspace*{1.2em}\kbvnum{(3)}\hspace{0.55em}ಕಡ್ಡಾಯವಾಗಿ ಕರ್ನಾಟಕದಲ್ಲಿರುವ ಕನ್ನಡವನ್ನು ಕನ್ನಡೇತರರು ಕಲಿಯಬೇಕು.\par
}
\clearpage

% === tex/0009.tex (Printed Page 7) ===
\renewcommand{\kbvchaptername}{ಅಕ್ಷರ ಮಾಲೆ}
\setcounter{page}{7}
\thispagestyle{kbvpage}
\kbvtypesetpage{\vspace{2.5pt plus 3pt minus 0.5pt}\noindent\hspace*{1.2em}\kbvnum{(10)}\hspace{0.55em}ಈ ಅಂಶಗಳನ್ನೊಳಗೊಂಡಂತೆ ಅರ್ಥವತ್ತಾದ ವಾಕ್ಯಗಳನ್ನು ರಚಿಸಿ :\par

\vspace{2.5pt plus 3pt minus 0.5pt}\noindent\hspace*{1.2em}\kbvnum{(೧)}\hspace{0.55em}ಕಬ್ಬು--ರಸ--ಸಿಹಿ\par

\vspace{2.5pt plus 3pt minus 0.5pt}\noindent\hspace*{1.2em}\kbvnum{(೨)}\hspace{0.55em}ರಾಮ--ರಾವಣ--ಸಂಹಾರ\par

\vspace{2.5pt plus 3pt minus 0.5pt}\noindent\hspace*{1.2em}\kbvnum{(೩)}\hspace{0.55em}ವ್ಯಾಸ--ಮಹಾಭಾರತ--ರಚನೆ\par

\vspace{2.5pt plus 3pt minus 0.5pt}\noindent\hspace*{1.2em}\kbvnum{(೪)}\hspace{0.55em}ಹಾಲು--ಮೊಸರು--ಬೆಣ್ಣೆ--ತುಪ್ಪ\par

\vspace{2.5pt plus 3pt minus 0.5pt}\noindent\hspace*{1.2em}\kbvnum{(೫)}\hspace{0.55em}ಐದು--ಮೂರು--ಎಂಟು\par

\kbvchapter{ಪಾಠ ೨}{ಅಕ್ಷರ ಮಾಲೆ}

\kbvsechead{ಕನ್ನಡ ಅಕ್ಷರ ಮಾಲೆ :}

\vspace{4pt plus 4pt minus 1pt}\begin{center}{\fontsize{11pt}{17pt}\selectfont \makebox[1.7em]{{\kbvsecfont ಅ}}\makebox[1.7em]{{\kbvsecfont ಆ}}\makebox[1.7em]{{\kbvsecfont ಇ}}\makebox[1.7em]{{\kbvsecfont ಈ}}\makebox[1.7em]{{\kbvsecfont ಉ}}\makebox[1.7em]{{\kbvsecfont ಊ}}\makebox[1.7em]{{\kbvsecfont ಋ}}\makebox[1.7em]{{\kbvsecfont ೠ}}\makebox[1.7em]{{\kbvsecfont ಎ}}\makebox[1.7em]{{\kbvsecfont ಏ}}\makebox[1.7em]{{\kbvsecfont ಐ}}\makebox[1.7em]{{\kbvsecfont ಒ}}\makebox[1.7em]{{\kbvsecfont ಓ}}\makebox[1.7em]{{\kbvsecfont ಔ}} \\
\makebox[1.7em]{ಂ}\makebox[1.7em]{(ಅಂ)}\makebox[1.7em]{;}\makebox[1.7em]{ಃ}\makebox[1.7em]{(ಅಃ)} \\
\makebox[1.7em]{ಕ}\makebox[1.7em]{ಖ}\makebox[1.7em]{ಗ}\makebox[1.7em]{ಘ}\makebox[1.7em]{ಙ} \\
\makebox[1.7em]{ಚ}\makebox[1.7em]{ಛ}\makebox[1.7em]{ಜ}\makebox[1.7em]{ಝ}\makebox[1.7em]{ಞ} \\
\makebox[1.7em]{ಟ}\makebox[1.7em]{ಠ}\makebox[1.7em]{ಡ}\makebox[1.7em]{ಢ}\makebox[1.7em]{ಣ} \\
\makebox[1.7em]{ತ}\makebox[1.7em]{ಥ}\makebox[1.7em]{ದ}\makebox[1.7em]{ಧ}\makebox[1.7em]{ನ} \\
\makebox[1.7em]{ಪ}\makebox[1.7em]{ಫ}\makebox[1.7em]{ಬ}\makebox[1.7em]{ಭ}\makebox[1.7em]{ಮ} \\
\makebox[1.7em]{ಯ}\makebox[1.7em]{ರ}\makebox[1.7em]{ಲ}\makebox[1.7em]{ವ}\makebox[1.7em]{ಶ}\makebox[1.7em]{ಷ}\makebox[1.7em]{ಸ}\makebox[1.7em]{ಹ}\makebox[1.7em]{ಳ}}\end{center}\vspace{4pt plus 4pt minus 1pt}

\begin{kbvrule}
\noindent\kbvbullet\hspace{0.5em}ಕನ್ನಡ ಅಕ್ಷರ ಮಾಲೆಯಲ್ಲಿರುವ ಬಿಡಿ ಅಕ್ಷರಗಳ ಸಂಖ್ಯೆ ಐವತ್ತು. ಅವುಗಳನ್ನು ಸ್ವರಗಳು, ವ್ಯಂಜನಗಳು ಮತ್ತು ಯೋಗವಾಹಗಳು ಎಂದು ಮೂರು ವಿಧವಾಗಿ ವಿಂಗಡಿಸಬಹುದು.
\end{kbvrule}

\kbvsechead{ಸ್ವರಗಳು :}

\begin{kbvrule}
\noindent\kbvbullet\hspace{0.5em}ಅ, ಆ, ಇ, ಈ, ಉ, ಊ, ಋ, ೠ, ಎ, ಏ, ಐ ಮತ್ತು ಒ, ಓ, ಔ ಸ್ವತಂತ್ರವಾಗಿ ಉಚ್ಚರಿಸಬಹುದಾದ ಅಕ್ಷರಗಳು ; ಇವುಗಳಿಗೆ ಸ್ವರಗಳೆಂದು ಹೆಸರು.
\end{kbvrule}
}
\clearpage

% === tex/0010.tex (Printed Page 8) ===
\renewcommand{\kbvchaptername}{ಅಕ್ಷರ ಮಾಲೆ}
\setcounter{page}{8}
\thispagestyle{kbvpage}
\kbvtypesetpage{\begin{kbvrule}
\noindent\kbvbullet\hspace{0.5em}ಅ-ಕಾರದಿಂದ ಔ-ಕಾರದವರೆಗಿನ ಸ್ವರಗಳ ಒಟ್ಟು ಸಂಖ್ಯೆ ಹದಿನಾಲ್ಕು. ಇವುಗಳಲ್ಲಿ ಅ, ಇ, ಉ, ಋ, ಎ, ಒಗಳು ಮೊಟಕಾಗಿ (ಅಲ್ಪಕಾಲ) ಉಚ್ಚರಿಸಲು ಬರುವ ಸ್ವರಗಳು. ಈ ಆರು ಸ್ವರಗಳಿಗೆ ಹ್ರಸ್ವ ಸ್ವರಗಳೆಂದು ಹೆಸರು.\par
\noindent\kbvbullet\hspace{0.5em}ಆ, ಈ, ಊ, ೠ, ಏ, ಐ, ಓ, ಔ- ಇವು ನೀಳವಾಗಿ (ದೀರ್ಘವಾಗಿ) ಉಚ್ಚರಿಸಲು ಬರುವ ಸ್ವರಗಳು. ಈ ಎಂಟು ಸ್ವರಗಳಿಗೆ ದೀರ್ಘ ಸ್ವರಗಳೆಂದು ಹೆಸರು.
\end{kbvrule}

\kbvsechead{ವ್ಯಂಜನಗಳು :}

\begin{kbvrule}
\noindent\kbvbullet\hspace{0.5em}ಕ-ಕಾರದಿಂದ ಳ-ಕಾರದವರೆಗಿನ ಅಕ್ಷರಗಳನ್ನು ಸ್ವತಂತ್ರವಾಗಿ ಉಚ್ಚರಿಸಲು ಸಾಧ್ಯವಿಲ್ಲ. ಸ್ವರಗಳು ಸೇರಿದಾಗ ಮಾತ್ರವೇ ಇವುಗಳನ್ನು ಉಚ್ಚರಿಸಲು ಸಾಧ್ಯ. ಈ ಜಾತಿಗೆ ಸೇರಿದ ಅಕ್ಷರಗಳಿಗೆ ವ್ಯಂಜನಗಳೆಂದು ಹೆಸರು. ಇವುಗಳ ಒಟ್ಟು ಸಂಖ್ಯೆ ಮೂವತ್ತನಾಲ್ಕು.\par
\noindent\kbvbullet\hspace{0.5em}ವ್ಯಂಜನಗಳು ಅರ್ಧಾಕ್ಷರಗಳು. ಉದಾಹರಣೆಗೆ, "ಕ" ವ್ಯಂಜನದ ಮೂಲ ರೂಪ "ಕ್"; ಇದರೊಂದಿಗೆ "ಅ" ಸೇರಿದಾಗ "ಕ" ಆಗುವುದು. ಕ್ + ಅ = ಕ. ಇದೇ ರೀತಿಯಾಗಿ ಚ್ + ಆ = ಚಾ, ಟ್ + ಇ = ಟಿ, ತ್ + ಈ = ತೀ, ಪ್ + ಉ = ಪು ಆಗುವುದು.\par
\noindent\kbvbullet\hspace{0.5em}ವ್ಯಂಜನಗಳಲ್ಲಿ ವರ್ಗೀಯ ವ್ಯಂಜನಗಳು ಮತ್ತು ಅವರ್ಗೀಯ ವ್ಯಂಜನಗಳು ಎಂದು ಎರಡು ವಿಧ.\par
\noindent\kbvbullet\hspace{0.5em}ವರ್ಗೀಯ ವ್ಯಂಜನಗಳು : ವ್ಯಂಜನಗಳಲ್ಲಿ ಮೊದಲ ಇಪ್ಪತ್ತೈದು ಅಕ್ಷರಗಳು ವರ್ಗೀಯ ವ್ಯಂಜನಗಳು. ಇವುಗಳನ್ನು ಐದು ವರ್ಗಗಳಾಗಿ ವಿಂಗಡಿಸಲಾಗಿದೆ; ಇಂತಹ ಪ್ರತಿಯೊಂದು ವರ್ಗದಲ್ಲಿರುವ ಅಕ್ಷರಗಳ ಸಂಖ್ಯೆ ಐದು.
\end{kbvrule}

\vspace{5pt plus 4pt minus 1pt}
\begin{center}
{\renewcommand{\arraystretch}{1.24}\setlength{\tabcolsep}{10pt}
\begin{tabular}{l c l}
ಕ ಖ ಗ ಘ ಙ & \ldots\ldots & ಕ-ವರ್ಗ \\
ಚ ಛ ಜ ಝ ಞ & \ldots\ldots & ಚ-ವರ್ಗ \\
ಟ ಠ ಡ ಢ ಣ & \ldots\ldots & ಟ-ವರ್ಗ \\
ತ ಥ ದ ಧ ನ & \ldots\ldots & ತ-ವರ್ಗ \\
ಪ ಫ ಬ ಭ ಮ & \ldots\ldots & ಪ-ವರ್ಗ \\
\end{tabular}}
\end{center}
\vspace{5pt plus 4pt minus 1pt}
}
\clearpage

% === tex/0011.tex (Printed Page 9) ===
\renewcommand{\kbvchaptername}{ಅಕ್ಷರ ಮಾಲೆ}
\setcounter{page}{9}
\thispagestyle{kbvpage}
\kbvtypesetpage{\begin{kbvrule}
\noindent\kbvbullet\hspace{0.5em}ವರ್ಗೀಯ ವ್ಯಂಜನಗಳನ್ನು ಉಚ್ಚಾರದ ಆಧಾರದ ಮೇಲೆ ಅಲ್ಪಪ್ರಾಣ, ಮಹಾಪ್ರಾಣ, ಅನುನಾಸಿಕ ಎಂಬುದಾಗಿ ಮೂರು ಬಗೆಯಾಗಿ ವಿಂಗಡಿಸಬಹುದು.
\end{kbvrule}

\vspace{2.5pt plus 3pt minus 0.5pt}\noindent\hspace*{1.2em}\kbvnum{(೧)}\hspace{0.55em}“ಕಚಟ ತಪ” ಮತ್ತು “ಗಜಡ ದಬ” ಇವು ತೇಲಿಸಿ-ಉಚ್ಚರಿಸ ಸುವ ವ್ಯಂಜನಗಳು ; ಇವುಗಳಿಗೆ “ಅಲ್ಪಪ್ರಾಣ” ಗಳೆಂದು ಹೆಸರು.\par

\vspace{2.5pt plus 3pt minus 0.5pt}\noindent\hspace*{1.2em}\kbvnum{(೨)}\hspace{0.55em}“ಖಛಠ ಥಫ” ಮತ್ತು “ಘ ಝ ಢ ಧ ಭ” - ಇವು ಹೆಚ್ಚು ಒತ್ತು ಕೊಟ್ಟು ಉಚ್ಚರಿಸುವ ವ್ಯಂಜನಗಳು ; ಇವುಗಳಿಗೆ “ಮಹಾಪ್ರಾಣ” ಗಳೆಂದು ಹೆಸರು.\par

\vspace{2.5pt plus 3pt minus 0.5pt}\noindent\hspace*{1.2em}\kbvnum{(೩)}\hspace{0.55em}ಙ ಞ ಣ ನ ಮ - ಇವು ಮೂಗಿನ ಸಹಾಯದಿಂದ ಉಚ್ಚರಿಸ ಸುವ ವ್ಯಂಜನಗಳು ; ಇವುಗಳಿಗೆ “ಅನುನಾಸಿಕ” ಗಳೆಂದು ಹೆಸರು.\par

\begin{kbvrule}
\noindent\kbvbullet\hspace{0.5em}ಅವರ್ಗೀಯ ವ್ಯಂಜನಗಳು : ಯ, ರ, ಲ, ವ, ಶ, ಷ, ಸ, ಹ, ಳ - ಇವು ಅವರ್ಗೀಯ ವ್ಯಂಜನಗಳು ; ಯ-ಕಾರದಿಂದ ಳ-ಕಾರದ ವರೆಗಿನ ಈ ವ್ಯಂಜನಗಳ ಒಟ್ಟು ಸಂಖ್ಯೆ ಒಂಭತ್ತು.
\end{kbvrule}

\kbvsechead{ಯೋಗವಾಹಗಳು :}

\indent ಸ್ವರಾಕ್ಷರದ ಬಲಬದಿಯಲ್ಲಿ ಒಂದು ಸೊನ್ನೆ ಬಂದರೆ “ಅನುಸ್ವಾರ” ವೆನಿಸುವುದು ; ಮತ್ತು ಒಂದರ ಮೇಲೊಂದರಂತೆ ಎರಡು ಸೊನ್ನೆಗಳು ಬಂದರೆ “ವಿಸರ್ಗ”ವೆನಿಸುವುದು. ಈ ಎರಡು ಪ್ರಭೇದಗಳನ್ನು ಯೋಗವಾಹಗಳೆನ್ನುವರು.\par

\indent ಅನುಸ್ವಾರ (ಂ) ; ವಿಸರ್ಗ (ಃ)\par

\indent ಉದಾಹರಣೆಗಳು : ಅಂ ; - ಅಃ\par

\indent ಯೋಗವಾಹಗಳನ್ನು ಸ್ವರಗಳ ಸಹಯೋಗದೊಂದಿಗೆ ಮಾತ್ರ ಉಚ್ಚರಿಸ ಬಹುದಾದ ಕಾರಣ ಅವುಗಳಿಗೆ ಆ ಹೆಸರು ಬಂದಿದೆ.\par

\kbvsechead{ಗುಣಿತಾಕ್ಷರಗಳು :}

\indent ಮೂಲ ವ್ಯಂಜನಗಳು ಅರ್ಧಾಕ್ಷರಗಳು. ಇವುಗಳ ಉಚ್ಚಾರ ಕಠಿಣ ಹಾಗೂ ಅಸ್ಪಷ್ಟ. ಈ ಅಪೂರ್ಣ ಅಕ್ಷರಗಳೊಂದಿಗೆ ಸ್ವರಗಳು ಸೇರಿದಾಗ ಸ್ಪಷ್ಟ ವಾಗಿ ಉಚ್ಚರಿಸಬಹುದಾದ ಹೊಸ ಅಕ್ಷರಗಳು ಲಭಿಸುತ್ತವೆ... ಹೀಗೆ ರೂಪ ತಳೆಯುವ ವ್ಯಂಜನಗಳಿಗೆ “ಗುಣಿತಾಕ್ಷರ” ಗಳೆಂದು ಹೆಸರು.\par
}
\clearpage

% === tex/0012.tex (Printed Page 10) ===
\renewcommand{\kbvchaptername}{ಅಕ್ಷರ ಮಾಲೆ}
\setcounter{page}{10}
\thispagestyle{kbvpage}
\kbvtypesetpage{\begin{kbvderiv}
\vspace{3pt plus 3pt minus 1pt}
\hspace*{1.8em}ಉದಾಹರಣೆ ; ಗ್‌ + ಅ = ಗ, ಗ್‌ + ಆ = ಗಾ, ಗ್‌ + ಇ = ಗಿ, \\
\hspace*{1.8em}ಗ್‌ + ಈ = ಗೀ ಇತ್ಯಾದಿ. ಇಲ್ಲಿ ಗ್‌ ಅಗ್ರ ಗುಣಿತಾಕ್ಷರಗಳು,
\vspace{3pt plus 3pt minus 1pt}
\end{kbvderiv}

\indent ಈ ರೀತಿಯಾಗಿ ಮೂವತ್ತನಾಲ್ಕು ವ್ಯಂಜನಗಳಿಗೂ ಸ್ವರಗಳನ್ನು ಸೇರಿಸಿ ಗುಣಿತಾಕ್ಷರಗಳನ್ನು ಬರೆಯಬಹುದು. ಬರವಣಿಗೆಯಲ್ಲಿ ಸ್ವರಗಳ ಬದಲು ಸಂಕೇತಗಳನ್ನು ವ್ಯಂಜನಗಳಿಗೆ ಸೇರಿಸಿ ಗುಣಿತಾಕ್ಷರಗಳನ್ನು ಬರೆಯುವುದು ಪದ್ಧತಿ. ಕೆಳಗೆ ಕೊಟ್ಟಿರುವ ಪಟ್ಟಿಯಲ್ಲಿ “ಕ್‌” ವ್ಯಂಜನದ ಗುಣಿತಾಕ್ಷರಗಳನ್ನು ಮಾದರಿಗಾಗಿ ಕೊಟ್ಟಿದೆ.\par

\kbvsechead{ಕಾಗುಣಿತ (ಕ್‌ ವ್ಯಂಜನದ ಗುಣಿತಾಕ್ಷರಗಳು) :}

\vspace{5pt plus 4pt minus 1pt}
\begin{center}
{\renewcommand{\arraystretch}{1.20}\setlength{\tabcolsep}{10pt}
\begin{tabular}{c | c | l}
\toprule
{\kbvsecfont ಸ್ವರ} & {\kbvsecfont ಸ್ವರ ಸಂಕೇತ} & {\kbvsecfont ಗುಣಿತಾಕ್ಷರ} \\
\midrule
ಅ & ಾ (ತಲೆಕಟ್ಟು) & ಕ್‌ + ಅ = ಕ \\
ಆ & ಾ & ಕ್‌ + ಆ = ಕಾ \\
ಇ & ಿ & ಕ್‌ + ಇ = ಕಿ \\
ಈ & ೀ & ಕ್‌ + ಈ = ಕೀ \\
ಉ & ು & ಕ್‌ + ಉ = ಕು \\
ಊ & ೂ & ಕ್‌ + ಊ = ಕೂ \\
ಋ & ೃ & ಕ್‌ + ಋ = ಕೃ \\
ೠ & ೄ & ಕ್‌ + ೠ = ಕೄ \\
ಎ & ೆ & ಕ್‌ + ಎ = ಕೆ \\
ಏ & ೇ & ಕ್‌ + ಏ = ಕೇ \\
ಐ & ೈ & ಕ್‌ + ಐ = ಕೈ \\
ಒ & ೊ & ಕ್‌ + ಒ = ಕೊ \\
ಓ & ೋ & ಕ್‌ + ಓ = ಕೋ \\
ಔ & ೌ & ಕ್‌ + ಔ = ಕೌ \\
ಅಂ & ಂ & ಕ್‌ + ಅಂ = ಕಂ \\
ಅಃ & ಃ & ಕ್‌ + ಅಃ = ಕಃ \\
\bottomrule
\end{tabular}}
\end{center}
\vspace{5pt plus 4pt minus 1pt}
}
\clearpage

% === tex/0013.tex (Printed Page 11) ===
\renewcommand{\kbvchaptername}{ಅಕ್ಷರ ಮಾಲೆ}
\setcounter{page}{11}
\thispagestyle{kbvpage}
\kbvtypesetpage{\kbvsechead{ಒತ್ತಕ್ಷರಗಳು :}

\indent ಎರಡು ಅಥವಾ ಮೂರು ವ್ಯಂಜನಗಳನ್ನು ಒತ್ತಿ ಕೂಡಿಸಿ ಒಂದಕ್ಷರವಾಗಿಸಿ ಉಚ್ಚರಿಸಬಹುದು. ಈ ಬಗೆಯ ಸಂಯುಕ್ತಾಕ್ಷರಗಳಿಗೆ ಒತ್ತಕ್ಷರಗಳೆಂದು ಹೆಸರು.\par

\indent ಉದಾಹರಣೆಗೆ,\par

\begin{kbvderiv}
\vspace{4pt plus 3pt minus 1pt}
\begin{center}\begin{adjustbox}{max width=\textwidth}
{\renewcommand{\arraystretch}{1.22}\setlength{\tabcolsep}{4pt}
\begin{tabular}{l c l}
ಟ್ + ಟ್ + (ಅ) & {\color{kbvgrey}$=$} & \kbvresult{ಟ್ಟ} \\
ಸ್ + ತ್ + ರ್ + (ಅ) & {\color{kbvgrey}$=$} & \kbvresult{ಸ್ತ್ರ} \\
\end{tabular}}
\end{adjustbox}\end{center}
\vspace{4pt plus 3pt minus 1pt}
\end{kbvderiv}

\begin{kbvrule}
\noindent\kbvbullet\hspace{0.5em}ಒಂದೇ ಜಾತಿಯ ವ್ಯಂಜನಗಳು ಸೇರಿ ಆಗುವ ಒತ್ತಕ್ಷರಗಳಿಗೆ ಸಜಾತೀಯ ಒತ್ತಕ್ಷರಗಳೆಂದು ಹೆಸರು.
\end{kbvrule}

\vspace{4pt plus 3pt minus 1pt}
\begin{center}
{\renewcommand{\arraystretch}{1.28}\setlength{\tabcolsep}{11pt}
\begin{tabular}{c c c c c}
ಕ್ಕ & ಖ್ಖ & ಗ್ಗ & ಘ್ಘ & ಙ್ಙ \\
ಚ್ಚ & ಛ್ಛ & ಜ್ಜ & ಝ್ಝ & ಞ್ಞ \\
ಟ್ಟ & ಠ್ಠ & ಡ್ಡ & ಢ್ಢ & ಣ್ಣ \\
ತ್ತ & ಥ್ಥ & ದ್ದ & ಧ್ಧ & ನ್ನ \\
ಪ್ಪ & ಫ್ಫ & ಬ್ಬ & ಭ್ಭ & ಮ್ಮ \\
\end{tabular}}
\end{center}
\vspace{2pt}

\begin{center}{\fontsize{11pt}{15pt}\selectfont ಯ್ಯ\hspace{1.1em}ರ್ರ\hspace{1.1em}ಲ್ಲ\hspace{1.1em}ವ್ವ\hspace{1.1em}ಶ್ಶ\hspace{1.1em}ಷ್ಷ\hspace{1.1em}ಸ್ಸ\hspace{1.1em}ಹ್ಹ\hspace{1.1em}ಳ್ಳ}\end{center}\vspace{3pt}

\noindent\hspace*{2.2em}—ಇವು ಸಜಾತೀಯ ಒತ್ತಕ್ಷರಗಳು.\par

\begin{kbvrule}
\noindent\kbvbullet\hspace{0.5em}ವಿಭಿನ್ನ ಜಾತಿಯ ವ್ಯಂಜನಗಳು ಸೇರಿ ಆಗುವ ಒತ್ತಕ್ಷರಗಳಿಗೆ ವಿಜಾತೀಯ ಒತ್ತಕ್ಷರಗಳೆಂದು ಹೆಸರು.
\end{kbvrule}

\vspace{3pt plus 2pt}\noindent\hspace*{1.2em}\kbvlabel{ಉದಾ}\par\vspace{1pt}

\begin{kbvderiv}
\vspace{4pt plus 3pt minus 1pt}
\noindent{\renewcommand{\arraystretch}{1.24}\setlength{\tabcolsep}{5pt}
\begin{tabularx}{\textwidth}{X X}
ಬ್ + ಜ್ = (ಬ್ಜ) & ಸ್ + ತ್ + ರ್ = (ಸ್ತ್ರ) \\
ಕ್ + ತ್ = (ಕ್ತ) & ಷ್ + ಟ್ + ರ್ = (ಷ್ಟ್ರ) \\
\end{tabularx}}
\vspace{4pt plus 3pt minus 1pt}
\end{kbvderiv}

\kbvexercise{ಅಭ್ಯಾಸ ಪತ್ರಿಕೆ ೨}

\vspace{2.5pt plus 3pt minus 0.5pt}\noindent\hspace*{1.2em}\kbvnum{(1)}\hspace{0.55em}ಕನ್ನಡ ಅಕ್ಷರಮಾಲೆಯಲ್ಲಿರುವ ಅಕ್ಷರಗಳ ಸಂಖ್ಯೆಯೆಷ್ಟು?\par

\vspace{2.5pt plus 3pt minus 0.5pt}\noindent\hspace*{1.2em}\kbvnum{(2)}\hspace{0.55em}ಸ್ವರ ಎಂದರೇನು? ಸ್ವರಗಳಲ್ಲಿ ಎಷ್ಟು ವಿಧ?\par

\vspace{2.5pt plus 3pt minus 0.5pt}\noindent\hspace*{1.2em}\kbvnum{(3)}\hspace{0.55em}ಅನುನಾಸಿಕ ಎಂದರೇನು? ಉದಾಹರಣೆಗಳ ಸಹಿತ ತಿಳಿಸಿ.\par
}
\clearpage

% === tex/0014.tex (Printed Page 12) ===
\renewcommand{\kbvchaptername}{ನಾಮಪದ}
\setcounter{page}{12}
\thispagestyle{kbvpage}
\kbvtypesetpage{\vspace{2.5pt plus 3pt minus 0.5pt}\noindent\hspace*{1.2em}\kbvnum{(4)}\hspace{0.55em}ಕೆಳಗೆ ಕೊಟ್ಟಿರುವ ಪದಗಳಲ್ಲಿ ಕಾಣಸಿಗುವ ವರ್ಗೀಯ ವ್ಯಂಜನ, ಅವರ್ಗೀಯ ವ್ಯಂಜನ ಮತ್ತು ಯೋಗವಾಹಗಳನ್ನು ಗುರ್ತಿಸಿ:

\vspace{5pt plus 4pt minus 1pt}
\noindent{\renewcommand{\arraystretch}{1.24}\setlength{\tabcolsep}{4pt}
\begin{tabularx}{\textwidth}{X X X X X}
ಅಗಸ & ಅಂಗಳ & ಆಸನ & ಆಭರಣ & ಔನಃ \\
ಉರಗ & ಓಲಗ & ಐದಳೆ & ಮದಗಜ & ಆಂತಃ \\
\end{tabularx}}
\vspace{5pt plus 4pt minus 1pt}

\vspace{2.5pt plus 3pt minus 0.5pt}\noindent\hspace*{1.2em}\kbvnum{(5)}\hspace{0.55em}ಗುಣಿತಾಕ್ಷರ ಎಂದರೇನು? “ಚ್” ವ್ಯಂಜನದ ಗುಣಿತಾಕ್ಷರಗಳ ಬರೆದು ತಿಳಿಸಿ.\par

\vspace{2.5pt plus 3pt minus 0.5pt}\noindent\hspace*{1.2em}\kbvnum{(6)}\hspace{0.55em}ಸಜಾತೀಯ ಒತ್ತಕ್ಷರವುಳ್ಳ ನಾಲ್ಕು ಪದಗಳನ್ನು ಬರೆಯಿರಿ.\par

\vspace{2.5pt plus 3pt minus 0.5pt}\noindent\hspace*{1.2em}\kbvnum{(7)}\hspace{0.55em}ವಿಜಾತೀಯ ಒತ್ತಕ್ಷರವುಳ್ಳ ನಾಲ್ಕು ಪದಗಳನ್ನು ಬರೆಯಿರಿ.\par

\vspace{2.5pt plus 3pt minus 0.5pt}\noindent\hspace*{1.2em}\kbvnum{(8)}\hspace{0.55em}ಗ್ಗ, ಜ್ರ, ನ್ಯ, ಶ್ಚ, ರೃ, ಸ್ಹ, ದೃ, ಜ್ವ—ಈ ಒತ್ತ ಗನ್ನು ಬಿಡಿಸಿ ಬರೆಯಿರಿ.\par

\vspace{2.5pt plus 3pt minus 0.5pt}\noindent\hspace*{1.2em}\kbvnum{(9)}\hspace{0.55em}“ಸೀತಾಪತಿ” -ಈ ಶಬ್ದದಲ್ಲಿರುವ ಸ್ವರ, ವ್ಯಂಜನಗಳ ಬಿಡಿಸಿ ಬರೆಯಿರಿ.\par

\vspace{2.5pt plus 3pt minus 0.5pt}\noindent\hspace*{1.2em}\kbvnum{(10)}\hspace{0.55em}ಕ ಕಾರದಿಂದ ಳ ಕಾರದವರೆಗಿನ ಎಲ್ಲ ವ್ಯಂಜನಗಳ ಗುಣಿತಾಕ್ಷರಗಳನ್ನೊಳಗೊಂಡ ಪಟ್ಟಿ ಸಿದ್ಧಪಡಿಸಿ (ಯಾವುದೇ ಅಕ್ಷರವನ್ನು ಹೆಸರಿಸುವಾಗ ಅದರ ಮುಂದೆ “ಕಾರ” ಎಂಬ ಪ್ರತ್ಯಯ ಸೇರಿಸುವುದು ರೂಢಿಯಲ್ಲಿದೆ.)\par

\kbvchapter{ಪಾಠ ೩}{ನಾಮಪದ}

\noindent\hspace*{2.2em}ಕೋತಿ ಲಾಗ ಹಾಕಿತು.\par

\noindent\hspace*{2.2em}ಮೇಜು ಬಹಳ ದೊಡ್ಡದು.\par

\noindent\hspace*{2.2em}ಬೆಂಗಳೂರು ದೊಡ್ಡ ವಾಣಿಜ್ಯ ಕೇಂದ್ರ.\par

\noindent\hspace*{2.2em}ರಾಮಯ್ಯ ಈ ಕಡೆ ಬರುತ್ತಿದ್ದಾನೆ.\par

\indent ಈ ವಾಕ್ಯಗಳಲ್ಲಿ ದಪ್ಪ ಅಕ್ಷರಗಳಲ್ಲಿರುವ ಪದಗಳನ್ನು ಚೆನ್ನಾಗಿ ಗಮನಿಸಿ-- ಕೋತಿ - ಒಂದು ಪ್ರಾಣಿ ; ಮೇಜು - ಒಂದು ವಸ್ತು ; ಬೆಂಗಳೂರು - ಒಂದು  ನಗರ ; ರಾಮಯ್ಯ - ಒಬ್ಬ ವ್ಯಕ್ತಿ. ಹೀಗೆ ಪ್ರಾಣಿ, ವಸ್ತು, ಸ್ಥಳ, ಮನುಷ್ಯ ಇವೇ ಮೊದಲಾದವುಗಳನ್ನು ಹೆಸರಿಸುವ ಪದಗಳಿಗೆ “ನಾಮ ಪ್ರಕೃತಿ”ಗಳೆಂದು  ಹೆಸರು.\par
}
\clearpage

% === tex/0015.tex (Printed Page 13) ===
\renewcommand{\kbvchaptername}{ನಾಮಪದ}
\setcounter{page}{13}
\thispagestyle{kbvpage}
\kbvtypesetpage{\kbvsechead{ನಾಮ ಪ್ರಕೃತಿ :}

\indent ಯಾವುದೇ ವ್ಯಕ್ತಿಯ ಹೆಸರನ್ನಾಗಲೀ, ವಸ್ತುವಿನ ಹೆಸರನ್ನಾಗಲೀ, ಸ್ಥಳದ ಹೆಸರನ್ನಾಗಲೀ, ಸಸ್ಯದ ಹೆಸರನ್ನಾಗಲೀ, ಪಶುಪಕ್ಷಿಯ ಹೆಸರನ್ನಾಗಲೀ ಸೂಚಿಸುವ ಪದಕ್ಕೆ “ನಾಮ ಪ್ರಕೃತಿ” ಎಂದು ಹೆಸರು.\par

\noindent\hspace*{1.2em}\kbvlabel{ಉದಾ}\hspace{0.6em}ಲಕ್ಷ್ಮಣ, ಮಾಲಿನಿ, ಕಬ್ಬಿಣ, ಕಲ್ಲು, ಕಾಡು, ಊರು, ಗುಲಾಬಿ, ಬಾಳೆ, ಎಮ್ಮೆ, ಪಾರಿವಾಳ ಇವು ನಾಮ ಪ್ರಕೃತಿಗಳು.\par

\kbvsechead{ಪ್ರತ್ಯಯ :}

\noindent\hspace*{2.2em}ಹೊಲವು ಊರಿಗೆ ಸಮೀಪದಲ್ಲಿದೆ.\par

\noindent\hspace*{2.2em}ಹೊಲದಲ್ಲಿ ಕಸಕಡ್ಡಿ ತುಂಬಿದೆ.\par

\noindent\hspace*{2.2em}ರೈತನು ಹೊಲವನ್ನು ಉಳುತ್ತಿದ್ದಾನೆ.\par

\noindent\hspace*{2.2em}ಹೊಲದ ಅಂಚಿನಲ್ಲಿ ಹುತ್ತವಿದೆ.\par

\indent ಈ ವಾಕ್ಯಗಳಲ್ಲಿ ದಪ್ಪ ಅಕ್ಷರಗಳಲ್ಲಿರುವ ಹೊಲವು, ಹೊಲದಲ್ಲಿ, ಹೊಲ\par

\indent ವನ್ನು, ಹೊಲದ-- ಪದಗಳನ್ನು ಗಮನಿಸಿ. ಇಲ್ಲಿ “ಹೊಲ” ಎಂಬ ನಾಮ ಪ್ರಕೃತಿಯನ್ನು ಸಂದರ್ಭೋಚಿತ ರೀತಿಯಲ್ಲಿ ಬಳಸಿರುವುದು ಕಂಡುಬರುತ್ತದೆ.\par

\noindent\hspace*{2.2em}ಈಗ ಆ ಪದಗಳನ್ನು ಬಿಡಿಸಿ ನೋಡೋಣ.\par

\begin{kbvderiv}
\vspace{4pt plus 3pt minus 1pt}
\begin{center}\begin{adjustbox}{max width=\textwidth}
{\renewcommand{\arraystretch}{1.22}\setlength{\tabcolsep}{4pt}
\begin{tabular}{l c l}
ಹೊಲ + ಉ & {\color{kbvgrey}$=$} & ಹೊಲವು (ವ್ + ಉ = ವು) \\
ಹೊಲ + ಅಲ್ಲಿ & {\color{kbvgrey}$=$} & \kbvresult{ಹೊಲದಲ್ಲಿ} \\
ಹೊಲ + ಅನ್ನು & {\color{kbvgrey}$=$} & \kbvresult{ಹೊಲವನ್ನು} \\
ಹೊಲ + ದ & {\color{kbvgrey}$=$} & \kbvresult{ಹೊಲದ} \\
\end{tabular}}
\end{adjustbox}\end{center}
\vspace{4pt plus 3pt minus 1pt}
\end{kbvderiv}

\indent ಮೇಲಿನ ನಾಲ್ಕು ಸಂದರ್ಭಗಳಲ್ಲಿಯೂ ಹೊಲ ಎಂಬ ನಾಮ ಪ್ರಕೃತಿ ಯೊಂದಿಗೆ ಕ್ರಮವಾಗಿ ಉ, ಅಲ್ಲಿ, ಅನ್ನು, ಆ ಸೇರಿರುವುದನ್ನು ಕಾಣಬಹುದು ; ಇವೇ “ಪ್ರತ್ಯಯಗಳು” ಎಂದು.\par

\begin{kbvrule}
\noindent\kbvbullet\hspace{0.5em}ಪ್ರಕೃತಿಯೊಂದಿಗೆ ಪ್ರತ್ಯಯ ಸೇರಿ ಪದವಾಗುವುದು. ಆದುದ ರಿಂದ ಹೊಲವು, ಹೊಲದಲ್ಲಿ, ಹೊಲವನ್ನು, ಹೊಲದ ಎಂಬ ಪದಗಳು ನಾಮ ಪದಗಳಾಗಿರುವುವು.
\end{kbvrule}

\begin{kbvrule}
\noindent{\kbvsecfont ವ್ಯಾಖ್ಯೆ :}\hspace{0.5em}ನಾಮ ಪ್ರಕೃತಿಗಳ ಮುಂದೆ ಪ್ರತ್ಯಯಗಳು ಸೇರಿ ನಾಮ ಪದಗಳೆನಿಸುವುವು.
\end{kbvrule}
}
\clearpage

% === tex/0016.tex (Printed Page 14) ===
\renewcommand{\kbvchaptername}{ವಾಕ್ಯರಚನೆ}
\setcounter{page}{14}
\thispagestyle{kbvpage}
\kbvtypesetpage{\indent ಉದಾಹರಣೆಗೆ,\par

\begin{kbvderiv}
\vspace{4pt plus 3pt minus 1pt}
\begin{center}\begin{adjustbox}{max width=\textwidth}
{\renewcommand{\arraystretch}{1.22}\setlength{\tabcolsep}{4pt}
\begin{tabular}{l c l}
(ನಾಮ ಪ್ರಕೃತಿ) + (ಪ್ರತ್ಯಯ) & {\color{kbvgrey}$=$} & \kbvresult{(ನಾಮಪದ)} \\
ಲಕ್ಷ್ಮಣ + ಉ & {\color{kbvgrey}$=$} & \kbvresult{ಲಕ್ಷ್ಮಣನು} \\
ರಾಮ + ಅ & {\color{kbvgrey}$=$} & \kbvresult{ರಾಮನ} \\
ಸೀತೆ + ಅನ್ನು & {\color{kbvgrey}$=$} & \kbvresult{ಸೀತೆಯನ್ನು} \\
ಅರಣ್ಯ + ಅಲ್ಲಿ & {\color{kbvgrey}$=$} & \kbvresult{ಅರಣ್ಯದಲ್ಲಿ} \\
\end{tabular}}
\end{adjustbox}\end{center}
\vspace{4pt plus 3pt minus 1pt}
\end{kbvderiv}

\kbvexercise{ಅಭ್ಯಾಸ ಪತ್ರಿಕೆ ೩}

\vspace{2.5pt plus 3pt minus 0.5pt}\noindent\hspace*{1.2em}\kbvnum{(1)}\hspace{0.55em}ನಾಮಪದ ಎಂದರೇನು ? ಉದಾಹರಣೆಗಳ ಸಹಿತ ತಿಳಿಸಿ.\par

\vspace{2.5pt plus 3pt minus 0.5pt}\noindent\hspace*{1.2em}\kbvnum{(2)}\hspace{0.55em}ನಾಮ ಪ್ರಕೃತಿಗೂ ನಾಮಪದಕ್ಕೂ ಇರುವ ವ್ಯತ್ಯಾಸವೇನು?\par

\vspace{2.5pt plus 3pt minus 0.5pt}\noindent\hspace*{1.2em}\kbvnum{(3)}\hspace{0.55em}ಕೆಳಗೆ ಕೊಟ್ಟಿರುವ ವಾಕ್ಯಗಳಲ್ಲಿ ನಾಮಪದಗಳನ್ನು ಗುರ್ತಿಸಿ ಮತ್ತು ಪ್ರತಿಯೊಂದು ನಾಮಪದವನ್ನೂ ಬಿಡಿಸಿ ಪ್ರಕೃತಿ, ಪ್ರತ್ಯಯ ತಿಳಿಸಿ\par

\vspace{2.5pt plus 3pt minus 0.5pt}\noindent\hspace*{1.2em}\kbvnum{(೧)}\hspace{0.55em}ಭೀಮನು ಕೀಚಕನನ್ನು ಕೊಂದನು.\par

\vspace{2.5pt plus 3pt minus 0.5pt}\noindent\hspace*{1.2em}\kbvnum{(೨)}\hspace{0.55em}ಕೋತಿಯು ಶಾಲಿನಿಯಿಂದ ಹಣ್ಣನ್ನು ಕಸಿಯಿತು.\par

\vspace{2.5pt plus 3pt minus 0.5pt}\noindent\hspace*{1.2em}\kbvnum{(೩)}\hspace{0.55em}ಕರ್ಣನು ಸರ್ಪಾಸ್ತ್ರವನ್ನು ಪ್ರಯೋಗಿಸಿದನು.\par

\vspace{2.5pt plus 3pt minus 0.5pt}\noindent\hspace*{1.2em}\kbvnum{(೪)}\hspace{0.55em}ಆಮೆಯು ನೀರಿನಲ್ಲಿ ವಾಸಿಸುತ್ತದೆ.\par

\vspace{2.5pt plus 3pt minus 0.5pt}\noindent\hspace*{1.2em}\kbvnum{(೫)}\hspace{0.55em}ಬೇಟೆಗಾರನು ಕಾಡಿಗೆ ಹೋದನು.\par

\kbvchapter{ಪಾಠ ೪}{ವಾಕ್ಯರಚನೆ}

\indent ವಸ್ತು ಸ್ಥಿತಿಯ ಅರಿವುಂಟುಮಾಡುವ ಶಬ್ದಗಳೇ ಪದಗಳು. ಪದಗಳ ರಚನೆಗೆ ಅಕ್ಷರಗಳು ಅತ್ಯಗತ್ಯ. ರೂಢಿಯಲ್ಲಿರುವ ಅರ್ಥಬರುವಂತೆ ಅಕ್ಷರಗಳನ್ನು ಕೂಡಿಸಿದಾಗಲೇ ಪದ ಉತ್ಪತ್ತಿಯಾಗುವುದು.\par

\indent ಹಲವು ಪದಗಳು ಸೇರಿ ವಾಕ್ಯವಾಗುವುದು. ವಾಕ್ಯವು ಅರ್ಥವತ್ತಾಗಿರ ಬೇಕಾದರೆ, ಆ ವಾಕ್ಯದಲ್ಲಿನ ಪದಗಳಲ್ಲಿ ಸೂಕ್ತವಾದ ಹೊಂದಾಣಿಕೆ ಇರಬೇಕು.\par

\indent ಉದಾಹರಣೆಗೆ ಕತ್ತೆಯು, ಭಾರವನ್ನು, ಹೊರುತ್ತದೆ ಎಂಬ ಪದ ಗಳನ್ನು ಗಮನಿಸೋಣ. ಇಲ್ಲಿ “ಕತ್ತೆಯು” ಎಂದಷ್ಟೇ ಹೇಳಿದರೆ ಮಾತಿನ ಅರ್ಥ ತಿಳಿದುಬರುವುದಿಲ್ಲ. ಹೀಗೆ, ಮೇಲೆ ತಿಳಿಸಿದ ಯಾವುದೊಂದು ಪದವನ್ನು\kbvcontinues
}
\clearpage

% === tex/0017.tex (Printed Page 15) ===
\renewcommand{\kbvchaptername}{ವಾಕ್ಯರಚನೆ}
\setcounter{page}{15}
\thispagestyle{kbvpage}
\kbvtypesetpage{\kbvcontinued ಹೇಳಿದರೂ ವಸ್ತುಸ್ಥಿತಿಯ ಅರಿವುಂಟಾಗುವುದಿಲ್ಲ. “ಕತ್ತೆಯು ಭಾರವನ್ನು” ಎಂದಾಗಲೀ, “ಕತ್ತೆಯು ಹೊರುತ್ತದೆ” ಎಂದಾಗಲೀ, “ಭಾರವನ್ನು ಹೊರು ತ್ತದೆ” ಎಂದಾಗಲೀ ಹೇಳಿದರೂ ಪೂರ್ಣವಾದ ಅರ್ಥ ಬರುವುದಿಲ್ಲ. ಅರ್ಥ ಪೂರ್ಣವಾಗಬೇಕಾದರೆ “ಕತ್ತೆಯು ಭಾರವನ್ನು ಹೊರುತ್ತದೆ” ಎಂಬು ದಾಗಿಯೇ ಹೇಳಬೇಕು. ಆದುದರಿಂದ ಸಮಂಜಸವಾದ ಅರ್ಥ ಬರುವ ವಾಕ್ಯಗಳ ರಚನೆಯಲ್ಲಿ ಪದಗಳ ಆಯ್ಕೆ ಹಾಗೂ ಜೋಡಣೆಯ ಕ್ರಮ ಬಹಳ ಮುಖ್ಯ. ಇದೇ ವಾಕ್ಯ ರಚನೆಯ ಒಳ ತಿರುಳು.\par

\kbvsechead{ವಾಕ್ಯ ವಿಭಜನೆ :}

\indent ಕತ್ತೆಯು ಭಾರವನ್ನು ಹೊರುತ್ತದೆ ಈ ವಾಕ್ಯದಲ್ಲಿ ಮೂರು ಪದ ಗಳಿವೆ. ಇವುಗಳಲ್ಲಿ “ಕತ್ತೆಯು” ಎಂಬುದು ನಾಮಪದ. ಈ ನಾಮಪದದ ಮೂಲರೂಪ “ಕತ್ತೆ” ; ಇದು ನಾಮ ಪ್ರಕೃತಿ. ಈ ನಾಮ ಪ್ರಕೃತಿಗೆ “ಉ” ಎಂಬ ಪ್ರತ್ಯಯ ಸೇರಿ “ಕತ್ತೆಯು” ಎಂಬ ನಾಮಪದವಾಗಿರುವುದು. ಆದರೆ, ಕತ್ತೆಯು ಎಂಬ ಪದ ಕ್ರಿಯೆಯ ಅರ್ಥವನ್ನು ಕೊಡುವುದಿಲ್ಲ. “ಕತ್ತೆಯು ಭಾರವನ್ನು ಹೊರುತ್ತದೆ” ಎನ್ನುವಾಗ ಕೆಲಸ ಮಾಡುವುದು ಕತ್ತೆ ಎಂಬುದಾಗಿ ತಿಳಿದುಬರುವುದು; ಇದು ಹೊರುವ ಕೆಲಸವನ್ನು ಮಾಡುವುದು,\par

\vspace{2.5pt plus 3pt minus 0.5pt}\noindent\hspace*{1.2em}\kbvnum{(೧)}\hspace{0.55em}ಕರ್ತೃಪದ ಬಡಗಿಯು ಮೇಜನ್ನು ಮಾಡಿದನು. ರೇವತಿಯು ಬಟ್ಟೆಗಳನ್ನು ಹೊಲಿಯುತ್ತಾಳೆ. ವಾಲ್ಮೀಕಿಯು ರಾಮಾಯಣವನ್ನು ಬರೆದನು.\par

\indent --ಈ ಮೂರು ವಾಕ್ಯಗಳಲ್ಲಿಯೂ ಕ್ರಮವಾಗಿ ಬಡಗಿಯು, ರೇವತಿಯು ವಾಲ್ಮೀಕಿಯು ಕೆಲಸವನ್ನು ಮಾಡುವವರ ಹೆಸರುಗಳನ್ನು ತಿಳಿಸುವ ನಾಮಪದ ಗಳಾಗಿವೆ. ಹೀಗೆ ಕೆಲಸವನ್ನು ಮಾಡುವವರ ಹೆಸರುಗಳನ್ನು ತಿಳಿಸುವ ನಾಮಪದಗಳಿಗೆ ಕರ್ತೃಪದಗಳೆಂದು ಹೆಸರು.\par

\vspace{2.5pt plus 3pt minus 0.5pt}\noindent\hspace*{1.2em}\kbvnum{(೨)}\hspace{0.55em}ಕ್ರಿಯಾಪದ “ಕತ್ತೆಯು ಭಾರವನ್ನು ಹೊರುತ್ತದೆ” ಎಂಬ ವಾಕ್ಯದಲ್ಲಿ ಹೊರುತ್ತದೆ ಎಂಬ ಪದವು ಕತ್ತೆ (ಕರ್ತೃ)ಯು ಮಾಡುವ ಕೆಲಸವನ್ನು ತಿಳಿಸುತ್ತದೆ. ಹೀಗೆ\kbvcontinues
}
\clearpage

% === tex/0018.tex (Printed Page 16) ===
\renewcommand{\kbvchaptername}{ವಾಕ್ಯರಚನೆ}
\setcounter{page}{16}
\thispagestyle{kbvpage}
\kbvtypesetpage{\kbvcontinued ಯಾವುದೇ ವಾಕ್ಯದಲ್ಲಿ ಕರ್ತೃವು ಮಾಡುವ ಕೆಲಸವನ್ನು ತಿಳಿಸುವ ಫಲವನ್ನು ಕ್ರಿಯಾನಾಮವಾಗಿರುವುದು. ಮೇಲಿನ ಮೂರು ಮುಖಗಳಲ್ಲಿ ಕ್ರಮವಾಗಿ ಮಾಡಿದನು, ಹೊಲಿಯುತ್ತಾಳೆ ಬರೆದನು - ಕ್ರಿಯಾನಾಮಗಳಾಗಿರುವುವು. ಕರ್ತೃವಿನಿಂದ ಆಗುವ ಕಾರ್ಯ ಈ ಪದಗಳಲ್ಲಿ ವ್ಯಕ್ತವಾಗಿದೆ.\par

\vspace{2.5pt plus 3pt minus 0.5pt}\noindent\hspace*{1.2em}\kbvnum{(3)}\hspace{0.55em}ಕರ್ಮನಾಮ\par

\indent ಕತ್ತೆಯು ಭಾರವನ್ನು ಹೊರುತ್ತದೆ - - ಎಂಬ ವಾಕ್ಯಕ್ಕೆ ಸಂಬಂಧಿಸಿದಂತೆ\par

\indent ಕೇಳಬಹುದಾದ ಮೂರು ಪ್ರಶ್ನೆಗಳು ಇಲ್ಲಿವೆ. ಈ ಪ್ರಶ್ನೆಗಳನ್ನು ಗಮನಿಸಿ,\par

\indent ({\kbvlatin i}) ಭಾರವನ್ನು ಹೊರುವುದು ಯಾವುದು ? --- ಕ್ರಿಯಾನಾಮವನ್ನು ಅವಲಂಬಿಸಿ ಕೇಳಲಾಗುವ ಈ ಪ್ರಶ್ನೆಗೆ “ಕತ್ತೆಯು” ಎಂಬ ಉತ್ತರ ಸಹಜವಾಗಿ ಬರುವುದು. ಈ ರೀತಿಯಾಗಿ ಪ್ರಶ್ನಿಸಿ ಕರ್ತೃಪದವನ್ನು ಸುಲಭವಾಗಿ ಕಂಡುಹಿಡಿಯಬಹುದು.\par

\indent ({\kbvlatin ii}) ಕತ್ತೆಯು ಮಾಡುವ ಕಾರ್ಯವೇನು ? --- ಕರ್ತೃ ಮಾಡುವ ಕೆಲಸವೇನೆಂದು ಕೇಳುವ ಈ ಪ್ರಶ್ನೆಗೆ “ಹೊರುವುದು” ಎಂಬ ಉತ್ತರ ಸಹಜವಾಗಿ ಬರುವುದು. ಈ ರೀತಿಯಾಗಿ ಪ್ರಶ್ನಿಸಿ ಕ್ರಿಯಾನಾಮವನ್ನು ಸುಲಭವಾಗಿ ಕಂಡುಹಿಡಿಯಬಹುದು.\par

\indent ({\kbvlatin iii}) ಕತ್ತೆಯು ಏನನ್ನು ಹೊರುವುದು ? --- ಕ್ರಿಯಾಪದವನ್ನೇ ಅವಲಂಬಿಸಿ ಕೇಳಲಾಗುವ ಈ ಪ್ರಶ್ನೆಗೆ “ಭಾರವನ್ನು” ಎಂಬ ಉತ್ತರ ಸಹಜವಾಗಿ ಬರುವುದು. ಕ್ರಿಯಾತ್ಮಕ ಕರ್ತೃವಿನ ಇಚ್ಛೆ ಅಥವಾ ಇಚ್ಛಿತ ಕಾರ್ಯ ಯಾವುದೆಂಬುದು ಈ ಉತ್ತರದಿಂದ ತಿಳಿದುಬರುವುದು.\par

\indent ಹೀಗೆ ಕ್ರಿಯಾಪದವನ್ನು ಅವಲಂಬಿಸಿ ಏನನ್ನು? ಅಥವಾ ಯಾರನ್ನು? ಅಥವಾ ಯಾವುದನ್ನು? ಎಂಬ ಪ್ರಶ್ನೆ ಹಾಕಿದಾಗ ಉತ್ತರವಾಗಿ ಬರುವ ಪದವನ್ನು ಕರ್ಮನಾಮ ಎನ್ನುವರು.\par

\indent ಮೇಲೆ ತಿಳಿಸಿರುವ ವಾಕ್ಯಗಳಲ್ಲಿ “ಮೇಜನ್ನು”, “ಬಟ್ಟೆಗಳನ್ನು”, “ರಾಮಾಯಣವನ್ನು”... ಕರ್ಮನಾಮಗಳಾಗಿರುವುವು.\par
}
\clearpage

% === tex/0019.tex (Printed Page 17) ===
\renewcommand{\kbvchaptername}{ವಾಕ್ಯರಚನೆ}
\setcounter{page}{17}
\thispagestyle{kbvpage}
\kbvtypesetpage{\begin{kbvrule}
\noindent\kbvbullet\hspace{0.5em}ವಾಕ್ಯವು ಕರ್ತೃ, ಕರ್ಮ, ಕ್ರಿಯಾಪದಗಳನ್ನೊಳಗೊಂಡ ಪದ ಸಮುದಾಯವೆಂಬುದನ್ನು ತಿಳಿಮುದಾಯಿತು. ಆದರೆ ಕೆಲವು ವಾಕ್ಯಗಳಲ್ಲಿ ಕರ್ಮಪದದ ಅಗತ್ಯವೇ ಇರುವುದಿಲ್ಲ. ಉದಾಹರಣೆಗೆ, ಗಿಣಿಯು ಹಾರುತ್ತಿದೆ. ಕತ್ತೆ ಒದೆಯಿತು. ರಾಮನು ಬರುತ್ತಾನೆ. ವಿಮಾನ ಇಳಿಯಿತು. ನಾಯಿಯು ಓಡಿತು. ಹೂವು ಅರಳಿತು.
\end{kbvrule}

\indent ಈ ವಾಕ್ಯಗಳಲ್ಲಿ ಕರ್ಮಪದಗಳಿಲ್ಲ ; ಇಲ್ಲಿರುವುದು ಕರ್ತೃ ಮತ್ತು ಕ್ರಿಯಾಪದಗಳು ಮಾತ್ರ... ಆದರೂ ಈ ವಾಕ್ಯಗಳು ಅರ್ಥಪೂರ್ಣವಾಗಿವೆ. {[}ಈ ಬಗ್ಗೆ ಹೆಚ್ಚಿನ ವಿವರಗಳನ್ನು ಮುಂದಿನ ಪಾಠದಲ್ಲಿ ತಿಳಿಯಬಹುದು.{]}\par

\kbvexercise{ಅಭ್ಯಾಸ ಪತ್ರಿಕೆ ೪}

\vspace{2.5pt plus 3pt minus 0.5pt}\noindent\hspace*{1.2em}\kbvnum{(1)}\hspace{0.55em}ಈ ವಾಕ್ಯಗಳಲ್ಲಿರುವ ಕರ್ತೃಪದ, ಕ್ರಿಯಾಪದ ಮತ್ತು ಕರ್ಮ\par

\kbvsechead{ಪದಗಳನ್ನು ಗುರ್ತಿಸಿ :}

\vspace{2.5pt plus 3pt minus 0.5pt}\noindent\hspace*{1.2em}\kbvnum{1.}\hspace{0.55em}ಚಿರತೆಯು ಹಸುವನ್ನು ಕೊಂದಿತು.\par

\vspace{2.5pt plus 3pt minus 0.5pt}\noindent\hspace*{1.2em}\kbvnum{2.}\hspace{0.55em}ಭೀಮಸೇನನು ಹಿಡಂಬಿಯನ್ನು ಮದುವೆಯಾದನು.\par

\vspace{2.5pt plus 3pt minus 0.5pt}\noindent\hspace*{1.2em}\kbvnum{3.}\hspace{0.55em}ಹೂವು ಅರಳಿತು ; ಸುಗಂಧ ಹರಡಿತು.\par

\vspace{2.5pt plus 3pt minus 0.5pt}\noindent\hspace*{1.2em}\kbvnum{4.}\hspace{0.55em}ರೈಲುಗಾಡಿ ಚಲಿಸಿತು.\par

\vspace{2.5pt plus 3pt minus 0.5pt}\noindent\hspace*{1.2em}\kbvnum{5.}\hspace{0.55em}ಕಳ್ಳನು ಸರವನ್ನು ಕದ್ದನು.\par

\vspace{2.5pt plus 3pt minus 0.5pt}\noindent\hspace*{1.2em}\kbvnum{6.}\hspace{0.55em}ಬಟ್ಟೆಯನ್ನು ಕತ್ತರಿಯಿಂದ ಕತ್ತರಿಸುತ್ತೇವೆ.\par

\vspace{2.5pt plus 3pt minus 0.5pt}\noindent\hspace*{1.2em}\kbvnum{7.}\hspace{0.55em}ಸಕ್ಕರೆಯಿಂದ ಮಿಠಾಯಿ ತಯಾರಿಸುತ್ತಾರೆ.\par

\vspace{2.5pt plus 3pt minus 0.5pt}\noindent\hspace*{1.2em}\kbvnum{8.}\hspace{0.55em}ಅರ್ಚಕನು ಹೂವುಗಳಿಂದ ದೇವರನ್ನು ಪೂಜಿಸಿದನು.\par

\vspace{2.5pt plus 3pt minus 0.5pt}\noindent\hspace*{1.2em}\kbvnum{9.}\hspace{0.55em}ಹಡಗು ನೀರಿನಲ್ಲಿ ಮುಳುಗಿತು.\par

\vspace{2.5pt plus 3pt minus 0.5pt}\noindent\hspace*{1.2em}\kbvnum{10.}\hspace{0.55em}ಗಾಳಿಯು ವೇಗವಾಗಿ ಬೀಸುತ್ತಿದೆ.\par

\vspace{2.5pt plus 3pt minus 0.5pt}\noindent\hspace*{1.2em}\kbvnum{(2)}\hspace{0.55em}ಕೆಳಗೆ ಕೊಟ್ಟಿರುವ ಕ್ರಿಯಾಪದಗಳಿಗೆ ಸರಿಹೊಂದುವ ಕರ್ತೃ,\par

\kbvsechead{ಕರ್ಮ ಪದಗಳನ್ನು ಆಯ್ಕೆಮಾಡಿ ವಾಕ್ಯಗಳನ್ನು ಬರೆಯಿರಿ :}

\vspace{4pt plus 4pt minus 1pt}\noindent{\kbvsecfont\fontsize{11pt}{14.5pt}\selectfont\color{kbvaccentdark} 1.\hspace{0.45em}ಕಚ್ಚಿತು (ಮಾದರಿ :}\hspace{0.5em}ಮುಂಗುಸಿಯು ಹಾವನ್ನು ಕಚ್ಚಿತು.)\par
}
\clearpage

% === tex/0020.tex (Printed Page 18) ===
\renewcommand{\kbvchaptername}{ಕ್ರಿಯಾಪದ}
\setcounter{page}{18}
\thispagestyle{kbvpage}
\kbvtypesetpage{\vspace{2.5pt plus 3pt minus 0.5pt}\noindent\hspace*{1.2em}\kbvnum{೨.}\hspace{0.55em}ತಿನ್ನುತ್ತಿದೆ ೪. ಕರಿಯುತ್ತಿದ್ದಾಳೆ\par

\vspace{2.5pt plus 3pt minus 0.5pt}\noindent\hspace*{1.2em}\kbvnum{೩.}\hspace{0.55em}ಮಾಡುತ್ತಿದ್ದಾನೆ ೫. ಆಡುತ್ತಿದ್ದಾರೆ.\par

\kbvchapter{ಪಾಠ ೫}{ಕ್ರಿಯಾಪದ}

\indent ಕ್ರಿಯೆಯ ಅರ್ಥವನ್ನು ಕೊಡುವ ಪದವೇ ಕ್ರಿಯಾಪದ; ಇದು ೪ ನಿರ್ವಹಿಸುವ ಕೆಲಸವನ್ನು ತಿಳಿಸುತ್ತದೆ.\par

\begin{kbvderiv}
\vspace{4pt plus 3pt minus 1pt}
\noindent{\renewcommand{\arraystretch}{1.24}\setlength{\tabcolsep}{5pt}
\begin{tabularx}{\textwidth}{X X X}
ಕುದುರೆಯು ಓಡುತ್ತಿದೆ & ಪದ್ಮಿನಿಯು ಹಾಡುತ್ತಾಳೆ & ನಾಯಿ \\
ರೊಟ್ಟಿಯನ್ನು ತಿನ್ನುತ್ತದೆ & ಮಂಜುನಾಥನು ಕಾಲ್ಚೆಂಡಿನ ಆಟ &  \\
\end{tabularx}}
\vspace{4pt plus 3pt minus 1pt}
\end{kbvderiv}

\indent ಆಡುತ್ತಾನೆ--ಈ ವಾಕ್ಯಗಳಲ್ಲಿರುವ ಕ್ರಿಯಾಪದಗಳೆಂದರೆ ಓಡು, ಹಾಡುತ್ತಾಳೆ, ತಿನ್ನುತ್ತದೆ, ಆಡುತ್ತಾನೆ. ಈ ಕ್ರಿಯಾಪದಗಳನ್ನು ೬\par

\begin{kbvderiv}
\vspace{4pt plus 3pt minus 1pt}
\begin{center}\begin{adjustbox}{max width=\textwidth}
{\renewcommand{\arraystretch}{1.22}\setlength{\tabcolsep}{4pt}
\begin{tabular}{l c l}
ಓಡುತ್ತಿದೆ & {\color{kbvgrey}$=$} & ಓಡು + ಉತ್ತ + ಇದೆ \\
ಹಾಡುತ್ತಾಳೆ & {\color{kbvgrey}$=$} & ಹಾಡು + ಉತ್ತ + ಆಳೆ \\
ತಿನ್ನುತ್ತದೆ & {\color{kbvgrey}$=$} & ತಿನ್ನು + ಉತ್ತ + ಅದೆ \\
ಆಡುತ್ತಾನೆ & {\color{kbvgrey}$=$} & ಆಡು + ಉತ್ತ + ಆನೆ \\
\end{tabular}}
\end{adjustbox}\end{center}
\vspace{4pt plus 3pt minus 1pt}
\end{kbvderiv}

\indent ಇಲ್ಲಿ “ಓಡು” ಎಂಬುದು “ಓಡುತ್ತಿದೆ” ಎಂಬ ಕ್ರಿಯಾಪದದ ವ. ರೂಪ; ಇದು ಸ್ವತಂತ್ರವಾದ ಅರ್ಥವನ್ನು ಹೊಂದಿರುವ ಶಬ್ದ. ಕ್ರಿಯೆ ಅರ್ಥವನ್ನು ಕೊಡುವ, ಆದರೆ ಪ್ರತ್ಯಯವನ್ನು ಹೊಂದಿರದ ಈ ಶಃ ಕ್ರಿಯಾಪ್ರಕೃತಿಯಾಗಿರುವುದು. ಕ್ರಿಯಾಪ್ರಕೃತಿಯನ್ನು “ಧಾತು” ಎ ದಾಗಿಯೂ ಕರೆಯುವುದುಂಟು. ಕ್ರಿಯಾಪ್ರಕೃತಿಯ ಮುಂದಿರುವ ಉತ್ತ, ಎಂಬುವು ಪ್ರತ್ಯಯಗಳು. “ಉತ್ತ” ಎಂಬುದನ್ನು “ಕಾಲ ಸೂ ಪ್ರತ್ಯಯ”ವೆಂತಲೂ, “ಇದೆ” ಎಂಬುದನ್ನು ಆಖ್ಯಾತ ಪ್ರತ್ಯಯವೇಂ ಹೇಳುತ್ತೇವೆ.\par

\begin{kbvrule}
\noindent{\kbvsecfont ವ್ಯಾಖ್ಯೆ :}\hspace{0.5em}ಕ್ರಿಯಾಪದದ ಮೂಲ ರೂಪಕ್ಕೆ ಕ್ರಿಯಾಪ್ರಕೃತಿ ೮2 ಧಾತು ಎಂದು ಹೆಸರು. ಸ್ವತಂತ್ರವಾಗಿರುವ ಈ ಧಾತುಗಳೇ ಮ ಧಾತುಗಳು,
\end{kbvrule}
}
\clearpage

% === tex/0021.tex (Printed Page 19) ===
\renewcommand{\kbvchaptername}{ಕ್ರಿಯಾಪದ}
\setcounter{page}{19}
\thispagestyle{kbvpage}
\kbvtypesetpage{\indent ಮೇಲೆ ಕೊಟ್ಟಿರುವ ಉದಾಹರಣೆಗಳಲ್ಲಿ ಓಡು, ಹಾಡು, ತಿನ್ನು, ಆಡು ಕ್ರಿಯಾಪ್ರಕೃತಿಗಳು ಅಥವಾ ಧಾತುಗಳು ; “ಉತ್ತ” ಎಂಬುದು ಕಾಲಸೂಚಕ ಪ್ರತ್ಯಯ ಮತ್ತು ಇದೆ, ಆಳೆ, ಅದೆ, ಆನೆ ಆಖ್ಯಾತ ಪ್ರತ್ಯಯಗಳು. ಸೂತ್ರ ರೂಪದಲ್ಲಿ ಹೇಳುವುದಾದರೆ,\par

\indent ಕ್ರಿಯಾ = ಕ್ರಿಯಾ + ಕಾಲಸೂಚಕ + ಆಖ್ಯಾತ\par

\indent ಪದ + ಪ್ರಕೃತಿ + ಪ್ರತ್ಯಯ + ಪ್ರತ್ಯಯ ವಿಶೇಷ ಸೂಚನೆ : ಕೆಲವು ನಾಮಪ್ರಕೃತಿಗಳನ್ನು ಮತ್ತು ಅನುಕರಣ ಶಬ್ದಗಳನ್ನು ಧಾತುಗಳಂತೆ ಬಳಸುವುದುಂಟು. ನಾಮಪ್ರಕೃತಿಯ ಮೇಲಾಗಲೀ ಅನುಕರಣ ಶಬ್ದದ ಮೇಲಾಗಲೀ “ಇಸು” ಪ್ರತ್ಯಯ ಸೇರಲು ಅದು ಧಾತು ವೆನಿಸುವುದು. ಅಂತಹ ಧಾತುಗಳಿಗೆ ಪ್ರತ್ಯಯಾಂತ ಧಾತುಗಳೆಂದು ಹೆಸರು.\par

\vspace{3pt plus 2pt}\noindent\hspace*{1.2em}\kbvlabel{ಉದಾ}\par\vspace{1pt}

\begin{kbvderiv}
\vspace{4pt plus 3pt minus 1pt}
\begin{center}\begin{adjustbox}{max width=\textwidth}
{\renewcommand{\arraystretch}{1.22}\setlength{\tabcolsep}{4pt}
\begin{tabular}{l c l}
ಕನ್ನಡ + ಇಸು & {\color{kbvgrey}$=$} & \kbvresult{ಕನ್ನಡಿಸು}{\color{kbvgrey} (ಪ್ರತ್ಯಯಾಂತ ಧಾತು)} \\
ಉತ್ತರ + ಇಸು & {\color{kbvgrey}$=$} & \kbvresult{ಉತ್ತರಿಸು}{\color{kbvgrey} ( ೨ )} \\
“ಪ್ರಯತ್ನ + ಇಸು & {\color{kbvgrey}$=$} & \kbvresult{ಪ್ರಯತ್ನಿಸು}{\color{kbvgrey} ( ೨ )} \\
ಸಿದ್ಧಿ + ಇಸು & {\color{kbvgrey}$=$} & \kbvresult{ಸಿದ್ಧಿಸು}{\color{kbvgrey} ( ೨ )} \\
ಥಳಥಳ + ಇಸು & {\color{kbvgrey}$=$} & \kbvresult{ಥಳಥಳಿಸು}{\color{kbvgrey} ( ೨ )} \\
ಘಮಘಮ + ಇಸು & {\color{kbvgrey}$=$} & \kbvresult{ಘಮಘಮಿಸು}{\color{kbvgrey} ( ೨ )} \\
\end{tabular}}
\end{adjustbox}\end{center}
\vspace{4pt plus 3pt minus 1pt}
\end{kbvderiv}

\indent ಪ್ರತ್ಯಯಾಂತ ಧಾತುಗಳು ಆಖ್ಯಾತ ಪ್ರತ್ಯಯಗಳನ್ನು ಹೊಂದಿ ಕ್ರಿಯಾಪದಗಳಾಗುವುವು.\par

\begin{kbvderiv}
\vspace{4pt plus 3pt minus 1pt}
\begin{center}\begin{adjustbox}{max width=\textwidth}
{\renewcommand{\arraystretch}{1.22}\setlength{\tabcolsep}{4pt}
\begin{tabular}{l c l}
ಉದಾ : ಕನ್ನಡಿಸು + ಉತ್ತ + ಆನೆ & {\color{kbvgrey}$=$} & \kbvresult{ಕನ್ನಡಿಸುತ್ತಾರೆ} \\
ಸಿದ್ಧಿಸು + ಇದೆ & {\color{kbvgrey}$=$} & \kbvresult{ಸಿದ್ಧಿಸಿತು} \\
ಥಳಥಳಿಸು + ಉದು & {\color{kbvgrey}$=$} & \kbvresult{ಥಳಥಳಿಸುವುದು} \\
\end{tabular}}
\end{adjustbox}\end{center}
\vspace{4pt plus 3pt minus 1pt}
\end{kbvderiv}

\kbvsechead{ಅಕರ್ಮಕ ಸಕರ್ಮಕ ಧಾತುಗಳು :}

\begin{kbvrule}
\noindent\kbvbullet\hspace{0.5em}ಧಾತುಗಳಲ್ಲಿ ಎರಡು ವಿಧ. ಆ ಎರಡು ವಿಧಗಳು ಯಾವುವೆಂದರೆ,
\end{kbvrule}

\vspace{2.5pt plus 3pt minus 0.5pt}\noindent\hspace*{1.2em}\kbvnum{(೧)}\hspace{0.55em}ಅಕರ್ಮಕ ಧಾತು ಮತ್ತು (೨) ಸಕರ್ಮಕ ಧಾತು.\par
}
\clearpage

% === tex/0022.tex (Printed Page 20) ===
\renewcommand{\kbvchaptername}{ಕ್ರಿಯಾಪದ}
\setcounter{page}{20}
\thispagestyle{kbvpage}
\kbvtypesetpage{\indent “ಕುಮಾರು ಓಡುತ್ತದೆ” ಎಂಬ ವಾಕ್ಯದಲ್ಲಿ ಕರ್ಮಪದವಿಲ್ಲ. ಎಂಬ ಧಾತುವಿಗೆ ಕರ್ಮಪದದ ಅಗತ್ಯ ಬೀಳುವುದಿಲ್ಲ... ಹೀಗೆ, ಕರ್ಮ ಅಗತ್ಯ ಇರದ ಧಾತುವೇ ಅಕರ್ಮಕ ಧಾತು\par

\indent “ನಾಯಿಯು ರೈತನು ತಿನ್ನುತ್ತದೆ”--ಇದು ಕರ್ಮಪದ ಹೊಂದಿರುವ ವಾಕ್ಯ. “ತಿನ್ನು” ಎಂಬ ಧಾತುವಿಗೆ ಕರ್ಮಪದದ ಅಗತ್ಯವ ಹೀಗೆ ಕರ್ಮಪದದ ಅಗತ್ಯವಿರುವ ಧಾತುವೇ ಸಕರ್ಮಕ ಧಾತು.\par

\indent ಉದಾಹರಣೆಗೆ ಬೀಳು, ಏಕ ಕುಣಿ, 'ಹೋಗು, ಬರು, ಅಕರ್ಮಕ ಧಾತುಗಳು ; ಕಾಣು, ಮಾಡು, ನೋಡು, 'ಹಾಡು 8 ಕೊಡು, ಸಕರ್ಮಕ ಧಾತುಗಳು. (1 ೫1 ಗ\par

\kbvexercise{ಅಭ್ಯಾಸ ಪತ್ರಿಕೆ ೫}

\vspace{2.5pt plus 3pt minus 0.5pt}\noindent\hspace*{1.2em}\kbvnum{1.}\hspace{0.55em}ಕ್ರಿಯಾಪದ ಎಂದರೇನು ? ಉದಾಹರಣೆ ಸಹಿತ ತಿಳಿಸಿ.\par

\vspace{2.5pt plus 3pt minus 0.5pt}\noindent\hspace*{1.2em}\kbvnum{2.}\hspace{0.55em}ಕುಣಿಯುತ್ತದೆ, ಮಾಡುತ್ತದೆ, ಬರೆಯುತ್ತಾಳೆ- ಈ ಕ್ರಿಯಾ ಗಳಲ್ಲಿ ಪ್ರಕೃತಿ, ಪ್ರತ್ಯಯಗಳನ್ನು ಬಿಡಿಸಿ ಬರೆಯಿರಿ.\par

\vspace{2.5pt plus 3pt minus 0.5pt}\noindent\hspace*{1.2em}\kbvnum{3.}\hspace{0.55em}ಧಾತು ಎಂದರೇನು ? ಎರಡು ಉದಾಹರಣೆಗಳನ್ನು ಕೊಡಿ.\par

\vspace{2.5pt plus 3pt minus 0.5pt}\noindent\hspace*{1.2em}\kbvnum{4.}\hspace{0.55em}ಆಖ್ಯಾತ ಪ್ರತ್ಯಯಗಳಿಗೆ ಎರಡು ಉದಾಹರಣೆಗಳನ್ನು ಕೊಡಿ\par

\vspace{2.5pt plus 3pt minus 0.5pt}\noindent\hspace*{1.2em}\kbvnum{5.}\hspace{0.55em}ಅಕರ್ಮಕ ಧಾತುವಿಗೂ ಸಕರ್ಮಕ ಧಾತುವಿಗೂ' ಇಃ ವ್ಯತ್ಯಾಸವೇನು ? ಉದಾಹರಣೆಗಳ ಸಹಿತ ತಿಳಿಸಿ.\par

\vspace{2.5pt plus 3pt minus 0.5pt}\noindent\hspace*{1.2em}\kbvnum{6.}\hspace{0.55em}ಕರ್ಮಪದವನ್ನು ಅಪೇಕ್ಷಿಸುವ ಎರಡು ಧಾತುಗಳನ್ನು ತಿಳಿಸಿ,\par

\vspace{2.5pt plus 3pt minus 0.5pt}\noindent\hspace*{1.2em}\kbvnum{7.}\hspace{0.55em}ಕರ್ಮಪದವನ್ನು ಅಪೇಕ್ಷಿಸದ ಎರಡು ಧಾತುಗಳನ್ನು ತಿಳಿಸಿ.\par

\vspace{2.5pt plus 3pt minus 0.5pt}\noindent\hspace*{1.2em}\kbvnum{8.}\hspace{0.55em}ಪ್ರತ್ಯಯಾಂತ ಧಾತುವಿಗೆ ಒಂದು ಉದಾಹರಣೆಯನ್ನು ಕೊಡಿ\par

\vspace{2.5pt plus 3pt minus 0.5pt}\noindent\hspace*{1.2em}\kbvnum{9.}\hspace{0.55em}“ದುಃಖಿಸುತ್ತಾನೆ” ಈ ಕ್ರಿಯಾಪದದಲ್ಲಿ ಮೂರು ಪ್ರತ್ಯಯ ಸೇರಿವೆ ; ಅವುಗಳನ್ನು ಬಿಡಿಸಿ ಹೇಳಿ.\par

\vspace{2.5pt plus 3pt minus 0.5pt}\noindent\hspace*{1.2em}\kbvnum{10.}\hspace{0.55em}ಪ್ರಕೃತಿ, ಪ್ರತ್ಯಯಗಳನ್ನು ಸೇರಿಸಿ ಪದಗಳನ್ನು ಬರೆಯಿರಿ.\par

\vspace{2.5pt plus 3pt minus 0.5pt}\noindent\hspace*{1.2em}\kbvnum{(೧)}\hspace{0.55em}ಕಲಿ + ಉತ್ತ + ಆರೆ (೨) ಮಾಡು + ದ + ಅನು\par

\vspace{2.5pt plus 3pt minus 0.5pt}\noindent\hspace*{1.2em}\kbvnum{(೩)}\hspace{0.55em}ನೋಡು + ದ + ಅಳು (೪) ಮಲಗು + ಇಸು\par
}
\clearpage

% === tex/0023.tex (Printed Page 21) ===
\renewcommand{\kbvchaptername}{ಕನ್ನಡ ಸಂಧಿಗಳು}
\setcounter{page}{21}
\thispagestyle{kbvpage}
\kbvtypesetpage{\kbvchapter{ಪಾಠ ೬}{ಕನ್ನಡ ಸಂಧಿಗಳು}

\indent “ನಾವೆಲ್ಲರೂ ಜೊತೆಗೂಡಿ ಊರನ್ನು ಸೇರೋಣ.” - ಈ ವಾಕ್ಯದ ಅಭಿಪ್ರಾಯವೇನು ಎಂಬುದನ್ನು ನೀವು ಗ್ರಹಿಸಬಲ್ಲಿರಿ. ಆದರೆ, ಇಲ್ಲಿ ಕೆಲವು ಪದಗಳನ್ನು ಬಿಡಿಬಿಡಿಯಾಗಿ ಹೇಳಿರುವುದು ಅಷ್ಟೊಂದು ತೃಪ್ತಿಕರವಾಗಿಲ್ಲ. ವಿಚಿತ್ರ ತೆರನಾಗಿರುವ ಈ ಕ್ರಮ ರೂಢಿಗೆ ವಿರುದ್ಧವೆಂತಲೇ ಹೇಳಬಹುದು. ಈ ವಾಕ್ಯದಲ್ಲಿ ಕೆಲವೊಂದು ಕಡೆ ಅಕ್ಷರಗಳನ್ನು ಒಟ್ಟುಗೂಡಿಸಿ ಹೇಳಿ:\par

\indent - ದಲ್ಲಿ ಕೇಳಲು ಹಿತವೆನಿಸುವುದು,\par

\indent - ನಾವು + ಎಲ್ಲರೂ = ನಾವೆಲ್ಲರೂ, ಜೊತೆ + ಕೂಡಿ = ಜೊತೆಗೂಡಿ,\par

\indent - ಊರು + ಅನ್ನು = ಊರನ್ನು ಎಂದಾಗಲಿ ಸುಧಾರಿತ ವಾಕ್ಯವು “ನಾವೆಲ್ಲರೊ ಜೊತೆಗೂಡಿ ಊರನ್ನು ಸೇರೋಣ” ಎಂದಾಗುವುದು.\par

\indent ನಾವೆಲ್ಲರೊ ಎಂಬ ಪದದಲ್ಲಿ ಪೂರ್ವಾರ್ಧ ಭಾಗವು “ನಾವು” -ಮತ್ತು ಉತ್ತರಾರ್ಧ ಭಾಗವು “ಎಲ್ಲರೊ” ಆಗಿರುವುದು. ಇಲ್ಲಿ ಪೂರ್ವಾರ್ಧದ ಕಡೆಯ ಅಕ್ಷರ “ವು” ಮತ್ತು ಉತ್ತರಾರ್ಧದ ಮೊದಲ ಅಕ್ಷರ “ಎ” ಇವೆರಡೂ ಸೇರಿ “ವೆ” ಆಗಿರುವುದನ್ನು ಕಾಣಬಹುದು. ಹೀಗೆ, ಉಚ್ಚರಿಸಲು ಕಾಲ ವಿಳಂಬವಾಗದ ರೀತಿಯಲ್ಲಿ ಎರಡು ಅಕ್ಷರಗಳು ನಡಬಿಡಿ ಒಂದಕ್ಕೊಂದು ಸೇರುವುದಕ್ಕೆ “ಸಂಧಿ” ಎಂದು ಹೆಸರು.\par

\begin{kbvrule}
\noindent\kbvbullet\hspace{0.5em}ಸಂಧಿಯಾಗುವಾಗ ಒಂದು ಅಕ್ಷರ ಬಿಟ್ಟು ಹೋಗಬಹುದು. ಒಂದು ಅಕ್ಷರ ಹೊಸದಾಗಿ ಬಂದು ಸೇರಬಹುದು; ಒಂದು ಅಕ್ಷರಕ್ಕೆ ಬದಲಾಗಿ ಮತ್ತೊಂದು ಅಕ್ಷರ ಬರಬಹುದು. ಈ ರೀತಿಯಾಗಿ ಉಂಟಾಗುವ ಬದಲಾವಣೆಗಳನ್ನು ಸಂಧಿಕಾರ್ಯವೆನ್ನುತ್ತೇವೆ.
\end{kbvrule}

\vspace{6pt plus 5pt minus 1pt}\begin{center}{\kbvsecfont\fontsize{12.5pt}{16pt}\selectfont\color{kbvaccentdark} ಕನ್ನಡ ಸಂಧಿಗಳು}\end{center}\vspace{3pt plus 3pt minus 1pt}

\indent ಕನ್ನಡ ಸಂಧಿಗಳಲ್ಲಿ ಮೂರು ಪ್ರಕಾರಗಳಿವೆ; ಅವು ಯಾವುವೆಂದರೆ\par

\vspace{2pt plus 2pt}\noindent\hspace*{1.5em}(೧) ಲೋಪ ಸಂಧಿ,\hfill (೨) ಆಗಮ ಸಂಧಿ,\hfill (೩) ಆದೇಶ ಸಂಧಿ\hspace*{1.5em}\par\vspace{2pt plus 2pt}

\kbvsechead{(೧) ಲೋಪ ಸಂಧಿ:}

\indent “ನಾವೆಲ್ಲರೊ ಜೊತೆಗೂಡಿ ಮನೆಗೆ ಹೋಗೋಣ” ಎಂಬ ವಾಕ್ಯದಲ್ಲಿ ನಾವೆಲ್ಲರೊ ಎಂಬ ಸಂಧಿ ಪದವನ್ನು ಬಿಡಿಸಿ ನೋಡೋಣ.\par
}
\clearpage

% === tex/0024.tex (Printed Page 22) ===
\renewcommand{\kbvchaptername}{ಕನ್ನಡ ಸಂಧಿಗಳು}
\setcounter{page}{22}
\thispagestyle{kbvpage}
\kbvtypesetpage{\begin{kbvderiv}
\vspace{4pt plus 3pt minus 1pt}
\begin{center}\begin{adjustbox}{max width=\textwidth}
{\renewcommand{\arraystretch}{1.22}\setlength{\tabcolsep}{4pt}
\begin{tabular}{l c l}
 & {\color{kbvgrey}$=$} & ನಾವು + ಎಲ್ಲರೂ \\
 & {\color{kbvgrey}$=$} & ನಾ + (ವ್ + ಉ) + ಎಲ್ಲರೂ {[}\ensuremath{\because} ವು = ವ್ + ಉ{]} \\
 & {\color{kbvgrey}$=$} & ನಾ + (ವ್ + ಎ) + ಎಲ್ಲರೂ.... {[}ಉ-ಕಾರ ಲೋಪ{]} \\
 & {\color{kbvgrey}$=$} & ನಾವೆಲ್ಲರೂ ..... {[}\ensuremath{\because} ವ್ + ಎ = ವೆ{]} \\
\end{tabular}}
\end{adjustbox}\end{center}
\vspace{4pt plus 3pt minus 1pt}
\end{kbvderiv}

\indent ನಾವು, ಎಲ್ಲರೂ ಎಂಬ ಪದಗಳಲ್ಲಿ ವು, ಎ, ಅಕ್ಷರಗಳು ಒಂದಕ್ಕೊಂದು ಸೇರಿ ಸಂಧಿಯಾದಾಗ ಉ-ಕಾರವು ಕಾಣದಾಗುವುದು. ಇದನ್ನು ಉ-ಕಾರ ಲೋಪವಾಯಿತು ಎಂಬುದಾಗಿ ಹೇಳುತ್ತೇವೆ. ಹೀಗೆ ಸಂಧಿಕಾರ್ಯವು ನಡೆಯುವಾಗ ಒಂದು ಅಕ್ಷರವು ಬಿಟ್ಟು ಹೋಗುವುದಕ್ಕೆ “ಲೋಪ ಸಂಧಿ” ಎಂದು ಹೆಸರು.\par

\begin{kbvnote}
\noindent{\kbvsecfont ವಿವರಣೆ :}\hspace{0.5em}ಮೇಲಿನ ಉದಾಹರಣೆಯಲ್ಲಿ “ಉ” ಮತ್ತು “ಎ” ಸ್ವರಾಕ್ಷರ ಗಳು, ಹೀಗೆ, ಎರಡು ಸ್ವರಾಕ್ಷರಗಳು ಸಂಧಿಸುವಾಗ ಪದದ ಅರ್ಥ ಕೆಡದ ರೀತಿಯಲ್ಲಿ ಮೊದಲನೆಯ ಸ್ವರಾಕ್ಷರವು ಲೋಪವಾಗುವುದು. ಯಾವ ಅಕ್ಷರ ಲೋಪವಾಗುವುದೋ ಆ ಅಕ್ಷರವನ್ನು ಸೇರಿಸಿ ಸಂಧಿಯನ್ನು ಹೆಸರಿಸುವುದು ವಾಡಿಕೆಯಾಗಿದೆ. ಇಲ್ಲಿ ಉ-ಕಾರ ಲೋಪವಾಗಿರುವುದರಿಂದ, ಇದು ಉ-ಕಾರ ಲೋಪ ಸಂಧಿ.
\end{kbvnote}

\vspace{3pt plus 2pt}\noindent\hspace*{1.2em}\kbvlabel{ಉದಾ}\par\vspace{1pt}

\begin{kbvderiv}
\vspace{4pt plus 3pt minus 1pt}
\begin{center}\begin{adjustbox}{max width=\textwidth}
{\renewcommand{\arraystretch}{1.22}\setlength{\tabcolsep}{4pt}
\begin{tabular}{l c l}
ಅಲ್ಲಿ + ಅಲ್ಲಿ & {\color{kbvgrey}$=$} & \kbvresult{ಅಲ್ಲಲ್ಲಿ,}{\color{kbvgrey} .... (ಇ-ಕಾರ ಲೋಪ ಸಂಧಿ)} \\
ನನ್ನ + ಅಕೆ & {\color{kbvgrey}$=$} & \kbvresult{ನನ್ನಾಕೆ,}{\color{kbvgrey} .... (ಅ-ಕಾರ ಲೋಪ ಸಂಧಿ)} \\
ನನಗೆ + ಅಲ್ಲಿ & {\color{kbvgrey}$=$} & \kbvresult{ನನಗೆಲ್ಲಿ,}{\color{kbvgrey} .... (ಎ-ಕಾರ ಲೋಪ ಸಂಧಿ)} \\
\end{tabular}}
\end{adjustbox}\end{center}
\vspace{4pt plus 3pt minus 1pt}
\end{kbvderiv}

\vspace{4pt plus 4pt minus 1pt}\noindent{\kbvsecfont\fontsize{11pt}{14.5pt}\selectfont\color{kbvaccentdark} (೨)\hspace{0.45em}ಆಗಮ ಸಂಧಿ:}\hspace{0.5em}“ಕಳ್ಳರು ಕುರಿಯನ್ನು ಕದ್ದರು” ಎಂಬ ವಾಕ್ಯದಲ್ಲಿ ಕುರಿಯನ್ನು ಎಂಬ\par

\begin{kbvderiv}
\vspace{4pt plus 3pt minus 1pt}
\begin{center}\begin{adjustbox}{max width=\textwidth}
{\renewcommand{\arraystretch}{1.22}\setlength{\tabcolsep}{4pt}
\begin{tabular}{l c l}
ಸಂಧಿಪದವನ್ನು ಬಿಡಿಸಿ ನೋಡೋಣ. ಕುರಿ + ಅನ್ನು & {\color{kbvgrey}$=$} & \kbvresult{ಕುರಿಯನ್ನು} \\
ಕುರಿ + ಅನ್ನು & {\color{kbvgrey}$=$} & ಕು + (ರ್ + ಇ) + (ಅ) + ನ್ನು \\
ಕುರಿಯನ್ನು & {\color{kbvgrey}$=$} & ಕು + (ರ್ + ಇ + ಯ್ + ಅ) + ನ್ನು \\
\end{tabular}}
\end{adjustbox}\end{center}
\vspace{4pt plus 3pt minus 1pt}
\end{kbvderiv}

\indent ಇಲ್ಲಿ ಸ್ವರದ ಮುಂದೆ ಸ್ವರ ಬಂದು ಸಂಧಿಕಾರ್ಯವಾಗುವುದನ್ನು ಕಾಣ\par

\indent ಬಹುದು. ಹಿಂದೆ ತಿಳಿಸಿರುವಂತೆ ಸಂಧಿಮಾಡಿದಲ್ಲಿ ಕುರಿ + ಅನ್ನು = ಕುರೆನ್ನು\kbvcontinues
}
\clearpage

% === tex/0025.tex (Printed Page 23) ===
\renewcommand{\kbvchaptername}{ಕನ್ನಡ ಸಂಧಿಗಳು}
\setcounter{page}{23}
\thispagestyle{kbvpage}
\kbvtypesetpage{\kbvcontinued ಆಗುವುದು. ಹೀಗಾದಲ್ಲಿ ಸಂಧಿಪದವು ಅರ್ಥರಹಿತವಾಗುವುದು. ಇಂತಹ ಸಂದರ್ಭದಲ್ಲಿ ಎರಡೂ ಸ್ವರಗಳ ಮಧ್ಯೆ ಹೊಸದೊಂದು ಅಕ್ಷರ ಸೇರಿ ಸಂಧಿ- ಕಾರ್ಯವಾಗುವುದು. ಮೇಲಿನ ಉದಾಹರಣೆಯಲ್ಲಿ “ಯ್” ಎಂಬ ಅಕ್ಷರವು ಹೊಸದಾಗಿ ಬಂದು ಸೇರಿರುವುದನ್ನು ನೋಡಬಹುದು. ಹೀಗೆ\par

\indent ಸಂಧಿಕಾರ್ಯ ನಡೆಯುವಾಗ ಒಂದು ಅಕ್ಷರವು ಹೊಸದಾಗಿ ಬಂದು ಸೇರುವುದಕ್ಕೆ “ಆಗಮ ಸಂಧಿ” ಎಂದು ಹೆಸರು.\par

\begin{kbvnote}
\noindent{\kbvsecfont ವಿವರಣೆ :}\hspace{0.5em}ಮೇಲಿನ ಉದಾಹರಣೆಯಲ್ಲಿ ಯ-ಕಾರವು ಹೊಸದಾಗಿ ಬಂದು ಸೇರಿದೆ; ಆದುದರಿಂದ ಇದನ್ನು ಯ-ಕಾರಾಗಮ ಸಂಧಿ ಎಂದು ಹೇಳುತ್ತೇವೆ.
\end{kbvnote}

\indent ಬಾಣ + ಅನ್ನು = ಬಾಣವನ್ನು ; ಇಲ್ಲಿ “ವ” ಕಾರವು ಹೊಸದಾಗಿ ಬಂದು ಸೇರಿ ಸಂಧಿಕಾರ್ಯವಾಗಿದೆ. ಇದು ವ-ಕಾರಾಗಮ ಸಂಧಿ.\par

\vspace{3pt plus 2pt}\noindent\hspace*{1.2em}\kbvlabel{ಉದಾ}\par\vspace{1pt}

\begin{kbvrule}
\noindent\kbvbullet\hspace{0.5em}ಯ-ಕಾರ ಆಗಮ ಸಂಧಿ \ensuremath{\square} ವ-ಕಾರ ಆಗಮ ಸಂಧಿ
\end{kbvrule}

\begin{kbvderiv}
\vspace{4pt plus 3pt minus 1pt}
\noindent{\renewcommand{\arraystretch}{1.24}\setlength{\tabcolsep}{5pt}
\begin{tabularx}{\textwidth}{X X}
ಹಣೆ + ಅಲ್ಲಿ = ಹಣೆಯಲ್ಲಿ & ಮಗು + ಇಗೆ = ಮಗುವಿಗೆ \\
ಕೊಠಡಿ + ಇಂದ = ಕೊಠಡಿಯಿಂದ & ಗೋ + ಇಂದ = ಗೋವಿಂದ \\
ಮೈ + ಅನ್ನು = ಮೈಯನ್ನು & ಹೆಣ + ಅನ್ನು = ಹೆಣವನ್ನು \\
ದಾರಿ + ಅನ್ನು = ದಾರಿಯನ್ನು & ಹೂ + ಇನ = ಹೂವಿನ \\
\end{tabularx}}
\vspace{4pt plus 3pt minus 1pt}
\end{kbvderiv}

\kbvsechead{(೩)\hspace{0.45em}ಆದೇಶ ಸಂಧಿ :}

\indent ನಾವೆಲ್ಲರೂ ಜೊತೆಗೂಡಿ ಹೊರಡೋಣ ಎಂಬ ವಾಕ್ಯದಲ್ಲಿ ಜೊತೆಗೂಡಿ ಎಂಬ ಸಂಧಿ ಪದವನ್ನು ಬಿಡಿಸಿ ನೋಡೋಣ.\par

\begin{kbvderiv}
\vspace{4pt plus 3pt minus 1pt}
\begin{center}\begin{adjustbox}{max width=\textwidth}
{\renewcommand{\arraystretch}{1.22}\setlength{\tabcolsep}{4pt}
\begin{tabular}{l c l}
ಜೊತೆ + ಕೂಡಿ & {\color{kbvgrey}$=$} & ಜೊ + (ತ್ + ಎ) + (ಕ್ + ಊ) + ಡಿ \\
ಜೊತೆಗೂಡಿ & {\color{kbvgrey}$=$} & ಜೊ + (ತ್ + ಎ + ಗ್ + ಊ) + ಡಿ \\
\end{tabular}}
\end{adjustbox}\end{center}
\vspace{4pt plus 3pt minus 1pt}
\end{kbvderiv}

\indent ಇಲ್ಲಿ ಸ್ವರದ ಮುಂದೆ ವ್ಯಂಜನ ಬಂದು ಸಂಧಿಕಾರ್ಯವಾಗುವುದನ್ನು ಕಾಣಬಹುದು. ಇಲ್ಲಿ ಉತ್ತರಾರ್ಥ ಭಾಗದ ಮೊದಲಿಗಿರುವ “ಕ್” ವ್ಯಂಜನಕ್ಕೆ ಬದಲಾಗಿ ಸಂಧಿ ಪದದಲ್ಲಿ “ಗ್” ವ್ಯಂಜನ ಸೇರಿ ಸಂಧಿಕಾರ್ಯವಾಗಿರುವುದು.\par
}
\clearpage

% === tex/0026.tex (Printed Page 24) ===
\renewcommand{\kbvchaptername}{ಕನ್ನಡ ಸಂಧಿಗಳು}
\setcounter{page}{24}
\thispagestyle{kbvpage}
\kbvtypesetpage{\indent ಹೀಗೆ, ಸಂಧಿಕಾರ್ಯವು ನಡೆಯುವಾಗ ಒಂದು ಅಪ್ರ\par

\kbvsechead{ಬದಲು ಮತ್ತೊಂದು ಅಕ್ಷರ ಬಂದು ಸೇರುವುದಕ್ಕೆ “ಆದೇಶ ಸಂ:}

\noindent\hspace*{2.2em}ಎಂದು ಹೆಸರು.\par

\vspace{3pt plus 2pt}\noindent\hspace*{1.2em}\kbvlabel{ಉದಾ}\par\vspace{1pt}

\begin{kbvderiv}
\vspace{4pt plus 3pt minus 1pt}
\begin{center}\begin{adjustbox}{max width=\textwidth}
{\renewcommand{\arraystretch}{1.22}\setlength{\tabcolsep}{4pt}
\begin{tabular}{l c l}
ಊರು + ಕೋಲು & {\color{kbvgrey}$=$} & \kbvresult{ಊರುಗೋಲು} \\
ಕೆಳ + ತುಟಿ & {\color{kbvgrey}$=$} & \kbvresult{ಕೆಳದುಟಿ.} \\
\end{tabular}}
\end{adjustbox}\end{center}
\vspace{4pt plus 3pt minus 1pt}
\end{kbvderiv}

\indent ಈ ಉದಾಹರಣೆಗಳಲ್ಲಿ ಕ್ರಮವಾಗಿ ಕ-ಕಾರಕ್ಕೆ ಬದಲಾಗಿ ಗ-ಕಾರವೂ, ತ-ಕಾರಕ್ಕೆ ಬದಲಾಗಿ ದ-ಕಾರವೂ, ಪ-ಕಾರಕ್ಕೆ ಬದಲಾಗಿ ಬ-ಕಾರವೂ ಆದೇಶವಾಗಿ ಬಂದಿರುವುದನ್ನು ಗಮನಿಸಿ ; ಇವುಗಳೆಲ್ಲವೂ ವ್ಯಂಜನಗಳು.\par

\kbvsechead{ನೆನಪಿಡಿ :}

\vspace{2.5pt plus 3pt minus 0.5pt}\noindent\hspace*{1.2em}\kbvnum{(೧)}\hspace{0.55em}ಸ್ವರಾಕ್ಷರದ ಮುಂಜಿ ಸ್ವರಾಕ್ಷರವು ಬಂದು ಸಂಧಿಕಾರ್ಯ ವುಂಟಾದಲ್ಲಿ, ಅದು ಸ್ವರ ಸಂಧಿ ಎನಿಸುವುದು. ಲೋಪ ಮತ್ತು ಆಗಮ ಸಂಧಿಗಳು ಸ್ವರ ಸಂಧಿಗಳು.\par

\vspace{2.5pt plus 3pt minus 0.5pt}\noindent\hspace*{1.2em}\kbvnum{(೨)}\hspace{0.55em}ಸ್ವರಾಕ್ಷರದ ಮುಂಜಿ ವ್ಯಂಜನವು ಬಂದು ಸಂಧಿಕಾರ್ಯ ವುಂಟಾದಲ್ಲಿ, ಅದು ವ್ಯಂಜನ ಸಂಧಿ ಎನಿಸುವುದು. ಆದೇಶ ಸಂಧಿಯು ವ್ಯಂಜನ ಸಂಧಿಯಾಗಿರುವುದು.\par

\vspace{2.5pt plus 3pt minus 0.5pt}\noindent\hspace*{1.2em}\kbvnum{(೩)}\hspace{0.55em}ಸ್ವರದ ಮುಂಜಿ ಸ್ವರವು ಬಂದರೂ ಸಂಧಿಕಾರ್ಯವಾಗದಿರುವ ಸಂದರ್ಭಗಳುಂಟು, ಇದನ್ನು ಪ್ರಕೃತಿಭಾನ ಎನ್ನುವರು. (ಸಂಧಿಕಾರ್ಯ ವಾಗಬೇಕಾದ ಸಂದರ್ಭಗಳಲ್ಲಿ ಪದಗಳು ಇದ್ದ ಹಾಗೆಯೇ ಇರುವುದೇ ಪ್ರಕೃತಿಭಾನ.)\par

\vspace{3pt plus 2pt}\noindent\hspace*{1.2em}\kbvlabel{ಉದಾ}\par\vspace{1pt}

\begin{kbvderiv}
\vspace{4pt plus 3pt minus 1pt}
\noindent{\renewcommand{\arraystretch}{1.24}\setlength{\tabcolsep}{5pt}
\begin{tabularx}{\textwidth}{X X X X}
ಆ ಅಗಸ & ಎಲಾ ಅದೇನು & ಬಾಣ ಇಲ್ಲಿ & ಮುದಿ ಆನೆ \\
ಲಕ್ಷ್ಮಣಾ ! ಅತ್ತ ನೋಡು & ಅಬ್ಬಾ ! ಅದೇನು ಸಂಭ್ರಮ &  &  \\
\end{tabularx}}
\vspace{4pt plus 3pt minus 1pt}
\end{kbvderiv}
}
\clearpage

% === tex/0027.tex (Printed Page 25) ===
\renewcommand{\kbvchaptername}{ಸಂಸ್ಕೃತ ಸಂಧಿಗಳು}
\setcounter{page}{25}
\thispagestyle{kbvpage}
\kbvtypesetpage{\kbvexercise{ಅಭ್ಯಾಸ ಪತ್ರಿಕೆ ೬}

\vspace{2.5pt plus 3pt minus 0.5pt}\noindent\hspace*{1.2em}\kbvnum{(1)}\hspace{0.55em}ಸಂಧಿ ಎಂದರೇನು?\par

\vspace{2.5pt plus 3pt minus 0.5pt}\noindent\hspace*{1.2em}\kbvnum{(2)}\hspace{0.55em}ಕನ್ನಡ ಸಂಧಿಗಳಲ್ಲಿ ಎಷ್ಟು ಪ್ರಕಾರಗಳಿವೆ?\par

\vspace{2.5pt plus 3pt minus 0.5pt}\noindent\hspace*{1.2em}\kbvnum{(3)}\hspace{0.55em}ಕನ್ನಡ ಸಂಧಿಗಳಲ್ಲಿ ಸ್ವರ ಸಂಧಿಗಳು ಯಾವುವು?\par

\vspace{2.5pt plus 3pt minus 0.5pt}\noindent\hspace*{1.2em}\kbvnum{(4)}\hspace{0.55em}ವ್ಯಂಜನ ಸಂಧಿ ಎಂದರೇನು? ಉದಾಹರಣೆ ಸಹಿತ ತಿಳಿಸಿ.\par

\vspace{2.5pt plus 3pt minus 0.5pt}\noindent\hspace*{1.2em}\kbvnum{(5)}\hspace{0.55em}ಲೋಪ ಸಂಧಿ ಎಂದರೇನು? ಉದಾಹರಣೆ ಸಹಿತ ವಿವರಿಸಿ.\par

\vspace{2.5pt plus 3pt minus 0.5pt}\noindent\hspace*{1.2em}\kbvnum{(6)}\hspace{0.55em}ಆಗಮ ಸಂಧಿ ಎಂದರೇನು? ಉದಾಹರಣೆ ಸಹಿತ ವಿವರಿಸಿ.\par

\vspace{2.5pt plus 3pt minus 0.5pt}\noindent\hspace*{1.2em}\kbvnum{(7)}\hspace{0.55em}ಆದೇಶ ಸಂಧಿ ಎಂದರೇನು? ಉದಾಹರಣೆ ಸಹಿತ ವಿವರಿಸಿ.\par

\vspace{2.5pt plus 3pt minus 0.5pt}\noindent\hspace*{1.2em}\kbvnum{(8)}\hspace{0.55em}ಯ-ಕಾರಾಗಮ ಸಂಧಿಗೆ ಎರಡು ಉದಾಹರಣೆಗಳನ್ನು ಕೊಡಿ.\par

\vspace{2.5pt plus 3pt minus 0.5pt}\noindent\hspace*{1.2em}\kbvnum{(9)}\hspace{0.55em}-- ಇವುಗಳನ್ನು ಕೂಡಿಸಿ ಬರೆದು, ಸಂಧಿಯ ಹೆಸರು ತಿಳಿಸಿ.\par

\begin{kbvderiv}
\vspace{4pt plus 3pt minus 1pt}
\noindent{\renewcommand{\arraystretch}{1.24}\setlength{\tabcolsep}{5pt}
\begin{tabularx}{\textwidth}{X X X}
ಹಸಿರು + ಎಲೆ & ಹಣೆ + ಕಣ್ಣು & ಕೆಂಪು + ತುಟಿ \\
ಈಗ + ಈಗ & ಗಿಡ + ಇದು & ತಂತಿ + ಇಂದ \\
ಮೈ + ತಡವು & ಬಾರೆ + ಅಮ್ಮ & ನನಗೆ + ಇದೆ \\
\end{tabularx}}
\vspace{4pt plus 3pt minus 1pt}
\end{kbvderiv}

\vspace{2.5pt plus 3pt minus 0.5pt}\noindent\hspace*{1.2em}\kbvnum{(10)}\hspace{0.55em}ಈ ಸಂಧಿ-ಪದಗಳನ್ನು ಬಿಡಿಸಿ ಬರೆದು, ಸಂಧಿಯ ಹೆಸರು ತಿಳಿಸಿ.\par

\begin{kbvderiv}
\vspace{4pt plus 3pt minus 1pt}
\noindent{\renewcommand{\arraystretch}{1.24}\setlength{\tabcolsep}{5pt}
\begin{tabularx}{\textwidth}{X X X}
ಊರನ್ನು & ಹುಲ್ಲುಗಾವಲು & ಮನೆಯನ್ನು \\
ದೇಶವನ್ನು & ಆಗಾಗ & ಮಳೆಗಾಲ \\
\end{tabularx}}
\vspace{4pt plus 3pt minus 1pt}
\end{kbvderiv}

\kbvchapter{ಪಾಠ ೭}{ಸಂಸ್ಕೃತ ಸಂಧಿಗಳು}

\indent ಕನ್ನಡ ಭಾಷೆಯಲ್ಲಿ ಬಹುತೇಕ ಸಂಸ್ಕೃತ ಶಬ್ದಗಳು ಸೇರಿಹೋಗಿವೆ. ಇಂತಹ ಸಂಸ್ಕೃತ ಶಬ್ದಗಳಲ್ಲಿ ಎರಡು ಶಬ್ದಗಳು ಸೇರಿ ಸಂಧಿಯಾಗಿರುವ ಪಕ್ಷಗಳು ಸಾಕಷ್ಟಿವೆ. ಇಲ್ಲಿ ಸಂಧಿಕಾರ್ಯವು ಸಂಸ್ಕೃತ ಭಾಷಾಶಾಸ್ತ್ರಕ್ಕೆ ಅನುಗುಣವಾಗಿ ನಡೆಯುವ ಕಾರಣ ಸಂಸ್ಕೃತ ಸಂಧಿಗಳನ್ನು ತಿಳಿಯಬೇಕಾದ್ದು.\par

\indent ಸಂಸ್ಕೃತ ಸಂಧಿಗಳನ್ನು ಸ್ವರ ಸಂಧಿಗಳು ಮತ್ತು ವ್ಯಂಜನ ಸಂಧಿಗಳು ಎಂದು ಎರಡು ಬಗೆಯಾಗಿ ವಿಂಗಡಿಸಬಹುದು. ಇವುಗಳಲ್ಲಿ (1) ಸವರ್ಣ ಸಂಧಿ, (2) ಗುಣ ಸಂಧಿ, (3) ವೃದ್ಧಿ ಸಂಧಿ, ಮತ್ತು (4) ಯಣ್‌\kbvcontinues
}
\clearpage

% === tex/0028.tex (Printed Page 26) ===
\renewcommand{\kbvchaptername}{ಸಂಸ್ಕೃತ ಸಂಧಿಗಳು}
\setcounter{page}{26}
\thispagestyle{kbvpage}
\kbvtypesetpage{\kbvcontinued ಸಂಧಿ-ಇವು ಸ್ವರ ಸಂಧಿಗಳು. ವ್ಯಂಜನ ಸಂಧಿಗಳಲ್ಲಿ ಮುಖ್ಯವಾದವುಗಳು\par

\vspace{2.5pt plus 3pt minus 0.5pt}\noindent\hspace*{1.2em}\kbvnum{(೧)}\hspace{0.55em}ಜಶ್ತ್ವ ಸಂಧಿ, (೨) ಶ್ಚುತ್ವ ಸಂಧಿ, ಮತ್ತು (೩) ಅನುನಾಸಿಕ ಸಂಧಿ. ಈ ಸಂಧಿಗಳಿಗೆ ಅನ್ವಯಿಸುವ ಸೂತ್ರಗಳನ್ನು ಉದಾಹರಣೆಗಳ ಸಹಿತ ತಿಳಿಯೋಣ.\par

\kbvsubtitle{ಸ್ವರ ಸಂಧಿಗಳು}

\vspace{4pt plus 4pt minus 1pt}\noindent{\kbvsecfont\fontsize{11pt}{14.5pt}\selectfont\color{kbvaccentdark} (೧)\hspace{0.45em}ಸವರ್ಣದೀರ್ಘ ಸಂಧಿ :}\hspace{0.5em}(ಅ, ಆ); (ಇ, ಈ); (ಉ, ಊ); (ಋ, ೠ); (ಎ, ಏ); (ಒ, ಓ)-ಇವು ಒಂದೇ ಜಾತಿಯ ಸ್ವರಗಳನ್ನುಳ್ಳ ಆರು ಜೊತೆ ಅಕ್ಷರಗಳು. ಇವುಗಳಿಗೆ ಸವರ್ಣ ಸ್ವರಗಳೆಂದು ಹೆಸರು. ಈ ಹಿನ್ನೆಲೆಯಲ್ಲಿ “ದೇವಾಸುರೆ” ಎಂಬ ಶಬ್ದದಲ್ಲಿನ ಸಂಧಿಕಾರ್ಯವನ್ನು ನೋಡೋಣ.\par

\begin{kbvderiv}
\vspace{4pt plus 3pt minus 1pt}
\begin{center}\begin{adjustbox}{max width=\textwidth}
{\renewcommand{\arraystretch}{1.22}\setlength{\tabcolsep}{4pt}
\begin{tabular}{l c l}
ದೇವಾಸುರೆ & {\color{kbvgrey}$=$} & ದೇವ + ಅಸುರ \\
 & {\color{kbvgrey}$=$} & ದೇ + ವ್‌ + (ಅ + ಅ) + ಸುರ \\
\end{tabular}}
\end{adjustbox}\end{center}
\vspace{4pt plus 3pt minus 1pt}
\end{kbvderiv}

\indent ಇಲ್ಲಿ “ಅಕಾರಕ್ಕೆ” ಅಕಾರ ಪರವಾಗಿದೆ ; ಇವು ಸವರ್ಣ ಸ್ವರಗಳು.\par

\indent  = ದೇ + ವ್‌ + (ಆ) + ಸುರ\par

\indent ಇಲ್ಲಿ ಅ, ಅಗಳ ಸ್ಥಾನದಲ್ಲಿ “ಆ” ಕಾರ ಆದೇಶವಾಗಿ ಬಂದಿದೆ.\par

\indent  = ದೇವಾಸುರ (ವ್‌ + ಆ = ವಾ)\par

\indent ಇಲ್ಲಿ ಎರಡು ಸವರ್ಣ ಸ್ವರಗಳಿಗೆ ಒಂದೇ ದೀರ್ಧ ಸ್ವರವು ಆದೇಶವಾಗಿ ಬಂದಿದೆ. ಇದನ್ನು ಸವರ್ಣದೀರ್ಧ ಸಂಧಿ ಎನ್ನುವರು.\par

\begin{kbvrule}
\noindent{\kbvsecfont ಸೂತ್ರ :}\hspace{0.5em}ಸವರ್ಣ ಸ್ವರಗಳು ಪರವಾದಾಗ (ಒಂದರ ಮುಂದೆ ಮತ್ತೊಂದು ಬಂದಾಗ) ಅವೆರಡರ ಸ್ಥಾನದಲ್ಲಿ ಅದೇ ಜಾತಿಯ ದೀರ್ಧ ಸ್ವರವು ಆದೇಶವಾಗಿ ಬರುವುದು. ಇದನ್ನು ಸವರ್ಣ ದೀರ್ಧ ಸಂಧಿಯೆನ್ನುವರು.
\end{kbvrule}

\vspace{3pt plus 2pt}\noindent\hspace*{1.2em}\kbvlabel{ಉದಾ}\par\vspace{1pt}

\begin{kbvderiv}
\vspace{4pt plus 3pt minus 1pt}
\begin{center}\begin{adjustbox}{max width=\textwidth}
{\renewcommand{\arraystretch}{1.22}\setlength{\tabcolsep}{4pt}
\begin{tabular}{l c l}
ಗುರು + ಉಪದೇಶ & {\color{kbvgrey}$=$} & ಗುರೂಪದೇಶ .... (ಉ + ಉ = ಊ) \\
ಸ್ವರ + ಅಕ್ಷರ & {\color{kbvgrey}$=$} & ಸ್ವರಾಕ್ಷರ .... (ಅ + ಅ = ಆ) \\
ಗಿರಿ + ಈಶ & {\color{kbvgrey}$=$} & ಗಿರೀಶ .... (ಇ + ಈ = ಈ) \\
\end{tabular}}
\end{adjustbox}\end{center}
\vspace{4pt plus 3pt minus 1pt}
\end{kbvderiv}
}
\clearpage

% === tex/0029.tex (Printed Page 27) ===
\renewcommand{\kbvchaptername}{ಸಂಸ್ಕೃತ ಸಂಧಿಗಳು}
\setcounter{page}{27}
\thispagestyle{kbvpage}
\kbvtypesetpage{\begin{kbvderiv}
\vspace{4pt plus 3pt minus 1pt}
\begin{center}\begin{adjustbox}{max width=\textwidth}
{\renewcommand{\arraystretch}{1.22}\setlength{\tabcolsep}{4pt}
\begin{tabular}{l c l}
ಗಿರಿ + ಇಂದ್ರ & {\color{kbvgrey}$=$} & ಗಿರೀಂದ್ರ .....(ಇ + ಇ = ಈ) \\
ಧ್ವಜ + ಆರೋಹಣ & {\color{kbvgrey}$=$} & ಧ್ವಜಾರೋಹಣ .....(ಅ + ಆ = ಆ) \\
ಲಕ್ಷ್ಮೀ + ಈಶ & {\color{kbvgrey}$=$} & ಲಕ್ಷ್ಮೀಶ .....(ಈ + ಈ = ಈ) \\
\end{tabular}}
\end{adjustbox}\end{center}
\vspace{4pt plus 3pt minus 1pt}
\end{kbvderiv}

\vspace{4pt plus 4pt minus 1pt}\noindent{\kbvsecfont\fontsize{11pt}{14.5pt}\selectfont\color{kbvaccentdark} (೨)\hspace{0.45em}ಗುಣ ಸಂಧಿ :}\hspace{0.5em}“ಗಣೇಶ” ಈ ಶಬ್ದವನ್ನು ಬಿಡಿಸಿ ನೋಡೋಣ.\par

\begin{kbvderiv}
\vspace{4pt plus 3pt minus 1pt}
\begin{center}\begin{adjustbox}{max width=\textwidth}
{\renewcommand{\arraystretch}{1.22}\setlength{\tabcolsep}{4pt}
\begin{tabular}{l c l}
ಗಣೇಶ & {\color{kbvgrey}$=$} & ಗಣ + ಈಶ \\
 & {\color{kbvgrey}$=$} & ಗ + ಣ್‌ + (ಅ + ಈ) + ಶ \\
\end{tabular}}
\end{adjustbox}\end{center}
\vspace{4pt plus 3pt minus 1pt}
\end{kbvderiv}

\noindent\hspace*{2.2em}ಇಲ್ಲಿ ಅ ಕಾರಕ್ಕೆ ಈ ಕಾರ ಪರವಾಗಿದೆ.\par

\indent  = ಗ + ಣ್‌ + (ಏ) + ಶ ಇಲ್ಲಿ ಅ, ಈ ಗಳ ಸ್ಥಾನದಲ್ಲಿ ಏ ಕಾರ ಆದೇಶವಾಗಿ ಬಂದಿದೆ.\par

\indent  = ಗಣೇಶ .....(\ensuremath{\because} ಣ್‌ + ಈ = ಣೇ) ತೀರ್ಮಾನ : ಅ ಕಾರಕ್ಕೆ ಈ ಕಾರ ಪರವಾದಾಗ ಏ ಕಾರ ಆದೇಶವಾಗಿ ಬರುವುದು.\par

\indent ಇದೇ ರೀತಿಯಾಗಿ,\par

\begin{kbvderiv}
\vspace{4pt plus 3pt minus 1pt}
\begin{center}\begin{adjustbox}{max width=\textwidth}
{\renewcommand{\arraystretch}{1.22}\setlength{\tabcolsep}{4pt}
\begin{tabular}{l c l}
ದೇವ + ಇಂದ್ರ & {\color{kbvgrey}$=$} & ದೇವೇಂದ್ರ .....(ಅ + ಇ = ಏ) \\
ಮಹಾ + ಈಶ್ವರೀ & {\color{kbvgrey}$=$} & ಮಹೇಶ್ವರೀ .....(ಆ + ಈ = ಏ) \\
ಜನ + ಉಪಕಾರಿ & {\color{kbvgrey}$=$} & ಜನೋಪಕಾರಿ .....(ಅ + ಉ = ಓ) \\
ಮಹಾ + ಉನ್ನತಿ & {\color{kbvgrey}$=$} & ಮಹೋನ್ನತಿ .....(ಅ + ಉ = ಓ) \\
ರಾಜ + ಋಷಿ & {\color{kbvgrey}$=$} & ರಾಜರ್ಷಿ .....(ಅ + ಋ = ಅರ್‌) \\
ದೇವ + ಋಷಿ & {\color{kbvgrey}$=$} & ದೇವರ್ಷಿ .....(ಅ + ಋ = ಅರ್‌) \\
\end{tabular}}
\end{adjustbox}\end{center}
\vspace{4pt plus 3pt minus 1pt}
\end{kbvderiv}

\indent ಈ ಉದಾಹರಣೆಗಳಲ್ಲಿ ಪೂರ್ವದ ಸ್ವರವು “ಅ” ಅಥವಾ “ಆ” ಆಗಿರುವುದು ; ಅವುಗಳಿಗೆ ಪರವಾಗಿರುವ ಸ್ವರಗಳು ಇ, ಈ, ಉ, ಋ ಆಗಿರುವುವು. ಅ, ಆಗಳಿಗೆ ಇ, ಈ ಗಳು ಪರವಾದಾಗ ಏಕಾರವೂ, ಅಕಾರಕ್ಕೆ ಉಕಾರ ಪರವಾದಾಗ ಓಕಾರವೂ, ಅಕಾರಕ್ಕೆ ಋಕಾರ ಪರವಾದಾಗ “ಅರ್‌” ಕಾರವೂ, ಆದೇಶಗಳಾಗಿ ಬಂದಿರುವುವು. ಈ ತೆರನಾದ ಸಂಧಿಯು ಗುಣ ಸಂಧಿಯೆನಿಸುವುದು.\par
}
\clearpage

% === tex/0030.tex (Printed Page 28) ===
\renewcommand{\kbvchaptername}{ಸಂಸ್ಕೃತ ಸಂಧಿಗಳು}
\setcounter{page}{28}
\thispagestyle{kbvpage}
\kbvtypesetpage{\begin{kbvrule}
\noindent{\kbvsecfont ಸೂತ್ರ :}\hspace{0.5em}ಸಂಧಿಕಾರ್ಯ ನಡೆಯುವಾಗ ಅ, ಆ ಕಾರಗಳಿಗೆ ಕರಳು ಪರವಾದರೆ “ಏ” ಕಾರವೂ, ಉ, ಊ ಕಾರಗಳು ಪರವಾದ ಕಾರವು ಋಕಾರವು ಪರವಾದರೆ “ಅರ್” ಕಾರವೂ ಆದೇಶಗಳಾಗಿ ಬರುವುವು. ಅದು ಗುಣಸಂಧಿ ಎನಿಸುವುದು.
\end{kbvrule}

\vspace{4pt plus 4pt minus 1pt}\noindent{\kbvsecfont\fontsize{11pt}{14.5pt}\selectfont\color{kbvaccentdark} (೩)\hspace{0.45em}ವೃದ್ಧಿ ಸಂಧಿ :}\hspace{0.5em}“ಏಕೈಕ” — ಈ ಶಬ್ದವನ್ನು ಬಿಡಿಸಿ ನೋಡೋಣ.\par

\begin{kbvderiv}
\vspace{4pt plus 3pt minus 1pt}
\begin{center}\begin{adjustbox}{max width=\textwidth}
{\renewcommand{\arraystretch}{1.22}\setlength{\tabcolsep}{4pt}
\begin{tabular}{l c l}
ಏಕೈಕ & {\color{kbvgrey}$=$} & ಏಕ + ಏಕ \\
 & {\color{kbvgrey}$=$} & ಏಕ್ + (ಅ + ಏ) + ಕ \\
\end{tabular}}
\end{adjustbox}\end{center}
\vspace{4pt plus 3pt minus 1pt}
\end{kbvderiv}

\indent ಇಲ್ಲಿ ಅ ಕಾರಕ್ಕೆ ಏ ಕಾರ ಪರವಾಗಿದೆ. ಇಲ್ಲಿ ಅ, ಏ ಗಳ ಸ್ಥಾನದಲ್ಲಿ “ಐ” ಕಾರ ಆದೇಶವಾಗಿ ಬಂದಿದೆ.\par

\indent  = ಏಕೈಕ (“ಕ್” — ಐಕ್)\par

\indent ತೀರ್ಮಾನ : ಅಕಾರಕ್ಕೆ ಏಕಾರ ಪರವಾದಾಗ ಐಕಾರ ಆದೇಶವಾಗಿ ಬರುವುದು.\par

\indent ಇದೇ ರೀತಿಯಾಗಿ,\par

\begin{kbvderiv}
\vspace{4pt plus 3pt minus 1pt}
\begin{center}\begin{adjustbox}{max width=\textwidth}
{\renewcommand{\arraystretch}{1.22}\setlength{\tabcolsep}{4pt}
\begin{tabular}{l c l}
ಅಷ್ಟ + ಐಶ್ವರ್ಯ & {\color{kbvgrey}$=$} & ಅಷ್ಟೈಶ್ವರ್ಯ (ಅ + ಐ = ಐ) \\
ವನ + ಓಷಧಿ & {\color{kbvgrey}$=$} & ವನೌಷಧಿ (ಅ + ಓ = ಔ) \\
ಮಹಾ + ಔದಾರ್ಯ & {\color{kbvgrey}$=$} & ಮಹೌದಾರ್ಯ (ಆ + ಔ = ಔ) \\
ಭಾವ + ಐಕ್ಯತೆ & {\color{kbvgrey}$=$} & ಭಾವೈಕ್ಯತೆ (ಅ + ಐ = ಐ) \\
\end{tabular}}
\end{adjustbox}\end{center}
\vspace{4pt plus 3pt minus 1pt}
\end{kbvderiv}

\kbvsechead{ಈ ಉದಾಹರಣೆಗಳಲ್ಲಿ ಪೂರ್ವದ ಸ್ವರವು ಅ ಅಥವಾ ಆ ಆಗಿರುವುದು :}

\indent ಅವುಗಳಿಗೆ ಪರವಾಗಿರುವ ಸ್ವರಗಳು ಏ, ಐ, ಓ, ಔ ಆಗಿರುವುವು. ಅ, ಆ ಗಳಿಗೆ ಏ, ಐ ಗಳು ಪರವಾದಾಗ “ಐ” ಕಾರವೂ ಓ, ಔ ಗಳು ಪರವಾದಾಗ “ಔ” ಕಾರವೂ ಆದೇಶಗಳಾಗಿ ಬಂದಿರುವುವು. ಈ ತೆರನಾದ ಸಂಧಿಯು ವೃದ್ಧಿ ಸಂಧಿಯೆನಿಸುವುದು.\par
}
\clearpage

% === tex/0031.tex (Printed Page 29) ===
\renewcommand{\kbvchaptername}{ಸಂಸ್ಕೃತ ಸಂಧಿಗಳು}
\setcounter{page}{29}
\thispagestyle{kbvpage}
\kbvtypesetpage{\indent ಬ್ರಹ್ಮಸೂತ್ರ : ಸಂಧಿಕಾರ್ಯ ನಡೆಯುವಾಗ ಅ, ಆ ಕಾರಗಳಲ್ಲಿ ಯಾವುದಾದರೊಂದು ಸ್ವರಕ್ಕೆ ಏ, ಐ ಕಾರಗಳು ಪರವಾದರೆ ಐ ಕಾರವೂ ಓ, ಔ ಕಾರಗಳು ಪರವಾದರೆ ಔಕಾರವೂ ಆದೇಶಗಳಾಗಿ ಬರುವುದಾದರೆ ಅದು ವೃದ್ಧಿಸಂಧಿಯೆನಿಸುವುದು.\par

\vspace{4pt plus 4pt minus 1pt}\noindent{\kbvsecfont\fontsize{11pt}{14.5pt}\selectfont\color{kbvaccentdark} {[}೪{]}\hspace{0.45em}ಯಣ್‌ ಸಂಧಿ :}\hspace{0.5em}ಇತ್ಯಾದಿ, ಗುರ್ವಾಜ್ಞೆ, ನಿತ್ರಾರ್ಜಿತ - ಈ ಶಬ್ದಗಳನ್ನು ಬಿಡಿಸಿ ನೋಡೋಣ...\par

\begin{kbvderiv}
\vspace{4pt plus 3pt minus 1pt}
\begin{center}\begin{adjustbox}{max width=\textwidth}
{\renewcommand{\arraystretch}{1.22}\setlength{\tabcolsep}{4pt}
\begin{tabular}{l c l}
{[}{\kbvlatin i}{]} ಇತ್ಯಾದಿ & {\color{kbvgrey}$=$} & ಇತಿ + ಆದಿ \\
 & {\color{kbvgrey}$=$} & ಇ + {[}ತ್ + ಆ + ದಿ{]} + ಆ \\
\end{tabular}}
\end{adjustbox}\end{center}
\vspace{4pt plus 3pt minus 1pt}
\end{kbvderiv}

\indent ಸ ಇಕಾರಕ್ಕೆ ಯ್‌ಕಾರಾದೇಶವಾಗಲು\par

\begin{kbvderiv}
\vspace{4pt plus 3pt minus 1pt}
\begin{center}\begin{adjustbox}{max width=\textwidth}
{\renewcommand{\arraystretch}{1.22}\setlength{\tabcolsep}{4pt}
\begin{tabular}{l c l}
 & {\color{kbvgrey}$=$} & ಇ + (ತ್ + ಆ + ದಿ) + ಆ \\
 & {\color{kbvgrey}$=$} & ಇತ್ಯಾದಿ.... {[}\ensuremath{\because} ತ್ + ಆ + ದಿ = ತ್ಯಾದಿ{]} \\
\end{tabular}}
\end{adjustbox}\end{center}
\vspace{4pt plus 3pt minus 1pt}
\end{kbvderiv}

\begin{kbvderiv}
\vspace{4pt plus 3pt minus 1pt}
\begin{center}\begin{adjustbox}{max width=\textwidth}
{\renewcommand{\arraystretch}{1.22}\setlength{\tabcolsep}{4pt}
\begin{tabular}{l c l}
{[}{\kbvlatin ii}{]} ಗುರ್ವಾಜ್ಞೆ & {\color{kbvgrey}$=$} & ಗುರು + ಆಜ್ಞೆ \\
 & {\color{kbvgrey}$=$} & ಗು + {[}ವ್ + ಆ + ಜ್ಞೆ{]} + ಆ \\
 & {\color{kbvgrey}$=$} & ಗುರ್ವಾಜ್ಞೆ.... {[}\ensuremath{\because} ವ್ + ಆ + ಜ್ಞೆ = ವಾಜ್ಞೆ{]} \\
\end{tabular}}
\end{adjustbox}\end{center}
\vspace{4pt plus 3pt minus 1pt}
\end{kbvderiv}

\begin{kbvderiv}
\vspace{4pt plus 3pt minus 1pt}
\begin{center}\begin{adjustbox}{max width=\textwidth}
{\renewcommand{\arraystretch}{1.22}\setlength{\tabcolsep}{4pt}
\begin{tabular}{l c l}
{[}{\kbvlatin iii}{]} ನಿತ್ರಾರ್ಜಿತ & {\color{kbvgrey}$=$} & ಪಿತೃ + ಆರ್ಜಿತ \\
 & {\color{kbvgrey}$=$} & ಪಿ + {[}ತ್ + ಋ + ಆ{]} + ರ್ಜಿತ \\
 & {\color{kbvgrey}$=$} & ಪಿ + (ತ್ + ಋ + ಆ) + ರ್ಜಿತ \\
 & {\color{kbvgrey}$=$} & ನಿತ್ರಾರ್ಜಿತ.... (\ensuremath{\because} ತ್ + ಋ + ಆ = ತ್ರಾ) \\
\end{tabular}}
\end{adjustbox}\end{center}
\vspace{4pt plus 3pt minus 1pt}
\end{kbvderiv}

\indent ಮೇಲಿನ ಉದಾಹರಣೆಗಳಲ್ಲಿ ಪೂರ್ವಾರ್ಧದಲ್ಲಿರುವ ಸ್ವರಾಕ್ಷರಗಳಾದ ಇ, ಉ, ಋ ಕಾರಗಳಿಗೆ ಕ್ರಮವಾಗಿ ಯ್‌, ವ್‌, ರ್‌ಗಳು ಆದೇಶವಾಗಿ ಬಂದಿರುವುದನ್ನು ನೋಡಬಹುದು ; ಇಲ್ಲಿ ಸವರ್ಣ ಸ್ವರ ಪರವಾಗಿ ಬಂದಿಲ್ಲದಿರುವುದು ಗಮನಾರ್ಹ ಅಂಶ. ಈ ತೆರನಾದ ಸಂಧಿಯು 'ಯಣ್‌ ಸಂಧಿ' ಎಂದು ಬಿಂದುಸುವುದು.\par
}
\clearpage

% === tex/0032.tex (Printed Page 30) ===
\renewcommand{\kbvchaptername}{ಸಂಸ್ಕೃತ ಸಂಧಿಗಳು}
\setcounter{page}{30}
\thispagestyle{kbvpage}
\kbvtypesetpage{\begin{kbvrule}
\noindent{\kbvsecfont ಸೂತ್ರ :}\hspace{0.5em}ಸಂಧಿ ಕಾರ್ಯದಲ್ಲಿ ಇ, ಈ, ಉ, ಊ, ಋ-ಕಾರಗಳು ಸವರ್ಣವಲ್ಲದ ಸ್ವರ ಪರವಾದರೆ ಇ, ಈ ಕಾರಗಳಿಗೆ ಯ ಕಾರವೂ ಉ, ಊ ಕಾರಗಳಿಗೆ ವ್‌-ಕಾರವೂ ಋ-ಕಾರಕ್ಕೆ ರ್‌-ಕಾರವೂ ಆದೇಶಗಳಾಗಿ ಬರುವು ದಾದರೆ ಅದು ಯಣ್‌ ಸಂಧಿ ಎನಿಸುವುದು.
\end{kbvrule}

\begin{kbvderiv}
\vspace{4pt plus 3pt minus 1pt}
\begin{center}\begin{adjustbox}{max width=\textwidth}
{\renewcommand{\arraystretch}{1.22}\setlength{\tabcolsep}{4pt}
\begin{tabular}{l c l}
ಗತಿ + ಅಂತರ & {\color{kbvgrey}$=$} & \kbvresult{ಗತ್ಯಂತರ} \\
ಅತಿ + ಅಂತ & {\color{kbvgrey}$=$} & \kbvresult{ಅತ್ಯಂತ} \\
ಆದಿ + ಅಂತ & {\color{kbvgrey}$=$} & \kbvresult{ಆದ್ಯಂತ} \\
ಮನು + ಅಂತರ & {\color{kbvgrey}$=$} & \kbvresult{ಮನ್ವಂತರ} \\
\end{tabular}}
\end{adjustbox}\end{center}
\vspace{4pt plus 3pt minus 1pt}
\end{kbvderiv}

\indent : — ವ್ಯಂಜನ ಸಂಧಿಗಳು\par

\vspace{4pt plus 4pt minus 1pt}\noindent{\kbvsecfont\fontsize{11pt}{14.5pt}\selectfont\color{kbvaccentdark} (೧)\hspace{0.45em}ಜಶ್ತ್ವ ಸಂಧಿ:}\hspace{0.5em}ಕೆಳಗೆ ಕೊಟ್ಟಿರುವ ಉದಾಹರಣೆಗಳನ್ನು ನೋಡೋಣ.\par

\indent ದಿಕ್‌ + ದರ್ಶನ = ದಿಗ್‌ + ದರ್ಶನ = ದಿಗ್ದರ್ಶನ (ಕಕಾರಕ್ಕೆ ಗಕಾರಾದೇಶ)\par

\begin{kbvderiv}
\vspace{4pt plus 3pt minus 1pt}
\begin{center}\begin{adjustbox}{max width=\textwidth}
{\renewcommand{\arraystretch}{1.22}\setlength{\tabcolsep}{4pt}
\begin{tabular}{l c l}
ಅಚ್‌ + ಅಂತ & {\color{kbvgrey}$=$} & ಅಜ್‌ + ಅಂತ = ಅಜಂತ \\
ಷಟ್‌ + ಆನನ & {\color{kbvgrey}$=$} & ಷಡ್‌ + ಆನನ = ಷಡಾನನ \\
\end{tabular}}
\end{adjustbox}\end{center}
\vspace{4pt plus 3pt minus 1pt}
\end{kbvderiv}

\indent (ಟಕಾರಕ್ಕೆ ದಕಾರಾದೇಶ)\par

\indent ಜಗತ್‌ + ಈಶ್ವರ = ಜಗದ್‌ + ಈಶ್ವರ = ಜಗದೀಶ್ವರ (ತಕಾರಕ್ಕೆ ದಕಾರಾದೇಶ)\par

\indent ಅಪ್‌ + ಜ = ಅಬ್‌ + ಜ = ಅಬ್ಜ (ಪಕಾರಕ್ಕೆ ಬಕಾರಾದೇಶ)\par

\indent ಇಲ್ಲಿ ಐದು ಉದಾಹರಣೆಗಳಿವೆ. ಈ ಉದಾಹರಣೆಗಳಲ್ಲಿ ಪೂರ್ವ ಶಬ್ದದ ಕೊನೆಯಲ್ಲಿರುವ ಕ್‌, ಚ್‌, ಟ್‌, ತ್‌, ಪ್‌ ಗಳಿಗೆ ಕ್ರಮವಾಗಿ ಗ್‌, ಜ್‌, ಡ್‌, ದ್‌, ಬ್‌ ಗಳು ಆದೇಶಗಳಾಗಿ ಬಂದು ಸಂಧಿಕಾರ್ಯ ವೇರ್ಪಟ್ಟಿದೆ. ಈ ತೆರನಾದ ಸಂಧಿಯು ಜಶ್ತ್ವ ಸಂಧಿ ಎನಿಸುವುದು.\par
}
\clearpage

% === tex/0033.tex (Printed Page 31) ===
\renewcommand{\kbvchaptername}{ಸಂಸ್ಕೃತ ಸಂಧಿಗಳು}
\setcounter{page}{31}
\thispagestyle{kbvpage}
\kbvtypesetpage{\begin{kbvrule}
\noindent{\kbvsecfont ಸೂತ್ರ :}\hspace{0.5em}ಸಂಧಿಕಾರ್ಯದಲ್ಲಿ ಪೂರ್ವಶಬ್ದದ ಅಂತ್ಯದಲ್ಲಿರುವ ಕ, ಚ, ಟ, ತ, ಪ ವ್ಯಂಜನಗಳಿಗೆ ಯಾವುದೇ ವರ್ಣ ಪರವಾದರೂ ಕ್ರಮವಾಗಿ ಗ, ಜ, ಡ, ದ, ಬ ವ್ಯಂಜನಗಳು ಆದೇಶಗಳಾಗಿ ಬರುವುದಾದರೆ ಅದು ಜಶ್ತ್ವ ಸಂಧಿ ಯೆನಿಸುವುದು.
\end{kbvrule}

\kbvsechead{(೧)\hspace{0.45em}ಶ್ಚುತ್ವ ಸಂಧಿ :}

\noindent\hspace*{2.2em}ಈ ಉದಾಹರಣೆಗಳನ್ನು ಗಮನಿಸಿ...\par

\indent ({\kbvlatin i}) ಮನಸ್‌ + ಚಂಚಲ = ಮನಶ್‌ + ಚಂಚಲ = ಮನಶ್ಚಂಚಲ (ಸಕಾರಕ್ಕೆ ಚಕಾರ ಪರವಾಗಿದ್ದು ಸಕಾರಕ್ಕೆ ಶಕಾರಾದೇಶವಾಗಿದೆ)\par

\indent ({\kbvlatin ii}) ತಕಾರ + ಚಕಾರ = ತಕಾರ + ಚಕಾರ = ತಕಾರಚಕಾರ (ತಕಾರಕ್ಕೆ ಚಕಾರ ಪರವಾಗಿದ್ದು ತಕಾರಕ್ಕೆ ಚಕಾರಾದೇಶವಾಗಿದೆ)\par

\indent ({\kbvlatin iii}) ಬೃಹತ್‌ + ಛತ್ರ = ಬೃಹಚ್‌ + ಛತ್ರ = ಬೃಹಚ್ಛತ್ರ (ತಕಾರಕ್ಕೆ ಛಕಾರ ಪರವಾಗಿದ್ದು ತಕಾರಕ್ಕೆ ಚಕಾರಾದೇಶವಾಗಿದೆ)\par

\indent ({\kbvlatin iv}) ಮನಸ್‌ + ಶುದ್ಧಿ = ಮನಶ್‌ + ಶುದ್ಧಿ = ಮನಶ್ಶುದ್ಧಿ (ಸಕಾರಕ್ಕೆ ಶಕಾರ ಪರವಾಗಿದ್ದು ಸಕಾರಕ್ಕೆ ಶಕಾರಾದೇಶವಾಗಿದೆ)\par

\indent —ಇವು ಶ್ಚುತ್ವ ಸಂಧಿಗೆ ಉದಾಹರಣೆಗಳು. ಶ್ಚುತ್ವ ಸಂಧಿ ಯಾಗುವ ಸಂದರ್ಭಗಳನ್ನು ಕೆಳಗೆ ಕೊಟ್ಟಿರುವ ಸೂತ್ರದಿಂದ ತಿಳಿಯಬಹುದು.\par

\begin{kbvrule}
\noindent{\kbvsecfont ಸೂತ್ರ :}\hspace{0.5em}ಪೂರ್ವ ಪದದ ಅಂತ್ಯದಲ್ಲಿ ಸಕಾರವಾಗಲೀ ತವರ್ಗದಲ್ಲಿನ ಯಾವುದಾದರೊಂದು ವರ್ಣವಾಗಲೀ ಇದ್ದು, ಆ ವರ್ಣಗಳಿಗೆ ಶಕಾರವಾಗಲೀ ಚವರ್ಗದಲ್ಲಿನ ಯಾವುದಾದರೊಂದು ವರ್ಣವಾಗಲೀ ಪರವಾದಾಗ, ಸಕಾರಕ್ಕೆ ಶಕಾರವೂ ತವರ್ಗಕ್ಕೆ ಚವರ್ಗವೂ ಆದೇಶಗಳಾಗಿ ಬರುವುದಾದರೆ ಅದು ಶ್ಚುತ್ವ ಸಂಧಿಯೆನಿಸುವುದು.
\end{kbvrule}

\kbvsechead{(೨)\hspace{0.45em}ಅನುನಾಸಿಕ ಸಂಧಿ :}

\indent ಈ ಉದಾಹರಣೆಗಳನ್ನು ಗಮನಿಸಿ—\par

\indent ವಾಕ್‌ + ಮಯ = ವಾಙ್‌ + ಮಯ = ವಾಙ್ಮಯ (ಕಕಾರಕ್ಕೆ ಮಕಾರ ಪರವಾಗಿದ್ದು ಙ ಕಾರ ಆದೇಶವಾಗಿದೆ)\par
}
\clearpage

% === tex/0034.tex (Printed Page 32) ===
\renewcommand{\kbvchaptername}{ಸಂಸ್ಕೃತ ಸಂಧಿಗಳು}
\setcounter{page}{32}
\thispagestyle{kbvpage}
\kbvtypesetpage{\indent ಷಟ್ + ಮಾಸ = ಷಣ್ + ಮಾಸ = ಷಣ್ಮಾಸ (ಟಕಾರಕ್ಕೆ ಮಕಾರ ಪರವಾಗಿದ್ದು ಣಕಾರ ಆದೇಶವಾಗಿದೆ\par

\indent ಚಿತ್ + ಮಯ = ಚಿನ್ + ಮಯ = ಚಿನ್ಮಯ (ತಕಾರಕ್ಕೆ ಮಕಾರ ಪರವಾಗಿದ್ದು ನಕಾರ ಆದೇಶವಾಗಿದೆ ಇವು ಅನುನಾಸಿಕ ಸಂಧಿಗೆ ಉದಾಹರಣೆಗಳು ; ಅನುನಾಸಿಕ ಸಂಧಿ ಯಾಗುವ ಸಂದರ್ಭಗಳನ್ನು ಕೆಳಗೆ ಕೊಟ್ಟಿರುವ ಸೂತ್ರದಿಂದ ತಿಳಿಯಬಹುದು\par

\begin{kbvrule}
\noindent{\kbvsecfont ಸೂತ್ರ :}\hspace{0.5em}ಪೂರ್ವ ಶಬ್ದದ ಅಂತ್ಯದಲ್ಲಿರುವ ಕ, ಚ, ಟ, ತ, ಪ ವ್ಯಂಜನಗಳಿಗೆ ಯಾವುದೇ ಅನುನಾಸಿಕಾಕ್ಷರ ಪರವಾದರೂ ಕ ಚ ಟ ತ ಪ ಗಳಿಗೆ ಙ, ಞ, ಣ, ನ, ಮ ಗಳು ಕ್ರಮವಾಗಿ ಆದೇಶಗಳಾಗಿ ಬರುವುದಾದರೆ ಅದು ಅನುನಾಸಿಕ ಸಂಧಿಯೆನಿಸುವುದು.
\end{kbvrule}

\kbvexercise{ಅಭ್ಯಾಸ ಪತ್ರಿಕೆ ೭}

\vspace{2.5pt plus 3pt minus 0.5pt}\noindent\hspace*{1.2em}\kbvnum{{[}1{]}}\hspace{0.55em}ಸ್ವರ ಸಂಧಿ ಎಂದರೇನು ? ಸಂಸ್ಕೃತ ಸ್ವರ ಸಂಧಿಗಳಲ್ಲಿ ಎಷ್ಟು ಪ್ರಕಾರಗಳಿವೆ ?\par

\vspace{2.5pt plus 3pt minus 0.5pt}\noindent\hspace*{1.2em}\kbvnum{{[}2{]}}\hspace{0.55em}ವ್ಯಂಜನ ಸಂಧಿ ಎಂದರೇನು ? ಮುಖ್ಯವಾದ ಮೂರು ಸಂಸ್ಕೃತ ವ್ಯಂಜನ ಸಂಧಿಗಳನ್ನು ತಿಳಿಸಿ.\par

\vspace{2.5pt plus 3pt minus 0.5pt}\noindent\hspace*{1.2em}\kbvnum{{[}3{]}}\hspace{0.55em}ಸವರ್ಣ ದೀರ್ಘ ಸಂಧಿ ಎಂದರೇನು? ಒಂದು ಉದಾಹರಣೆ ಸಹಿತ ವಿವರಿಸಿ.\par

\vspace{2.5pt plus 3pt minus 0.5pt}\noindent\hspace*{1.2em}\kbvnum{{[}4{]}}\hspace{0.55em}ಗುಣ ಸಂಧಿ ಎಂದರೇನು ? ಉದಾಹರಣೆಗಳ ಸಹಿತ ತಿಳಿಸಿ.\par

\vspace{2.5pt plus 3pt minus 0.5pt}\noindent\hspace*{1.2em}\kbvnum{{[}5{]}}\hspace{0.55em}ವೃದ್ಧಿ ಸಂಧಿ ಎಂದರೇನು ? ಉದಾಹರಣೆಗಳ ಸಹಿತ ತಿಳಿಸಿ.\par

\vspace{2.5pt plus 3pt minus 0.5pt}\noindent\hspace*{1.2em}\kbvnum{{[}6{]}}\hspace{0.55em}ಯಣ್ ಸಂಧಿಯ ಸೂತ್ರವನ್ನು ನಿರೂಪಿಸಿ, ಉದಾಹರಣೆಗಳ ಮೂಲಕ ಸಮರ್ಥಿಸಿ?\par

\vspace{2.5pt plus 3pt minus 0.5pt}\noindent\hspace*{1.2em}\kbvnum{{[}7{]}}\hspace{0.55em}ಜಶ್ತ್ವ ಸಂಧಿಗೆ ಎರಡು ಉದಾಹರಣೆಗಳನ್ನು ಕೊಡಿ.\par

\vspace{2.5pt plus 3pt minus 0.5pt}\noindent\hspace*{1.2em}\kbvnum{{[}8{]}}\hspace{0.55em}ಶ್ಚುತ್ವ ಸಂಧಿಯ ಸೂತ್ರವನ್ನು ನಿರೂಪಿಸಿ.\par

\vspace{2.5pt plus 3pt minus 0.5pt}\noindent\hspace*{1.2em}\kbvnum{{[}9{]}}\hspace{0.55em}ಅನುನಾಸಿಕ ಸಂಧಿ ಎಂದರೇನು ? ಉದಾಹರಣೆ ಸಹಿತ ತಿಳಿಸಿ.\par

\vspace{2.5pt plus 3pt minus 0.5pt}\noindent\hspace*{1.2em}\kbvnum{{[}10{]}}\hspace{0.55em}ಇವುಗಳನ್ನು ಕೂಡಿಸಿ ಬರೆದು ಸಂಧಿಯನ್ನು ತಿಳಿಸಿರಿ.\par

\indent ಅಷ್ಟ್ + ಮಾಸಂ; ಮಾಸ್ + ಮಾನ; ಮಾಸಂ + ಮಾಸಂ\par
}
\clearpage

% === tex/0035.tex (Printed Page 33) ===
\renewcommand{\kbvchaptername}{ನಾಮಪದದ ವಿಭೇದಗಳು}
\setcounter{page}{33}
\thispagestyle{kbvpage}
\kbvtypesetpage{\begin{kbvderiv}
\vspace{4pt plus 3pt minus 1pt}
\noindent{\renewcommand{\arraystretch}{1.24}\setlength{\tabcolsep}{5pt}
\begin{tabularx}{\textwidth}{X X X}
ಷಟ್‌ + ಚಕ್ರ & ವಾಕ್‌ + ದೇವಿ & ಪೂರ್ವ + ಉತ್ತರ \\
ದುರ್ಗ + ಅಷ್ಟಮಿ & ಅತಿ + ಆಸೆ & ಯಶಸ್‌ + ಚಂದ್ರಿಕೆ \\
\end{tabularx}}
\vspace{4pt plus 3pt minus 1pt}
\end{kbvderiv}

\vspace{2.5pt plus 3pt minus 0.5pt}\noindent\hspace*{1.2em}\kbvnum{(11)}\hspace{0.55em}ಇವುಗಳನ್ನು ಬಿಡಿಸಿ ಬರೆದು ಸಂಧಿಯನ್ನು ತಿಳಿಸಿ. ಧನಾರ್ಜನೆ, ಮಹೇಶ್ವರ, ಸಿದ್ಧೌಷಧಿ, ಪ್ರತ್ಯುತ್ತರ... ಜಗದ್ರಕ್ಷಕ, ಷಡ್ದರ್ಶನ, ಚಿದಾನಂದ, ಮನಶ್ಶಾಂತಿ ಕಠಶೂಕ್ತಿ, ವೀರೇಂದ್ರ, ಮಹೌನ್ನತ್ಯ, ಷಡ್ಭುಜ - ಸಚ್ಚರಿತ್ರ, ಅಬ್ಧಿ, ದಿಗಂತ, ಅತ್ಯವಸರ\par

\vspace{2.5pt plus 3pt minus 0.5pt}\noindent\hspace*{1.2em}\kbvnum{(12)}\hspace{0.55em}. ಇವುಗಳನ್ನು ಬಿಡಿಸಿ ಬರೆಯಿರಿ ; ಬ್ರಹ್ಮರ್ಷಿ, ವೀರ್ಯಾಣು, ಅತುಳೈಶ್ವರ್ಯ, ದೇವಾಲಯ ಷಡ್ದಳ ಜನೈಕ್ಯತೆ, ಮನಶ್ಚಾಪಲ್ಯ, . ವಿಜಯೋತ್ಸವ\par

\kbvchapter{ಪಾಠ ೮}{ನಾಮಪದದ ವಿಭೇದಗಳು}

\indent “ರವಿಯು ತೋಟಕ್ಕೆ ನುಗ್ಗಿದ ಎಮ್ಮೆಯನ್ನು ಕೋಲಿನಿಂದ ಥಳಿಸಿದನು.”\par

\indent —ಈ ವಾಕ್ಯದಲ್ಲಿ ರವಿಯು, ತೋಟಕ್ಕೆ, ಎಮ್ಮೆಯನ್ನು, ಕೋಲಿನಿಂದ ಎಂಬ ಪದಗಳು ನಾಮಪದಗಳಾಗಿರುವುವು... ಈ ನಾಮಪದಗಳಲ್ಲಿ, ರವಿ, ತೋಟ ಎಮ್ಮೆ , ಕೋಲು. ಇವು ನಾಮಪ್ರಕೃತಿಗಳು ' ಮನುಷ್ಯ, ಸ್ಥಳ ಪ್ರಾಣಿ, ಪದಾರ್ಥ ಮುಂತಾದವುಗಳ ಹೆಸರುಗಳನ್ನು ತಿಳಿಸುವ ಈ ನಾಮಪ್ರಕೃತಿಗಳಿಗೆ “ನಾಮವಾಚಕ ಪ್ರಕೃತಿ” ಗಳೆಂದೂ ಹೆಸರು... ನಾಮವಾಚಕ ಪ್ರಕೃತಿಗಳಲ್ಲಿ ಹಲವು ಪ್ರಕಾರಗಳುಂಟು ; ಅವುಗಳಲ್ಲಿ ಮುಖ್ಯವಾದವುಗಳೆಂದರೆ—\par

\vspace{4pt plus 3pt minus 1pt}
\begin{center}
{\renewcommand{\arraystretch}{1.22}\setlength{\tabcolsep}{5pt}
\begin{tabular}{r l c l}
(೧) & ವಸ್ತುವಾಚಕಗಳು & ಉದಾ : & ಮರ, ಗೌರಿ \\
(೨) & ಗುಣವಾಚಕಗಳು & ಉದಾ : & `: : ದೊಡ್ಡ, ಹಳೆಯ \\
(೩) & ಸಂಖ್ಯಾವಾಚಕಗಳು & ಉದಾ : & ಒಂದು ಹತ್ತು \\
\end{tabular}}
\end{center}
\vspace{4pt plus 3pt minus 1pt}

\vspace{2.5pt plus 3pt minus 0.5pt}\noindent\hspace*{1.2em}\kbvnum{(೪)}\hspace{0.55em}. ಸಂಖ್ಯೇಯವಾಚಕಗಳು. ಉದಾ : ಐವರು, ಹತ್ತನೆಯ\par

\vspace{4pt plus 3pt minus 1pt}
\begin{center}
{\renewcommand{\arraystretch}{1.22}\setlength{\tabcolsep}{5pt}
\begin{tabular}{r l c l}
(೫) & ಪರಿಮಾಣವಾಚಕಗಳು & ಉದಾ : & ಅಷ್ಟು, ಇಷ್ಟು \\
(೬) & ಪ್ರಕಾರವಾಚಕಗಳು & ಉದಾ : & ಅಂತಹ, ಇಂತಹ \\
(೭) & ಸರ್ವನಾಮಗಳು & ಉದಾ : & ಅವನು ಇವನು \\
\end{tabular}}
\end{center}
\vspace{4pt plus 3pt minus 1pt}
}
\clearpage

% === tex/0036.tex (Printed Page 34) ===
\renewcommand{\kbvchaptername}{ನಾಮಪದದ ವಿಭೇದಗಳು}
\setcounter{page}{34}
\thispagestyle{kbvpage}
\kbvtypesetpage{\vspace{2.5pt plus 3pt minus 0.5pt}\noindent\hspace*{1.2em}\kbvnum{(೧)}\hspace{0.55em}ವಸ್ತುವಾಚಕಗಳು\par

\indent ಪಾರ್ವತಿ, ನಂದಿನಿ, ಶೇಷಾಚಲ, ವೆಂಕಟೇಶ, ಸೂರ್ಯ, ಚಂದ್ರ, ರಾಮನಗರ, ದಾವಣಗೆರೆ, ಹೇಮಾವತಿ, ಅಣ್ಣ, ತಂಗಿ, ನದಿ, ಬೆಟ್ಟ, ಹುಲಿ, ಇಲಿ, ಬಿಂದಿಗೆ, ಪುಸ್ತಕ, ಬಳಪ, ಕುರುಡ, ಕುಂಟ, ಮೂಗ, ಹೆಳವ, ಕಿವುಡ, ಅರ್ಚಕ, ನೀರುಗಂಟ, ಪುರೋಹಿತ ಇವೇ ಮೊದಲಾದ ಶಬ್ದಗಳು ವಸ್ತುವಾಚಕಗಳು.\par

\begin{kbvrule}
\noindent{\kbvsecfont ವ್ಯಾಖ್ಯೆ :}\hspace{0.5em}ಮನುಷ್ಯ, ಸ್ಥಳ, ಪ್ರಾಣಿ, ವಸ್ತು ಇವುಗಳ ಹೆಸರನ್ನು ತಿಳಿಸುವ ನಾಮಪದಗಳಿಗೆ ವಸ್ತುವಾಚಕಗಳೆಂದು ಹೆಸರು. ವಸ್ತುವಾಚಕಗಳನ್ನು ಮೂರು ಪ್ರಕಾರವಾಗಿ ವಿಂಗಡಿಸಬಹುದು.
\end{kbvrule}

\indent ({\kbvlatin i}) ರೂಢನಾಮಗಳು : ಅಣ್ಣ, ತಂಗಿ, ನದಿ, ಬೆಟ್ಟ, ಹುಲಿ, ಇಲಿ, ಬಿಂದಿಗೆ, ಪುಸ್ತಕ, ಬಳಪ — ಇವು ರೂಢಿಯಲ್ಲಿರುವ ಸಾಮಾನ್ಯವಸ್ತುವಾಚಕ ಶಬ್ದಗಳು. ಇಲ್ಲಿ, ನದಿ ಎಂದರೆ ಯಾವುದೇ ನದಿಯಾಗಬಹುದು; ಪುಸ್ತಕ ಎಂದರೆ ಯಾವುದೇ ಪುಸ್ತಕವಾಗಬಹುದು. ಇಂತಹ ವಸ್ತುವಾಚಕಗಳಿಗೆ ರೂಢನಾಮಗಳೆಂದು ಹೆಸರು.\par

\indent ({\kbvlatin ii}) ಅಂಕಿತನಾಮಗಳು : ಪಾರ್ವತಿ, ನಂದಿನಿ, ಶೇಷಾಚಲ, ವೆಂಕಟೇಶ, ಸೂರ್ಯ, ಚಂದ್ರ, ರಾಮನಗರ, ದಾವಣಗೆರೆ, ಹೇಮಾವತಿ — ಇವು ನಿರ್ದಿಷ್ಟ ವ್ಯಕ್ತಿಯನ್ನೋ ಸ್ಥಳವನ್ನೋ ಬೆಟ್ಟವನ್ನೋ ನದಿಯನ್ನೋ ಗುರ್ತಿಸುವ ಸಲುವಾಗಿ ಜನರು ತಮ್ಮ ಇಷ್ಟಾನುಸಾರ ಇಟ್ಟುಕೊಂಡಿರುವ ಹೆಸರುಗಳು. ಇಂತಹ ವಸ್ತುವಾಚಕಗಳಿಗೆ ಅಂಕಿತನಾಮಗಳೆಂದು ಹೆಸರು.\par

\indent ({\kbvlatin iii}) ಅನ್ವರ್ಥನಾಮಗಳು : ಕುರುಡ, ಕುಂಟ, ಮೂಗ, ಹೆಳವ, ಕಿವುಡ, ಅರ್ಚಕ, ನೀರುಗಂಟ, ಪುರೋಹಿತ — ಇವು ವೃತ್ತಿಗೆ ಅನುಗುಣವಾಗಿ, ಅರ್ಥಕ್ಕೆ ಅನುಗುಣವಾಗಿ ವಾಡಿಕೆಯಲ್ಲಿ ಬಂದಿರುವ ಅನ್ವರ್ಥಕನಾಮಗಳು. ಇಂತಹ ವಸ್ತುವಾಚಕಗಳಿಗೆ ಅನ್ವರ್ಥನಾಮಗಳೆಂದು ಹೆಸರು.\par

\begin{kbvnote}
\noindent{\kbvsecfont ವಿವರಣೆ :}\hspace{0.5em}ಕುರುಡ-ಕಣ್ಣು ಕಾಣದವನು.
\end{kbvnote}

\noindent\hspace*{2.2em}ನೀರುಗಂಟ-ನೀರು ಹಾಯಿಸುವವನು-ಇತ್ಯಾದಿ.\par
}
\clearpage

% === tex/0037.tex (Printed Page 35) ===
\renewcommand{\kbvchaptername}{ನಾಮಪದದ ವಿಭೇದಗಳು}
\setcounter{page}{35}
\thispagestyle{kbvpage}
\kbvtypesetpage{\kbvsechead{(೨)\hspace{0.45em}ಗುಣವಾಚಕಗಳು :}

\vspace{2.5pt plus 3pt minus 0.5pt}\noindent\hspace*{1.2em}\kbvnum{(1)}\hspace{0.55em}ಮೈಸೂರಿನ ತಾತಯ್ಯ ಒಳ್ಳೆಯ ವ್ಯಕ್ತಿಯಾಗಿದ್ದರು.\par

\vspace{2.5pt plus 3pt minus 0.5pt}\noindent\hspace*{1.2em}\kbvnum{(2)}\hspace{0.55em}ಇದು ಹಳೆಯ ಮನೆ.\par

\vspace{2.5pt plus 3pt minus 0.5pt}\noindent\hspace*{1.2em}\kbvnum{(3)}\hspace{0.55em}ಕಪಿಗೆ ಉದ್ದವಾದ ಬಾಲವುಂಟು.\par

\vspace{2.5pt plus 3pt minus 0.5pt}\noindent\hspace*{1.2em}\kbvnum{(4)}\hspace{0.55em}ನನ್ನ ಮೇಜು ಚಿಕ್ಕದು.\par

\vspace{2.5pt plus 3pt minus 0.5pt}\noindent\hspace*{1.2em}\kbvnum{(5)}\hspace{0.55em}ದುಷ್ಟರಾದ ಜನರಿಂದ ದೂರವಿರಬೇಕು.\par

\indent ಮೇಲಿನ ವಾಕ್ಯಗಳಲ್ಲಿ ದಪ್ಪಕ್ಷರಗಳಲ್ಲಿರುವ ಪದಗಳನ್ನು ಗಮನಿಸಿ, ಒಳ್ಳೆಯ, ಹಳೆಯ, ಉದ್ದವಾದ, ಚಿಕ್ಕದು, ದುಷ್ಟರಾದ—ಇವು ಮನುಷ್ಯ, ಸ್ಥಳ, ಪ್ರಾಣಿ ಹಾಗೂ ವಸ್ತುಗಳ ಗುಣ, ರೀತಿ, ಸ್ವಭಾವಗಳನ್ನು ತಿಳಿಸುವ ; ಪದಗಳಾಗಿವೆ ; ಇವು ವಿಶೇಷಣಗಳು.\par

\indent “ಮೈಸೂರಿನ ತಾತಯ್ಯ ಒಳ್ಳೆಯ ವ್ಯಕ್ತಿಯಾಗಿದ್ದರು” ಎಂಬ ವಾಕ್ಯವನ್ನು ಆಧಾರವಾಗಿಟ್ಟು ಕೊಂಡು, “ಅವರು ಎಂತಹ ವ್ಯಕ್ತಿಯಾಗಿದ್ದರು?”, ಎಂಬುದಾಗಿ ಪ್ರಶ್ನೆ ಹಾಕಿದರೆ ಉತ್ತರ ರೂಪವಾಗಿ ಬರುವ ಪದ “ಒಳ್ಳೆಯ”, ಈ ಪದವು ತಾತಯ್ಯನವರ ಗುಣ ಅಥವಾ ಸ್ವಭಾವವನ್ನು ತಿಳಿಸುವುದು. ಇದೇ ಕ್ರಮವನ್ನು ಅನುಸರಿಸಿ, ಪ್ರತಿಯೊಂದು ವಾಕ್ಯದಲ್ಲೂ ನಾಮಪದದ ಮೇಲೆ “ಎಂತಹ” ಎಂಬ ಪ್ರಶ್ನೆ ಹಾಕಿ ಉತ್ತರ ಪಡೆಯಿರಿ ; ನಿಮಗೆ ಲಭಿಸುವ ಉತ್ತರವು ನಾಮಪದದ ವಿಶೇಷಣವೇ ಆಗಿರುವುದು.\par

\begin{kbvrule}
\noindent{\kbvsecfont ವ್ಯಾಖ್ಯೆ :}\hspace{0.5em}ಮನುಷ್ಯ, ಸ್ಥಳ, ಪ್ರಾಣಿ ಹಾಗೂ ವಸ್ತುಗಳ ಗುಣ, ರೀತಿ, ಸ್ವಭಾವಗಳನ್ನು ತಿಳಿಸುವ ವಿಶೇಷಣಗಳಿಗೆ ಗುಣವಾಚಕಗಳೆಂದು ಹೆಸರು.
\end{kbvrule}

\kbvsechead{(೩)\hspace{0.45em}ಸಂಖ್ಯಾವಾಚಕಗಳು :}

\indent ಒಂದು, ಎರಡು, ಮೂರು, ಹತ್ತು, ನೂರು ಇವೇ ಮೊದಲಾದ ನಾಮಪದಗಳು ಸಂಖ್ಯೆಗಳನ್ನು ಹೆಸರಿಸುತ್ತವೆ ; ಈ ನಾಮಪದಗಳಿಗೆ ಸಂಖ್ಯಾ ವಾಚಕಗಳೆಂದು ಹೆಸರು.\par

\vspace{2.5pt plus 3pt minus 0.5pt}\noindent\hspace*{1.2em}\kbvnum{(1)}\hspace{0.55em}ನನ್ನ ತೂಕ ಐವತ್ತು ಕಿಲೋಗ್ರಾಂ...\par

\vspace{2.5pt plus 3pt minus 0.5pt}\noindent\hspace*{1.2em}\kbvnum{(2)}\hspace{0.55em}ಈ ಪುಸ್ತಕದ ಬೆಲೆ ಎಂಟು ರೂಪಾಯಿಗಳು.\par
}
\clearpage

% === tex/0038.tex (Printed Page 36) ===
\renewcommand{\kbvchaptername}{ನಾಮಪದದ ವಿಭೇದಗಳು}
\setcounter{page}{36}
\thispagestyle{kbvpage}
\kbvtypesetpage{\vspace{2.5pt plus 3pt minus 0.5pt}\noindent\hspace*{1.2em}\kbvnum{(3)}\hspace{0.55em}ನಾನು ಪ್ರತಿ ದಿನ ರಾತ್ರಿ ಹತ್ತು ಗಂಟೆಗೆ ಮಲಗುತ್ತೇನೆ. ಈ ವಾಕ್ಯಗಳಲ್ಲಿ ದಪ್ಪಕ್ಷರಗಳಲ್ಲಿರುವ ಪದಗಳು ಸಂಖ್ಯಾವಾಚಕಗಳು.\par

\vspace{4pt plus 4pt minus 1pt}\noindent{\kbvsecfont\fontsize{11pt}{14.5pt}\selectfont\color{kbvaccentdark} (೪)\hspace{0.45em}ಸಂಖ್ಯೇಯ ವಾಚಕಗಳು :}\hspace{0.5em}"ಸಂಖ್ಯಾತ್ಮಕ ನಾಮಪದಗಳಿಗೆ ಸಂಖ್ಯೇಯ ವಾಚಕಗಳೆಂದು ಹೆಸರು. ಉದಾಹರಣೆಗೆ, “ದ್ರೌಪದಿಯು ಐವರು ಮಡದಿಯಾಗಿದ್ದಳು” ; ಈ ವಾಕ್ಯ ದಲ್ಲಿ ಐದು ಜನರು ಎಂಬ ಅರ್ಥದಲ್ಲಿ “ಐವರು” ಎಂಬ ಪದ ಬಳಸಲಾಗಿದೆ. ಇದು ಸಂಖ್ಯೇಯವಾಚಕ.\par

\vspace{2.5pt plus 3pt minus 0.5pt}\noindent\hspace*{1.2em}\kbvnum{(1)}\hspace{0.55em}ಓಟದ ಪಂದ್ಯದಲ್ಲಿ ಗೆದ್ದವರಲ್ಲಿ ನಾನು ಎರಡನೆಯವನು.\par

\vspace{2.5pt plus 3pt minus 0.5pt}\noindent\hspace*{1.2em}\kbvnum{(2)}\hspace{0.55em}ಕನ್ನಡದ ಶ್ರೇಷ್ಠ ಸಾಹಿತಿಗಳಲ್ಲಿ ಕುವೆಂಪು ಒಬ್ಬರು. ಈ ವಾಕ್ಯಗಳಲ್ಲಿ ದಪ್ಪಕ್ಷರಗಳಲ್ಲಿರುವ ಪದಗಳು ಸಂಖ್ಯೇಯವಾಚಕಗಳು.\par

\vspace{4pt plus 4pt minus 1pt}\noindent{\kbvsecfont\fontsize{11pt}{14.5pt}\selectfont\color{kbvaccentdark} (೫)\hspace{0.45em}ಪರಿಮಾಣವಾಚಕಗಳು :}\hspace{0.5em}ಅಂದಾಜಿನ ಮೇರೆಗೆ ವಸ್ತುಗಳ ಪರಿಮಾಣ ಸೂಚಿಸುವ ನಾಮಪದ ಗಳಿಗೆ ಪರಿಮಾಣ ವಾಚಕಗಳೆಂದು ಹೆಸರು.\par

\noindent\hspace*{1.2em}\kbvlabel{ಉದಾ}\hspace{0.6em}ಇಷ್ಟು, ಅಷ್ಟು, ರವೆಯಷ್ಟು, ಬೆಟ್ಟದಷ್ಟು, ಕೆಲವು, ಹಲವು—ಇವು ಪರಿಮಾಣವಾಚಕಗಳು.\par

\vspace{4pt plus 4pt minus 1pt}\noindent{\kbvsecfont\fontsize{11pt}{14.5pt}\selectfont\color{kbvaccentdark} (೬)\hspace{0.45em}ಪ್ರಕಾರವಾಚಕಗಳು :}\hspace{0.5em}ಹೋಲಿಕೆಯ ಅರ್ಥ `ಬರುವಂತೆ' ಮನುಷ್ಯ, ಪ್ರಾಣಿ, ಸ್ಥಳ, ವಸ್ತು ಇವುಗಳ ರೀತಿ ಸ್ವಭಾವ, ಸ್ಥಿತಿ ಯಾವ ಪ್ರಕಾರವಾಗಿರುವುದೆಂಬುದನ್ನು ತಿಳಿಸುವ ನಾಮಪದಗಳಿಗೆ ಪ್ರಕಾರವಾಚಕಗಳೆಂದು ಹೆಸರು. 'ಅಂಥ, ಇಂಥ, ಅಂತಹ, ಇಂತಹ ಮೊದಲಾದ ಪದಗಳನ್ನು ಪ್ರಕಾರವಾಚಕಗಳಾಗಿ ಬಳಸುವುದುಂಟು.\par

\noindent\hspace*{1.2em}\kbvlabel{ಉದಾ}\hspace{0.6em}ಅಂಥ ವ್ಯಕ್ತಿ ಮತ್ತೊಬ್ಬನಿಲ್ಲ.\par

\noindent\hspace*{2.2em}ಇಂಥ ಜನಕ್ಕೆ ಈ ಪಾಡು ಬರಬಾರದಿತ್ತು.\par

\noindent\hspace*{2.2em}ಅಂತಹವನು ಹುಟ್ಟುವುದು ಅಪರೂಪ.\par

\noindent\hspace*{2.2em}ನೀನು ಇಂತಹ ಕೆಲಸ ಮಾಡಿದ್ದು ತಪ್ಪು.\par
}
\clearpage

% === tex/0039.tex (Printed Page 37) ===
\renewcommand{\kbvchaptername}{ನಾಮಪದದ ವಿಭೇದಗಳು}
\setcounter{page}{37}
\thispagestyle{kbvpage}
\kbvtypesetpage{\kbvsechead{(2)\hspace{0.45em}ಸರ್ವನಾಮಗಳು :}

\indent “ದೇವರು ದೊಡ್ಡವನು ; ದೇವರು ಸರ್ವಶಕ್ತ” — ಈ ವಾಕ್ಯಗಳನ್ನು ಓದಿ ನೋಡಿ. ಎರಡನೆಯ ವಾಕ್ಯದಲ್ಲಿ “ದೇವರು” ಎಂಬ ನಾಮಪದ ಪುನರಾವೃತ್ತವಾಗಿದೆ. ಈ ಪದಕ್ಕೆ ಬದಲಾಗಿ ಅವನು ಎಂಬ ಪದ ಬಳಸಿ ದಲ್ಲಿ ಓದಲು ಹಿತವೆನಿಸುವುದು. “ದೇವರು ದೊಡ್ಡವನು; ಅವನು ಸರ್ವಶಕ್ತ.”\par

\noindent\hspace*{2.2em}ಇದೇ ರೀತಿಯ ಹಲವು ಉದಾಹರಣೆಗಳನ್ನು ಪರಿಶೀಲಿಸೋಣ.\par

\vspace{2.5pt plus 3pt minus 0.5pt}\noindent\hspace*{1.2em}\kbvnum{(1)}\hspace{0.55em}ಪಾರ್ವತಿಯು ಶಿವನ ಹೆಂಡತಿ; ಅವಳು ಪರ್ವತೇಶ್ವರನ ಮಗಳು.\par

\vspace{2.5pt plus 3pt minus 0.5pt}\noindent\hspace*{1.2em}\kbvnum{(2)}\hspace{0.55em}ಬೆಂಗಳೂರು ದೊಡ್ಡ ಪಟ್ಟಣ; ಅದು ಕರ್ನಾಟಕದ ರಾಜಧಾನಿ.\par

\vspace{2.5pt plus 3pt minus 0.5pt}\noindent\hspace*{1.2em}\kbvnum{(3)}\hspace{0.55em}ನಾನು ಇಲ್ಲಿಯೇ ಕಾದಿರುತ್ತೇನೆ ; ನೀನು ಆದಷ್ಟು ಬೇಗ ಬಾ.\par

\vspace{2.5pt plus 3pt minus 0.5pt}\noindent\hspace*{1.2em}\kbvnum{(4)}\hspace{0.55em}ಇವರು ಯಾರು ?...{[}ಈ ಪ್ರಶ್ನೆಗೆ ಉತ್ತರವಾಗಿ “ಇವರು ಗಾಂಧೀಜಿ” ಎಂದು ಹೇಳುವುದಾದರೆ, ಗಾಂಧೀಜಿ. ಎಂಬ ನಾಮಪದಕ್ಕೆ ಬದ ಲಾಗಿ ಯಾರು ಎಂಬ ಪ್ರಶ್ನಾರ್ಥಕ ಪದ ಬಳಸಿರುವುದು ವ್ಯಕ್ತವಾಗುವುದು.{]}\par

\indent ಮೇಲಿನ ಉದಾಹರಣೆಗಳಲ್ಲಿ ನಾಮಪದಗಳಿಗೆ ಬದಲಾಗಿ ಅವನು ಅವಳು, ಅದು, ನಾನು, ನೀನು, ಯಾರು ಎಂಬ ಪದಗಳನ್ನು ಬಳಸಿರುವುದು ಕಂಡುಬರುತ್ತದೆ. ಹೀಗೆ ನಾಮಪದಗಳಿಗೆ ಪ್ರತಿಯಾಗಿ ಅವುಗಳ ಸ್ಥಾನದಲ್ಲಿ ಬಳಸುವ ಪದಗಳಿಗೆ “ಸರ್ವನಾಮ” ಗಳೆಂದು ಹೆಸರು.\par

\kbvsechead{ಸರ್ವನಾಮಗಳಲ್ಲಿ ಮೂರು ವಿಧ :}

\vspace{2.5pt plus 3pt minus 0.5pt}\noindent\hspace*{1.2em}\kbvnum{(1)}\hspace{0.55em}ಪುರುಷಾರ್ಥಕ ಸರ್ವನಾಮಗಳು\par

\vspace{2.5pt plus 3pt minus 0.5pt}\noindent\hspace*{1.2em}\kbvnum{(2)}\hspace{0.55em}ಆತ್ಮಾರ್ಥಕ ಸರ್ವನಾಮಗಳು\par

\vspace{2.5pt plus 3pt minus 0.5pt}\noindent\hspace*{1.2em}\kbvnum{(3)}\hspace{0.55em}ಪ್ರಶ್ನಾರ್ಥಕ ಸರ್ವನಾಮಗಳು\par

\kbvsechead{(1)\hspace{0.45em}ಪುರುಷಾರ್ಥಕ ಸರ್ವನಾಮಗಳು :}

\indent ಉದಾಹರಣೆಗೆ ನಾನು, ನೀನು, ಅವನು, ‘ಇವನು’, ‘ಅದು’, ‘ಇದು’ ಇತ್ಯಾದಿ ಪುರುಷಾರ್ಥಕ ಸರ್ವನಾಮಗಳು. ಈ ಸರ್ವನಾಮಗಳನ್ನು ಮೂರು ಗುಂಪುಗಳಾಗಿ ವಿಂಗಡಿಸಬಹುದು.\par
}
\clearpage

% === tex/0040.tex (Printed Page 38) ===
\renewcommand{\kbvchaptername}{ನಾಮಪದದ ವಿಭೇದಗಳು}
\setcounter{page}{38}
\thispagestyle{kbvpage}
\kbvtypesetpage{\indent ({\kbvlatin i}) ನಾನು, ನಾವು—ಇವು “ಉತ್ತಮ ಪುರುಷ” ಸರ್ವನಾಮಗಳು. ({\kbvlatin ii}) ನೀನು, ನೀವು—ಇವು “ಮಧ್ಯಮ ಪುರುಷ” ಸರ್ವನಾಮಗಳು. ({\kbvlatin iii}) ಅವನು, ಅವಳು, ಅವರು, ಇವನು, ಇವಳು, ಇವರು, ಅದು, ಅವು, ಇದು, ಇವು—ಇವು “ಪ್ರಥಮ ಪುರುಷ” ಸರ್ವನಾಮಗಳು.\par

\vspace{4pt plus 4pt minus 1pt}\noindent{\kbvsecfont\fontsize{11pt}{14.5pt}\selectfont\color{kbvaccentdark} (1)\hspace{0.45em}ಆತ್ಮಾರ್ಥಕ ಸರ್ವನಾಮಗಳು :}\hspace{0.5em}ಉದಾಹರಣೆಗೆ ತಾನು, ತಾವು ಇತ್ಯಾದಿ ಆತ್ಮಾರ್ಥಕ ಸರ್ವನಾಮಗಳು.\par

\vspace{4pt plus 4pt minus 1pt}\noindent{\kbvsecfont\fontsize{11pt}{14.5pt}\selectfont\color{kbvaccentdark} (2)\hspace{0.45em}ಪ್ರಶ್ನಾರ್ಥಕ ಸರ್ವನಾಮಗಳು :}\hspace{0.5em}ಉದಾಹರಣೆಗೆ ಯಾರು, ಎಲ್ಲಿ, ಯಾವುದು, ಏನು ಇತ್ಯಾದಿ ಪ್ರಶ್ನಾರ್ಥಕ ಸರ್ವನಾಮಗಳು.\par

\kbvexercise{ಅಭ್ಯಾಸ ಪತ್ರಿಕೆ ೮}

\vspace{2.5pt plus 3pt minus 0.5pt}\noindent\hspace*{1.2em}\kbvnum{(1)}\hspace{0.55em}ವಸ್ತುವಾಚಕಗಳು ಎಂದರೇನು ? ಉದಾಹರಣೆಗಳ ಸಹಿತ ತಿಳಿಸಿ.\par

\vspace{2.5pt plus 3pt minus 0.5pt}\noindent\hspace*{1.2em}\kbvnum{(2)}\hspace{0.55em}ವಸ್ತುವಾಚಕದ ಮೂರು ಪ್ರಕಾರಗಳನ್ನು ತಿಳಿಸಿ.\par

\vspace{2.5pt plus 3pt minus 0.5pt}\noindent\hspace*{1.2em}\kbvnum{(3)}\hspace{0.55em}ರೂಢನಾಮಗಳಿಗೂ ಅಂಕಿತನಾಮಗಳಿಗೂ ಇರುವ ವ್ಯತ್ಯಾಸವೇನು ? ಉದಾಹರಣೆಗಳ ಸಹಿತ ತಿಳಿಸಿ.\par

\vspace{2.5pt plus 3pt minus 0.5pt}\noindent\hspace*{1.2em}\kbvnum{(4)}\hspace{0.55em}ಗುಣವಾಚಕ ಎಂದರೇನು ? ಉದಾಹರಣೆಗಳ ಸಹಿತ ವಿವರಿಸಿ.\par

\vspace{2.5pt plus 3pt minus 0.5pt}\noindent\hspace*{1.2em}\kbvnum{(5)}\hspace{0.55em}ಸುಯೋಧನ, ದಶರಥ, ಸುಬಾಹು, ದ್ರೌಪದಿ—ಈ ಪುರಾಣ ಪ್ರಸಿದ್ಧ ಹೆಸರುಗಳು ಯಾವ ಪ್ರಕಾರದ ನಾಮವಾಚಕಗಳು?\par

\vspace{2.5pt plus 3pt minus 0.5pt}\noindent\hspace*{1.2em}\kbvnum{(6)}\hspace{0.55em}ಸಂಖ್ಯಾವಾಚಕಗಳಿಗೂ ಸಂಖ್ಯೇಯವಾಚಕಗಳಿಗೂ ಇರುವ ವ್ಯತ್ಯಾಸವನ್ನು ಉದಾಹರಣೆಗಳ ಮೂಲಕ ಸ್ಪಷ್ಟಪಡಿಸಿ.\par

\vspace{2.5pt plus 3pt minus 0.5pt}\noindent\hspace*{1.2em}\kbvnum{(7)}\hspace{0.55em}ಈ ಪದಗಳನ್ನು ಬಳಸಿಕೊಂಡು ಸ್ವಂತ ವಾಕ್ಯಗಳನ್ನು ರಚಿಸಿ : ಅಷ್ಟೊಂದು, ಇಂತಿಷ್ಟು, ಇಂಥಾದ್ದು, ಎಂತಹ\par

\vspace{2.5pt plus 3pt minus 0.5pt}\noindent\hspace*{1.2em}\kbvnum{(8)}\hspace{0.55em}ಇವುಗಳಿಗೆ ಎರಡೆರಡು ಉದಾಹರಣೆಗಳನ್ನು ಕೊಡಿ : ರೂಢನಾಮಗಳು, ಪರಿಮಾಣವಾಚಕಗಳು ಅಂಕಿತನಾಮಗಳು, ಪ್ರಕಾರವಾಚಕಗಳು ಅನ್ವರ್ಥನಾಮಗಳು, ಪ್ರಥಮ ಪುರುಷ ಸರ್ವನಾಮಗಳು\par
}
\clearpage

% === tex/0041.tex (Printed Page 39) ===
\renewcommand{\kbvchaptername}{ನಾಮಪದಗಳ ರೂಪ ಬದಲಾವಣೆ ({\kbvlatin i})}
\setcounter{page}{39}
\thispagestyle{kbvpage}
\kbvtypesetpage{\vspace{2.5pt plus 3pt minus 0.5pt}\noindent\hspace*{1.2em}\kbvnum{(9)}\hspace{0.55em}ಸರ್ವನಾಮಗಳು ಎಂದರೇನು ? ಉದಾಹರಿಸಿ.\par

\vspace{2.5pt plus 3pt minus 0.5pt}\noindent\hspace*{1.2em}\kbvnum{(10)}\hspace{0.55em}ರಾಮನು ದಶರಥನ ಮಗ. ರಾಮನು ತಂದೆಯ ಮಾತಿನಂತೆ ಕಾಡಿಗೆ ಹೋದನು. ರಾಮನೊಂದಿಗೆ ಸೀತೆಯೂ ಕಾಡಿಗೆ ಹೋದಳು. ಸೀತಾರಾಮರ ಜೊತೆಯಲ್ಲಿ ಲಕ್ಷ್ಮಣನೂ ಹೊರಟನು.\par

\indent — ಈ ವಾಕ್ಯಗಳಲ್ಲಿ ಅಗತ್ಯವಿದ್ದೆಡೆ ಸರ್ವನಾಮಗಳನ್ನು ಸೇರಿಸಿ ಬರೆಯಿರಿ.\par

\vspace{2.5pt plus 3pt minus 0.5pt}\noindent\hspace*{1.2em}\kbvnum{(11)}\hspace{0.55em}ಉತ್ತಮ ಪುರುಷ ಸರ್ವನಾಮಗಳಿಗೂ ಮಧ್ಯಮ ಪುರುಷ ಸರ್ವನಾಮಗಳಿಗೂ ಇರುವ ವ್ಯತ್ಯಾಸವೇನು ? ಉದಾಹರಿಸಿ.\par

\vspace{2.5pt plus 3pt minus 0.5pt}\noindent\hspace*{1.2em}\kbvnum{(12)}\hspace{0.55em}ನಿನ್ನ ವಾಸಸ್ಥಳ ಯಾವುದು? ನೀನು ಎಲ್ಲಿಂದ ಬಂದೆ? ನಿನ್ನನ್ನು ಕಳುಹಿಸಿದವರು ಯಾರು? ನೀನು ಬಂದ ಉದ್ದೇಶವೇನು?\par

\indent — ಇಲ್ಲಿ ಕಾಣಿಸಿರುವ ಪ್ರಶ್ನಾರ್ಥಕ ಸರ್ವನಾಮಗಳನ್ನು ಗುರ್ತಿಸಿ.\par

\kbvchapter{ಪಾಠ ೯}{ನಾಮಪದಗಳ ರೂಪ ಬದಲಾವಣೆ ({\kbvlatin i})}

\kbvsubtitle{ಲಿಂಗ}

\noindent\hspace*{2.2em}ಭೀಮಣ್ಣ ಓಡುತ್ತಿದ್ದಾನೆ.\par

\noindent\hspace*{2.2em}ಆ ಕಮಲೆ ಹಾಡುತ್ತಿದ್ದಾಳೆ.\par

\noindent\hspace*{2.2em}ನವಿಲು ಕುಣಿಯುತ್ತಿದೆ.\par

\noindent\hspace*{2.2em}ಗಡಿಯಾರ ನಿಂತುಹೋಗಿದೆ.\par

\indent ಈ ವಾಕ್ಯಗಳಲ್ಲಿ ದಪ್ಪ ಅಕ್ಷರಗಳಲ್ಲಿರುವ ಪದಗಳನ್ನು ಗಮನಿಸಿ. “ಭೀಮಣ್ಣ” ಎಂಬ ಶಬ್ದ ಕಿವಿಗೆ ಬಿದ್ದೊಡನೆ ಇವನು ಪುರುಷ {[}ಗಂಡಸು{]} ಎಂದು ಹೇಳಬಹುದು. ಈ ಕಾರಣದಿಂದ “ಭೀಮಣ್ಣ” ಎನ್ನುವುದು ಪುರುಷಲಿಂಗ ಶಬ್ದ ಎನ್ನುತ್ತೇವೆ. “ಕಮಲೆ” - ಎಂಬ ಶಬ್ದ ಕಿವಿಗೆ ಬಿದ್ದೊಡನೆ ಇವಳು ಸ್ತ್ರೀ {[}ಹೆಂಗಸು{]} ಎಂದು ಹೇಳಬಹುದು. ಈ ಕಾರಣದಿಂದ “ಕಮಲೆ” ಎನ್ನುವುದು ಸ್ತ್ರೀಲಿಂಗ ಶಬ್ದ ಎನ್ನುತ್ತೇವೆ. ನವಿಲು, ಗಡಿಯಾರ - ಈ ಶಬ್ದಗಳು ಮನಸ್ಸಿನ ಮೇಲೆ ಗಂಡು, ಎಂಬ ಅಭಿಪ್ರಾಯವನ್ನಾಗಲೀ ಹೆಣ್ಣು ಎಂಬ ಅಭಿಪ್ರಾಯವನ್ನಾಗಲೀ ಉಂಟುಮಾಡುವುದಿಲ್ಲ. ಆದುದರಿಂದ “ನವಿಲು”, “ಗಡಿಯಾರ” ಎಂಬ ಶಬ್ದಗಳನ್ನು ನಪುಂಸಕ ಲಿಂಗ ಶಬ್ದಗಳೆನ್ನುತ್ತೇವೆ.\par
}
\clearpage

% === tex/0042.tex (Printed Page 40) ===
\renewcommand{\kbvchaptername}{ನಾಮಪದಗಳ ರೂಪ ಬದಲಾವಣೆ ({\kbvlatin i})}
\setcounter{page}{40}
\thispagestyle{kbvpage}
\kbvtypesetpage{\indent ಹೀಗೆ ಮನುಷ್ಯ, ಪ್ರಾಣಿ, ಸ್ಥಳ, ವಸ್ತು—ಇವುಗಳನ್ನು ಹೆಸರಿಸುವ ನಾಮಪದಗಳನ್ನು ಪುಲ್ಲಿಂಗ ಶಬ್ದಗಳು, ಸ್ತ್ರೀಲಿಂಗ ಶಬ್ದಗಳು ಮತ್ತು ನಪುಂಸಕ ಲಿಂಗ ಶಬ್ದಗಳು ಎಂಬುದಾಗಿ ಮೂರು ಬಗೆಯಾಗಿ ವಿಂಗಡಿಸಿರುವುದನ್ನು ಕನ್ನಡ ವ್ಯಾಕರಣ ಶಾಸ್ತ್ರದಲ್ಲಿ ಕಾಣಬಹುದು.\par

\vspace{2.5pt plus 3pt minus 0.5pt}\noindent\hspace*{1.2em}\kbvnum{(೧)}\hspace{0.55em}ಯಾವ ಶಬ್ದವು ಗಂಡಸು ಎಂಬ ಅರ್ಥವನ್ನು ಉಂಟುಮಾಡುವುದು. ಅದು ಪುಲ್ಲಿಂಗ ಶಬ್ದವಾಗಿರುವುದು.\par

\noindent\hspace*{1.2em}\kbvlabel{ಉದಾ}\hspace{0.6em}ಕೃಷ್ಣ, ಅಜ್ಜ, ಗಂಡ, ಅಣ್ಣ, ಹುಡುಗ, ಗೌಡ ಒಕ್ಕಲಿಗ, ಅವನು—ಇತ್ಯಾದಿ.\par

\vspace{2.5pt plus 3pt minus 0.5pt}\noindent\hspace*{1.2em}\kbvnum{(೨)}\hspace{0.55em}ಯಾವ ಶಬ್ದವು ಹೆಣ್ಣು ಎಂಬ ಅರ್ಥವನ್ನು ಉಂಟುಮಾಡುವುದು. ಅದು ಸ್ತ್ರೀಲಿಂಗ ಶಬ್ದವಾಗಿರುವುದು.\par

\noindent\hspace*{1.2em}\kbvlabel{ಉದಾ}\hspace{0.6em}ರುಕ್ಮಿಣಿ, ಅಜ್ಜಿ, ಹೆಂಡತಿ, ಅಕ್ಕ, ಹುಡುಗಿ, ಗೌಡಿತಿ ಒಕ್ಕಲಗಿತ್ತಿ, ಅವಳು—ಇತ್ಯಾದಿ.\par

\vspace{2.5pt plus 3pt minus 0.5pt}\noindent\hspace*{1.2em}\kbvnum{(೩)}\hspace{0.55em}ಯಾವ ಶಬ್ದವು ಗಂಡಸು ಎಂಬ ಅರ್ಥವನ್ನು ಹೆಣ್ಣು ಎಂಬ ಅರ್ಥವನ್ನು ಉಂಟುಮಾಡುವುದಿಲ್ಲವೋ ಅದು ನಪುಂಸಕ ಲಿಂಗ ಶಬ್ದವಾಗಿರುವುದು. {[}ಪ್ರಾಣಿಗಳಲ್ಲಿ, ಪಕ್ಷಿಗಳಲ್ಲಿ ಮತ್ತು ಕೆಲವು ಜಾತಿಯ ಸಸ್ಯಗಳಲ್ಲಿ ಗಂಡು ಹೆಣ್ಣು ಎಂಬ ಪ್ರಭೇದಗಳುಂಟು. ಆದಾಗ್ಯೂ ಅವುಗಳು ವನ್ನೂ ನಪುಂಸಕ ಲಿಂಗಕ್ಕೆ ಸೇರಿದವುಗಳೆಂತಲೇ ಹೇಳುವುದು ರೂಢಿ ಯಾಗಿದೆ{]}.\par

\noindent\hspace*{1.2em}\kbvlabel{ಉದಾ}\hspace{0.6em}ಕೋಣ, ಎಮ್ಮೆ, ಹುಂಜ, ಕೋಳಿ, ಮರ, ಮನೆ, ಹಾಲು, ನದಿ, ಕಲ್ಲು, ಕತ್ತರಿ, ಮಗು, ಗಂಡಾನೆ, ಹೆಣ್ಣಾನೆ—ಇತ್ಯಾದಿ.\par

\indent ಸೂಚನೆ: ಕನ್ನಡ ವ್ಯಾಕರಣ ಶಾಸ್ತ್ರದಲ್ಲಿ ಪುಲ್ಲಿಂಗ, ಸ್ತ್ರೀಲಿಂಗ ವಿಭೇದಗಳನ್ನು ಮನುಷ್ಯವರ್ಗಕ್ಕೆ ಸೀಮಿತಗೊಳಿಸಿರುವುದು ಗಮನಾರ್ಹ ಅಂಶ.\par

\kbvsechead{ವಿಶೇಷ ಸಂದರ್ಭಗಳು:}

\begin{kbvrule}
\noindent\kbvbullet\hspace{0.5em}ಕೆಲವು ಶಬ್ದಗಳನ್ನು ಪುಲ್ಲಿಂಗದಲ್ಲಿ ಹಾಗೂ ನಪುಂಸಕ ಲಿಂಗದಲ್ಲಿ ಬಳಸುವ ವಾಡಿಕೆಯುಂಟು. ಅಂತಹ ಶಬ್ದಗಳು ಪುನ್ನಪುಂಸಕ ಲಿಂಗಕ್ಕೆ ಸೇರಿದವುಗಳೆಂದು ಹೇಳುತ್ತೇವೆ.
\end{kbvrule}
}
\clearpage

% === tex/0043.tex (Printed Page 41) ===
\renewcommand{\kbvchaptername}{ನಾಮಪದಗಳ ರೂಪ ಬದಲಾವಣೆ ({\kbvlatin i})}
\setcounter{page}{41}
\thispagestyle{kbvpage}
\kbvtypesetpage{\noindent\hspace*{1.2em}\kbvlabel{ಉದಾ}\hspace{0.6em}ಸೂರ್ಯನು ಚೆನ್ನಾಗಿ ಪ್ರಕಾಶಿಸುತ್ತಿದ್ದಾನೆ. (ಪುಲ್ಲಿಂಗ)\par

\indent ಬೆಳಿಗ್ಗೆ ಆರು ಗಂಟೆಗೆ ಸೂರ್ಯೋದಯವಾಯಿತು. (ನಪುಂಸಕ ಲಿಂಗ)\par

\begin{kbvrule}
\noindent\kbvbullet\hspace{0.5em}ಕೆಲವು ಶಬ್ದಗಳನ್ನು ಸ್ತ್ರೀಲಿಂಗದಲ್ಲಿ ಹಾಗೂ ನಪುಂಸಕ ಲಿಂಗದಲ್ಲಿ ಬಳಸುವ ವಾಡಿಕೆಯುಂಟು. ಅಂತಹ ಶಬ್ದಗಳು ಸ್ತ್ರೀನಪುಂಸಕ ಲಿಂಗಕ್ಕೆ ಸೇರಿದವುಗಳೆಂದು ಹೇಳುತ್ತೇವೆ.
\end{kbvrule}

\noindent\hspace*{1.2em}\kbvlabel{ಉದಾ}\hspace{0.6em}ಅಣ್ಣಮ್ಮದೇವಿ ವರ ಕೊಟ್ಟಳು. (ಸ್ತ್ರೀಲಿಂಗ)\par

\begin{kbvrule}
\noindent\kbvbullet\hspace{0.5em}ಅಣ್ಣಮ್ಮದೇವಿ ವರ ಕೊಟ್ಟಿತು. (ನಪುಂಸಕ ಲಿಂಗ)\par
\noindent\kbvbullet\hspace{0.5em}ಕೆಲವು ಶಬ್ದಗಳನ್ನು ಮೂರು ಲಿಂಗಗಳಲ್ಲೂ ಬಳಸುವ ಸಾಧ್ಯತೆ ಯುಂಟು. ಅಂತಹ ಶಬ್ದಗಳು ವಿಶೇಷ್ಯಾಧೀನ ಲಿಂಗಕ್ಕೆ ಅಥವಾ ವಾಚ್ಯ ಲಿಂಗಕ್ಕೆ ಸೇರಿದವುಗಳೆಂದು ಹೇಳುತ್ತೇವೆ.
\end{kbvrule}

\noindent\hspace*{1.2em}\kbvlabel{ಉದಾ}\hspace{0.6em}. ಅವನು ಜೀತದ ಆಳು. (ಪು.ಲಿಂ.) ಆ\par

\indent ಮನೆಗೆಲಸದ ಆಳು ಈ ದಿನ ಬಂದಿದ್ದಳು. (ಸ್ತ್ರೀ ಲಿಂ.) ಊರಿನಿಂದ ಆಳು ಬಂದಿತ್ತು. (ನ.ಲಿಂ.) ನಾನು ನಿನ್ನ ತಂದೆ. (ಪು.ಲಿಂ.) ನಾನು ನಿನ್ನ ತಾಯಿ. (ಸ್ತ್ರೀ.ಲಿಂ.) ನಾನು ನಿಮ್ಮ ಮಗು. (ನ.ಲಿಂ.)\par

\indent ವಿ. ಸೂ. : ಪ್ರಧಾನಿ, ಮಂತ್ರಿ, ವಿಜ್ಞಾನಿ, ಪ್ರಜೆ, ಕಾರ್ಯದರ್ಶಿ ಈ ಮೊದಲಾದ ಶಬ್ದಗಳು ವಾಚ್ಯಲಿಂಗ ಶಬ್ದಗಳಾಗಿರುವುವು.\par

\begin{kbvrule}
\noindent\kbvbullet\hspace{0.5em}ಕೆಲವು ಶಬ್ದಗಳನ್ನು ನಪುಂಸಕ ಲಿಂಗದಲ್ಲಿ ಮಾತ್ರ ಬಳಸುವುದು - ವಾಡಿಕೆ :: ಅವು ಪುಲ್ಲಿಂಗವಾದರೂ ಆಗಿರಬಹುದು. ಸ್ತ್ರೀಲಿಂಗವಾದರೂ ಆಗಿರಬಹುದು. ಅಂತಹ ಶಬ್ದಗಳು ನಿತ್ಯನಪುಂಸಕ ಲಿಂಗಕ್ಕೆ ಸೇರಿದವು - ಗಳೆಂದು ಹೇಳುತ್ತೇವೆ.
\end{kbvrule}

\noindent\hspace*{1.2em}\kbvlabel{ಉದಾ}\hspace{0.6em}ಮಗು, ಕೂಸು, ಶಿಶು, ಕಂದ, ಹಸುಳೆ, ದಂಡು, ಸೇನೆ, ದಳ, ಜನ ಇತ್ಯಾದಿ. \par
}
\clearpage

% === tex/0044.tex (Printed Page 42) ===
\renewcommand{\kbvchaptername}{ನಾಮಪದಗಳ ರೂಪ ಬದಲಾವಣೆ ({\kbvlatin i})}
\setcounter{page}{42}
\thispagestyle{kbvpage}
\kbvtypesetpage{\kbvsechead{ವಿಶೇಷ ರೀತಿಯಾಗಿ ಸ್ತ್ರೀಲಿಂಗ ರೂಪ ಶಬ್ದಗಳ ಉತ್ಪತ್ತಿ:}

\indent ಅಪ್ಪ \ensuremath{\times} ಅಮ್ಮ, ತಂದೆ \ensuremath{\times} ತಾಯಿ, ಚಿಕ್ಕಪ್ಪ \ensuremath{\times} ಚಿಕ್ಕಮ್ಮ, ತಮ್ಮ \ensuremath{\times} ತಂಗಿ, ವಿಧುರ \ensuremath{\times} ವಿಧವೆ, ರಾಜ \ensuremath{\times} ರಾಣಿ -- ಇವೇ ಮೊದಲಾದವುಗಳು ಸಾಮಾನ್ಯ ರೀತಿಯ ಪುಲ್ಲಿಂಗ-ಸ್ತ್ರೀಲಿಂಗ ಶಬ್ದಗಳು. ಆದರೆ ಕೆಲವು ವಿಶೇಷ ಸಂದರ್ಭಗಳಲ್ಲಿ ಪುಲ್ಲಿಂಗ ನಾಮಪ್ರಕೃತಿಗಳ ಮೇಲೆ ಎ, ಇ, ಅಳು, ಇತ್ತಿ , ಗಿತ್ತಿ, ಇತ್ತಿ ಮೊದಲಾದ ಪ್ರತ್ಯಯಗಳು ಸೇರಿ ಸ್ತ್ರೀಲಿಂಗ ರೂಪ ಶಬ್ದಗಳು ಉತ್ಪತ್ತಿಯಾಗುವುವು.\par

\vspace{3pt plus 2pt}\noindent\hspace*{1.2em}\kbvlabel{ಉದಾ}\par\vspace{1pt}

\begin{kbvderiv}
\vspace{4pt plus 3pt minus 1pt}
\begin{center}\begin{adjustbox}{max width=\textwidth}
{\renewcommand{\arraystretch}{1.22}\setlength{\tabcolsep}{5pt}
\begin{tabular}{l c l l}
ಗುಣವಂತ & $\times$ & ಗುಣವಂತೆ & \ldots\ (ಗುಣವಂತ + ಎ) \\
ಶಿಕ್ಷಕ & $\times$ & ಶಿಕ್ಷಕಿ & \ldots\ (ಶಿಕ್ಷಕ + ಇ) \\
ಮಗ & $\times$ & ಮಗಳು & \ldots\ (ಮಗ + ಅಳು) \\
ಬೀಯಗ & $\times$ & ಬೀಯಗಿತಿ & \ldots\ (ಬೀಯಗ + ಇತಿ) \\
ಕುಂಬಾರ & $\times$ & ಕುಂಬಾರಗಿತ್ತಿ & \ldots\ (ಕುಂಬಾರ + ಗಿತ್ತಿ) \\
ಹೂವಾಡಿಗ & $\times$ & ಹೂವಾಡಿಗಿತ್ತಿ & \ldots\ (ಹೂವಾಡಿಗ + ಇತ್ತಿ) \\
\end{tabular}}
\end{adjustbox}\end{center}
\vspace{4pt plus 3pt minus 1pt}
\end{kbvderiv}

\kbvexercise{ಅಭ್ಯಾಸ ಪತ್ರಿಕೆ ೯}

\vspace{2.5pt plus 3pt minus 0.5pt}\noindent\hspace*{1.2em}\kbvnum{(1)}\hspace{0.55em}ಲಿಂಗಗಳಲ್ಲಿ ಮುಖ್ಯವಾಗಿ ಎಷ್ಟು ಪ್ರಭೇದಗಳಿವೆ? ಅವು ಯಾವುವು? ಉದಾಹರಣೆಗಳ ಸಹಿತ ತಿಳಿಸಿ.\par

\vspace{2.5pt plus 3pt minus 0.5pt}\noindent\hspace*{1.2em}\kbvnum{(2)}\hspace{0.55em}ಈ ಪದಗಳ ಪುಲ್ಲಿಂಗ ರೂಪಗಳನ್ನು ತಿಳಿಸಿ:\par

\indent ಭಿಕ್ಷುಕಿ, ಸಂನ್ಯಾಸಿನಿ, ಗೆಳತಿ, ಕುಂಟಿ, ಯುವತಿ, ಕಿವುಡ, ಮಾಟಗಾತಿ, ಮೂಗಿ, ಹೆಳವ, ಕನ್ಯೆ.\par

\vspace{2.5pt plus 3pt minus 0.5pt}\noindent\hspace*{1.2em}\kbvnum{(3)}\hspace{0.55em}ಈ ಪದಗಳ ಸ್ತ್ರೀಲಿಂಗ ರೂಪಗಳನ್ನು ತಿಳಿಸಿ:\par

\indent ಗಾಣಿಗ, ಸೇವಕ, ದಾಸ, ಸೊಗಸುಗಾರ, ಪೂಜಾರಿ, ಮುದುಕ, ಹಣವಂತ, ಅಯ್ಯ, ಸಹೋದರ, ವರ.\par

\vspace{2.5pt plus 3pt minus 0.5pt}\noindent\hspace*{1.2em}\kbvnum{(4)}\hspace{0.55em}ಈ ಪದಗಳು ಯಾವ ಲಿಂಗಕ್ಕೆ ಸೇರಿದವುಗಳು?\par

\indent ಗಂಡುಮಗು, ದೇವರು, ಪೂಜಾರಿ, ಸಲಗ, ಅಮ್ಮ, ಗುಮಾಸ್ತೆ, ಚಂದ್ರ, ಗುಲಾಬಿ, ಹೋರಿಗಳು.\par
}
\clearpage

% === tex/0045.tex (Printed Page 43) ===
\renewcommand{\kbvchaptername}{ನಾಮಪದಗಳ ರೂಪ ಬದಲಾವಣೆ ({\kbvlatin ii})}
\setcounter{page}{43}
\thispagestyle{kbvpage}
\kbvtypesetpage{\vspace{2.5pt plus 3pt minus 0.5pt}\noindent\hspace*{1.2em}\kbvnum{(5)}\hspace{0.55em}“ನೀನು?” ಎಂಬ ಸರ್ವನಾಮವನ್ನು ಪುಲ್ಲಿಂಗ, ಸ್ತ್ರೀಲಿಂಗ, ಹಾಗೂ ನಪುಂಸಕ ಲಿಂಗಗಳಲ್ಲಿ ಬಳಸಬಹುದೆಂಬುದನ್ನು ಉದಾಹರಣೆಗಳ ಮೂಲಕ ಸಮರ್ಥಿಸಿ.\par

\vspace{2.5pt plus 3pt minus 0.5pt}\noindent\hspace*{1.2em}\kbvnum{(6)}\hspace{0.55em}ನಪುಂಸಕ ಲಿಂಗಕ್ಕೆ ಒಂದು ಉದಾಹರಣೆಯನ್ನು ಕೊಡಿ.\par

\vspace{2.5pt plus 3pt minus 0.5pt}\noindent\hspace*{1.2em}\kbvnum{(7)}\hspace{0.55em}ನಿತ್ಯ ನಪುಂಸಕ ಲಿಂಗ ಎಂದರೇನು ? ಎರಡು ಉದಾಹರಣೆಗಳನ್ನು ಕೊಡಿ.\par

\vspace{2.5pt plus 3pt minus 0.5pt}\noindent\hspace*{1.2em}\kbvnum{(8)}\hspace{0.55em}“ನಾನು ಈ ದೇಶದ ಪ್ರಜೆ.”— ಈ ವಾಕ್ಯದಲ್ಲಿರುವ ನಾಮಪದಗಳನ್ನು ಗುರ್ತಿಸಿ.ಲಿಂಗವನ್ನು ತಿಳಿಸಿ.\par

\kbvchapter{ಪಾಠ ೧೦}{ನಾಮಪದಗಳ ರೂಪ ಬದಲಾವಣೆ ({\kbvlatin ii})}

\kbvsubtitle{ವಚನ}

\indent ಒಂದು ಗಾಡಿ ಮುಂದೆ ಹೋಗುತ್ತಿದೆ ; ಆದರೆ ಹಿಂದೆ ಅನೇಕ ಗಾಡಿಗಳು ಬರುತ್ತಿವೆ.\par

\indent ಈ ವಾಕ್ಯಗಳಲ್ಲಿ ದಪ್ಪಕ್ಷರಗಳಲ್ಲಿರುವ “ಗಾಡಿ” ಗಾಡಿಗಳು” ಎಂಬ ಶಬ್ದಗಳನ್ನು ಗಮನಿಸಿ. ಇಲ್ಲಿ “ಗಾಡಿ” ಎಂಬ ಶಬ್ದ ಒಂದು ವಸ್ತುವನ್ನು ಮಾತ್ರ ಸೂಚಿಸುತ್ತದೆ ; ಆದರೆ “ಗಾಡಿಗಳು” ಎಂಬ ಶಬ್ದ ಒಂದಕ್ಕಿಂತ ಹೆಚ್ಚು ವಸ್ತುಗಳು ಎಂಬುದನ್ನು ಸೂಚಿಸುತ್ತದೆ.\par

\indent ಹೀಗೆ, ಒಂದು ವಸ್ತುವನ್ನು ಕುರಿತಾದ ಶಬ್ದವನ್ನು ಏಕವಚನ ಶಬ್ದ ಎಂದು ಒಂದಕ್ಕಿಂತ ಹೆಚ್ಚು ವಸ್ತುಗಳನ್ನು ಕುರಿತಾದ ಶಬ್ದವನ್ನು ಬಹುವಚನ ಶಬ್ದವೆಂದು ಹೇಳುತ್ತೇವೆ. ಕೆಳಗಿನ ಉದಾಹರಣೆಗಳನ್ನು ನೋಡಿ...\par

\vspace{5pt plus 4pt minus 1pt}
\begin{center}
{\renewcommand{\arraystretch}{1.24}\setlength{\tabcolsep}{10pt}\rowcolors{2}{kbvrow}{white}
\begin{tabularx}{0.88\textwidth}{X X}
{\kbvsecfont\color{kbvaccentdark}ಏಕವಚನ} & {\kbvsecfont\color{kbvaccentdark}ಬಹುವಚನ} \\[1pt]
\arrayrulecolor{kbvline}\specialrule{0.5pt}{0pt}{2pt}\arrayrulecolor{black}
ಆಳು & ಆಳುಗಳು \\
ಹಸು & ಹಸುಗಳು \\
ರೂಪಾಯಿ & ರೂಪಾಯಿಗಳು \\
ಗಿಣಿ & ಗಿಣಿಗಳು \\
\end{tabularx}}
\end{center}
\vspace{5pt plus 4pt minus 1pt}
}
\clearpage

% === tex/0046.tex (Printed Page 44) ===
\renewcommand{\kbvchaptername}{ನಾಮಪದಗಳ ರೂಪ ಬದಲಾವಣೆ ({\kbvlatin ii})}
\setcounter{page}{44}
\thispagestyle{kbvpage}
\kbvtypesetpage{\indent ಈ ಉದಾಹರಣೆಗಳಲ್ಲಿ ಏಕವಚನ ಶಬ್ದಗಳಿಗೆ “ಗಳು” ಎಂಬ ಆಗಮ ಸೇರಿ ಬಹುವಚನ ಶಬ್ದಗಳಾಗಿರುವುದು ಕಂಡುಬರುತ್ತದೆ. ಆಳು, ಹಸು, ರೂಪಾಯಿ, ಗಿಣಿ—ಇವು ನಪುಂಸಕ ಲಿಂಗ ಶಬ್ದಗಳು. ಆದುದರಿಂದ ನಪುಂಸಕ ಲಿಂಗ ಶಬ್ದಗಳಿಗೆ ಬಹುವಚನದಲ್ಲಿ “ಗಳು” ಎಂಬ ಆಗಮ ಸೇರುತ್ತದೆಂದು ತಿಳಿಯಬಹುದು.\par

\indent ವಿಶೇಷ ಸಂದರ್ಭ : ಅಕಾರಾಂತವಲ್ಲದ ಪುಲ್ಲಿಂಗ ಹಾಗೂ ಸ್ತ್ರೀಲಿಂಗ ಶಬ್ದಗಳಿಗೆ “ಗಳು” ಎಂಬ ಆಗಮ ಸೇರಿ ಬಹುವಚನ ರೂಪಗಳಾಗುವುದು.\par

\vspace{3pt plus 2pt}\noindent\hspace*{1.2em}\kbvlabel{ಉದಾ}\par\vspace{1pt}

\begin{kbvderiv}
\vspace{4pt plus 3pt minus 1pt}
\begin{center}\begin{adjustbox}{max width=\textwidth}
{\renewcommand{\arraystretch}{1.22}\setlength{\tabcolsep}{4pt}
\begin{tabular}{l c l}
ಗಂಡು + ಗಳು & {\color{kbvgrey}$=$} & \kbvresult{ಗಂಡುಗಳು} \\
ಹೆಣ್ಣು + ಗಳು & {\color{kbvgrey}$=$} & \kbvresult{ಹೆಣ್ಣುಗಳು} \\
\end{tabular}}
\end{adjustbox}\end{center}
\vspace{4pt plus 3pt minus 1pt}
\end{kbvderiv}

\begin{kbvrule}
\noindent\kbvbullet\hspace{0.5em}ಏಕವಚನದಲ್ಲಿರುವ ಪುಲ್ಲಿಂಗ, ಸ್ತ್ರೀಲಿಂಗ ಶಬ್ದಗಳಿಗೆ ಸಾಮಾನ್ಯ ವಾಗಿ “ಅರು” ಎಂಬ ಆಗಮ ಸೇರಲು ಬಹುವಚನ ರೂಪವುಂಟಾಗುವುದು.
\end{kbvrule}

\vspace{3pt plus 2pt}\noindent\hspace*{1.2em}\kbvlabel{ಉದಾ}\par\vspace{1pt}

\begin{kbvderiv}
\vspace{4pt plus 3pt minus 1pt}
\begin{center}\begin{adjustbox}{max width=\textwidth}
{\renewcommand{\arraystretch}{1.22}\setlength{\tabcolsep}{4pt}
\begin{tabular}{l c l}
ಅರಸ + ಅರು & {\color{kbvgrey}$=$} & \kbvresult{ಅರಸರು}{\color{kbvgrey} (ಬ.ವ.)} \\
ಹುಡುಗ + ಅರು & {\color{kbvgrey}$=$} & \kbvresult{ಹುಡುಗರು} \\
ಮುದುಕಿ + ಅರು & {\color{kbvgrey}$=$} & \kbvresult{ಮುದುಕಿಯರು} \\
ತಂಗಿ + ಅರು & {\color{kbvgrey}$=$} & \kbvresult{ತಂಗಿಯರು} \\
\end{tabular}}
\end{adjustbox}\end{center}
\vspace{4pt plus 3pt minus 1pt}
\end{kbvderiv}

\begin{kbvrule}
\noindent\kbvbullet\hspace{0.5em}ಏಕವಚನದಲ್ಲಿರುವ ಕೆಲವು ಪುಲ್ಲಿಂಗ, ಸ್ತ್ರೀಲಿಂಗ ಅಕಾರಾಂತ ಬಂಧುವಾಚಕ ಶಬ್ದಗಳಿಗೆ “ಅಂದಿರು” ಎಂಬ ಆಗಮ ಸೇರಲು ಬಹುವಚನ ರೂಪವುಂಟಾಗುವುದು.
\end{kbvrule}

\vspace{3pt plus 2pt}\noindent\hspace*{1.2em}\kbvlabel{ಉದಾ}\par\vspace{1pt}

\begin{kbvderiv}
\vspace{4pt plus 3pt minus 1pt}
\begin{center}\begin{adjustbox}{max width=\textwidth}
{\renewcommand{\arraystretch}{1.22}\setlength{\tabcolsep}{4pt}
\begin{tabular}{l c l}
ಅಣ್ಣ + ಅಂದಿರು & {\color{kbvgrey}$=$} & \kbvresult{ಅಣ್ಣಂದಿರು}{\color{kbvgrey} (ಬ.ವ.)} \\
ಅಳಿಯ + ಅಂದಿರು & {\color{kbvgrey}$=$} & \kbvresult{ಅಳಿಯಂದಿರು} \\
ಅಕ್ಕ + ಅಂದಿರು & {\color{kbvgrey}$=$} & \kbvresult{ಅಕ್ಕಂದಿರು} \\
ತಾಯಿ + ಅಂದಿರು & {\color{kbvgrey}$=$} & \kbvresult{ತಾಯಂದಿರು} \\
\end{tabular}}
\end{adjustbox}\end{center}
\vspace{4pt plus 3pt minus 1pt}
\end{kbvderiv}

\begin{kbvrule}
\noindent\kbvbullet\hspace{0.5em}ಕೆಲವು ಏಕವಚನ ಶಬ್ದಗಳಿಗೆ ಕಳು, ವು ಮೊದಲಾದ ಆಗಮಗಳು ಸೇರಿ ಬಹುವಚನ ರೂಪಗಳಾಗುವುದು.
\end{kbvrule}
}
\clearpage

% === tex/0047.tex (Printed Page 45) ===
\renewcommand{\kbvchaptername}{ನಾಮಪದಗಳ ರೂಪ ಬದಲಾವಣೆ ({\kbvlatin ii})}
\setcounter{page}{45}
\thispagestyle{kbvpage}
\kbvtypesetpage{\vspace{3pt plus 2pt}\noindent\hspace*{1.2em}\kbvlabel{ಉದಾ}\par\vspace{1pt}

\begin{kbvderiv}
\vspace{4pt plus 3pt minus 1pt}
\begin{center}\begin{adjustbox}{max width=\textwidth}
{\renewcommand{\arraystretch}{1.22}\setlength{\tabcolsep}{4pt}
\begin{tabular}{l c l}
ಮಗು + ಕಳು & {\color{kbvgrey}$=$} & \kbvresult{ಮಕ್ಕಳು}{\color{kbvgrey} (ಬ.ವ.)} \\
ಮಗ + ಕಳು & {\color{kbvgrey}$=$} & \kbvresult{ಮಕ್ಕಳು} \\
ಮಗಳು + ಕಳು & {\color{kbvgrey}$=$} & \kbvresult{ಮಕ್ಕಳು} \\
ದೊಡ್ಡದು + ವು & {\color{kbvgrey}$=$} & \kbvresult{ದೊಡ್ಡವು} \\
ಇದು + ವು & {\color{kbvgrey}$=$} & \kbvresult{ಇವು} \\
ನಾನು + ವು & {\color{kbvgrey}$=$} & \kbvresult{ನಾವು} \\
\end{tabular}}
\end{adjustbox}\end{center}
\vspace{4pt plus 3pt minus 1pt}
\end{kbvderiv}

\kbvexercise{ಅಭ್ಯಾಸ ಪತ್ರಿಕೆ ೧೦}

\vspace{2.5pt plus 3pt minus 0.5pt}\noindent\hspace*{1.2em}\kbvnum{(1)}\hspace{0.55em}ವಚನಗಳಲ್ಲಿ ಎಷ್ಟು ಪ್ರಕಾರಗಳಿವೆ? ಅವು ಯಾವುವು?\par

\vspace{2.5pt plus 3pt minus 0.5pt}\noindent\hspace*{1.2em}\kbvnum{(2)}\hspace{0.55em}ಏಕವಚನಕ್ಕೂ ಬಹುವಚನಕ್ಕೂ ಇರುವ ವ್ಯತ್ಯಾಸವೇನು?\par

\vspace{2.5pt plus 3pt minus 0.5pt}\noindent\hspace*{1.2em}\kbvnum{(3)}\hspace{0.55em}ಗಳು, ಅರು, ಅಂದಿರು - ಈ ಆಗಮಗಳು ಸೇರಿ ಆಗುವ ಬಹುವಚನ ಶಬ್ದಗಳಿಗೆ ಒಂದೊಂದು ಉದಾಹರಣೆಯನ್ನು ಕೊಡಿ.\par

\vspace{2.5pt plus 3pt minus 0.5pt}\noindent\hspace*{1.2em}\kbvnum{(4)}\hspace{0.55em}ಈ ಶಬ್ದಗಳ ಬಹುವಚನ ರೂಪಗಳನ್ನು ತಿಳಿಸಿ - ಕುದುರೆ, ಸೈನಿಕ, ಅಜ್ಜ, ರಾಣಿ, ಅದು, ಅವನು, ಹೆಂಡತಿ.\par

\vspace{2.5pt plus 3pt minus 0.5pt}\noindent\hspace*{1.2em}\kbvnum{(5)}\hspace{0.55em}ಈ ಶಬ್ದಗಳ ಏಕವಚನ ರೂಪಗಳನ್ನು ತಿಳಿಸಿ - ಮುದುಕರು, ತಮ್ಮಂದಿರು, ಪಕ್ಷಿಗಳು, ಕಳ್ಳರು, ಚಿಕ್ಕವು, ಇವರು.\par

\vspace{2.5pt plus 3pt minus 0.5pt}\noindent\hspace*{1.2em}\kbvnum{(6)}\hspace{0.55em}“ನಿಮ್ಮ ಅಣ್ಣಂದಿರು ರಾಮನ ಪಾತ್ರವನ್ನು ಚೆನ್ನಾಗಿ ಅಭಿನಯಿಸಿದರು” - ಈ ವಾಕ್ಯದಲ್ಲಿ ಅಣ್ಣಂದಿರು ಎಂಬ ಶಬ್ದವನ್ನು ಯಾವ ಅರ್ಥದಲ್ಲಿ ಬಳಸಲಾಗಿದೆ?\par

\vspace{2.5pt plus 3pt minus 0.5pt}\noindent\hspace*{1.2em}\kbvnum{(7)}\hspace{0.55em}ತಂದೆಯವರು, ಸುಬ್ಬಮ್ಮನವರು, ರಾಮಾಶಾಸ್ತ್ರಿಗಳು, ದಿವಾನರು - ಈ ನಾಮಪದಗಳಲ್ಲಿ ಕಂಡುಬರುವ ವಿಶೇಷವೇನು?\par

\vspace{2.5pt plus 3pt minus 0.5pt}\noindent\hspace*{1.2em}\kbvnum{(8)}\hspace{0.55em}ಈ ಬಹುವಚನ ಶಬ್ದಗಳಿಂದ ಆಗಮಗಳನ್ನು ಪ್ರತ್ಯೇಕಿಸಿ ಕ್ರಮವಾಗಿ ಬರೆಯಿರಿ - ಮಾವಂದಿರು, ಹಾವಾಡಿಗರು, ಇಲಿಗಳು, ಜವಾನರು.\par
}
\clearpage

% === tex/0048.tex (Printed Page 46) ===
\renewcommand{\kbvchaptername}{ನಾಮಪದಗಳ ರೂಪಬದಲಾವಣೆ ({\kbvlatin iii})}
\setcounter{page}{46}
\thispagestyle{kbvpage}
\kbvtypesetpage{\kbvchapter{ಪಾಠ ೧೧}{ನಾಮಪದಗಳ ರೂಪಬದಲಾವಣೆ ({\kbvlatin iii})}

\kbvsubtitle{ವಿಭಕ್ತಿ ಪ್ರತ್ಯಯಗಳು}

\kbvsechead{ಈ ವಾಕ್ಯಗಳನ್ನು ಓದಿ ನೋಡಿ :}

\vspace{2.5pt plus 3pt minus 0.5pt}\noindent\hspace*{1.2em}\kbvnum{(೧)}\hspace{0.55em}ರಾಮನು ದಶರಥನ ಮಗ.\par

\vspace{2.5pt plus 3pt minus 0.5pt}\noindent\hspace*{1.2em}\kbvnum{(೨)}\hspace{0.55em}ರಾಮನನ್ನು ಸೀತೆಯು ಮದುವೆಯಾದಳು.\par

\vspace{2.5pt plus 3pt minus 0.5pt}\noindent\hspace*{1.2em}\kbvnum{(೩)}\hspace{0.55em}ರಾಮನಿಂದ ಮಾರೀಚನ ಸಂಹಾರವಾಯಿತು.\par

\vspace{2.5pt plus 3pt minus 0.5pt}\noindent\hspace*{1.2em}\kbvnum{(೪)}\hspace{0.55em}ರಾಮನಿಗೆ ಬಿಲ್ಲುವಿದ್ಯೆ ಚೆನ್ನಾಗಿ ತಿಳಿದಿತ್ತು.\par

\vspace{2.5pt plus 3pt minus 0.5pt}\noindent\hspace*{1.2em}\kbvnum{(೫)}\hspace{0.55em}ರಾಮನ ದೆಸೆಯಿಂದ ಸುಗ್ರೀವನಿಗೆ ರಾಜ್ಯ ದಕ್ಕಿತು.\par

\vspace{2.5pt plus 3pt minus 0.5pt}\noindent\hspace*{1.2em}\kbvnum{(೬)}\hspace{0.55em}ರಾಮನ ಭಂಟ ಹನುಮಂತ.\par

\vspace{2.5pt plus 3pt minus 0.5pt}\noindent\hspace*{1.2em}\kbvnum{(೭)}\hspace{0.55em}ರಾಮನಲ್ಲಿ ಪ್ರಜಾವಾತ್ಸಲ್ಯ ತುಂಬಿತ್ತು.\par

\vspace{2.5pt plus 3pt minus 0.5pt}\noindent\hspace*{1.2em}\kbvnum{(೮)}\hspace{0.55em}ರಾಮನೇ !... ನೀನೇ ನಮಗೆ ದಿಕ್ಕು.\par

\indent ಮೇಲಿನ ವಾಕ್ಯಗಳಲ್ಲಿ ದಪ್ಪಕ್ಷರಗಳಲ್ಲಿರುವ ಪದಗಳನ್ನು ಗಮನಿಸಿ. ಇಲ್ಲಿ “ರಾಮ” ಎಂಬ ನಾಮಪ್ರಕೃತಿಯ ಮೇಲೆ ಪ್ರತ್ಯಯಗಳು ಸೇರಿ ನಾಮಪದಗಳಾಗಿರುವುದನ್ನು ಕಾಣಬಹುದು.\par

\begin{kbvderiv}
\vspace{4pt plus 3pt minus 1pt}
\begin{center}\begin{adjustbox}{max width=\textwidth}
{\renewcommand{\arraystretch}{1.22}\setlength{\tabcolsep}{4pt}
\begin{tabular}{l c l}
ರಾಮನು & {\color{kbvgrey}$=$} & ರಾಮ + (ಉ) \\
ರಾಮನನ್ನು & {\color{kbvgrey}$=$} & ರಾಮ + (ಅನ್ನು) \\
ರಾಮನಿಂದ & {\color{kbvgrey}$=$} & ರಾಮ + (ಇಂದ) \\
ರಾಮನಿಗೆ & {\color{kbvgrey}$=$} & ರಾಮ + (ಇಗೆ) \\
ರಾಮನ ದೆಸೆಯಿಂದ & {\color{kbvgrey}$=$} & ರಾಮನ + (ದೆಸೆಯಿಂದ) \\
ರಾಮನ & {\color{kbvgrey}$=$} & ರಾಮ + (ಅ) \\
ರಾಮನಲ್ಲಿ & {\color{kbvgrey}$=$} & ರಾಮ + (ಅಲ್ಲಿ) \\
ರಾಮನೇ & {\color{kbvgrey}$=$} & ರಾಮ + (ಏ) \\
\end{tabular}}
\end{adjustbox}\end{center}
\vspace{4pt plus 3pt minus 1pt}
\end{kbvderiv}

\indent ಆವರಣಗಳಲ್ಲಿರುವ ಪ್ರತ್ಯಯಗಳು ಕ್ರಮವಾಗಿ ಮೇಲಿನ ವಾಕ್ಯಗಳಲ್ಲಿ ಕರ್ತೃ, ಕರ್ಮ, ಕ್ರಿಯಾಪದಗಳಿಗಿರುವ ಸಂಬಂಧವನ್ನು ಅರ್ಥವತ್ತಾಗಿ ತಿಳಿಸುತ್ತವೆ. ಈ ರೀತಿಯಾಗಿ, ನಾಮಪ್ರಕೃತಿಯ ಮುಂದೆ ಸೇರುವ ಪ್ರತ್ಯಯಗಳಿಗೆ “ವಿಭಕ್ತಿ ಪ್ರತ್ಯಯಗಳೆಂದು ಹೆಸರು.\par
}
\clearpage

% === tex/0049.tex (Printed Page 47) ===
\renewcommand{\kbvchaptername}{ನಾಮಪದಗಳ ರೂಪಬದಲಾವಣೆ ({\kbvlatin iii})}
\setcounter{page}{47}
\thispagestyle{kbvpage}
\kbvtypesetpage{\indent ವಿಭಕ್ತಿ ಪ್ರತ್ಯಯಗಳನ್ನು ಎಂಟು ವಿಧಗಳಾಗಿ ವಿಂಗಡಿಸಬಹುದು. ಅವುಗಳು ಯಾವುವೆಂದರೆ—\par

\vspace{5pt plus 4pt minus 1pt}
\noindent{\renewcommand{\arraystretch}{1.24}\setlength{\tabcolsep}{4pt}\rowcolors{2}{kbvrow}{white}
\begin{tabularx}{\textwidth}{X X X}
{\kbvsecfont\color{kbvaccentdark}ವಿಭಕ್ತಿ} & {\kbvsecfont\color{kbvaccentdark}ವಿಭಕ್ತಿ ಪ್ರತ್ಯಯ} &  \\[1pt]
\arrayrulecolor{kbvline}\specialrule{0.5pt}{0pt}{2pt}\arrayrulecolor{black}
(೧) ಪ್ರಥಮಾ ವಿಭಕ್ತಿ & ... & ಉ \\
(೨) ದ್ವಿತೀಯಾ ವಿಭಕ್ತಿ & ... & ಅನ್ನು \\
(೩) ತೃತೀಯಾ ವಿಭಕ್ತಿ & ... & ಇಂದ \\
(೪) ಚತುರ್ಥಿ ವಿಭಕ್ತಿ & ... & ಇಗೆ, ಗೆ, ಕ್ಕೆ, ಅಕ್ಕೆ \\
(೫) ಪಂಚಮೀ ವಿಭಕ್ತಿ & ... & ದೆಸೆಯಿಂದ \\
(೬) ಷಷ್ಠಿ ವಿಭಕ್ತಿ & ... & ಎಂ, ಅ \\
(೭) ಸಪ್ತಮೀ ವಿಭಕ್ತಿ & ... & ಅಲ್ಲಿ \\
(೮) ಸಂಬೋಧನಾ ವಿಭಕ್ತಿ & ... & ಎ, ಏ—ಇತ್ಯಾದಿ \\
\end{tabularx}}
\vspace{5pt plus 4pt minus 1pt}

\kbvsechead{\textbf{ವಿಭಕ್ತಿ ರೂಪಗಳು}}

\vspace{4pt plus 4pt minus 1pt}\noindent{\kbvsecfont\fontsize{11pt}{14.5pt}\selectfont\color{kbvaccentdark} (೧)\hspace{0.45em}ಪ್ರಥಮಾ ವಿಭಕ್ತಿ:}\hspace{0.5em}ಪ್ರಥಮಾ ವಿಭಕ್ತಿ ಪ್ರತ್ಯಯ “ಉ” ಆಗಿದೆ. ಈ ವಿಭಕ್ತಿ ಪ್ರತ್ಯಯವು ಸಾಮಾನ್ಯವಾಗಿ ಕರ್ತೃಪದದ ಮೇಲೆ ಬರುವುದು.\par

\vspace{3pt plus 2pt}\noindent\hspace*{1.2em}\kbvlabel{ಉದಾ}\par\vspace{1pt}

\vspace{2.5pt plus 3pt minus 0.5pt}\noindent\hspace*{1.2em}\kbvnum{1.}\hspace{0.55em}ಸೂರ್ಯನು ಪೂರ್ವದಲ್ಲಿ ಉದಯಿಸಿದನು.\par

\indent (ಇಲ್ಲಿ ಸೂರ್ಯ- + ನ್ + ಉ = ಸೂರ್ಯನು)\par

\vspace{2.5pt plus 3pt minus 0.5pt}\noindent\hspace*{1.2em}\kbvnum{2.}\hspace{0.55em}ಎಮ್ಮೆಯು ಹಾಲನ್ನು ಕೊಡುತ್ತದೆ.\par

\begin{kbvrule}
\noindent\kbvbullet\hspace{0.5em}ಹಲವು ಸಂದರ್ಭಗಳಲ್ಲಿ ಪ್ರಥಮಾ ವಿಭಕ್ತಿಯ ಏಕವಚನ ರೂಪವನ್ನು ನಾಮಪ್ರಕೃತಿ ರೂಪದಲ್ಲೇ ಹೇಳುವುದು ವಾಡಿಕೆ. ಅಂದರೆ ನಾಮಪ್ರಕೃತಿಯ ರೂಪವೇ ಪ್ರಥಮಾ ವಿಭಕ್ತಿಯ ರೂಪವಾಗಿರುವುದು.
\end{kbvrule}

\vspace{3pt plus 2pt}\noindent\hspace*{1.2em}\kbvlabel{ಉದಾ}\par\vspace{1pt}

\vspace{2.5pt plus 3pt minus 0.5pt}\noindent\hspace*{1.2em}\kbvnum{1.}\hspace{0.55em}ಚಂದ್ರ ಆಕಾಶದಲ್ಲಿ ಹೊಳೆಯುತ್ತಾನೆ.\par

\vspace{2.5pt plus 3pt minus 0.5pt}\noindent\hspace*{1.2em}\kbvnum{2.}\hspace{0.55em}ಗಂಧಕ ನೀರಲ್ಲಿ ಕರಗುವುದಿಲ್ಲ.\par

\vspace{4pt plus 4pt minus 1pt}\noindent{\kbvsecfont\fontsize{11pt}{14.5pt}\selectfont\color{kbvaccentdark} (೨)\hspace{0.45em}ದ್ವಿತೀಯಾ ವಿಭಕ್ತಿ:}\hspace{0.5em}ದ್ವಿತೀಯಾ ವಿಭಕ್ತಿ ಪ್ರತ್ಯಯ “ಅನ್ನು” ಆಗಿದೆ. ಈ ವಿಭಕ್ತಿ ಪ್ರತ್ಯಯವು ಸಾಮಾನ್ಯವಾಗಿ ಕರ್ಮಪದದ ಮೇಲೆ ಬರುವುದು.\par
}
\clearpage

% === tex/0050.tex (Printed Page 48) ===
\renewcommand{\kbvchaptername}{ನಾಮಪದಗಳ ರೂಪಬದಲಾವಣೆ ({\kbvlatin iii})}
\setcounter{page}{48}
\thispagestyle{kbvpage}
\kbvtypesetpage{\noindent\hspace*{1.2em}\kbvlabel{ಉದಾ}\hspace{0.6em}1. ಬೆಕ್ಕು ಇಲಿಯನ್ನು ಕೊಂದಿತು.\par

\indent {[}ಇಲಿ + ಅನ್ನು = ಇಲಿಯನ್ನು ; ಇದು ಯಕಾರಾಗಮಸಂಧಿ{]}\par

\vspace{2.5pt plus 3pt minus 0.5pt}\noindent\hspace*{1.2em}\kbvnum{2.}\hspace{0.55em}ಕೋತಿಯು ಹಣ್ಣನ್ನು ತಿಂದಿತು.\par

\vspace{4pt plus 4pt minus 1pt}\noindent{\kbvsecfont\fontsize{11pt}{14.5pt}\selectfont\color{kbvaccentdark} (3)\hspace{0.45em}ತೃತೀಯಾ ವಿಭಕ್ತಿ :}\hspace{0.5em}ತೃತೀಯಾ ವಿಭಕ್ತಿ ಪ್ರತ್ಯಯ “ಇಂದ” ಆಗಿದೆ.\par

\noindent\hspace*{1.2em}\kbvlabel{ಉದಾ}\hspace{0.6em}1. ಸಕ್ಕರೆಯಿಂದ ಮಿಠಾಯಿ ತಯಾರಿಸುತ್ತಾರೆ.\par

\vspace{2.5pt plus 3pt minus 0.5pt}\noindent\hspace*{1.2em}\kbvnum{2.}\hspace{0.55em}ಹಣದಿಂದ ಗುಣ ಸಂಪಾದಿಸಲು ಸಾಧ್ಯವಿಲ್ಲ.\par

\vspace{2.5pt plus 3pt minus 0.5pt}\noindent\hspace*{1.2em}\kbvnum{3.}\hspace{0.55em}ಅಭ್ಯಾಸ ಮಾಡಿ ಮುಂದುವರೆದಿದೆ.\par

\vspace{2.5pt plus 3pt minus 0.5pt}\noindent\hspace*{1.2em}\kbvnum{4.}\hspace{0.55em}ಮೂರಕ್ಕೆ ಐದು ಸೇರಿದರೆ ಎಂಟಾಗುವುದು.\par

\vspace{4pt plus 4pt minus 1pt}\noindent{\kbvsecfont\fontsize{11pt}{14.5pt}\selectfont\color{kbvaccentdark} (4)\hspace{0.45em}ಚತುರ್ಥಿ ವಿಭಕ್ತಿ :}\hspace{0.5em}ಚತುರ್ಥಿ ವಿಭಕ್ತಿ ಪ್ರತ್ಯಯಗಳು “ಇಗೆ, ಗೆ, ಕ್ಕೆ, ಅಕ್ಕೆ” ಆಗಿವೆ.\par

\noindent\hspace*{1.2em}\kbvlabel{ಉದಾ}\hspace{0.6em}1. ರಾಮನಿಗೆ ಕಾಲು ಉಳುಕಿದೆ.\par

\vspace{2.5pt plus 3pt minus 0.5pt}\noindent\hspace*{1.2em}\kbvnum{2.}\hspace{0.55em}ಆನೆಗೆ ನಾಲ್ಕು ಕಾಲುಗಳುಂಟು.\par

\vspace{4pt plus 4pt minus 1pt}\noindent{\kbvsecfont\fontsize{11pt}{14.5pt}\selectfont\color{kbvaccentdark} (5)\hspace{0.45em}ಪಂಚಮೀ ವಿಭಕ್ತಿ :}\hspace{0.5em}ಪಂಚಮೀ ವಿಭಕ್ತಿ ಪ್ರತ್ಯಯ “ದೆಸೆಯಿಂದ” ಎಂಬುದಾಗಿದೆ.\par

\noindent\hspace*{1.2em}\kbvlabel{ಉದಾ}\hspace{0.6em}1. ಹಣದ ದೆಸೆಯಿಂದ ಸ್ನೇಹ ಮುರಿದು ಬಿತ್ತು.\par

\vspace{2.5pt plus 3pt minus 0.5pt}\noindent\hspace*{1.2em}\kbvnum{2.}\hspace{0.55em}ಕೌರವರ ದೆಸೆಯಿಂದ ಯುದ್ಧ ಸಂಭವಿಸಿತು.\par

\vspace{4pt plus 4pt minus 1pt}\noindent{\kbvsecfont\fontsize{11pt}{14.5pt}\selectfont\color{kbvaccentdark} (6)\hspace{0.45em}ಷಷ್ಠಿ ವಿಭಕ್ತಿ :}\hspace{0.5em}ಷಷ್ಠಿ ವಿಭಕ್ತಿ ಪ್ರತ್ಯಯ “ಅ” ಆಗಿದೆ.\par

\noindent\hspace*{1.2em}\kbvlabel{ಉದಾ}\hspace{0.6em}1. ರಾಜನ ಹೆಂಡತಿ ರಾಣಿ.\par

\vspace{2.5pt plus 3pt minus 0.5pt}\noindent\hspace*{1.2em}\kbvnum{2.}\hspace{0.55em}ದುರ್ಜನರ ಸಂಗ ಮಾಡಬಾರದು.\par

\vspace{4pt plus 4pt minus 1pt}\noindent{\kbvsecfont\fontsize{11pt}{14.5pt}\selectfont\color{kbvaccentdark} (7)\hspace{0.45em}ಸಪ್ತಮೀ ವಿಭಕ್ತಿ :}\hspace{0.5em}ಸಪ್ತಮೀ ವಿಭಕ್ತಿ ಪ್ರತ್ಯಯ “ಅಲ್ಲಿ” ಆಗಿದೆ.\par

\noindent\hspace*{1.2em}\kbvlabel{ಉದಾ}\hspace{0.6em}1. ನನ್ನ ಜೇಬಿನಲ್ಲಿ ಹಣವಿಲ್ಲ.\par

\vspace{2.5pt plus 3pt minus 0.5pt}\noindent\hspace*{1.2em}\kbvnum{2.}\hspace{0.55em}ದೇವರಲ್ಲಿ ಭಕ್ತಿಯಿಡು.\par
}
\clearpage

% === tex/0051.tex (Printed Page 49) ===
\renewcommand{\kbvchaptername}{ನಾಮಪದಗಳ ರೂಪಬದಲಾವಣೆ ({\kbvlatin iii})}
\setcounter{page}{49}
\thispagestyle{kbvpage}
\kbvtypesetpage{\vspace{4pt plus 4pt minus 1pt}\noindent{\kbvsecfont\fontsize{11pt}{14.5pt}\selectfont\color{kbvaccentdark} ೮.\hspace{0.45em}ಸಂಬೋಧನಾ ವಿಭಕ್ತಿ :}\hspace{0.5em}ಮತ್ತೊಬ್ಬರನ್ನು ಸಮೀಪಕ್ಕೆ ಕರೆಯುವಾಗ, ಮತ್ತೊಬ್ಬರಲ್ಲಿ ಮೊರೆಯಿಡುವಾಗ, ಮತ್ತೊಬ್ಬರೊಂದಿಗೆ ಏನನ್ನಾದರೂ ಹೇಳುವಾಗ ಈ ವಿಭಕ್ತಿ ರೂಪವನ್ನು ಬಳಸುತ್ತೇವೆ. ಸಂಬೋಧನಾ ವಿಭಕ್ತಿ ಪ್ರತ್ಯಯಗಳು ಎ, ಏ ಆಗಿವೆ.\par

\begin{kbvderiv}
\vspace{4pt plus 3pt minus 1pt}
\begin{center}\begin{adjustbox}{max width=\textwidth}
{\renewcommand{\arraystretch}{1.22}\setlength{\tabcolsep}{4pt}
\begin{tabular}{l c l}
ಉದಾ : ತಾಯೆ ! ನನ್ನನ್ನು ಹರಸು. (ತಾಯಿ + ಎ & {\color{kbvgrey}$=$} & \kbvresult{ತಾಯೆ)} \\
ತಾಯೇ ! ನನ್ನನ್ನು ರಕ್ಷಿಸು. (ತಾಯಿ + ಏ & {\color{kbvgrey}$=$} & \kbvresult{ತಾಯೇ)} \\
\end{tabular}}
\end{adjustbox}\end{center}
\vspace{4pt plus 3pt minus 1pt}
\end{kbvderiv}

\begin{kbvrule}
\noindent\kbvbullet\hspace{0.5em}ಕೆಲವು ಸಂದರ್ಭಗಳಲ್ಲಿ ನಾಮಪ್ರಕೃತಿ ರೂಪವನ್ನೇ ಸಂಬೋಧನಾ ವಿಭಕ್ತಿ ರೂಪವಾಗಿ ಬಳಸುವುದುಂಟು.
\end{kbvrule}

\noindent\hspace*{1.2em}\kbvlabel{ಉದಾ}\hspace{0.6em}ಅಕ್ಕಾ ! ಇಲ್ಲಿ ಬಾ.\par

\indent ಈ ವಾಕ್ಯದಲ್ಲಿ “ಅಕ್ಕಾ” ಶಬ್ದವನ್ನು ಸಂಬೋಧನಾ ವಿಭಕ್ತಿ ರೂಪದಲ್ಲಿ ಬಳಸಲಾಗಿದೆ. “ಅಕ್ಕಾ” ಎನ್ನುವಲ್ಲಿ ನಾಮಪ್ರಕೃತಿಯ ಕಡೆಯ ಸ್ವರ ದೀರ್ಘವಾಗಿರುವುದು ಗಮನಾರ್ಹ. (ಇಲ್ಲಿ “ಅಕ್ಕಾ” ಎನ್ನುವುದು “ಆ” ಕಾರಾಂತ, ಸ್ತ್ರೀಲಿಂಗ ಶಬ್ದ.)\par

\indent ಕೆಲವೊಮ್ಮೆ ನಾಮಪ್ರಕೃತಿಯ ಕಡೆಯ ಸ್ವರ ಹ್ರಸ್ವವಾಗಿರುವುದೂ ಉಂಟು. ಉದಾಹರಣೆಗೆ “ಮಗು, ಈ ಚೆಂಡು ತೆಗೆದುಕೊ.” (ಇಲ್ಲಿ “ಮಗು” ಎನ್ನುವುದು ಉಕಾರಾಂತ, ನಪುಂಸಕ ಲಿಂಗ ಶಬ್ದ.)\par

\begin{kbvrule}
\noindent\kbvbullet\hspace{0.5em}ಪುಲ್ಲಿಂಗ, ಸ್ತ್ರೀಲಿಂಗ, ಮತ್ತು ನಪುಂಸಕ ಲಿಂಗಗಳಲ್ಲಿನ ನಾಮಪ್ರಕೃತಿಗಳಿಗೆ ವಿಭಕ್ತಿ ಪ್ರತ್ಯಯಗಳು ಸೇರಿ ಆಗುವ ಶಬ್ದಗಳನ್ನು ಏಕವಚನ ಹಾಗೂ ಬಹುವಚನಗಳಲ್ಲಿ ಹೇಳಬಹುದು.
\end{kbvrule}

\indent ಉದಾಹರಣೆಗೆ ಅಗಸ, ರಾಣಿ, ದೇವರು — ಈ ಶಬ್ದಗಳ ವಿಭಕ್ತಿ ರೂಪಗಳನ್ನು ನೋಡೋಣ.\par

\kbvsechead{ಅಕಾರಾಂತ ಪುಲ್ಲಿಂಗ ಶಬ್ದ “ಅಗಸ” :}

\vspace{5pt plus 4pt minus 1pt}
\noindent{\renewcommand{\arraystretch}{1.24}\setlength{\tabcolsep}{4pt}
\begin{tabularx}{\textwidth}{X X X X}
{\kbvsecfont\color{kbvaccentdark}ವಿ.} & {\kbvsecfont\color{kbvaccentdark}ಪ್ರ.} & {\kbvsecfont\color{kbvaccentdark}ಏಕವಚನ} & {\kbvsecfont\color{kbvaccentdark}ಬಹುವಚನ} \\[1pt]
\arrayrulecolor{kbvline}\specialrule{0.5pt}{0pt}{2pt}\arrayrulecolor{black}
ಪ್ರ. & ವಿ. & ಅಗಸನು & ಅಗಸರು \\
ದ್ವಿ. & ವಿ. & ಅಗಸನನ್ನು & ಅಗಸರನ್ನು \\
\end{tabularx}}
\vspace{5pt plus 4pt minus 1pt}
}
\clearpage

% === tex/0052.tex (Printed Page 50) ===
\renewcommand{\kbvchaptername}{ನಾಮಪದಗಳ ರೂಪಬದಲಾವಣೆ ({\kbvlatin iii})}
\setcounter{page}{50}
\thispagestyle{kbvpage}
\kbvtypesetpage{\vspace{5pt plus 4pt minus 1pt}
\noindent{\renewcommand{\arraystretch}{1.24}\setlength{\tabcolsep}{4pt}\rowcolors{1}{kbvrow}{white}
\begin{tabularx}{\textwidth}{X X X}
ತೃ. ವಿ. & ಅಗಸನಿಂದ & ಅಗಸರಿಂದ \\
ಚ. ವಿ. & ಅಗಸನಿಗೆ & ಅಗಸರಿಗೆ \\
ಪಂ. ವಿ. & ಅಗಸನ ದೆಸೆಯಿಂದ & ಅಗಸರ ದೆಸೆಯಿಂದ \\
ಷ. ವಿ. & ಅಗಸನ & ಅಗಸರ \\
ಸ. ವಿ. & ಅಗಸನಲ್ಲಿ & ಅಗಸರಲ್ಲಿ \\
ಸಂ. ವಿ. & ಅಗಸನೇ ! & ಅಗಸರೇ ! \\
\end{tabularx}}
\vspace{5pt plus 4pt minus 1pt}

\kbvsechead{ಇಕಾರಾಂತ ಸ್ತ್ರೀಲಿಂಗ ಶಬ್ದ “ರಾಣಿ” :}

\vspace{5pt plus 4pt minus 1pt}
\noindent{\renewcommand{\arraystretch}{1.24}\setlength{\tabcolsep}{4pt}\rowcolors{2}{kbvrow}{white}
\begin{tabularx}{\textwidth}{X X X}
{\kbvsecfont\color{kbvaccentdark}ವಿ. ಪ್ರ.} & {\kbvsecfont\color{kbvaccentdark}ಏ. ವ.} & {\kbvsecfont\color{kbvaccentdark}ಬ. ವ.} \\[1pt]
\arrayrulecolor{kbvline}\specialrule{0.5pt}{0pt}{2pt}\arrayrulecolor{black}
ಪ್ರ. ವಿ. & ರಾಣಿಯು & ರಾಣಿಯರು \\
ದ್ವಿ. ವಿ. & ರಾಣಿಯನ್ನು & ರಾಣಿಯರನ್ನು \\
ತೃ. ವಿ. & ರಾಣಿಯಿಂದ & ರಾಣಿಯರಿಂದ \\
ಚ. ವಿ. & ರಾಣಿಗೆ & ರಾಣಿಯರಿಗೆ \\
ಪಂ. ವಿ. & ರಾಣಿಯ ದೆಸೆಯಿಂದ & ರಾಣಿಯರ ದೆಸೆಯಿಂದ \\
ಷ. ವಿ. & ರಾಣಿಯ & ರಾಣಿಯರ \\
ಸ. ವಿ. & ರಾಣಿಯಲ್ಲಿ & ರಾಣಿಯರಲ್ಲಿ \\
ಸಂ. ವಿ. & ರಾಣಿಯೇ ! & ರಾಣಿಗಳೇ ! \\
\end{tabularx}}
\vspace{5pt plus 4pt minus 1pt}

\kbvsechead{ಉಕಾರಾಂತ ನಪುಂಸಕ ಲಿಂಗ ಶಬ್ದ “ದೇವರು” :}

\vspace{5pt plus 4pt minus 1pt}
\noindent{\renewcommand{\arraystretch}{1.24}\setlength{\tabcolsep}{4pt}\rowcolors{2}{kbvrow}{white}
\begin{tabularx}{\textwidth}{X X X}
{\kbvsecfont\color{kbvaccentdark}ವಿ. ಪ್ರ.} & {\kbvsecfont\color{kbvaccentdark}ಏ. ವ.} & {\kbvsecfont\color{kbvaccentdark}ಬ. ವ.} \\[1pt]
\arrayrulecolor{kbvline}\specialrule{0.5pt}{0pt}{2pt}\arrayrulecolor{black}
ಪ್ರ. ವಿ. & ದೇವರು & ದೇವರುಗಳು \\
ದ್ವಿ. ವಿ. & ದೇವರನ್ನು & ದೇವರುಗಳನ್ನು \\
ತೃ. ವಿ. & ದೇವರಿಂದ & ದೇವರುಗಳಿಂದ \\
ಚ. ವಿ. & ದೇವರಿಗೆ & ದೇವರುಗಳಿಗೆ \\
ಪಂ. ವಿ. & ದೇವರ ದೆಸೆಯಿಂದ & ದೇವರುಗಳ ದೆಸೆಯಿಂದ \\
ಷ. ವಿ. & ದೇವರ & ದೇವರುಗಳ \\
ಸ. ವಿ. & ದೇವರಲ್ಲಿ & ದೇವರುಗಳಲ್ಲಿ \\
ಸಂ. ವಿ. & ದೇವರೇ ! & ದೇವರುಗಳೇ ! \\
\end{tabularx}}
\vspace{5pt plus 4pt minus 1pt}
}
\clearpage

% === tex/0053.tex (Printed Page 51) ===
\renewcommand{\kbvchaptername}{ನಾಮಪದಗಳ ರೂಪಬದಲಾವಣೆ ({\kbvlatin iii})}
\setcounter{page}{51}
\thispagestyle{kbvpage}
\kbvtypesetpage{\kbvexercise{ಅಭ್ಯಾಸ ಪತ್ರಿಕೆ ೧೧}

\vspace{2.5pt plus 3pt minus 0.5pt}\noindent\hspace*{1.2em}\kbvnum{(1)}\hspace{0.55em}ವಿಭಕ್ತಿ ಪ್ರತ್ಯಯಗಳು ಯಾವುವು? ಅವುಗಳ ಉಪಯುಕ್ತತೆಯೇನು?\par

\vspace{2.5pt plus 3pt minus 0.5pt}\noindent\hspace*{1.2em}\kbvnum{(2)}\hspace{0.55em}ವಿಭಕ್ತಿ ಪ್ರತ್ಯಯಗಳ ವಿವಿಧ ಪ್ರಕಾರಗಳನ್ನು ಕ್ರಮಾನುಸಾರವಾಗಿ ಬರೆದು ತಿಳಿಸಿ.\par

\vspace{2.5pt plus 3pt minus 0.5pt}\noindent\hspace*{1.2em}\kbvnum{(3)}\hspace{0.55em}ಸಂಬೋಧನಾ ವಿಭಕ್ತಿರೂಪಕ್ಕೆ ಒಂದು ಉದಾಹರಣೆಯನ್ನು ಕೊಡಿ.\par

\vspace{2.5pt plus 3pt minus 0.5pt}\noindent\hspace*{1.2em}\kbvnum{(4)}\hspace{0.55em}“ಮರ” ಶಬ್ದದ ಏಕವಚನ ರೂಪ ವಿಭಕ್ತಿ ಪ್ರತ್ಯಯಗಳನ್ನು ಬರೆಯಿರಿ.\par

\vspace{2.5pt plus 3pt minus 0.5pt}\noindent\hspace*{1.2em}\kbvnum{(5)}\hspace{0.55em}“ಹಾವು” ಶಬ್ದದ ಬಹುವಚನ ರೂಪ ವಿಭಕ್ತಿ ಪ್ರತ್ಯಯಗಳನ್ನು ಬರೆಯಿರಿ.\par

\vspace{2.5pt plus 3pt minus 0.5pt}\noindent\hspace*{1.2em}\kbvnum{(6)}\hspace{0.55em}“ಹುಡುಗನು ನಂದಿಬೆಟ್ಟವನ್ನು ಹತ್ತುತ್ತಿದ್ದಾನೆ” ಈ ವಾಕ್ಯವನ್ನು ಬಹುವಚನ ರೂಪದಲ್ಲಿ ಬರೆಯಿರಿ.\par

\vspace{2.5pt plus 3pt minus 0.5pt}\noindent\hspace*{1.2em}\kbvnum{(7)}\hspace{0.55em}“ಆಳುಗಳು ತೋಟದಲ್ಲಿ ಹಣ್ಣಿನ ಮರಗಳನ್ನು ಗರಗಸದಿಂದ ಕೊಯ್ದು ನೆಲಕ್ಕೆ ಉರುಳಿಸಿದರು.” ಈ ವಾಕ್ಯದಲ್ಲಿರುವ ನಾಮಪದಗಳನ್ನು ಗುರ್ತಿಸಿ, ಅವುಗಳ ವಿಭಕ್ತಿ ರೂಪಗಳನ್ನು ತಿಳಿಸಿರಿ.\par

\vspace{2.5pt plus 3pt minus 0.5pt}\noindent\hspace*{1.2em}\kbvnum{(8)}\hspace{0.55em}ಈ ನಾಮಪದಗಳ ವಿಭಕ್ತಿ ರೂಪವನ್ನು ತಿಳಿಸಿ :\par

\noindent\hspace*{1.8em}ಕಾಡಿನಲ್ಲಿ,\hspace{0.9em}ಹುಲಿಯನ್ನು,\hspace{0.9em}ಪಿಸ್ತೂಲಿನಿಂದ,\hspace{0.9em}ಕಬ್ಬಿಣಕ್ಕೆ,\hspace{0.9em}ಚೇಳು ರೂಪಾಯಿಗೆ,\hspace{0.9em}ಹೊಲದ,\hspace{0.9em}ಸ್ವಾಮಿಯೇ!\par\vspace{3pt plus 2pt}

\vspace{2.5pt plus 3pt minus 0.5pt}\noindent\hspace*{1.2em}\kbvnum{(9)}\hspace{0.55em}ಈ ನಾಮಪದಗಳಲ್ಲಿ ಪ್ರಕೃತಿ ಪ್ರತ್ಯಯಗಳನ್ನು ಬಿಡಿಸಿ ಬರೆಯಿರಿ.\par

\noindent\hspace*{1.8em}ಸಿಂಹವು ಜನರನ್ನು,\hspace{0.9em}ಮೇಲಕ್ಕೆ,\hspace{0.9em}ಹುಡುಗರಿಗೆ,\hspace{0.9em}ಕಲ್ಲಿನಿಂದ,\hspace{0.9em}ಚಿನ್ನದಲ್ಲಿ,\hspace{0.9em}ಮುಂದೆ,\hspace{0.9em}ಮಕ್ಕಳೇ,\hspace{0.9em}ಮಕ್ಕಳಿರಾ,\hspace{0.9em}ಭೀಮನ.\par\vspace{3pt plus 2pt}

\vspace{2.5pt plus 3pt minus 0.5pt}\noindent\hspace*{1.2em}\kbvnum{(10)}\hspace{0.55em}“ಅವನು” ಶಬ್ದದ ವಿಭಕ್ತಿ ಪ್ರತ್ಯಯಗಳನ್ನು ಏಕವಚನ ಹಾಗೂ ಬಹುವಚನ ರೂಪಗಳಲ್ಲಿ ಬರೆಯಿರಿ.\par
}
\clearpage

% === tex/0054.tex (Printed Page 52) ===
\renewcommand{\kbvchaptername}{ಕ್ರಿಯಾಪದ-ಲಿಂಗ, ವಚನ, ಕಾಲ}
\setcounter{page}{52}
\thispagestyle{kbvpage}
\kbvtypesetpage{\kbvchapter{ಪಾಠ ೧೨}{ಕ್ರಿಯಾಪದ-ಲಿಂಗ, ವಚನ, ಕಾಲ}

\noindent\hspace*{2.2em}ಕೆಳಗೆ ಕೊಟ್ಟಿರುವ ವಾಕ್ಯಗಳನ್ನು ನೋಡಿರಿ.\par

\indent - ಹುಡುಗನು ರೊಟ್ಟಿ ತಿನ್ನುತ್ತಿದ್ದಾನೆ. (ಪು. ಏ.)\par

\indent - ಹುಡುಗರು ರೊಟ್ಟಿ ತಿನ್ನುತ್ತಿದ್ದಾರೆ. (ಪು. ಬ.)\par

\indent - ಹುಡುಗಿಯು ರೊಟ್ಟಿ ತಿನ್ನುತ್ತಿದ್ದಾಳೆ. (ಸ್ತ್ರೀ.)\par

\indent - ಹುಡುಗಿಯರು ರೊಟ್ಟಿ ತಿನ್ನುತ್ತಿದ್ದಾರೆ. (ಸ್ತ್ರೀ. ಬ.)\par

\indent - ನಾಯಿಯು ರೊಟ್ಟಿ ತಿನ್ನುತ್ತಿದೆ. (ನ. ಏ.)\par

\indent - ನಾಯಿಗಳು ರೊಟ್ಟಿ ತಿನ್ನುತ್ತಿವೆ. (ನ. ಬ.)\par

\indent - ಮಕ್ಕಳು ರೊಟ್ಟಿ ತಿನ್ನುತ್ತಿದ್ದಾರೆ. (ನ. ಬ.)\par

\indent ಈ ವಾಕ್ಯಗಳಲ್ಲಿ “ತಿನ್ನು” ಧಾತುವಿನ ಮೇಲೆ ಲಿಂಗ ವಚನಗಳಿಗೆ ಅನುಸಾರವಾಗಿ ಬೇರೆಬೇರೆ ಆಖ್ಯಾತ ಪ್ರತ್ಯಯಗಳು ಸೇರಿ ಕ್ರಿಯಾಪದಗಳಾಗಿರುವುದನ್ನು ಕಾಣಬಹುದು.\par

\indent ಮೊದಲನೆಯ ವಾಕ್ಯದಲ್ಲಿ “ಹುಡುಗನು” ಎಂಬ ಕರ್ತೃಪದವು ಪುಲ್ಲಿಂಗ ಏಕವಚನದಲ್ಲಿರುವುದು ; ಆದುದರಿಂದ ಕ್ರಿಯಾಪದವೂ ಪುಲ್ಲಿಂಗ ಏಕವಚನದಲ್ಲಿರುವುದು.\par

\indent ಹೀಗೆ, ಪ್ರತಿಯೊಂದು ವಾಕ್ಯದಲ್ಲೂ ಕರ್ತೃಪದವು ಯಾವ ಲಿಂಗದಲ್ಲಿರುವುದೋ ಕ್ರಿಯಾಪದವೂ ಅದೇ ಲಿಂಗದಲ್ಲಿರುವುದು ; ಅಲ್ಲದೆ, ಕರ್ತೃಪದವು ಯಾವ ವಚನದಲ್ಲಿರುವುದೋ ಕ್ರಿಯಾಪದವೂ ಅದೇ ವಚನದಲ್ಲಿರುವುದು. ಕ್ರಿಯಾಪದವು ಕರ್ತೃವಿನ ಕೆಲಸವನ್ನು ತಿಳಿಸುವುದರಿಂದ, ಕರ್ತೃ ಹಾಗೂ ಕ್ರಿಯಾಪದಗಳಲ್ಲಿ ಲಿಂಗ ವಚನಗಳು ಏಕರೀತಿಯಾಗಿರಬೇಕಾದ್ದು ಅತ್ಯಾವಶ್ಯಕ.\par

\indent ಕಾಲ ಸೂಚಕ ಪ್ರತ್ಯಯಗಳು\par

\vspace{4pt plus 4pt minus 1pt}\noindent{\kbvsecfont\fontsize{11pt}{14.5pt}\selectfont\color{kbvaccentdark} (1)\hspace{0.45em}ವರ್ತಮಾನ ಕಾಲದ ಕ್ರಿಯಾಪದ ರೂಪಗಳು :}\hspace{0.5em}ಕ್ರಿಯೆ ಈಗ ನಡೆಯುತ್ತಿರುವುದಾಗಿ ತಿಳಿಸುವ ಕ್ರಿಯಾಪದವು. ವರ್ತಮಾನ ಕಾಲದ ಕ್ರಿಯಾಪದವೆನಿಸುವುದು.\par
}
\clearpage

% === tex/0055.tex (Printed Page 53) ===
\renewcommand{\kbvchaptername}{ಕ್ರಿಯಾಪದ-ಲಿಂಗ, ವಚನ, ಕಾಲ}
\setcounter{page}{53}
\thispagestyle{kbvpage}
\kbvtypesetpage{\indent ಇಲ್ಲಿ ಕೊಟ್ಟಿರುವ ವಾಕ್ಯಗಳಲ್ಲಿ ಕಂಡುಬರುವ ಕ್ರಿಯಾಪದಗಳನ್ನು ಗಮನಿಸಿರಿ.\par

\indent (ಪ್ರ. ಪು. ಏ.) ... ಅವನು ಓಡುತ್ತಾನೆ - ಓಡು + ಉತ್ತ + ಆನೆ\par

\indent (ಪ್ರ. ಸ್ತ್ರೀ. ಏ.) ... ಅವಳು ಓಡುತ್ತಾಳೆ - ಓಡು + ಉತ್ತ + ಆಳೆ\par

\indent (ಪ್ರ. ನ. ಏ.) ... ಅದು ಓಡುತ್ತದೆ - ಓಡು + ಉತ್ತ + ಅದೆ\par

\indent (ಮ. ವಾ. ಏ.) ... ನೀನು ಓಡುತ್ತೀಯೆ - ಓಡು + ಉತ್ತ + ಈಯೆ\par

\indent (ಉ. ವಾ. ಏ.) ... ನಾನು ಓಡುತ್ತೇನೆ - ಓಡು + ಉತ್ತ + ಏನೆ\par

\indent (ಪ್ರ. ವಾ. ಬ.) ... ಅವರು ಓಡುತ್ತಾರೆ - ಓಡು + ಉತ್ತ + ಆರೆ\par

\indent (ಪ್ರ. ನ. ಬ.) ... ಅವು ಓಡುತ್ತವೆ - ಓಡು + ಉತ್ತ + ಅವೆ\par

\indent (ಮ. ವಾ. ಬ.) ... ನೀವು ಓಡುತ್ತೀರಿ - ಓಡು + ಉತ್ತ + ಈರಿ\par

\indent (ಉ. ವಾ. ಬ.) ... ನಾವು ಓಡುತ್ತೇವೆ - ಓಡು + ಉತ್ತ + ಏವೆ\par

\indent ಇಲ್ಲಿ ಕ್ರಿಯಾಪದಗಳಿವೆ. ಈ ಕ್ರಿಯಾಪದಗಳಲ್ಲಿ “ಓಡು” ಧಾತುವಿನ ಮುಂದೆ ಎಲ್ಲ ಕಡೆಗೂ “ಉತ್ತ” ಎಂಬ ಪ್ರತ್ಯಯ ಬಂದಿದೆ ; ಈ ಪ್ರತ್ಯಯದ ಮುಂದೆ ಏಕವಚನದಲ್ಲಿ ಆನೆ, ಆಳೆ, ಅದೆ, ಈಯೆ, ಏನೆ ಎಂಬ ಆಖ್ಯಾತ ಪ್ರತ್ಯಯಗಳೂ ಬಹುವಚನದಲ್ಲಿ ಆರೆ, ಅವೆ, ಈರಿ, ಏವೆ ಎಂಬ ಆಖ್ಯಾತ ಪ್ರತ್ಯಯಗಳೂ ಬಂದಿವೆ.\par

\indent ಮೂಲ ಧಾತುವಿಗೂ ಆಖ್ಯಾತ ಪ್ರತ್ಯಯಕ್ಕೂ ನಡುವೆ ಬಂದಿರುವ “ಉತ್ತ” ಎಂಬುದು ಕಾಲ ಸೂಚಕ ಪ್ರತ್ಯಯ. ಇದು ವರ್ತಮಾನ ಕಾಲ ವನ್ನು ಸೂಚಿಸುವುದು. ಅಂದರೆ, ಕ್ರಿಯಾಪದವು ವರ್ತಮಾನ ಕಾಲದಲ್ಲಿರುವು ದೆಂಬುದನ್ನು ತಿಳಿಸುವುದು.\par

\kbvsechead{ಸಂಕೇತಾಕ್ಷರಗಳ ನಿರ್ಧಾರ :}

\begin{kbvderiv}
\vspace{4pt plus 3pt minus 1pt}
\noindent{\renewcommand{\arraystretch}{1.24}\setlength{\tabcolsep}{5pt}
\begin{tabularx}{\textwidth}{X X X}
ಪ್ರ. - ಪ್ರಥಮ ಪುರುಷ & ಪು - ಪುರುಷ & ಏ. - ಏಕವಚನ \\
ಮ. - ಮಧ್ಯಮ ಪುರುಷ & ಸ್ತ್ರೀ. - ಸ್ತ್ರೀಲಿಂಗ & ಬ. - ಬಹುವಚನ \\
ಉ. - ಉತ್ತಮ ಪುರುಷ & ವಾ. - ವಾಚ್ಯಲಿಂಗ &  \\
\end{tabularx}}
\vspace{4pt plus 3pt minus 1pt}
\end{kbvderiv}
}
\clearpage

% === tex/0056.tex (Printed Page 54) ===
\renewcommand{\kbvchaptername}{ಕ್ರಿಯಾಪದ-ಲಿಂಗ, ವಚನ, ಕಾಲ}
\setcounter{page}{54}
\thispagestyle{kbvpage}
\kbvtypesetpage{\indent ಮೇಲಿನ ವಾಕ್ಯಗಳಲ್ಲಿ ಕಂಡುಬರುವ ಕ್ರಿಯಾಪದಗಳಿಂದ ಕ್ರಿಯೆ ಈಗ ನಡೆಯುತ್ತಿರುವುದು ಎಂದು ತಿಳಿದು ಬರುವ ಕಾರಣ, ಈ ಕ್ರಿಯಾಪದಗಳಿಗೆ ವರ್ತಮಾನ ಕಾಲದ ಕ್ರಿಯಾಪದಗಳೆಂದು ಹೆಸರು. ಆದುದರಿಂದ ಮೇಲಿನ ಕ್ರಿಯಾಪದಗಳಲ್ಲಿರುವ ಆನೆ, ಆಳೆ, ಅದೆ, ಈಯೆ, ಏನೆ, ಆರೆ, ಅವೆ, ಈರಿ, ಏವೆ—ಎಂಬುವು ವರ್ತಮಾನ ಕಾಲದ ಆಖ್ಯಾತ ಪ್ರತ್ಯಯಗಳು\par

\vspace{4pt plus 4pt minus 1pt}\noindent{\kbvsecfont\fontsize{11pt}{14.5pt}\selectfont\color{kbvaccentdark} (2)\hspace{0.45em}ಭೂತ ಕಾಲದ ಕ್ರಿಯಾಪದ ರೂಪಗಳು :}\hspace{0.5em}ಗತಿಸಿಹೋದ ಕಾಲದಲ್ಲಿ ಕ್ರಿಯೆ ನಡೆಯಿತೆಂದು ತಿಳಿಸುವ ಕ್ರಿಯಾಪದವು. ಭೂತಕಾಲದ ಕ್ರಿಯಾಪದವೆನಿಸುವುದು.\par

\indent ಇಲ್ಲಿ ಕೊಟ್ಟಿರುವ ವಾಕ್ಯಗಳಲ್ಲಿ ಕಂಡುಬರುವ ಕ್ರಿಯಾಪದಗಳನ್ನು ಗಮನಿಸಿರಿ.\par

\indent (ಪ್ರ. ಪು. ಏ.)... ಅವನು ಓಡಿದನು -- ಓಡು + ದ + ಅನು\par

\indent (ಪ್ರ. ಸ್ತ್ರೀ. ಏ.)... ಅವಳು ಓಡಿದಳು -- ಓಡು + ದ + ಅಳು\par

\indent (ಪ್ರ. ನ. ಏ)... ಅದು ಓಡಿತು... ಓಡು + ದ + ಇತು\par

\indent (ಮ. ವಾ. ಏ)... ನೀನು ಓಡಿದೆ... ಓಡು + ದ + ಎ\par

\indent (ಉ. ವಾ. ಏ.)... ನಾನು ಓಡಿದೆನು... ಓಡು + ದ + ಎನು\par

\indent (ಪ್ರ. ವಾ. ಬ)... ಅವರು ಓಡಿದರು... ಓಡು + ದ + ಅರು\par

\indent (ಪ್ರ. ನ. ಬ)... ಅವು ಓಡಿದವು... ಓಡು + ದ + ಅವು\par

\indent (ಮ. ವಾ. ಬ.)... ನೀವು ಓಡಿದಿರಿ... ಓಡು + ದ + ಇರಿ\par

\indent (ಉ. ವಾ. ಬ)... ನಾವು ಓಡಿದೆವು... ಓಡು + ದ + ಎವು\par

\indent ಇಲ್ಲಿ ಒಂಬತ್ತು ಕ್ರಿಯಾಪದಗಳಿವೆ. ಈ ಕ್ರಿಯಾಪದಗಳಲ್ಲಿ “ಓಡು” ಧಾತುವಿನ ಮುಂದೆ ಎಲ್ಲ ಕಡೆಗೂ "ದ" ಎಂಬ ಪ್ರತ್ಯಯ ಬಂದಿದೆ ; ಪ್ರತ್ಯಯದ ಮುಂದೆ ಏಕವಚನದಲ್ಲಿ ಅನು, ಅಳು, ಇತು, ಎ, ಎನು, ಎಂಬ ಆಖ್ಯಾತ ಪ್ರತ್ಯಯಗಳು ಬಹುವಚನದಲ್ಲಿ ಅರು, ಅವು, ಇರಿ, ಎವು, ಎಂಬ ಆಖ್ಯಾತ ಪ್ರತ್ಯಯಗಳು ಬಂದಿವೆ. ಮೂಲ ಧಾತುವಿಗೂ ಆಖ್ಯಾತ ಪ್ರತ್ಯಯಕ್ಕೂ ನಡುವೆ ಬಂದಿರುವ "ದ" ಎಂಬುದು ಕಾಲಸೂಚಕ ಪ್ರತ್ಯಯ. ಇದು ಭೂತಕಾಲವನ್ನು ಸೂಚಿಸುವುದು. ಅಂದರೆ, ಕ್ರಿಯಾಪದವು ಭೂತಕಾಲದಲ್ಲಿರುವುದೆಂಬುದನ್ನು ತಿಳಿಸುವುದು.\par
}
\clearpage

% === tex/0057.tex (Printed Page 55) ===
\renewcommand{\kbvchaptername}{ಕ್ರಿಯಾಪದ-ಲಿಂಗ, ವಚನ, ಕಾಲ}
\setcounter{page}{55}
\thispagestyle{kbvpage}
\kbvtypesetpage{\indent ಮೇಲಿನ ವಾಕ್ಯಗಳಲ್ಲಿ ಕಂಡುಬರುವ ಕ್ರಿಯಾಪದಗಳಿಂದ ಕ್ರಿಯೆ ಹಿಂದೆ ನಡೆದಿರುವುದಾಗಿ ತಿಳಿದು ಬರುವ ಕಾರಣ, ಈ ಕ್ರಿಯಾಪದಗಳಿಗೆ ಭೂತಕಾಲದ ಕ್ರಿಯಾಪದಗಳೆಂದು ಹೆಸರು. ಆದುದರಿಂದ ಮೇಲಿನ ಕ್ರಿಯಾ ಪದಗಳಲ್ಲಿರುವ ಅನು, ಅಳು, ಇತು, ಎ, ಎನು, ಅರು, ಅವು, ಇರಿ, ಎವು— ಎಂಬುವು ಭೂತಕಾಲದ ಆಖ್ಯಾತ ಪ್ರತ್ಯಯಗಳು.\par

\vspace{4pt plus 4pt minus 1pt}\noindent{\kbvsecfont\fontsize{11pt}{14.5pt}\selectfont\color{kbvaccentdark} (3)\hspace{0.45em}ಭವಿಷ್ಯತ್‌ ಕಾಲದ ಕ್ರಿಯಾಪದ ರೂಪಗಳು :}\hspace{0.5em}ಮುಂದೆ ಬರುವ ಕಾಲದ ಕ್ರಿಯೆ ನಡೆಯುವುದನ್ನು ತಿಳಿಸುವ ಕ್ರಿಯಾಪದವು ಭವಿಷ್ಯತ್‌ ಕಾಲದ ಕ್ರಿಯಾಪದವೆನಿಸುವುದು.\par

\indent ಈ ಕೆಳಗೆ ಕೊಟ್ಟಿರುವ ವಾಕ್ಯಗಳಲ್ಲಿ ಕಂಡುಬರುವ ಕ್ರಿಯಾಪದಗಳನ್ನು ಗಮನಿಸಿರಿ.\par

\indent (ಪ್ರ.ಪು.ಏ.).... ಅವನು ಓಡುವನು — ಓಡು + ವ + ಅನು\par

\indent (ಪ್ರ.ಸ್ತ್ರೀ.ಏ.)... ಅವಳು ಓಡುವಳು — ಓಡು + ವ + ಅಳು\par

\indent (ಪ್ರ.ನ.ಏ. ).... ಅದು ಓಡುವುದು — ಓಡು + ವ + ಉದು\par

\indent (ಮ.ವಾ.ಏ.).... ನೀನು ಓಡುವೆ — ಓಡು + ವ + ಎ\par

\indent (ಉ.ವಾ.ಏ.).... ನಾನು ಓಡುವೆನು — ಓಡು + ವ + ಎನು\par

\indent (ಪ್ರ.ವಾ ಬ.).... ಅವರು ಓಡುವರು — ಓಡು + ವ + ಅರು\par

\indent (ಪ್ರ.ನ.ಬ. ).... ಅವು ಓಡುವುವು — ಓಡು + ವ + ಉವು\par

\indent (ಮ.ವಾ.ಬ.)..., ನೀವು ಓಡುವಿರಿ — ಓಡು + ವ + ಇರಿ\par

\indent (ಉ.ವಾ.ಬ).... ನಾವು ಓಡುವೆವು — ಓಡು + ವ + ಎವು\par

\indent ಇಲ್ಲಿ ಒಂಭತ್ತು ಕ್ರಿಯಾಪದಗಳಿವೆ. ಈ ಕ್ರಿಯಾಪದಗಳಲ್ಲಿ “ಓಡು” ಧಾತುವಿನ ಮುಂದೆ ಎಲ್ಲ ಕಡೆಗೂ “ವ” ಎಂಬ ಪ್ರತ್ಯಯ ಬಂದಿದೆ ; ಈ ಪ್ರತ್ಯಯದ ಮುಂದೆ ಏಕವಚನದಲ್ಲಿ ಅನು, ಅಳು, ಉದು, ಎ, ಎನು—ಎಂಬ ಆಖ್ಯಾತ ಪ್ರತ್ಯಯಗಳೂ ಬಹುವಚನದಲ್ಲಿ ಅರು, ಉವು, ಇರಿ, ಎವು— ಎಂಬ ಆಖ್ಯಾತ ಪ್ರತ್ಯಯಗಳೂ ಬಂದಿವೆ.\par

\indent ಮೂಲ ಧಾತುವಿಗೂ ಆಖ್ಯಾತ ಪ್ರತ್ಯಯಕ್ಕೂ ನಡುವೆ ಬಂದಿರುವ “ವ” ಎಂಬುದು ಕಾಲಸೂಚಕ ಪ್ರತ್ಯಯ. ಇದು ಭವಿಷ್ಯತ್‌ ಕಾಲವನ್ನು ಸೂಚಿಸು.\par
}
\clearpage

% === tex/0058.tex (Printed Page 56) ===
\renewcommand{\kbvchaptername}{ಕ್ರಿಯಾಪದ-ಲಿಂಗ, ವಚನ, ಕಾಲ}
\setcounter{page}{56}
\thispagestyle{kbvpage}
\kbvtypesetpage{\indent ವುದು. ಅಂದರೆ, ಕ್ರಿಯಾಪದವು ಭವಿಷ್ಯತ್‌ ಕಾಲದಲ್ಲಿರುವುದೆಂಬುದನ್ನು ತಿಳಿಸುವುದು.\par

\indent ಮೇಲಿನ ವಾಕ್ಯಗಳಲ್ಲಿ ಕಂಡುಬರುವ ಕ್ರಿಯಾಪದಗಳಿಂದ ಇವು ಮುಂದೆ ನಡೆಯುವ ಕ್ರಿಯೆಯನ್ನು ತಿಳಿಸುವ ಕ್ರಿಯಾಪದಗಳೆಂದು ತಿಳಿದುಬರುವುದು ಆದುದರಿಂದ ಈ ಕ್ರಿಯಾಪದಗಳಿಗೆ ಭವಿಷ್ಯತ್‌ ಕಾಲದ ಕ್ರಿಯಾಪದಗಳೆಂದು ಹೆಸರು. ಅದರಂತೆಯೇ ಮೇಲಿನ ಕ್ರಿಯಾಪದಗಳಲ್ಲಿರುವ ಅನು, ಅಳು, ಉದು, ಎ, ಎನು, ಅರು, ಉವು, ಇರಿ, ಎವು—ಎಂಬುವು ಭವಿಷ್ಯತ್‌ ಕಾಲದ ಆಖ್ಯಾತ ಪ್ರತ್ಯಯಗಳು.\par

\vspace{2.5pt plus 3pt minus 0.5pt}\noindent\hspace*{1.2em}\kbvnum{(4)}\hspace{0.55em}ವಿಧ್ಯರ್ಥಕ ಕ್ರಿಯಾಪದಗಳು\par

\vspace{3pt plus 2pt}\noindent\hspace*{1.2em}\kbvlabel{ಉದಾ}\par\vspace{1pt}

\indent ಅವರು ಮನೆಗೆ ಬರಲಿ .... (ಬರು + ಅಲಿ)\par

\indent ಸೀತೆಯು ಶಾಲೆಗೆ ಹೋಗಲಿ .... (ಹೋಗು + ಅಲಿ)\par

\indent ನೀವು ಸುಖವಾಗಿ ಬಾಳಿರಿ .... (ಬಾಳು + ಇರಿ)\par

\indent ನಾನು ಊಟ ಮಾಡುವೆ .... (ಮಾಡು + ಎ)\par

\indent ನಾವು ಹಾಡು ಹೇಳೋಣ .... (ಹೇಳು + ಓಣ)\par

\indent ಮೇಲಿನ ಉದಾಹರಣೆಗಳಲ್ಲಿ ಧಾತುಗಳೊಂದಿಗೆ ಅಲಿ, ಇರಿ, ಎ, ಓಣ\par

\indent —ಈ ಆಖ್ಯಾತ ಪ್ರತ್ಯಯಗಳು ಸೇರಿ ಕ್ರಿಯಾಪದಗಳಾಗಿರುವುವು. ಅಪ್ಪಣೆ, ಆಜ್ಞೆ, ಆಶೀರ್ವಾದ, ಹಾರೈಕೆ, ಸಮ್ಮತಿ ಕೊಡುವ ಸಂದರ್ಭಗಳಲ್ಲಿ ಬಳಸುವ ಈ ಕ್ರಿಯಾಪದಗಳಿಗೆ ವಿಧ್ಯರ್ಥಕ ಕ್ರಿಯಾಪದಗಳೆಂದು ಹೆಸರು. ಈ ಕ್ರಿಯಾ ಪದಗಳಲ್ಲಿ ಕಾಲಸೂಚಕ ಪ್ರತ್ಯಯ ಸೇರಿರುವುದಿಲ್ಲ.\par

\vspace{2.5pt plus 3pt minus 0.5pt}\noindent\hspace*{1.2em}\kbvnum{(5)}\hspace{0.55em}ನಿಷೇಧಾರ್ಥಕ ಕ್ರಿಯಾಪದಗಳು\par

\noindent\hspace*{1.2em}\kbvlabel{ಉದಾ}\hspace{0.6em}“ಮಾಡು” ಧಾತು.\par

\vspace{5pt plus 4pt minus 1pt}
\begin{center}
{\renewcommand{\arraystretch}{1.24}\setlength{\tabcolsep}{10pt}\rowcolors{2}{kbvrow}{white}
\begin{tabularx}{0.88\textwidth}{X X}
{\kbvsecfont\color{kbvaccentdark}ಏಕವಚನ} & {\kbvsecfont\color{kbvaccentdark}ಬಹುವಚನ} \\[1pt]
\arrayrulecolor{kbvline}\specialrule{0.5pt}{0pt}{2pt}\arrayrulecolor{black}
ಮಾಡು + ಅನು = ಮಾಡನು & ಮಾಡು + ಅರು = ಮಾಡರು \\
ಮಾಡು + ಅಳು = ಮಾಡಳು & ಮಾಡು + ಅರು = ಮಾಡರು \\
\end{tabularx}}
\end{center}
\vspace{5pt plus 4pt minus 1pt}
}
\clearpage

% === tex/0059.tex (Printed Page 57) ===
\renewcommand{\kbvchaptername}{ಕ್ರಿಯಾಪದ-ಲಿಂಗ, ವಚನ, ಕಾಲ}
\setcounter{page}{57}
\thispagestyle{kbvpage}
\kbvtypesetpage{\begin{kbvderiv}
\vspace{4pt plus 3pt minus 1pt}
\noindent{\renewcommand{\arraystretch}{1.24}\setlength{\tabcolsep}{5pt}
\begin{tabularx}{\textwidth}{X X}
ಮಾಡು + ಅದು = ಮಾಡದು & ಮಾಡು + ಅವು = ಮಾಡವು \\
ಮಾಡು + ಎ = ಮಾಡೆ & ಮಾಡು + ಅರಿ = ಮಾಡಿರಿ \\
ಮಾಡು + ಎನು = ಮಾಡೆನು & ಮಾಡು + ಎವು = ಮಾಡೆವು \\
\end{tabularx}}
\vspace{4pt plus 3pt minus 1pt}
\end{kbvderiv}

\indent ಇವು “ಮಾಡು” ಧಾತುವಿನ ಮೇಲೆ ಹಲವು ಪ್ರತ್ಯಯಗಳು ಸೇರಿ “ಕ್ರಿಯೆಯು ನಡೆಯದು” ಎಂಬ ಅರ್ಥ ಕೊಡುವ ಕ್ರಿಯಾಪದಗಳಾಗಿರುವುವು. ಹೀಗೆ, ಕ್ರಿಯೆಯು ನಡೆಯದು ಎಂಬರ್ಥ ಕೊಡುವ ಕ್ರಿಯಾಪದಗಳಿಗೆ ನಿಷೇಧಾರ್ಥಕ ಕ್ರಿಯಾಪದಗಳೆಂದು ಹೆಸರು.\par

\indent : ಧಾತುಗಳ ಮೇಲೆ ಏಕವಚನದಲ್ಲಿ ಅನು, ಅಳು, ಅದು, ಎ, ಎನು ಎಂಬ ಆಖ್ಯಾತ ಪ್ರತ್ಯಯಗಳೂ ಬಹುವಚನದಲ್ಲಿ ಅರು, ಅವು, ಅರಿ, ಎವು ಎಂಬ ಆಖ್ಯಾತ ಪ್ರತ್ಯಯಗಳೂ ಸೇರಿ ನಿಷೇಧಾರ್ಥಕ ಕ್ರಿಯಾಪದಗಳಾಗುವುವು.\par

\vspace{2.5pt plus 3pt minus 0.5pt}\noindent\hspace*{1.2em}\kbvnum{(6)}\hspace{0.55em}ಸಂಭಾವನಾರ್ಥಕ ಕ್ರಿಯಾಪದಗಳು\par

\noindent\hspace*{1.2em}\kbvlabel{ಉದಾ}\hspace{0.6em}“ಕೊಡು” ಧಾತು\par

\vspace{5pt plus 4pt minus 1pt}
\begin{center}
{\renewcommand{\arraystretch}{1.24}\setlength{\tabcolsep}{10pt}\rowcolors{2}{kbvrow}{white}
\begin{tabularx}{0.88\textwidth}{X X}
{\kbvsecfont\color{kbvaccentdark}ಏಕವಚನ} & {\kbvsecfont\color{kbvaccentdark}ಬಹುವಚನ} \\[1pt]
\arrayrulecolor{kbvline}\specialrule{0.5pt}{0pt}{2pt}\arrayrulecolor{black}
ಕೊಡು + ಆನು = ಕೊಟ್ಟಾನು & ಕೊಡು + ಆರು = ಕೊಟ್ಟಾರು \\
ಕೊಡು + ಆಳು = ಕೊಟ್ಟಾಳು & ಕೊಡು + ಆರು = ಕೊಟ್ಟಾರು \\
ಕೊಡು + ಈತು = ಕೊಟ್ಟೀತು & ಕೊಡು + ಆವು = ಕೊಟ್ಟಾವು \\
ಕೊಡು + ಈಯೇ = ಕೊಟ್ಟೀಯೆ & ಕೊಡು + ಈರಿ = ಕೊಟ್ಟೀರಿ \\
ಕೊಡು + ಏನು = ಕೊಟ್ಟೇನು & ಕೊಡು + ಏವು = ಕೊಟ್ಟೇವು \\
\end{tabularx}}
\end{center}
\vspace{5pt plus 4pt minus 1pt}

\indent {[}ಗಮನಿಸಿ : “ಕೊಡು” ಧಾತುವಿನ ಮೇಲೆ ಸಂಭಾವನಾರ್ಥಕ ಆಖ್ಯಾತ ಪ್ರತ್ಯಯಗಳು ಸೇರಿದಾಗ ಧಾತುವಿನ ಕೊನೆಯ “ಡು” ಕಾರಕ್ಕೆ “ಟ್ಟ” ಕಾರ ಆದೇಶವಾಗಿ ಬರುವುದು.{]}\par

\indent ಇಲ್ಲಿ “ಕೊಡು” ಧಾತುವಿನ ಮೇಲೆ ಹಲವು ಪ್ರತ್ಯಯಗಳು ಸೇರಿ “ಕ್ರಿಯೆಯು ಸಂಭವಿಸುವಲ್ಲಿ ಸಂಶಯತೋರುವ” ಕ್ರಿಯಾಪದಗಳಾಗಿರುವುವು.\par
}
\clearpage

% === tex/0060.tex (Printed Page 58) ===
\renewcommand{\kbvchaptername}{ಕ್ರಿಯಾಪದ-ಲಿಂಗ, ವಚನ, ಕಾಲ}
\setcounter{page}{58}
\thispagestyle{kbvpage}
\kbvtypesetpage{\indent ಹೀಗೆ, ಕ್ರಿಯೆ ಸಂಭವಿಸುವಲ್ಲಿ ಸಂಶಯ ಅಥವಾ ಸಂದೇಹ ಪ್ರಕಟಿಸುವ ಕ್ರಿಯಾಪದಗಳಿಗೆ ಸಂಭಾವನಾರ್ಥಕ ಕ್ರಿಯಾಪದಗಳೆಂದು ಹೆಸರು. ಧಾತುಗಳ ಮೇಲೆ ಏಕವಚನದಲ್ಲಿ ಆನು, ಆಳು, ಈತು, ಈಯೆ, ಏನು ಎಂಬ ಆಖ್ಯಾತ ಪ್ರತ್ಯಯಗಳೂ ಬಹುವಚನದಲ್ಲಿ ಆರು, ಆವು, ಈರಿ, ಏವು—ಎಂಬ ಆಖ್ಯಾತ ಪ್ರತ್ಯಯಗಳೂ ಸೇರಿ ಸಂಭಾವನಾರ್ಥಕ ಕ್ರಿಯಾಪದ ಗಳಾಗುವುವು.\par

\begin{kbvrule}
\noindent\kbvbullet\hspace{0.5em}ಸಾಕು, ಬೇಕು, ಇಲ್ಲ, ಹೌದು, ಬೇಡ, ಅಲ್ಲ, ಉಂಟು ಇತ್ಯಾದಿ ಅವ್ಯಯಗಳನ್ನು ಕ್ರಿಯಾರ್ಥಕ ಪದಗಳೆನ್ನಬಹುದು. ಕ್ರಿಯಾಪದಗಳ ಅರ್ಥ ದಲ್ಲಿ ಬಳಸುವ ಈ ಅವ್ಯಯಗಳನ್ನು ಕ್ರಿಯಾರ್ಥಕಾವ್ಯಯಗಳೆನ್ನುವರು.
\end{kbvrule}

\begin{kbvderiv}
\vspace{4pt plus 3pt minus 1pt}
\noindent{\renewcommand{\arraystretch}{1.24}\setlength{\tabcolsep}{5pt}
\begin{tabularx}{\textwidth}{X X}
ಉದಾ : {[}೧{]} ಮಾತು ಸಾಕು, ಕೆಲಸ ಬೇಕು & {[}೨{]} ನನ್ನ ಜೇಬಿ \\
ನಲ್ಲಿ ಹಣ ಇಲ್ಲ & {[}೩{]} ನಿನ್ನ ಸಹವಾಸ ಬೇಡ. {[}೪{]} ಇದು ರೇಷ್ಮೆ ಸೀರೆ \\
\end{tabularx}}
\vspace{4pt plus 3pt minus 1pt}
\end{kbvderiv}

\indent ಅಲ್ಲ. (೫) ಅದು ಧ್ರುವ ನಕ್ಷತ್ರವೇ ಹೌದು. (೬) ಆಕೆಯಲ್ಲಿ ನನಗೆ ಪ್ರೀತಿ ಉಂಟು.\par

\kbvsubtitle{ಕರ್ತರಿ ಮತ್ತು ಕರ್ಮಣಿ ಪ್ರಯೋಗಗಳು}

\indent ವಾಕ್ಯಗಳ ಎರಡು ಶ್ರೇಣಿಗಳನ್ನು ಇಲ್ಲಿ ಕೊಟ್ಟಿದೆ. ಪ್ರತಿಯೊಂದು ಶ್ರೇಣಿಯಲ್ಲೂ ಮೂರು ವಾಕ್ಯಗಳಿವೆ, ಆ ವಾಕ್ಯಗಳನ್ನೂ ಅವುಗಳಲ್ಲಿರುವ ಕ್ರಿಯಾಪದಗಳನ್ನು ಗಮನಿಸಿರಿ.\par

\kbvsechead{ಮೊದಲನೆಯ ಶ್ರೇಣಿ :}

\indent ರಾಮನು ಹೊಲವನ್ನು ಮಾರುತ್ತಾನೆ. ...(ಮಾರು + ಉತ್ತ + ಆನೆ)\par

\indent ಸೀತೆಯು ಹೊಲವನ್ನು ಮಾರಿದಳು ....(ಮಾರು + ದ + ಅಳು)\par

\indent ಅವರು ಹೊಲವನ್ನು ಮಾರುವರು ....(ಮಾರು + ವ + ಅರು)\par

\kbvsechead{ಎರಡನೆಯ ಶ್ರೇಣಿ :}

\noindent\hspace*{2.2em}ರಾಮನಿಂದ ಹೊಲವು ಮಾರಲ್ಪಡುತ್ತದೆ.\par

\indent ...(ಮಾರು + ಅಲ್ಪಡು + ಉತ್ತ + ಅದೆ) ಸೀತೆಯಿಂದ ಹೊಲವು ಮಾರಲ್ಪಟ್ಟಿತು.\par

\indent ...(ಮಾರು + ಅಲ್ಪಡು + ದ + ಇತು)\par
}
\clearpage

% === tex/0061.tex (Printed Page 59) ===
\renewcommand{\kbvchaptername}{ಕ್ರಿಯಾಪದ-ಲಿಂಗ, ವಚನ, ಕಾಲ}
\setcounter{page}{59}
\thispagestyle{kbvpage}
\kbvtypesetpage{\indent ಅವರಿಂದ ಹೊಲವು ಮಾರಲ್ಪಡುವುದು\par

\indent (ಮಾರು + ಅಲ್ಪಡು + ಉ + ವದು)\par

\indent ಮೊದಲನೆಯ ಶ್ರೇಣಿಯಲ್ಲಿರುವ ವಾಕ್ಯಗಳಲ್ಲಿ ಕರ್ತೃಪದವು ಪ್ರಥಮಾ ವಿಭಕ್ತ್ಯಂತವಾಗಿರುವುದು ಮತ್ತು ಕರ್ಮಪದವು ದ್ವಿತೀಯಾ ವಿಭಕ್ತ್ಯಂತವಾಗಿರುವುದು. ಅಲ್ಲದೆ ಈ ವಾಕ್ಯಗಳಲ್ಲಿ ಕ್ರಿಯಾಪದದ ಲಿಂಗ ವಚನಗಳು ಕರ್ತೃಪದದ ಲಿಂಗ ವಚನಗಳಂತೆಯೇ ಇರುವುವು. ಉದಾಹರಣೆಗೆ, “ರಾಮನು” ಎಂಬ ಕರ್ತೃಪದವು ಪುಲ್ಲಿಂಗ ಏಕವಚನವಾಗಿರುವುದರಿಂದ ಕ್ರಿಯಾಪದವೂ ಪುಲ್ಲಿಂಗ ಏಕವಚನ ರೂಪವನ್ನು ಹೊಂದಿ “ಮಾರುತ್ತಾನೆ” ಎಂದಾಗಿರುವುದು. ಈ ಲಕ್ಷಣಗಳುಳ್ಳ ಪ್ರಯೋಗಗಳನ್ನು ಕರ್ತರಿ ಪ್ರಯೋಗಗಳೆನ್ನುವರು.\par

\indent ಎರಡನೆಯ ಶ್ರೇಣಿಯಲ್ಲಿರುವ ವಾಕ್ಯಗಳಲ್ಲಿ ಕರ್ತೃಪದವು ತೃತೀಯಾ ವಿಭಕ್ತ್ಯಂತವಾಗಿರುವುದು ಮತ್ತು ಕರ್ಮಪದವು ಪ್ರಥಮಾ ವಿಭಕ್ತ್ಯಂತವಾಗಿರುವುದು. ಅಲ್ಲದೆ, ಕರ್ಮಪದದ ಲಿಂಗ ವಚನಗಳಂತೆ ಕ್ರಿಯಾಪದವೂ ಇರುವುದು. ಅಲ್ಲದೆ, ವಿಶೇಷವೇನೆಂದರೆ, ಕ್ರಿಯಾಪದಗಳಲ್ಲಿ ಧಾತುವಿನ ಮುಂದೆ “ಅಲ್ಪಡು” ಎಂಬ ಪ್ರತ್ಯಯ ಸೇರಿರುವುದು. ಈ ಲಕ್ಷಣಗಳುಳ್ಳ ಪ್ರಯೋಗಗಳನ್ನು ಕರ್ಮಣಿ ಪ್ರಯೋಗಗಳೆನ್ನುವರು.\par

\vspace{3pt plus 2pt}\noindent\hspace*{1.2em}\kbvlabel{ಉದಾ}\par\vspace{1pt}

\vspace{2.5pt plus 3pt minus 0.5pt}\noindent\hspace*{1.2em}\kbvnum{1.}\hspace{0.55em}ರೈತನು ಭೂಮಿಯನ್ನು ಉಳುತ್ತಾನೆ. {[}ಕರ್ತರಿ ಪ್ರಯೋಗ{]} ರೈತನಿಂದ ಭೂಮಿಯು ಉಳಲ್ಪಡುತ್ತದೆ. {[}ಕರ್ಮಣಿ ಪ್ರಯೋಗ{]}\par

\vspace{2.5pt plus 3pt minus 0.5pt}\noindent\hspace*{1.2em}\kbvnum{2.}\hspace{0.55em}ಬೆಕ್ಕು ಇಲಿಯನ್ನು ಕೊಂದಿತು. {[}ಕರ್ತರಿ ಪ್ರಯೋಗ{]} ಬೆಕ್ಕಿನಿಂದ ಇಲಿಯು ಕೊಲ್ಲಲ್ಪಟ್ಟಿತು. {[}ಕರ್ಮಣಿ ಪ್ರಯೋಗ{]}\par

\vspace{2.5pt plus 3pt minus 0.5pt}\noindent\hspace*{1.2em}\kbvnum{3.}\hspace{0.55em}ಅಗಸರು ಬಟ್ಟೆಗಳನ್ನು ಒಗೆಯುವರು. {[}ಕರ್ತರಿ ಪ್ರಯೋಗ{]} ಅಗಸರಿಂದ ಬಟ್ಟೆಗಳು ಒಗೆಯಲ್ಪಡುವುವು. {[}ಕರ್ಮಣಿ ಪ್ರಯೋಗ{]}\par

\kbvsechead{ನೆನಪಿಡಿ :}

\vspace{2.5pt plus 3pt minus 0.5pt}\noindent\hspace*{1.2em}\kbvnum{(1)}\hspace{0.55em}ಸಕರ್ಮಕ ಧಾತುವಿನ ಮುಂದೆ “ಅಲ್ಪಡು” ಪ್ರತ್ಯಯ ಬಂದ ನಂತರದಲ್ಲಿ ಕಾಲಸೂಚಕ ಪ್ರತ್ಯಯವೂ ಆಖ್ಯಾತ ಪ್ರತ್ಯಯವೂ ಸೇರಿ ಕರ್ಮಣಿ ಪ್ರಯೋಗವಾಗುವುದು.\par
}
\clearpage

% === tex/0062.tex (Printed Page 60) ===
\renewcommand{\kbvchaptername}{ಕ್ರಿಯಾಪದ-ಲಿಂಗ, ವಚನ, ಕಾಲ}
\setcounter{page}{60}
\thispagestyle{kbvpage}
\kbvtypesetpage{\vspace{2.5pt plus 3pt minus 0.5pt}\noindent\hspace*{1.2em}\kbvnum{{[}2{]}}\hspace{0.55em}ಕರ್ತರಿ ಪ್ರಯೋಗದಲ್ಲಿ ಕ್ರಿಯಾಪದಗಳಿಗೆ ಕರ್ತೃವಿನ ಲಿಂಗವಚನಗಳೂ, ಕರ್ಮಣಿ ಪ್ರಯೋಗದಲ್ಲಿ ಕರ್ಮಪದದ ಲಿಂಗವಚನಗಳೂ ಬರುತ್ತವೆ.\par

\noindent\hspace*{1.2em}\kbvlabel{ಉದಾ}\hspace{0.6em}ಸುಜಾತ ಹಣ್ಣನ್ನು ತಿಂದಳು.\par

\noindent\hspace*{2.2em}ಚಿತ್ರವು ರವಿಯಿಂದ ಬರೆಯಲ್ಪಟ್ಟಿತು.\par

\vspace{2.5pt plus 3pt minus 0.5pt}\noindent\hspace*{1.2em}\kbvnum{{[}3{]}}\hspace{0.55em}ಒಂದೇ ಕ್ರಿಯೆಗೆ ವಿಭಿನ್ನ ಲಿಂಗದ ಅನೇಕ ಕರ್ತೃಪದಗಳಿರುವಾಗ ಕಟ್ಟಕಡೆಯ ಕರ್ತೃಪದದ ಲಿಂಗವೇ ಕ್ರಿಯಾಪದಕ್ಕೆ ಬರುವುದು.\par

\noindent\hspace*{1.2em}\kbvlabel{ಉದಾ}\hspace{0.6em}ಸರ್ಕಸ್ಸಿನಲ್ಲಿ ಪ್ರಾಣಿಗಳೂ ಜನರೂ ಸಾಹಸದ ಕೆಲಸಗಳನ್ನು ಮಾಡಿದರು. ರಾಮನೂ ಲಕ್ಷ್ಮಣನೂ ಕಪಿ ಸೈನ್ಯವೂ ರಾವಣನ ಮೇಲೆ ಯುದ್ಧ ಮಾಡಿತು.\par

\vspace{2.5pt plus 3pt minus 0.5pt}\noindent\hspace*{1.2em}\kbvnum{{[}4{]}}\hspace{0.55em}ಒಂದೇ ಕ್ರಿಯೆಗೆ ಉತ್ತಮ ಪುರುಷ ಸರ್ವನಾಮವೂ, ಇನ್ನಿತರ ಸರ್ವನಾಮಗಳೂ ಕರ್ತೃಗಳಾಗಿದ್ದಲ್ಲಿ ಕ್ರಿಯಾಪದವು ಉತ್ತಮ ಪುರುಷದಲ್ಲೇ ಇರುವುದು.\par

\noindent\hspace*{1.2em}\kbvlabel{ಉದಾ}\hspace{0.6em}ನಾನೂ ನೀನೂ ಅವನೂ ಊರಿಗೆ ಹೋದೆವು.\par

\noindent\hspace*{2.2em}ನಾನೂ ಅವನೂ ಸೇತಿಯಿಂದ ಬಂದೆವು.\par

\noindent\hspace*{2.2em}ನಾನೂ ನೀನೂ ಕೂಡಿ ಬಾಳೋಣ.\par

\vspace{2.5pt plus 3pt minus 0.5pt}\noindent\hspace*{1.2em}\kbvnum{{[}5{]}}\hspace{0.55em}ಒಂದೇ ಕ್ರಿಯೆಗೆ ಮಧ್ಯಮ ಪುರುಷ ಸರ್ವನಾಮವೂ, ಪ್ರಥಮ ಪುರುಷ ಸರ್ವನಾಮವೂ ಕರ್ತೃಗಳಾಗಿದ್ದಲ್ಲಿ ಕ್ರಿಯಾಪದವು ಮಧ್ಯಮ ಪುರುಷದಲ್ಲೇ ಇರುವುದು.\par

\noindent\hspace*{1.2em}\kbvlabel{ಉದಾ}\hspace{0.6em}ನೀನೂ ಅವನೂ ಹಣ್ಣುಗಳನ್ನು ತಿಂದಿರಿ.\par

\noindent\hspace*{2.2em}ಅವಳೂ ನೀನೂ ಸೇರಿ ಹೊಯ್ದ ಕಿತ್ತಿರಿ.\par

\kbvexercise{ಅಭ್ಯಾಸ ಪತ್ರಿಕೆ ೧೨}

\vspace{2.5pt plus 3pt minus 0.5pt}\noindent\hspace*{1.2em}\kbvnum{{[}1{]}}\hspace{0.55em}ಉತ್ತಮ ಪುರುಷ ಪುಲ್ಲಿಂಗ ಏಕವಚನದಲ್ಲಿ “ಬರು” ಧಾತು; ವರ್ತಮಾನ ಕಾಲದ ಕ್ರಿಯಾಪದ ರೂಪವನ್ನು ತಿಳಿಸಿರಿ.\par
}
\clearpage

% === tex/0063.tex (Printed Page 61) ===
\renewcommand{\kbvchaptername}{ಕ್ರಿಯಾಪದ-ಲಿಂಗ, ವಚನ, ಕಾಲ}
\setcounter{page}{61}
\thispagestyle{kbvpage}
\kbvtypesetpage{\vspace{2.5pt plus 3pt minus 0.5pt}\noindent\hspace*{1.2em}\kbvnum{(2)}\hspace{0.55em}ಈ ವಾಕ್ಯಗಳಲ್ಲಿನ ಕ್ರಿಯಾಪದಗಳ ಪುರುಷ, ಲಿಂಗ, ವಚನ, ಕಾಲಗಳನ್ನು ತಿಳಿಸಿರಿ.\par

\vspace{2.5pt plus 3pt minus 0.5pt}\noindent\hspace*{1.2em}\kbvnum{{[}ಅ{]}}\hspace{0.55em}ಬ್ಯಾಂಕರನು ಹಣವನ್ನು ಸಾಲವಾಗಿ ಕೊಡುತ್ತಾನೆ.\par

\vspace{2.5pt plus 3pt minus 0.5pt}\noindent\hspace*{1.2em}\kbvnum{{[}ಬ{]}}\hspace{0.55em}ನಾವು ಈ ದಿನ ಊರಿನಿಂದ ಬಂದೆವು.\par

\vspace{2.5pt plus 3pt minus 0.5pt}\noindent\hspace*{1.2em}\kbvnum{{[}ಕ{]}}\hspace{0.55em}ಬೆಕ್ಕು ಇಲಿಯನ್ನು ತಿಂದಿತು.\par

\vspace{2.5pt plus 3pt minus 0.5pt}\noindent\hspace*{1.2em}\kbvnum{{[}ಡ{]}}\hspace{0.55em}ಗೌರಿಯು ಜಾತ್ರೆಗೆ ಹೋಗುವಳು.\par

\vspace{2.5pt plus 3pt minus 0.5pt}\noindent\hspace*{1.2em}\kbvnum{(3)}\hspace{0.55em}“ಜನರು ಭಕ್ತಿಯಿಂದ ಗಣಪತಿಯನ್ನು ಪೂಜೆಮಾಡುತ್ತಾರೆ.”\par

\indent — ಈ ವಾಕ್ಯದಲ್ಲಿ “ಪೂಜೆ” ಎಂಬ ನಾಮಪ್ರಕೃತಿಯನ್ನು ಕ್ರಿಯಾಪದವಾಗಿ ಬಳಸಿರಿ.\par

\vspace{2.5pt plus 3pt minus 0.5pt}\noindent\hspace*{1.2em}\kbvnum{(4)}\hspace{0.55em}ಈ ಕ್ರಿಯಾಪದಗಳನ್ನು ನಿಮ್ಮ ಸ್ವಂತ ವಾಕ್ಯಗಳಲ್ಲಿ ಬಳಸಿರಿ.\par

\indent ಮಾಡೆವು ಜಪಿಸಲಿ ಹೋದಾನು\par

\indent - ಬರೆಯಲಿ - ತಿಂದಿರಿ ಕೊಟ್ಟಿ\par

\vspace{2.5pt plus 3pt minus 0.5pt}\noindent\hspace*{1.2em}\kbvnum{(5)}\hspace{0.55em}“ಏಳು” ಧಾತುವಿನ ಭೂತಕಾಲ ಕ್ರಿಯಾಪದ ರೂಪಗಳನ್ನು ತಿಳಿಸಿರಿ.\par

\vspace{2.5pt plus 3pt minus 0.5pt}\noindent\hspace*{1.2em}\kbvnum{(6)}\hspace{0.55em}“ಗೋಪಾಲ ರೆಡ್ಡಿಯು ನಾಳೆಯ ದಿನ ಹಣವನ್ನು ಕೊಡುತ್ತಾನೆ.”—ಈ ವಾಕ್ಯದಲ್ಲಿರುವ ಕ್ರಿಯಾಪದದ ವಿಶೇಷವೇನು?\par

\vspace{2.5pt plus 3pt minus 0.5pt}\noindent\hspace*{1.2em}\kbvnum{(7)}\hspace{0.55em}ಕಾಲಸೂಚಕ ಪ್ರತ್ಯಯಗಳು ಯಾವುವು? ಉದಾಹರಣೆಗಳ ಸಹಿತ ತಿಳಿಸಿ.\par

\vspace{2.5pt plus 3pt minus 0.5pt}\noindent\hspace*{1.2em}\kbvnum{(8)}\hspace{0.55em}ವಿಧ್ಯರ್ಥಕ ಕ್ರಿಯಾಪದ ಎಂದರೇನು? ಒಂದು ಉದಾಹರಣೆ ಸಹಿತ ತಿಳಿಸಿ.\par

\vspace{2.5pt plus 3pt minus 0.5pt}\noindent\hspace*{1.2em}\kbvnum{(9)}\hspace{0.55em}ನಿಷೇಧಾರ್ಥಕ ಕ್ರಿಯಾಪದ ಎಂದರೇನು? ಉದಾಹರಣೆಗಳ ಸಹಿತ ತಿಳಿಸಿ.\par

\vspace{2.5pt plus 3pt minus 0.5pt}\noindent\hspace*{1.2em}\kbvnum{(10)}\hspace{0.55em}“ತರು” ಧಾತುವಿನ ಸಂಭಾವನಾರ್ಥಕ ಕ್ರಿಯಾಪದಗಳ ಏಕವಚನ ರೂಪಗಳನ್ನು ಬರೆದು ತಿಳಿಸಿ.\par

\vspace{2.5pt plus 3pt minus 0.5pt}\noindent\hspace*{1.2em}\kbvnum{(11)}\hspace{0.55em}“ಮತಿ - ಉಂಟು, ಗತಿ - ಇಲ್ಲ”-ಇಲ್ಲಿರುವ ಕ್ರಿಯಾಪದಗಳನ್ನು ಗುರ್ತಿಸಿ, ಅವುಗಳ ವೈಶಿಷ್ಟ್ಯ ನಿರೂಪಣೆಯನ್ನು ತಿಳಿಸಿರಿ.\par
}
\clearpage

% === tex/0064.tex (Printed Page 62) ===
\renewcommand{\kbvchaptername}{ಸಮಾಸಗಳು}
\setcounter{page}{62}
\thispagestyle{kbvpage}
\kbvtypesetpage{\vspace{2.5pt plus 3pt minus 0.5pt}\noindent\hspace*{1.2em}\kbvnum{(12)}\hspace{0.55em}ಕರ್ಮಣಿ ಪ್ರಯೋಗದಲ್ಲಿ ಧಾತುವಿನ ಮುಂದೆ ಬರುವ ಪ್ರತ್ಯಯ ಯಾವುದು?\par

\vspace{2.5pt plus 3pt minus 0.5pt}\noindent\hspace*{1.2em}\kbvnum{(13)}\hspace{0.55em}“ಬಾಗಿಲನ್ನು ತೆರೆ,”-ಈ ವಾಕ್ಯವನ್ನು ಕರ್ಮಣಿ ಪ್ರಯೋಗ ದಲ್ಲಿ ಬರೆಯಿರಿ.\par

\vspace{2.5pt plus 3pt minus 0.5pt}\noindent\hspace*{1.2em}\kbvnum{(14)}\hspace{0.55em}“ರಾಮಾಯಣವು ವಾಲ್ಮೀಕಿಯಿಂದ ರಚಿಸಲ್ಪಟ್ಟಿತು.” ಈ ವಾಕ್ಯದ ಕರ್ತರಿ ಪ್ರಯೋಗವನ್ನು ತಿಳಿಸಿ.\par

\vspace{2.5pt plus 3pt minus 0.5pt}\noindent\hspace*{1.2em}\kbvnum{(15)}\hspace{0.55em}“ಬರು” ಧಾತುವಿನ ಭೂತಕಾಲದ ರೂಪಗಳನ್ನು, ಕೆಳಗೆ ಕೊಟ್ಟಿರುವ ವಾಕ್ಯಗಳಲ್ಲಿ ತೆರವಿನ ಎಡೆ ತುಂಬಿರಿ.\par

\vspace{2.5pt plus 3pt minus 0.5pt}\noindent\hspace*{1.2em}\kbvnum{{[}ಅ{]}}\hspace{0.55em}ನೀನೂ ನಾನು ಶಾಲೆಯಿಂದ……\par

\vspace{2.5pt plus 3pt minus 0.5pt}\noindent\hspace*{1.2em}\kbvnum{{[}ಬ{]}}\hspace{0.55em}ನೀನೂ ಅವನೂ ಶಾಲೆಯಿಂದ……\par

\kbvchapter{ಪಾಠ ೧೩}{ಸಮಾಸಗಳು}

\indent “ಅರಸರ ಮನೆ ಸುಂದರವಾಗಿದೆ.”- ಈ ವಾಕ್ಯದಲ್ಲಿ ದಪ್ಪ ಅಕ್ಷರಗಳಲ್ಲಿರುವ ಎರಡು ಪದಗಳಿಗೆ ಬದಲಾಗಿ ಅರಮನೆ ಎಂಬ ಒಂದೇ ಪದವನ್ನು ಬಳಸಬಹುದು. ರೂಢಿಯಲ್ಲಿ ಅರಸರ ಮನೆಯನ್ನು ಅರಮನೆಯೆಂತಲೇ ಕರೆಯುವರು. ಆದುದ ರಿಂದ, “ಅರಸರ ಮನೆ ಸುಂದರವಾಗಿದೆ” ಎಂದು ಹೇಳುವ ಬದಲು “ಅರಮನೆ ಸುಂದರವಾಗಿದೆ” ಎಂದು ಹೇಳಿದರೆ ತಪ್ಪಾಗಲಾರದು. ಅರಸರ, ಮನೆ-ಇವೆರಡು ಪದಗಳನ್ನೂ ಸೇರಿಸಿ ಒಂದೇ ಪದ ಮಾಡಿ ಹೇಳುವುದರಿಂದ ಕಾಲ ಉಳಿತಾಯದ ಜೊತೆಗೆ ಶ್ರಮವೂ ಕಡಿಮೆಯಾಗುವುದು.\par

\indent ಹೀಗೆ, ಎರಡು ಅಥವಾ ಹೆಚ್ಚು ಪದಗಳು ಅರ್ಥ ಬದಲಾವಣೆಯಾಗದ ರೀತಿಯಲ್ಲಿ ಒಂದಾಗಿ ಸೇರಿ ಒಂದು ಪದವಾಗುವುದನ್ನು ಸಮಾಸವೆಂದು ಹೇಳುತ್ತೇವೆ ; ಪದಗಳು ಒಂದಾಗಿ ಸೇರುವಾಗ ಅವುಗಳ ಮಧ್ಯದಲ್ಲಿರುವ ವಿಭಕ್ತಿ ಪ್ರತ್ಯಯ ಲೋಪವಾಗುವುದೆಂಬುದು ಗಮನಿಸಬೇಕಾದ ಅಂಶ.\par

\indent ಉದಾಹರಣೆಗೆ, ತಾವರೆಯ + ಕೊಳ = ತಾವರೆಕೊಳ ಇಲ್ಲಿ “ತಾವರೆಯ” ಎಂಬುದು ಪೂರ್ವಪದವಾಗಿರುವುದು. ಮತ್ತು “ಕೊಳ” ಎಂಬುದು ಉತ್ತರಪದವಾಗಿರುವುದು. ಪೂರ್ವ ಮತ್ತು ಉತ್ತರ\kbvcontinues
}
\clearpage

% === tex/0065.tex (Printed Page 63) ===
\renewcommand{\kbvchaptername}{ಸಮಾಸಗಳು}
\setcounter{page}{63}
\thispagestyle{kbvpage}
\kbvtypesetpage{\kbvcontinued ಪದಗಳು ಸೇರಿ ಸಮಸ್ತಪದವಾಗುವುದು. ಈ ಉದಾಹರಣೆಯಲ್ಲಿ “ತಾವರೆಕೊಳ” ಎಂಬುದು ಸಮಸ್ತಪದ. ಈ ಸಮಸ್ತಪದವೇ ಸಮಾಸ ಪದ.\par

\indent ಪೂರ್ವ ಪದ + ಉತ್ತರ ಪದ = ಸಮಸ್ತ ಪದ (ಸಮಾಸ)\par

\indent ಮೇಲಿನ ಉದಾಹರಣೆಯಲ್ಲಿ ಪೂರ್ವಪದವು ಷಷ್ಠೀ ವಿಭಕ್ತಿ ರೂಪದಲ್ಲಿದೆ; ಆದರೆ ಸಮಸ್ತಪದದಲ್ಲಿ ಈ ರೂಪ ಉಳಿದಿಲ್ಲ. ಆದುದರಿಂದ ಸಮಾಸವಾದಾಗ ವಿಭಕ್ತಿ ಪ್ರತ್ಯಯ ಲೋಪವಾದಂತಾಯಿತು.\par

\begin{kbvrule}
\noindent\kbvbullet\hspace{0.5em}ಸಮಸ್ತಪದವನ್ನು ಪೂರ್ವ ಮತ್ತು ಉತ್ತರಪದಗಳಾಗಿ ಬಿಡಿಸಿ ಬರೆಯುವ ಕ್ರಮಕ್ಕೆ ವಿಗ್ರಹ ವಾಕ್ಯ ಎಂದು ಹೆಸರು.
\end{kbvrule}

\vspace{3pt plus 2pt}\noindent\hspace*{1.2em}\kbvlabel{ಉದಾ}\par\vspace{1pt}

\begin{kbvderiv}
\vspace{4pt plus 3pt minus 1pt}
\begin{center}\begin{adjustbox}{max width=\textwidth}
{\renewcommand{\arraystretch}{1.22}\setlength{\tabcolsep}{4pt}
\begin{tabular}{l c l}
ಸಮಸ್ತಪದ & {\color{kbvgrey}$=$} & ಪೂರ್ವಪದ + ಉತ್ತರಪದ \\
ಮುಳ್ಳುಹಾಸಿಗೆ & {\color{kbvgrey}$=$} & ಮುಳ್ಳಿನ + ಹಾಸಿಗೆ \\
ದೊಡ್ಡತಲೆ & {\color{kbvgrey}$=$} & ದೊಡ್ಡದಾದ + ತಲೆ \\
ಹೊಸಗನ್ನಡ & {\color{kbvgrey}$=$} & ಹೊಸದು + ಕನ್ನಡ \\
\end{tabular}}
\end{adjustbox}\end{center}
\vspace{4pt plus 3pt minus 1pt}
\end{kbvderiv}

\begin{kbvrule}
\noindent\kbvbullet\hspace{0.5em}ಸಮಾಸಗಳಲ್ಲಿ ಹಲವು ಪ್ರಕಾರಗಳುಂಟು. ಅವುಗಳು ಯಾವುವೆಂಬುದನ್ನು ತಿಳಿಯೋಣ.
\end{kbvrule}

\vspace{4pt plus 4pt minus 1pt}\noindent{\kbvsecfont\fontsize{11pt}{14.5pt}\selectfont\color{kbvaccentdark} (೧)\hspace{0.45em}ತತ್ಪುರುಷ ಸಮಾಸ :}\hspace{0.5em}ಎರಡು ನಾಮಪದಗಳು ಸೇರಿ ಸಮಾಸವಾದಾಗ ಉತ್ತರ ಪದದ ಅರ್ಥವೇ ಪ್ರಧಾನವಾಗಿ ಉಳ್ಳ ಸಮಾಸವು ತತ್ಪುರುಷ ಸಮಾಸವೆನಿಸುವುದು. ಉದಾಹರಣೆಗೆ, “ತಲೆನೋವು” ಎಂಬ ಸಮಸ್ತಪದವನ್ನು ವಿಗ್ರಹವಾಕ್ಯ ಮಾಡಿನೋಡೋಣ.\par

\indent ತಲೆಯಲ್ಲಿ + ನೋವು = ತಲೆನೋವು ಇಲ್ಲಿ ಪೂರ್ವಪದ ಮತ್ತು ಉತ್ತರಪದಗಳಲ್ಲಿ ಉತ್ತರಪದದ ಅರ್ಥವೇ ಪ್ರಧಾನವಾಗಿದೆ; ಅಂದರೆ, ದೇಹದಲ್ಲಾಗುವ ನೋವು ಪ್ರಧಾನವೇ ಹೊರತು. ದೇಹದ ಯಾವ ಭಾಗದಲ್ಲಿ ನೋವಾಗುತ್ತಿರುವುದೆಂಬುದು ಅಷ್ಟು ಪ್ರಧಾನವಲ್ಲ, ಆದುದರಿಂದ, ಇಲ್ಲಿನ ಸಮಾಸವು ತತ್ಪುರುಷ ಸಮಾಸವಾಗಿರುವುದು.\par

\begin{kbvrule}
\noindent\kbvbullet\hspace{0.5em}ತತ್ಪುರುಷ ಸಮಾಸದಲ್ಲಿ ಪೂರ್ವಪದವು ತೃತೀಯಾ ವಿಭಕ್ತಿಯಿಂದ ಮೊದಲ್ಗೊಂಡು ಸಪ್ತಮೀ ವಿಭಕ್ತಿಯವರೆಗೆ ಯಾವುದಾದರೊಂದು
\end{kbvrule}
}
\clearpage

% === tex/0066.tex (Printed Page 64) ===
\renewcommand{\kbvchaptername}{ಸಮಾಸಗಳು}
\setcounter{page}{64}
\thispagestyle{kbvpage}
\kbvtypesetpage{\indent ವಿಭಕ್ತ್ಯಂತವಾಗಿರುವುದು ; ಹಾಗೆ ಪೂರ್ವಪದವು ಯಾವ ವಿಭಕ್ತಿಯಲ್ಲಿದೆ ಯೋ ಆ ವಿಭಕ್ತಿಯ ಹೆಸರಿನೊಂದಿಗೆ ಸಮಾಸ ಹೇಳುವುದು ರೂಢಿಯಾಗಿದೆ. ಮೇಲಿನ ಉದಾಹರಣೆಯಲ್ಲಿ “ತಲೆಯಲ್ಲಿ” ಎಂಬ ಪದವು ಸಪ್ತಮೀ ವಿಭಕ್ತಿಯಲ್ಲಿದೆ. ಆದುದರಿಂದ “ತಲೆನೋವು” ಎಂಬ ಶಬ್ದವು ಸಪ್ತಮೀ ತತ್ಪುರುಷ ಸಮಾಸ (ಸ.ತ.ಸ.) ವಾಗಿರುವುದು.\par

\vspace{3pt plus 2pt}\noindent\hspace*{1.2em}\kbvlabel{ಉದಾ}\par\vspace{1pt}

\vspace{5pt plus 4pt minus 1pt}
\noindent{\renewcommand{\arraystretch}{1.24}\setlength{\tabcolsep}{4pt}\rowcolors{1}{kbvrow}{white}
\begin{tabularx}{\textwidth}{X X X X X X}
ಮುಳ್ಳಿನ &  +  & ಹಾಸಿಗೆ &  =  & ಮುಳ್ಳುಹಾಸಿಗೆ & (ಸ.ತ.ಸ.) \\
ತಲೆಯ &  +  & ಕೂದಲು &  =  & ತಲೆಗೂದಲು & (ಸ.ತ.ಸ.) \\
ಮಳೆಯ &  +  & ಕಾಲ &  =  & ಮಳೆಗಾಲ & (ಸ.ತ.ಸ.) \\
ದೇವರ &  +  & ಮಂದಿರ &  =  & ದೇವಮಂದಿರ & (ಷ.ತ.ಸ.) \\
ವಯಸ್ಸಿನಿಂದ &  +  & ವೃದ್ಧ &  =  & ವಯೋವೃದ್ಧ & (ತೃ.ತ.ಸ.) \\
ವ್ಯಾಘ್ರದ &  +  & ದೆಸೆಯಿಂದ &  =  & ವ್ಯಾಘ್ರಭಯ & (ಪಂ.ತ.ಸ.) \\
ನರರಲ್ಲಿ &  +  & ಉತ್ತಮ &  =  & ನರೋತ್ತಮ & (ಸ.ತ.ಸ.) \\
\end{tabularx}}
\vspace{5pt plus 4pt minus 1pt}

\vspace{4pt plus 4pt minus 1pt}\noindent{\kbvsecfont\fontsize{11pt}{14.5pt}\selectfont\color{kbvaccentdark} (೨)\hspace{0.45em}ಕರ್ಮಧಾರಯ ಸಮಾಸ :}\hspace{0.5em}ಈ ವಿಗ್ರಹ ವಾಕ್ಯಗಳನ್ನು ನೋಡೋಣ.\par

\vspace{5pt plus 4pt minus 1pt}
\noindent{\renewcommand{\arraystretch}{1.24}\setlength{\tabcolsep}{4pt}
\begin{tabularx}{\textwidth}{X X X X X X}
ದೊಡ್ಡವನು &  +  & ಅಪ್ಪನು &  =  & ದೊಡ್ಡಪ್ಪನು & (ಪು.ಬ.) \\
ಹೆತ್ತವಳು &  +  & ಅಮ್ಮನು &  =  & ಹೆತ್ತಮ್ಮನು & (ಸ್ತ್ರೀ.ಬ.) \\
ಚಿಕ್ಕವು &  +  & ಮಕ್ಕಳು &  =  & ಚಿಕ್ಕಮಕ್ಕಳು & (ನ.ಬ.) \\
\end{tabularx}}
\vspace{5pt plus 4pt minus 1pt}

\indent ಈ ಮೂರು ಉದಾಹರಣೆಗಳಲ್ಲಿಯೂ ಪೂರ್ವ ಮತ್ತು ಉತ್ತರಪದ ಗಳಲ್ಲಿ ಲಿಂಗ ಭೇದವಾಗಲೀ ವಚನ ಭೇದವಾಗಲೀ ವಿಭಕ್ತಿರೂಪ ಭೇದವಾಗಲೀ ಕಂಡುಬರುವುದಿಲ್ಲ. ಅಪ್ಪನು, ಅಮ್ಮನು, ಮಕ್ಕಳು ಉತ್ತರಪದಗಳಾಗಿದ್ದು\par

\indent ಅವುಗಳ ಅರ್ಥವೇ ಪ್ರಧಾನವಾಗಿರುವುದು. ಅದೇ ಅಲ್ಲದೇ ದೊಡ್ಡವನು, ಹೆತ್ತವಳು, ಚಿಕ್ಕವು ಎಂಬ ಪದಗಳು ಕ್ರಮವಾಗಿ ಅಪ್ಪನು ಅಮ್ಮನು, ಮಕ್ಕಳು ಎಂಬ ಪದಗಳಿಗೆ ವಿಶೇಷಣಗಳಾಗಿವೆ. ಅಂದರೆ, ಮೇಲೆ ಕೊಟ್ಟಿರುವ ಪ್ರತಿಯೊಂದು ಉದಾಹರಣೆಯಲ್ಲಿ ಪೂರ್ವಪದವು ವಿಶೇಷಣ ಮತ್ತು ಉತ್ತರ\par

\indent ಪದವು ವಿಶೇಷ್ಯ ಆಗಿರುವುದು. ಇವು ಕರ್ಮಧಾರಯ ಸಮಾಸದ ಲಕ್ಷಣಗಳು.\par
}
\clearpage

% === tex/0067.tex (Printed Page 65) ===
\renewcommand{\kbvchaptername}{ಸಮಾಸಗಳು}
\setcounter{page}{65}
\thispagestyle{kbvpage}
\kbvtypesetpage{\begin{kbvrule}
\noindent{\kbvsecfont ವ್ಯಾಖ್ಯೆ :}\hspace{0.5em}ಉತ್ತರಪದದ ಅರ್ಥವೇ ಪ್ರಧಾನವಾಗಿದ್ದು, ಪೂರ್ವ ಹಾಗೂ ಉತ್ತರಪದಗಳು ಲಿಂಗ ವಚನ ವಿಭಕ್ತಿಗಳಿಂದ ಸಮಾನವಾಗಿದ್ದು, ಅವೆರಡೂ ಪದಗಳು ವಿಶೇಷಣ ವಿಶೇಷ್ಯ ಸಂಬಂಧದಿಂದ ಕೂಡಿ ಆಗುವ ಸಮಾಸಕ್ಕೆ ಕರ್ಮಧಾರಯ ಸಮಾಸ ಎಂದು ಹೆಸರು.
\end{kbvrule}

\vspace{3pt plus 2pt}\noindent\hspace*{1.2em}\kbvlabel{ಉದಾ}\par\vspace{1pt}

\begin{kbvderiv}
\vspace{4pt plus 3pt minus 1pt}
\begin{center}\begin{adjustbox}{max width=\textwidth}
{\renewcommand{\arraystretch}{1.22}\setlength{\tabcolsep}{4pt}
\begin{tabular}{l c l}
ಹಿರಿದು + ಬಾಗಿಲು & {\color{kbvgrey}$=$} & \kbvresult{ಹೆಬ್ಬಾಗಿಲು}{\color{kbvgrey} {[}ವಿ.ಪೂ.ಕ{]}} \\
ಹಿರಿದು + ಜೇನು & {\color{kbvgrey}$=$} & \kbvresult{ಹೆಜ್ಜೇನು}{\color{kbvgrey} {[}ವಿ.ಪೂ.ಕ{]}} \\
ತಣ್ಣನೆ + ಗಾಳಿ & {\color{kbvgrey}$=$} & \kbvresult{ತಂಗಾಳಿ}{\color{kbvgrey} {[}ವಿ.ಪೂ.ಕ{]}} \\
ಸಿಹಿಯಾದ + ನೀರು & {\color{kbvgrey}$=$} & \kbvresult{ಸಿಹಿನೀರು}{\color{kbvgrey} {[}ವಿ.ಪೂ.ಕ{]}} \\
ಹಣ್ಣಾದ + ಎಲೆ & {\color{kbvgrey}$=$} & \kbvresult{ಹಣ್ಣೆಲೆ}{\color{kbvgrey} {[}ವಿ.ಪೂ.ಕ{]}} \\
ಎಳದು + ಕರು & {\color{kbvgrey}$=$} & \kbvresult{ಎಳಗರು}{\color{kbvgrey} {[}ವಿ.ಪೂ.ಕ{]}} \\
ಕಿರಿದು + ದಾರಿ & {\color{kbvgrey}$=$} & \kbvresult{ಕಿರುದಾರಿ}{\color{kbvgrey} {[}ವಿ.ಪೂ.ಕ{]}} \\
ಮೆಲ್ಲಿತು + ನುಡಿ & {\color{kbvgrey}$=$} & \kbvresult{ಮೆಲ್ನುಡಿ}{\color{kbvgrey} {[}ವಿ.ಪೂ.ಕ{]}} \\
ಕೆಚ್ಚನೆ + ಚೇಳು & {\color{kbvgrey}$=$} & \kbvresult{ಕೆಂಚೇಳು}{\color{kbvgrey} {[}ವಿ.ಪೂ.ಕ{]}} \\
ಇನಿದು + ಸರ & {\color{kbvgrey}$=$} & \kbvresult{ಇಂಚರ}{\color{kbvgrey} {[}ವಿ.ಪೂ.ಕ{]}} \\
ನೀಲವಾದ + ಆಕಾಶ & {\color{kbvgrey}$=$} & \kbvresult{ನೀಲಾಕಾಶ}{\color{kbvgrey} {[}ವಿ.ಪೂ.ಕ{]}} \\
ಶ್ವೇತವಾದ + ಅಂಬರ & {\color{kbvgrey}$=$} & \kbvresult{ಶ್ವೇತಾಂಬರ}{\color{kbvgrey} {[}ವಿ.ಪೂ.ಕ{]}} \\
\end{tabular}}
\end{adjustbox}\end{center}
\vspace{4pt plus 3pt minus 1pt}
\end{kbvderiv}

\indent {[}ಸೂಚನೆ : ಈ ಸಮಾಸಗಳಲ್ಲಿ ಪೂರ್ವಪದಗಳು ವಿಶೇಷಣಗಳಾಗಿಯೂ ಉತ್ತರಪದಗಳು ವಿಶೇಷ್ಯಗಳಾಗಿಯೂ ಇರುವುದರಿಂದ, ಇವುಗಳಿಗೆ ವಿಶೇಷಣ ಪೂರ್ವಪದ ಕರ್ಮಧಾರಯ ಸಮಾಸ {[}ವಿ.ಪೂ.ಕ{]} ವೆಂದು ಹೇಳುವುದು ವಾಡಿಕೆ.{]}\par

\kbvsechead{(೩)\hspace{0.45em}ದ್ವಿಗು ಸಮಾಸ :}

\indent - ಮೂರು {[}ಪೂ.ಪ{]} + ಲೋಕ {[}ಉ.ಪ{]} = ಮೂಲೋಕ ಇಲ್ಲಿ “ಮೂರು” ಎಂಬ ಸಂಖ್ಯಾವಾಚಕವು “ಲೋಕ” ಎಂಬ ನಾಮಪದ ದೊಂದಿಗೆ ಸೇರಿ ಸಮಸ್ತಪದವಾಗಿರುವುದು - ಇದು ದ್ವಿಗು ಸಮಾಸ ಎನಿಸುವುದು.\par
}
\clearpage

% === tex/0068.tex (Printed Page 66) ===
\renewcommand{\kbvchaptername}{ಸಮಾಸಗಳು}
\setcounter{page}{66}
\thispagestyle{kbvpage}
\kbvtypesetpage{\begin{kbvrule}
\noindent{\kbvsecfont ವ್ಯಾಖ್ಯೆ :}\hspace{0.5em}ಪೂರ್ವಪದವು ಸಂಖ್ಯಾ ವಾಚಕವೂ ಉತ್ತರಪದವೂ ನಾಮಪದವೂ ಆಗಿರುವಾಗ, ಪೂರ್ವಪದವು ಉತ್ತರಪದದೊಂದಿಗೆ ಸೇರಿ ಆಗುವ ಸಮಾಸಕ್ಕೆ ದ್ವಿಗು ಸಮಾಸವೆಂದು ಹೆಸರು.
\end{kbvrule}

\begin{kbvderiv}
\vspace{4pt plus 3pt minus 1pt}
\begin{center}\begin{adjustbox}{max width=\textwidth}
{\renewcommand{\arraystretch}{1.22}\setlength{\tabcolsep}{4pt}
\begin{tabular}{l c l}
ಉದಾ : ನಾಲ್ಕು + ಮಡಿ & {\color{kbvgrey}$=$} & \kbvresult{ನಾಲ್ವಡಿ}{\color{kbvgrey} {[}ದ್ವಿ.ಸ.{]}} \\
ಎಂಟು + ದೆಸೆ & {\color{kbvgrey}$=$} & \kbvresult{ಎಣ್ದೆಸೆ}{\color{kbvgrey} {[}ದ್ವಿ.ಸ.{]}} \\
ಎರಡು + ಕೆಲ & {\color{kbvgrey}$=$} & \kbvresult{ಇರ್ಕೆಲ}{\color{kbvgrey} {[}ದ್ವಿ.ಸ.{]}} \\
ಒಂದು + ಕಟ್ಟು & {\color{kbvgrey}$=$} & \kbvresult{ಒಗ್ಗಟ್ಟು}{\color{kbvgrey} {[}ದ್ವಿ.ಸ.{]}} \\
ಒಂದು + ಕಣ್ಣು & {\color{kbvgrey}$=$} & \kbvresult{ಒಕ್ಕಣ್ಣು}{\color{kbvgrey} {[}ದ್ವಿ.ಸ.{]}} \\
\end{tabular}}
\end{adjustbox}\end{center}
\vspace{4pt plus 3pt minus 1pt}
\end{kbvderiv}

\vspace{4pt plus 4pt minus 1pt}\noindent{\kbvsecfont\fontsize{11pt}{14.5pt}\selectfont\color{kbvaccentdark} (೪)\hspace{0.45em}ಅಂಶಿ ಸಮಾಸ :}\hspace{0.5em}“ತುದಿಮೂಗು” ಎಂಬ ಸಮಸ್ತ ಪದವನ್ನು ವಿಗ್ರಹ ವಾಕ್ಯ ಮಾಡಿ\par

\indent ನೋಡೋಣ. ಮೂಗಿನ + ತುದಿ = ತುದಿಮೂಗು ಇಲ್ಲಿ ಪ್ರಧಾನವಾಗಿರುವುದು ಮೂಗಿನ ತುದಿಯೇ ಹೊರತು ಮೂಗಲ್ಲ “ತುದಿ” ಎಂಬುದು ಮೂಗಿನ ಒಂದು ಅಂಶ ಮಾತ್ರ. ಇಲ್ಲಿ ಮೂಗು ಅಂಶಿಯಾಗಿರುವುದು-\par

\indent - ಬ ಖಿ ಪ್ರ ವ ವಿಗ್ರಹ ವಾಕ್ಯದಲ್ಲಿ “ಮೂಗಿನ” ಎಂಬ ಪದ ಪೂರ್ವಪದವಾಗಿದ್ದರೂ, ಪೂರ್ವಪದ ಮತ್ತು ಉತ್ತರಪದಗಳು ಸೇರಿ ಸಮಾಸವಾದಾಗ ಸಮಸ್ತ ಪದದಲ್ಲಿ ಪೂರ್ವ ಮತ್ತು ಉತ್ತರಪದಗಳು ಅದಲು ಬದಲಾಗಿರುವುದನ್ನು ಕಾಣುತ್ತೇವೆ. ಆದುದರಿಂದ ಸಮಸ್ತಪದದಲ್ಲಿನ ಪೂರ್ವಪದವು ಪ್ರಧಾನವಾಗುವುದು. ಹೀಗೆ ಪೂರ್ವಪದದ ಅರ್ಥ ಪ್ರಧಾನವಾಗಿ ಉಳ್ಳ ಸಮಾಸಕ್ಕೆ ಅಂಶಿಸಮಾಸವೆಂದು ಹೆಸರು.\par

\begin{kbvrule}
\noindent{\kbvsecfont ವ್ಯಾಖ್ಯೆ :}\hspace{0.5em}ಪೂರ್ವ ಮತ್ತು ಉತ್ತರಪದಗಳು ಅಂಶಾಂಶಿಭಾವ ಸಂಬಂಧದಿಂದ ಸೇರಿ ಸಮಾಸವಾದಾಗ, ಪೂರ್ವಪದದ ಅರ್ಥವು ಪ್ರಧಾನ ವಾಗಿದ್ದರೂ ಅದು ಅಂಶಿ ಸಮಾಸವೆನಿಸುವುದು.
\end{kbvrule}

\indent ಅಂಶ--ಪೂರ್ಣದ ಒಂದು ಭಾಗ ; ಅಂಶಿ-ಅಂಶಗಳನ್ನು ಒಳಗೊಂಡಿರುವುದು.\par

\noindent\hspace*{2.2em}ಅಂಶಾಂಶಿಭಾವ--ಅಂಶಕ್ಕೂ ಅಂಶಿಗೂ ಇರುವ ಸಂಬಂಧ.\par
}
\clearpage

% === tex/0069.tex (Printed Page 67) ===
\renewcommand{\kbvchaptername}{ಸಮಾಸಗಳು}
\setcounter{page}{67}
\thispagestyle{kbvpage}
\kbvtypesetpage{\indent \textbf{ಉದಾ :}\par

\begin{kbvderiv}
\vspace{4pt plus 3pt minus 1pt}
\begin{center}\begin{adjustbox}{max width=\textwidth}
{\renewcommand{\arraystretch}{1.22}\setlength{\tabcolsep}{4pt}
\begin{tabular}{l c l}
ತಲೆಯ + ಮುಂದು & {\color{kbvgrey}$=$} & \kbvresult{ಮುಂದಲೆ} \\
ಕೈಯ + ಅಡಿ & {\color{kbvgrey}$=$} & \kbvresult{ಅಂಗೈ} \\
ಕಾಲು + ಅಡಿ & {\color{kbvgrey}$=$} & \kbvresult{ಅಂಗಾಲು} \\
ಕಣ್ಣ + ಕಡೆ & {\color{kbvgrey}$=$} & \kbvresult{ಕಡೆಗಣ್ಣು} \\
ಮೈಯ + ಒಳಗು & {\color{kbvgrey}$=$} & \kbvresult{ಒಳಮೈ} \\
ಇರುಳಿನ + ನಡು & {\color{kbvgrey}$=$} & \kbvresult{ನಟ್ಟಿರುಳು} \\
ತುಟಿಯ + ಕೆಳಗು & {\color{kbvgrey}$=$} & \kbvresult{ಕೆಳದುಟಿ} \\
\end{tabular}}
\end{adjustbox}\end{center}
\vspace{4pt plus 3pt minus 1pt}
\end{kbvderiv}

\vspace{4pt plus 4pt minus 1pt}\noindent{\kbvsecfont\fontsize{11pt}{14.5pt}\selectfont\color{kbvaccentdark} (೫)\hspace{0.45em}ದ್ವಂದ್ವ ಸಮಾಸ :}\hspace{0.5em}“ಮಳೆ ಬಂದು ಕೆರೆಕಟ್ಟೆ ತುಂಬಿತು.”— ಈ ವಾಕ್ಯದಲ್ಲಿ ದಪ್ಪಕ್ಷರಗಳಲ್ಲಿರುವ “ಕೆರೆಕಟ್ಟೆ” ಎಂಬ ಪದವನ್ನು ಕೆರೆಯೂ ಕಟ್ಟೆಯೂ ಎಂದು ಬಿಡಿಸಿ ಬರೆಯಬಹುದು. ಕೆರೆಯೂ ಕಟ್ಟೆಯೂ – ಕೆರೆಕಟ್ಟೆ ಇಲ್ಲಿ ಕೆರೆಯೂ ಮತ್ತು ಕಟ್ಟೆಯೂ ಒಂದಕ್ಕೊಂದು ಪೂರಕವಾಗಿರುವ ಅಥವಾ ಸಹಯೋಗ ಹೊಂದಿರುವ ಎರಡು ಪದಗಳು. ಇವುಗಳಲ್ಲಿ ಯಾವು ದೊಂದು ಪದದ ಅರ್ಥವು ಗೌಣವಲ್ಲ, ಎರಡು ಪದಗಳ ಅರ್ಥವೂ ಪ್ರಧಾನವೇ ಸರಿ. ಇಂತಹ ಪದಗಳು ಸೇರಿ ಆಗುವ ಸಮಾಸಕ್ಕೆ ದ್ವಂದ್ವ ಸಮಾಸವೆನ್ನುವರು.\par

\begin{kbvrule}
\noindent{\kbvsecfont ವ್ಯಾಖ್ಯೆ :}\hspace{0.5em}- ಸಹಯೋಗವಿರುವ ಎರಡು ಅಥವಾ ಹೆಚ್ಚು ನಾಮ ಪದಗಳು ಒಂದಕ್ಕೊಂದು ಸೇರಿ ಎಲ್ಲ ಪದಗಳ ಅರ್ಥವೂ ಪ್ರಧಾನವಾಗಿಯುಳ್ಳ ಸಮಾಸಕ್ಕೆ ದ್ವಂದ್ವ ಸಮಾಸವೆಂದು ಹೆಸರು.
\end{kbvrule}

\begin{kbvderiv}
\vspace{4pt plus 3pt minus 1pt}
\begin{center}\begin{adjustbox}{max width=\textwidth}
{\renewcommand{\arraystretch}{1.22}\setlength{\tabcolsep}{4pt}
\begin{tabular}{l c l}
ಉದಾ : ತಾಯಿಯೂ-ತಂದೆಯೂ & {\color{kbvgrey}$=$} & \kbvresult{ತಾಯಿತಂದೆ} \\
ಹಣ್ಣನ್ನೂ-ಕಾಯಿಯನ್ನೂ & {\color{kbvgrey}$=$} & \kbvresult{ಹಣ್ಣುಕಾಯಿ} \\
ಹಸುವನ್ನೂ-ಕರುವನ್ನೂ & {\color{kbvgrey}$=$} & \kbvresult{ಹಸುಕರು} \\
\end{tabular}}
\end{adjustbox}\end{center}
\vspace{4pt plus 3pt minus 1pt}
\end{kbvderiv}
}
\clearpage

% === tex/0070.tex (Printed Page 68) ===
\renewcommand{\kbvchaptername}{ಸಮಾಸಗಳು}
\setcounter{page}{68}
\thispagestyle{kbvpage}
\kbvtypesetpage{\indent ಇದೇ ರೀತಿಯಾಗಿ ಅರಿಶಿನ ಕುಂಕುಮ, ಫಲ ಪುಷ್ಪ, ಗಿಡ ಮರ ಬಳ್ಳಿ ಕೈಕಾಲು ಮುಖ, ಸೂರ್ಯ ಚಂದ್ರ ನಕ್ಷತ್ರ, ಗುಡುಗು ಸಿಡಿಲು ಮಿಂಚು ಶಬ್ದ ಸ್ಪರ್ಶ ರೂಪ ರಸ ಗಂಧ-ದ್ವಂದ್ವ ಸಮಾಸ ಪದಗಳು.\par

\kbvsechead{(೬)\hspace{0.45em}ಬಹುವ್ರೀಹಿ ಸಮಾಸ :}

\indent “ಹಣೆಗಣ್ಣ” ಎಂಬ ಪದವನ್ನು ಹಣೆಯಲ್ಲಿ + ಕಣ್ಣುಳ್ಳವನು ಎಂದು\par

\indent “ಫಾಲನೇತ್ರ” ಎಂಬ ಪದವನ್ನು ಫಾಲದಲ್ಲಿ + ನೇತ್ರವುಳ್ಳವನು ಎಂದು ಬಿಡಿಸಿ ಬರೆಯಬಹುದು.\par

\indent ಹಣೆಯಲ್ಲಿ + ಕಣ್ಣುಳ್ಳವನು = ಹಣೆಗಣ್ಣ\par

\indent ಈ ಉದಾಹರಣೆಯಲ್ಲಿ ಪೂರ್ವಪದದ ಅರ್ಥವಾಗಲೀ ಉತ್ತರಪದದ ಅರ್ಥವಾಗಲೀ ಪ್ರಧಾನವಲ್ಲ; ಹಣೆಯಲ್ಲಿ ಕಣ್ಣುಳ್ಳವನು ಯಾವನೋ ಅವನು ಪ್ರಧಾನ. ಇದರಿಂದ ಸಮಾಸದ ಯಾವುದೊಂದು ಪದವೂ ಪ್ರಧಾನವಾಗಿ ಅನ್ಯಪದದ ಅರ್ಥ ಪ್ರಧಾನವಾಗಿರುವುದೆಂದು ತಿಳಿದು ಬರುವುದು. ಲಕ್ಷಣವುಳ್ಳ ಸಮಾಸಕ್ಕೆ ಬಹುವ್ರೀಹಿ ಸಮಾಸವೆಂದು ಹೆಸರು.\par

\begin{kbvrule}
\noindent{\kbvsecfont ವ್ಯಾಖ್ಯೆ :}\hspace{0.5em}ಎರಡು ನಾಮಪದಗಳು ಸೇರಿ ಸಮಾಸವಾದಾಗ ಆ ಪದದ ಅರ್ಥವು ಪ್ರಧಾನವಾಗಿ ಉಳ್ಳ ಸಮಾಸಕ್ಕೆ ಬಹುವ್ರೀಹಿ ಸಮಾಸವೆಂದು.
\end{kbvrule}

\vspace{3pt plus 2pt}\noindent\hspace*{1.2em}\kbvlabel{ಉದಾ}\par\vspace{1pt}

\begin{kbvderiv}
\vspace{4pt plus 3pt minus 1pt}
\begin{center}\begin{adjustbox}{max width=\textwidth}
{\renewcommand{\arraystretch}{1.22}\setlength{\tabcolsep}{4pt}
\begin{tabular}{l c l}
ಒಂದು + ಕಣ್ಣುಳ್ಳವ & {\color{kbvgrey}$=$} & \kbvresult{ಒಕ್ಕಣ್ಣ} \\
ನಾಲ್ಕು + ಮೊಗವುಳ್ಳವ & {\color{kbvgrey}$=$} & \kbvresult{ನಾಲ್ಮೊಗ} \\
ಕೆಂಪು + ಕಣ್ಣುಳ್ಳವ & {\color{kbvgrey}$=$} & \kbvresult{ಕೆಂಗಣ್ಣ} \\
ಚಂದ್ರನಂತೆ + ಮುಖವುಳ್ಳವಳು & {\color{kbvgrey}$=$} & \kbvresult{ಚಂದ್ರಮುಖಿ} \\
ವಜ್ರದಂತೆ + ಶರೀರವುಳ್ಳವ & {\color{kbvgrey}$=$} & \kbvresult{ವಜ್ರಶರೀರಿ} \\
\end{tabular}}
\end{adjustbox}\end{center}
\vspace{4pt plus 3pt minus 1pt}
\end{kbvderiv}

\kbvsechead{(೭)\hspace{0.45em}ಕ್ರಿಯಾ ಸಮಾಸ:}

\begin{kbvrule}
\noindent{\kbvsecfont ವ್ಯಾಖ್ಯೆ :}\hspace{0.5em}ಪೂರ್ವಪದವು ಸಾಮಾನ್ಯವಾಗಿ ದ್ವಿತೀಯಾ ವಿಭಕ್ತಿ ನಾಮಪದವಾಗಿದ್ದು ಕ್ರಿಯಾಸೂಚಕ ಉತ್ತರಪದದೊಂದಿಗೆ ಸೇರಿ ಆಗುವ ಸಮಾಸಕ್ಕೆ ಕ್ರಿಯಾ ಸಮಾಸವೆಂದು ಹೆಸರು.
\end{kbvrule}
}
\clearpage

% === tex/0071.tex (Printed Page 69) ===
\renewcommand{\kbvchaptername}{ಸಮಾಸಗಳು}
\setcounter{page}{69}
\thispagestyle{kbvpage}
\kbvtypesetpage{\vspace{3pt plus 2pt}\noindent\hspace*{1.2em}\kbvlabel{ಉದಾ}\par\vspace{1pt}

\begin{kbvderiv}
\vspace{4pt plus 3pt minus 1pt}
\begin{center}\begin{adjustbox}{max width=\textwidth}
{\renewcommand{\arraystretch}{1.22}\setlength{\tabcolsep}{4pt}
\begin{tabular}{l c l}
ಕೈಯನ್ನು + ಮುಟ್ಟಿ & {\color{kbvgrey}$=$} & \kbvresult{ಕೈಮುಟ್ಟಿ} \\
ಕಣ್ಣನ್ನು + ತೆರೆದು & {\color{kbvgrey}$=$} & \kbvresult{ಕಣ್ಣೆರೆದು}{\color{kbvgrey} (ಕಣ್ಣೆರೆದು)} \\
ಮೈಯನ್ನು + ಮುಚ್ಚಿ & {\color{kbvgrey}$=$} & \kbvresult{ಮೈಮುಚ್ಚಿ} \\
ಕಾಲನ್ನು + ತೆಗೆದು & {\color{kbvgrey}$=$} & \kbvresult{ಕಾಲ್ದೆಗೆದು}{\color{kbvgrey} (ಕಾಲ್ದೆಗೆದು)} \\
\end{tabular}}
\end{adjustbox}\end{center}
\vspace{4pt plus 3pt minus 1pt}
\end{kbvderiv}

\kbvsechead{(8)\hspace{0.45em}ಗಮಕ ಸಮಾಸ :}

\begin{kbvrule}
\noindent{\kbvsecfont ವ್ಯಾಖ್ಯೆ :}\hspace{0.5em}ಪೂರ್ವಪದವು ಸರ್ವನಾಮವೋ ಕೃದಂತವೋ ಆಗಿದ್ದು ಉತ್ತರಪದದಲ್ಲಿರುವ ನಾಮಪದದೊಂದಿಗೆ ಸೇರಿ, ಆಗುವ ಸಮಾಸಕ್ಕೆ ಗಮಕ ಸಮಾಸವೆಂದು ಹೆಸರು.
\end{kbvrule}

\vspace{3pt plus 2pt}\noindent\hspace*{1.2em}\kbvlabel{ಉದಾ}\par\vspace{1pt}

\begin{kbvderiv}
\vspace{4pt plus 3pt minus 1pt}
\begin{center}\begin{adjustbox}{max width=\textwidth}
{\renewcommand{\arraystretch}{1.22}\setlength{\tabcolsep}{4pt}
\begin{tabular}{l c l}
ಅದು + ಬೆಟ್ಟ & {\color{kbvgrey}$=$} & \kbvresult{ಆ ಬೆಟ್ಟ}{\color{kbvgrey} (ಪೂ.ಪ. -- ಸರ್ವನಾಮ)} \\
ಇವನು + ಹುಡುಗ & {\color{kbvgrey}$=$} & \kbvresult{ಈ ಹುಡುಗ}{\color{kbvgrey} (ಪೂ.ಪ. -- ಸರ್ವನಾಮ)} \\
ಅವಳು + ನಟಿ & {\color{kbvgrey}$=$} & \kbvresult{ಆ ನಟಿ}{\color{kbvgrey} (ಪೂ.ಪ. -- ಸರ್ವನಾಮ)} \\
ತಿಂದುದು + ಕೂಳು & {\color{kbvgrey}$=$} & \kbvresult{ತಿಂದಕೂಳು}{\color{kbvgrey} (ಪೂ.ಪ. -- ಕೃದಂತ)} \\
ಬೆಂದುದು + ಬೇಳೆ & {\color{kbvgrey}$=$} & \kbvresult{ಬೆಂದಬೇಳೆ}{\color{kbvgrey} (ಪೂ.ಪ. -- ಕೃದಂತ)} \\
ಮಾಡಿದುದು + ಅಡಿಗೆ & {\color{kbvgrey}$=$} & \kbvresult{ಮಾಡಿದಡಿಗೆ}{\color{kbvgrey} (ಪೂ.ಪ. -- ಕೃದಂತ)} \\
\end{tabular}}
\end{adjustbox}\end{center}
\vspace{4pt plus 3pt minus 1pt}
\end{kbvderiv}

\kbvexercise{ಅಭ್ಯಾಸ ಪತ್ರಿಕೆ ೧೩}

\vspace{2.5pt plus 3pt minus 0.5pt}\noindent\hspace*{1.2em}\kbvnum{(1)}\hspace{0.55em}ಸಂಧಿಗೂ ಸಮಾಸಕ್ಕೂ ಇರುವ ವ್ಯತ್ಯಾಸವೇನು ?\par

\vspace{2.5pt plus 3pt minus 0.5pt}\noindent\hspace*{1.2em}\kbvnum{(2)}\hspace{0.55em}“ಮರಗಾಲು” -- ಈ ಸಂಧಿಪದದಲ್ಲಿ ಸಂಧಿಯನ್ನು ಬಿಡಿಸಿ, ಸಂಧಿಯ ಹೆಸರು ಬರೆಯಿರಿ.\par

\vspace{2.5pt plus 3pt minus 0.5pt}\noindent\hspace*{1.2em}\kbvnum{(3)}\hspace{0.55em}“ಮರಗಾಲು” -- ಈ ಪದದ ವಿಗ್ರಹವಾಕ್ಯ ಬರೆದು, ಸಮಾಸದ ಹೆಸರು ತಿಳಿಸಿ.\par

\vspace{2.5pt plus 3pt minus 0.5pt}\noindent\hspace*{1.2em}\kbvnum{(4)}\hspace{0.55em}ತತ್ಪುರುಷ ಸಮಾಸವೆಂದರೇನು ? ಉದಾಹರಣೆ ಸಹಿತ ತಿಳಿಸಿ.\par

\vspace{2.5pt plus 3pt minus 0.5pt}\noindent\hspace*{1.2em}\kbvnum{(5)}\hspace{0.55em}ಕರ್ಮಧಾರಯ ಸಮಾಸವೆಂದರೇನು ? ಉದಾಹರಣೆ ಸಹಿತ ವಿವರಿಸಿರಿ.\par
}
\clearpage

% === tex/0072.tex (Printed Page 70) ===
\renewcommand{\kbvchaptername}{ಕೃದಂತಗಳು}
\setcounter{page}{70}
\thispagestyle{kbvpage}
\kbvtypesetpage{\vspace{2.5pt plus 3pt minus 0.5pt}\noindent\hspace*{1.2em}\kbvnum{(6)}\hspace{0.55em}ತತ್ಪುರುಷ ಸಮಾಸಕ್ಕೂ ಕರ್ಮಧಾರಯ ಸಮಾಸಕ್ಕೂ ಇರುವ ಮುಖ್ಯ ವ್ಯತ್ಯಾಸವೇನು?\par

\vspace{2.5pt plus 3pt minus 0.5pt}\noindent\hspace*{1.2em}\kbvnum{(7)}\hspace{0.55em}ಅಂಶಿ ಸಮಾಸ ಎಂದರೇನು? ಉದಾಹರಣೆ ಸಹಿತ ವಿವರಿಸಿ.\par

\vspace{2.5pt plus 3pt minus 0.5pt}\noindent\hspace*{1.2em}\kbvnum{(8)}\hspace{0.55em}ದ್ವಿಗು ಸಮಾಸವೆಂದರೇನು? ಎರಡು ಉದಾಹರಣೆಗಳನ್ನು ಕೊಡಿ.\par

\vspace{2.5pt plus 3pt minus 0.5pt}\noindent\hspace*{1.2em}\kbvnum{(9)}\hspace{0.55em}ಬಹುವ್ರೀಹಿ ಸಮಾಸವೆಂದರೇನು? ಉದಾಹರಣೆ ಸಹಿತ ತಿಳಿಸಿ.\par

\vspace{2.5pt plus 3pt minus 0.5pt}\noindent\hspace*{1.2em}\kbvnum{(10)}\hspace{0.55em}ಕ್ರಿಯಾ ಸಮಾಸವೆಂದರೇನು? ಉದಾಹರಣೆ ಸಹಿತ ತಿಳಿಸಿ.\par

\vspace{2.5pt plus 3pt minus 0.5pt}\noindent\hspace*{1.2em}\kbvnum{(11)}\hspace{0.55em}ಗಮಕ ಸಮಾಸಕ್ಕೆ ಎರಡು ಉದಾಹರಣೆಗಳನ್ನು ಕೊಡಿ.\par

\vspace{2.5pt plus 3pt minus 0.5pt}\noindent\hspace*{1.2em}\kbvnum{(12)}\hspace{0.55em}ಇವುಗಳ ವಿಗ್ರಹವಾಕ್ಯ ಬರೆದು, ಸಮಾಸದ ಹೆಸರು ತಿಳಿಸಿ.\par

\vspace{5pt plus 4pt minus 1pt}
\noindent{\renewcommand{\arraystretch}{1.24}\setlength{\tabcolsep}{4pt}\rowcolors{1}{kbvrow}{white}
\begin{tabularx}{\textwidth}{X X X X}
ಮನೆದೇವರು & ಕೆಂದಾವರೆ & ಬಿರುಗಾಳಿ & ಹಿಂದಲೆ \\
ತಲೆದೆಗೆದು & ಚಾಕುಚೂರಿ & ಹುರಿಗಾಳು & ಕಣ್ಮುಚ್ಚಿ \\
ಕುಲದೇವತೆ & ಕಲ್ಲು & ಬಗ & ಮಳೆಬೆಳೆ \\
ಹಾಸಿಗೆಹೊದ್ದಿಕೆ & ಬಿಸಿನೀರು & ಬಯಲಾಟ & ಹೊರಮೈ \\
ಬಿಳಿದಾವರೆ & ಕಮಲಾಕ್ಷಿ & ಒಣಮರ & ಉಂಡನ್ನು \\
ಕೋದಂಡಪಾಣಿ & ಪೀತಾಂಬರ & ನಡುಮನೆ & ತಣ್ಣೀರು \\
ಕರಿತುರಗರಥ & ಕೋಳಿಮರಿ & ಒಳಗೆರೆ & ಮುಂಗೈ \\
ದಿವ್ಯಪ್ರಕಾಶ & ರೋಗಭಯ & ಮುಮ್ಮುಡಿ & ಹೆಮ್ಮಾರಿ \\
ದಾರಿಗಾಣದೆ & ಡಿಮುದು & ಕಿಬ್ಬೊಟ್ಟೆ & ಹೆದ್ದಾರಿ \\
ಮುಂಗುರುಳು & ಮೈದೋರು & ಮುಕ್ಕಣ್ಣ & ದಪ್ಪಕ್ಕೆ \\
ಮೈದಡವಿ & ಕೈಕೊಡು & ಇಮ್ಮಡಿ & ಶಶಿಧರ \\
\end{tabularx}}
\vspace{5pt plus 4pt minus 1pt}

\kbvchapter{ಪಾಠ ೧೪}{ಕೃದಂತಗಳು}

\indent ನಾಮ ಪ್ರಕೃತಿಗಳಲ್ಲಿ ಹಲವು ವಿಭೇದಗಳುಂಟು. ಇವುಗಳಲ್ಲಿ ಕೃನ್ನಾಮ ಪ್ರಕೃತಿ ಎಂಬುದು ಒಂದು ವಿಭೇದ. ಇದನ್ನು ಕೃದಂತ ನಾಮ ಪ್ರಕೃತಿ ಯೆಂತಲೂ ಕರೆಯುವುದುಂಟು.\par
}
\clearpage

% === tex/0073.tex (Printed Page 71) ===
\renewcommand{\kbvchaptername}{ಕೃದಂತಗಳು}
\setcounter{page}{71}
\thispagestyle{kbvpage}
\kbvtypesetpage{\indent “ಅವನು ಊಟ ಮಾಡಿದನು” ಎಂಬ ವಾಕ್ಯದಲ್ಲಿ “ಮಾಡಿದನು” ಎಂಬ ಶಬ್ದವನ್ನು ನೋಡೋಣ. ಇದು ಭೂತಕಾಲದ ಕ್ರಿಯೆಯನ್ನು ಗೊತ್ತುಪಡಿಸುವ ಕ್ರಿಯಾಪದ.\par

\indent ಮಾಡಿದನು = ಮಾಡು + ದ + ಅನು; ಇಲ್ಲಿ “ಮಾಡು” ಎಂಬ ಧಾತುವಿಗೆ “ಅನು” ಎಂಬ ಆಖ್ಯಾತ ಪ್ರತ್ಯಯ ಸೇರಿ, ಮಾಡಿದನು ಆಗಿರುವುದು. ಆದರೆ, ಧಾತುವಿಗೂ ಆಖ್ಯಾತ ಪ್ರತ್ಯಯಕ್ಕೂ ನಡುವೆ ಬಂದಿರುವ ‘ದ’ ಎಂಬ ಪ್ರತ್ಯಯವು ಕಾಲ ಸೂಚಕ ಪ್ರತ್ಯಯವಾಗಿರುವುದು. ಇದನ್ನು ಹಿಂದಿನ ಪ್ರಕರಣದಲ್ಲಿ ತಿಳಿದಿದ್ದೇವೆ. “ಮಾಡು” ಧಾತುವಿನ ಮತ್ತೊಂದು ರೂಪ “ಮಾಡಿದ” ಎಂಬ ಶಬ್ದ.\par

\indent ಮಾಡಿದ = ಮಾಡು + ದ + ಅ. ಇಲ್ಲಿ ಅಂತ್ಯದಲ್ಲಿರುವ ಪ್ರತ್ಯಯ “ಅ” ಮತ್ತು, ನಡುವೆ ಬಂದಿರುವ ಕಾಲ ಸೂಚಕ ಪ್ರತ್ಯಯ “ದ”. ಇದೇ ರೀತಿಯ ಹಲವು ಉದಾಹರಣೆಗಳನ್ನು ನೋಡೋಣ—\par

\begin{kbvderiv}
\vspace{4pt plus 3pt minus 1pt}
\begin{center}\begin{adjustbox}{max width=\textwidth}
{\renewcommand{\arraystretch}{1.22}\setlength{\tabcolsep}{4pt}
\begin{tabular}{l c l}
({\kbvlatin i}) ನೋಡಿದ & {\color{kbvgrey}$=$} & ನೋಡು + ದ + ಅ \\
({\kbvlatin ii}) ಹಾಡಿದ & {\color{kbvgrey}$=$} & ಹಾಡು + ದ + ಅ \\
\end{tabular}}
\end{adjustbox}\end{center}
\vspace{4pt plus 3pt minus 1pt}
\end{kbvderiv}

\indent (ಇಲ್ಲಿ “ದ” ಎಂಬುದು ಭೂತಕಾಲ ಸೂಚಕ ಪ್ರತ್ಯಯ)\par

\begin{kbvderiv}
\vspace{4pt plus 3pt minus 1pt}
\begin{center}\begin{adjustbox}{max width=\textwidth}
{\renewcommand{\arraystretch}{1.22}\setlength{\tabcolsep}{4pt}
\begin{tabular}{l c l}
({\kbvlatin iii}) ಓಡುವ & {\color{kbvgrey}$=$} & ಓಡು + ವ + ಅ \\
({\kbvlatin iv}) ತಿನ್ನುವ & {\color{kbvgrey}$=$} & ತಿನ್ನು + ವ + ಅ \\
\end{tabular}}
\end{adjustbox}\end{center}
\vspace{4pt plus 3pt minus 1pt}
\end{kbvderiv}

\indent (ಇಲ್ಲಿ “ವ” ಎಂಬುದು ವರ್ತಮಾನ ಕಾಲ ಸೂಚಕ ಪ್ರತ್ಯಯ)\par

\begin{kbvderiv}
\vspace{4pt plus 3pt minus 1pt}
\begin{center}\begin{adjustbox}{max width=\textwidth}
{\renewcommand{\arraystretch}{1.22}\setlength{\tabcolsep}{4pt}
\begin{tabular}{l c l}
({\kbvlatin v}) ನಡೆಯುವ & {\color{kbvgrey}$=$} & ನಡೆ + ಉವ + ಅ \\
({\kbvlatin vi}) ಕುಣಿಯುವ & {\color{kbvgrey}$=$} & ಕುಣಿ + ಉವ + ಅ \\
\end{tabular}}
\end{adjustbox}\end{center}
\vspace{4pt plus 3pt minus 1pt}
\end{kbvderiv}

\indent (ಇಲ್ಲಿ “ಉವ” ಎಂಬುದು ಭವಿಷ್ಯತ್ಕಾಲ ಸೂಚಕ ಪ್ರತ್ಯಯ) ಮೇಲಿನ ಆರು ಉದಾಹರಣೆಗಳಲ್ಲೂ ಪ್ರತಿಯೊಂದು ಶಬ್ದದ ಅಂತ್ಯದಲ್ಲಿ “ಅ” ಎಂಬ ಪ್ರತ್ಯಯವಿರುವುದು. ಈ ಪ್ರತ್ಯಯವನ್ನು ಕೃತ್‌ ಪ್ರತ್ಯಯ ಎನ್ನುತ್ತೇವೆ. ಹೀಗೆ ಕೃತ್‌ ಪ್ರತ್ಯಯವನ್ನು ಅಂತ್ಯದಲ್ಲಿ ಉಳ್ಳದ್ದೇ ಕೃದಂತ.\par

\noindent\hspace*{2.2em}{[}ಕೃತ್‌ + ಅಂತ = ಕೃದಂತ{]}.\par
}
\clearpage

% === tex/0074.tex (Printed Page 72) ===
\renewcommand{\kbvchaptername}{ಕೃದಂತಗಳು}
\setcounter{page}{72}
\thispagestyle{kbvpage}
\kbvtypesetpage{\kbvsechead{ಇದು ನಾಮಪದ:}

\indent ಧಾತುಗಳಿಗೆ ಕೃತ್‌ ಪ್ರತ್ಯಯ ಸೇರಿ ಆಗುವ ಶಬ್ದವೇ ಕೃದಂತ, ಇದನ್ನು ಕೃದಂತ ನಾಮಪ್ರಕೃತಿ ಅಥವಾ ಕೃನ್ನಾಮ ಪ್ರಕೃತಿ ಎಂದು ಕರೆಯುವುದು ವಾಡಿಕೆ. ಈ ರೀತಿಯಾಗಿ ನಾಮಪ್ರಕೃತಿಗಳಿಗೆ ನಾಮ ವಿಭಕ್ತಿ ಪ್ರತ್ಯಯಗಳು ಸೇರಿ ನಾಮಪದಗಳಾಗುವುದು; ಇವೇ ಕೃದಂತ ನಾಮಪದಗಳು.\par

\vspace{3pt plus 2pt}\noindent\hspace*{1.2em}\kbvlabel{ಉದಾ}\par\vspace{1pt}

\begin{kbvderiv}
\vspace{4pt plus 3pt minus 1pt}
\begin{center}\begin{adjustbox}{max width=\textwidth}
{\renewcommand{\arraystretch}{1.22}\setlength{\tabcolsep}{4pt}
\begin{tabular}{l c l}
ಓಡುವ + ಅವನು + ಉ & {\color{kbvgrey}$=$} & \kbvresult{ಓಡುವವನು} \\
ಓಡುವ + ಅವನು + ಅನ್ನು & {\color{kbvgrey}$=$} & \kbvresult{ಓಡುವವನನ್ನು} \\
ಓಡುವ + ಅವನು + ಇಂದ & {\color{kbvgrey}$=$} & \kbvresult{ಓಡುವವನಿಂದ} \\
ಓಡುವ + ಅವನು + ಇಗೆ & {\color{kbvgrey}$=$} & \kbvresult{ಓಡುವವನಿಗೆ} \\
ಓಡುವ + ಅವನು + ದೆಸೆಯಿಂದ & {\color{kbvgrey}$=$} & \kbvresult{ಓಡುವವನದೆಸೆಯಿಂದ} \\
ಓಡುವ + ಅವನು + ಅ & {\color{kbvgrey}$=$} & \kbvresult{ಓಡುವವನ} \\
ಓಡುವ + ಅವನು + ಅಲ್ಲಿ & {\color{kbvgrey}$=$} & \kbvresult{ಓಡುವವನಲ್ಲಿ} \\
ಓಡುವ + ಅವನು + ಏ & {\color{kbvgrey}$=$} & \kbvresult{ಓಡುವವನೇ} \\
\end{tabular}}
\end{adjustbox}\end{center}
\vspace{4pt plus 3pt minus 1pt}
\end{kbvderiv}

\indent ಮೇಲಿನ ಉದಾಹರಣೆಯಲ್ಲಿ ಕೃದಂತ ನಾಮಪ್ರಕೃತಿಗೂ ನಾಮ ವಿಭಕ್ತಿ ಪ್ರತ್ಯಯಕ್ಕೂ ನಡುವೆ “ಅವನು” ಎಂಬ ಸರ್ವನಾಮ ಆಗಮವಾದುದು. ಇದೇ ರೀತಿಯಾಗಿ, ಸ್ತ್ರೀಲಿಂಗ ಏಕವಚನದಲ್ಲಿ “ಅವಳು” ಮತ್ತು ನಪುಂಸಕಲಿಂಗ ಏಕವಚನದಲ್ಲಿ “ಅದು” ಎಂಬ ಸರ್ವನಾಮಗಳು ಬರುವುದು.\par

\begin{kbvrule}
\noindent\kbvbullet\hspace{0.5em}ಇದರ ಸಹಾಯದಿಂದ ನಾಮಪದಗಳನ್ನು ವಿಶೇಷಿಸುತ್ತಾರೆ.
\end{kbvrule}

\vspace{3pt plus 2pt}\noindent\hspace*{1.2em}\kbvlabel{ಉದಾ}\par\vspace{1pt}

\noindent\hspace*{1.8em}ಮಾ ಅಸ... (ಭೂತ ಕೃದಂತ)\hspace{0.9em}ಹಾ ಬಾರಿ... (ಭವಿಷ್ಯತ್‌ ಕೃದಂತ)\hspace{0.9em}ಸೂಚಕೆ; ಓಡ (ನಿಷೇಧ ಕೃದಂತ)\par\vspace{3pt plus 2pt}

\noindent\hspace*{2.2em}ನಿಷೇಧಾರ್ಥ ಸೂಚನೆ: “ಓಡದ” ಇಲ್ಲಿ “ಅದ” ಎಂಬುದು.\par
}
\clearpage

% === tex/0075.tex (Printed Page 73) ===
\renewcommand{\kbvchaptername}{ಕೃದಂತಗಳು}
\setcounter{page}{73}
\thispagestyle{kbvpage}
\kbvtypesetpage{\kbvsechead{ಕೃದಂತ ಭಾವನಾಮ :}

\indent ಕ್ರಿಯೆಯ ಭಾವವನ್ನು ತಿಳಿಸುವ ಶಬ್ದಗಳಿಗೆ ಕೃದಂತ ಭಾವನಾಮಗಳೆಂದು ಹೆಸರು; ಧಾತುಗಳ ಮೇಲೆ ಭಾವಾರ್ಥದಲ್ಲಿ ಕೃತ್‌ ಪ್ರತ್ಯಯಗಳು ಸೇರಿ ಕೃದಂತ ಭಾವನಾಮಗಳಾಗುವುದು.\par

\indent ಉದಾಹರಣೆಗೆ, “ಕುಣಿ” ಎಂಬ ಧಾತುವಿನ ಮೇಲೆ “ತ” ಎಂಬ ಪ್ರತ್ಯಯ ಸೇರಲು “ಕುಣಿತ” ಎಂದಾಗುವುದು. “ಕುಣಿತ” ಎಂಬ ಶಬ್ದರೂಪವು ಕುಣಿಯುವ ಭಾವವನ್ನು ತಿಳಿಸುವುದು. ಆದುದರಿಂದ ಇದು ಕೃದಂತ ಭಾವನಾಮವೆನಿಸುವುದು.\par

\begin{kbvrule}
\noindent\kbvbullet\hspace{0.5em}ಕೆಳಗೆ ಕೊಟ್ಟಿರುವ ಉದಾಹರಣೆಗಳಿಂದ ಕೃದಂತ ಭಾವನಾಮಗಳ ಅಥವಾ ಭಾವ ಕೃದಂತಗಳ ಅರ್ಥ ಮತ್ತಷ್ಟು ಸ್ಪಷ್ಟವಾಗಿ ತಿಳಿದುಬರುವುದು.
\end{kbvrule}

\vspace{5pt plus 4pt minus 1pt}
\noindent{\renewcommand{\arraystretch}{1.24}\setlength{\tabcolsep}{4pt}\rowcolors{2}{kbvrow}{white}
\begin{tabularx}{\textwidth}{X X X}
{\kbvsecfont\color{kbvaccentdark}ಮೂಲ} & {\kbvsecfont\color{kbvaccentdark}ಭಾವಾರ್ಥದಲ್ಲಿ} & {\kbvsecfont\color{kbvaccentdark}ಭಾವ} \\[1pt]
\arrayrulecolor{kbvline}\specialrule{0.5pt}{0pt}{2pt}\arrayrulecolor{black}
ಧಾತು & ಕೃತ್‌ ಪ್ರತ್ಯಯ & ಕೃದಂತ \\
ನಡುಗು &  + ಉಕ & ನಡುಕ \\
ನಗು &  + ಎ & ನಗೆ \\
ನೋಡು &  + ಟ & ನೋಟ \\
ಸೆಳೆ &  + ತ & ಸೆಳೆತ \\
ತೊಡು &  + ಇಗೆ & ತೊಡಿಗೆ \\
ತೊಡು &  + ಗೆ & ತೊಡುಗೆ \\
ನಾಚು &  + ಅಕೆ & ನಾಚಿಕೆ \\
ಉಳು &  + ವುದು & ಉಳುವುದು \\
\end{tabularx}}
\vspace{5pt plus 4pt minus 1pt}

\kbvsechead{ಕೃದಂತಾವ್ಯಯ :}

\indent ಧಾತುವಿನಿಂದ ಹುಟ್ಟಿ ಅವ್ಯಯದ ಲಕ್ಷಣಗಳನ್ನು ಹೊಂದಿರುವ ಶಬ್ದರೂಪಗಳಿಗೆ ಕೃದಂತಾವ್ಯಯಗಳು ಎಂದು ಹೆಸರು.\par

\indent ಧಾತುಗಳ ಮೇಲೆ ಅಲು, ಅಲಿಕ್ಕೆ, ದು, ಇ, ದರೆ, ಉತ್ತ್ಯ, ಅದೆ, ಅ, ಉತ್ತ್ಯ, ಇರಲು-ಮೊದಲಾದ ಪ್ರತ್ಯಯಗಳು ಸೇರಲು ಕೃದಂತಾವ್ಯಯಗಳು.\par
}
\clearpage

% === tex/0076.tex (Printed Page 74) ===
\renewcommand{\kbvchaptername}{ಕೃದಂತಗಳು}
\setcounter{page}{74}
\thispagestyle{kbvpage}
\kbvtypesetpage{\indent ಸಿದ್ಧಿಸುವುವು ; ಇವು ಲಿಂಗ ವಚನ, ವಿಭಕ್ತಿಗಳಿಂದ ವ್ಯತ್ಯಾಸವಾಗದ ರೂಪಗಳು.\par

\vspace{3pt plus 2pt}\noindent\hspace*{1.2em}\kbvlabel{ಉದಾ}\par\vspace{1pt}

\begin{kbvderiv}
\vspace{4pt plus 3pt minus 1pt}
\begin{center}\begin{adjustbox}{max width=\textwidth}
{\renewcommand{\arraystretch}{1.22}\setlength{\tabcolsep}{4pt}
\begin{tabular}{l c l}
೧. ಬೆಳಕು ಬರಲು ಕತ್ತಲೆ ಸರಿಯಿತು. (ಬರಲು & {\color{kbvgrey}$=$} & ಬರು + ಅಲು) \\
೨. ಈ ಹಣ್ಣು ತಿನ್ನಲಿಕ್ಕೆ ರುಚಿಯಾಗಿದೆ. (ತಿನ್ನಲಿಕ್ಕೆ & {\color{kbvgrey}$=$} & ತಿನ್ನು + ಅಲಿಕ್ಕೆ) \\
೩. ನೀನು ರೊಟ್ಟಿ ತಿಂದು ಹೋಗು. (ತಿಂದು & {\color{kbvgrey}$=$} & ತಿನ್ನು + ದು) \\
೪. ಮಾಡಿ ನೋಡು ಮದ್ದಿನ ಗುಣ. (ಮಾಡಿ & {\color{kbvgrey}$=$} & ಮಾಡು + ಇ) \\
೫. ಊರಿಗೆ ಹೋದರೆ ಹಿಂದಿರುಗಲಾರೆ. (ಹೋದರೆ & {\color{kbvgrey}$=$} & ಹೋಗು + ದರೆ) \\
೬. ಅವಳು ನಗುತ್ತೆ ಮನೆಗೆ ಹೋದಳು. (ನಗುತ್ತೆ & {\color{kbvgrey}$=$} & ನಗು + ಉತ್ತೆ) \\
೭. ನಾನು ನಿನ್ನನ್ನು ನೋಡದೆ ಇರಲಾರೆ. (ನೋಡದೆ & {\color{kbvgrey}$=$} & ನೋಡು + ಅದೆ) \\
೮. ಕಮಲೆಯು ನನ್ನನ್ನು ಬರೆ ಹೇಳಿದ್ದಾಳೆ. (ಬರೆ & {\color{kbvgrey}$=$} & ಬರು + ಅ) \\
೯. ಕಿಂದರಿ ಜೋಗಿ ಹಾಡುತೆ ಬರುತಿಹನು. (ಹಾಡುತೆ & {\color{kbvgrey}$=$} & ಹಾಡು + ಉತ್ತೆ) \\
\end{tabular}}
\end{adjustbox}\end{center}
\vspace{4pt plus 3pt minus 1pt}
\end{kbvderiv}

\vspace{2.5pt plus 3pt minus 0.5pt}\noindent\hspace*{1.2em}\kbvnum{೧೦.}\hspace{0.55em}ರಾಷ್ಟ್ರಾಧ್ಯಕ್ಷರು ಬರುತ್ತಿರಲು ಮಕ್ಕಳು ಚಪ್ಪಾಳೆ ತಟ್ಟಿದರು.\par

\indent (ಬರುತ್ತಿರಲು = ಬರು + ಉತ್ತೆ + ಅರಲು)\par

\indent ನೆನಪಿಡಿ : ಕೃದಂತಾವ್ಯಯಗಳು ಪೂರ್ಣ ಕ್ರಿಯಾರ್ಥವನ್ನು ಕೊಡುವುದಿಲ್ಲವಾದ ಕಾರಣ ಮತ್ತೊಂದು ಕ್ರಿಯೆಯನ್ನು ಅಪೇಕ್ಷಿಸುವುದು. ಇವುಗಳ ಲಕ್ಷಣ, ಉದಾ: ನಗುತ್ತ ಇದ್ದನು; ನೋಡದೆ ಇರಲಾರೆ; ಮಾಡಿ ನೋಡು.\par

\kbvexercise{ಅಭ್ಯಾಸ ಪತ್ರಿಕೆ ೧೪}

\vspace{2.5pt plus 3pt minus 0.5pt}\noindent\hspace*{1.2em}\kbvnum{1.}\hspace{0.55em}“ಬರೆಯುವನು” ಈ ಕ್ರಿಯಾಪದದಲ್ಲಿ ಸೇರಿರುವ ಕಾಲಸೂಚಕ ಪ್ರತ್ಯಯ ಯಾವುದು? ಮತ್ತು ಆಖ್ಯಾತ ಪ್ರತ್ಯಯ ಯಾವುದು?\par

\vspace{2.5pt plus 3pt minus 0.5pt}\noindent\hspace*{1.2em}\kbvnum{2.}\hspace{0.55em}ಸ್ತ್ರೀಲಿಂಗ ಏಕವಚನದಲ್ಲಿ “ಕೊಡು” ಧಾತುವಿನ ಭವಿಷ್ಯಕಾಲದ ಕ್ರಿಯಾಪದ ರೂಪವನ್ನು ಬರೆಯಿರಿ.\par

\vspace{2.5pt plus 3pt minus 0.5pt}\noindent\hspace*{1.2em}\kbvnum{3.}\hspace{0.55em}“ಹೋದನು” ಈ ಕ್ರಿಯಾಪದವನ್ನು ಬಿಡಿಸಿ ಬರೆಯಿರಿ.\par

\vspace{2.5pt plus 3pt minus 0.5pt}\noindent\hspace*{1.2em}\kbvnum{4.}\hspace{0.55em}“ಕಾಣುವ” ಈ ಶಬ್ದದಲ್ಲಿ ಸೇರಿರುವ ಪ್ರತ್ಯಯಗಳು ಯಾವುವು? ಅವುಗಳ ವಿಶೇಷವೇನು?\par

\vspace{2.5pt plus 3pt minus 0.5pt}\noindent\hspace*{1.2em}\kbvnum{5.}\hspace{0.55em}ಕೃದಂತ ನಾಮ ಎಂದರೇನು? ಉದಾಹರಣೆಗಳ ಸಹಿತ ತಿಳಿಸಿ.\par
}
\clearpage

% === tex/0077.tex (Printed Page 75) ===
\renewcommand{\kbvchaptername}{ತದ್ಧಿತಾಂತಗಳು}
\setcounter{page}{75}
\thispagestyle{kbvpage}
\kbvtypesetpage{\vspace{2.5pt plus 3pt minus 0.5pt}\noindent\hspace*{1.2em}\kbvnum{6.}\hspace{0.55em}ಕೃದಂತ ನಾಮಪದಗಳೆಂದರೇನು? ಉದಾಹರಿಸಿ.\par

\vspace{2.5pt plus 3pt minus 0.5pt}\noindent\hspace*{1.2em}\kbvnum{7.}\hspace{0.55em}ಕೃದಂತ ನಾಮಗಳು ನಾಮಪದಗಳನ್ನು ವಿಶೇಷಿಸುತ್ತಿರುವುದು ಎಂಬುದನ್ನು ಸಮರ್ಥಿಸಲು ಎರಡು ಉದಾಹರಣೆಗಳನ್ನು ಕೊಡಿ.\par

\vspace{2.5pt plus 3pt minus 0.5pt}\noindent\hspace*{1.2em}\kbvnum{8.}\hspace{0.55em}ಕೃದಂತ ಭಾವನಾಮ ಎಂದರೇನು? ಉದಾಹರಣೆ ಸಹಿತ ತಿಳಿಸಿ.\par

\vspace{2.5pt plus 3pt minus 0.5pt}\noindent\hspace*{1.2em}\kbvnum{9.}\hspace{0.55em}ಮಾಡು, ಕೊರೆ, ಓಡು, ನಂಬು, ಅಳೆ—ಈ ಧಾತುಗಳಿಂದ ಹುಟ್ಟುವ ಕೃದಂತ ಭಾವನಾಮಗಳನ್ನು ತಿಳಿಸಿರಿ.\par

\vspace{2.5pt plus 3pt minus 0.5pt}\noindent\hspace*{1.2em}\kbvnum{10.}\hspace{0.55em}“ಹೊಗಳುವ” ಈ ಕೃದಂತ ನಾಮಪ್ರಕೃತಿಗೆ ವಿಭಕ್ತಿ ಪ್ರತ್ಯಯಗಳನ್ನು ಸೇರಿಸಿ ಬರೆಯಿರಿ.\par

\vspace{2.5pt plus 3pt minus 0.5pt}\noindent\hspace*{1.2em}\kbvnum{11.}\hspace{0.55em}ಇರಿತ, ಉಡಿಗೆ—ಈ ಶಬ್ದಗಳನ್ನು ನಿಮ್ಮ ಸ್ವಂತ ವಾಕ್ಯಗಳಲ್ಲಿ ಬಳಸಿರಿ.\par

\vspace{2.5pt plus 3pt minus 0.5pt}\noindent\hspace*{1.2em}\kbvnum{12.}\hspace{0.55em}ಕೃದಂತಾವ್ಯಯ ಎಂದರೇನು? ಉದಾಹರಣೆಗಳ ಸಹಿತ ತಿಳಿಸಿ.\par

\vspace{2.5pt plus 3pt minus 0.5pt}\noindent\hspace*{1.2em}\kbvnum{13.}\hspace{0.55em}ಕೃದಂತಾವ್ಯಯಗಳ ಲಕ್ಷಣಗಳೇನು?\par

\vspace{2.5pt plus 3pt minus 0.5pt}\noindent\hspace*{1.2em}\kbvnum{14.}\hspace{0.55em}“ಬೀದಿಯಲ್ಲಿ ಬರುತ್ತ ಮುಗ್ಗರಿಸಿ ಬಿದ್ದನು.” ಈ ವಾಕ್ಯದಲ್ಲಿ ಕೃದಂತಾವ್ಯಯ ಯಾವುದು?\par

\vspace{2.5pt plus 3pt minus 0.5pt}\noindent\hspace*{1.2em}\kbvnum{15.}\hspace{0.55em}“ಮಾಡಲಿಕ್ಕೆ” ಈ ಶಬ್ದದಲ್ಲಿ ಧಾತು ಯಾವುದು? ಮತ್ತು ಪ್ರತ್ಯಯ ಯಾವುದು?\par

\vspace{2.5pt plus 3pt minus 0.5pt}\noindent\hspace*{1.2em}\kbvnum{16.}\hspace{0.55em}ಓಟ, ಹೇಸಿಗೆ—ಈ ಪದಗಳಲ್ಲಿರುವ ಮೂಲ ಧಾತುಗಳನ್ನು ಗುರ್ತಿಸಿ.\par

\vspace{2.5pt plus 3pt minus 0.5pt}\noindent\hspace*{1.2em}\kbvnum{17.}\hspace{0.55em}ಭೂತ ಕೃದಂತ, ಭವಿಷ್ಯತ್‌ ಕೃದಂತ, ನಿಷೇಧ ಕೃದಂತ—ಇವು ಗಳಿಗೆ ಒಂದೊಂದು ಉದಾಹರಣೆಯನ್ನು ಕೊಡಿ.\par

\kbvchapter{ಪಾಠ ೧೫}{ತದ್ಧಿತಾಂತಗಳು}

\indent ಆಟವನ್ನು ಆಡುವವನು ಆಟಗಾರ; ಇಲ್ಲಿ “ಆಟ” ಎಂಬ ನಾಮ ಪ್ರಕೃತಿಯೊಂದಿಗೆ “ಗಾರ” ಎಂಬ ಪ್ರತ್ಯಯ ಸೇರಿ “ಆಟಗಾರ” ಆಗಿರುವುದು. “ಗಾರ” ಪ್ರತ್ಯಯವನ್ನು “ಆಡುವವನು” ಎಂಬ ಅರ್ಥದಲ್ಲಿ ಬಳಸಲಾಗಿದೆ.\par

\indent ಆಟವನ್ನು + ಗಾರ = ಆಟ + ಅನ್ನು + ಗಾರ; ಆದರೆ, ಒಂದೇ ನಾಮ ಪ್ರಕೃತಿಯ ಮೇಲೆ ಎರಡು ಪ್ರತ್ಯಯಗಳು ಬರುವುದು ಅಸಂಭವ. ಈ ಕಾರಣ\kbvcontinues
}
\clearpage

% === tex/0078.tex (Printed Page 76) ===
\renewcommand{\kbvchaptername}{ತದ್ಧಿತಾಂತಗಳು}
\setcounter{page}{76}
\thispagestyle{kbvpage}
\kbvtypesetpage{\kbvcontinued ದಿಂದ “ಅನ್ನು” ಎಂಬ ಪ್ರತ್ಯಯ ಲೋಪವಾಗಿ, ಆಟ + ಗಾರ = ಆಟಗಾರ ಎಂದಾಗುವುದು. ಇಲ್ಲಿ ಬಳಸಿರುವ “ಗಾರ” ಪ್ರತ್ಯಯವನ್ನು ತದ್ಧಿತಪ್ರತ್ಯಯ ವೆನ್ನುತ್ತೇವೆ. ಈ ರೀತಿಯಾದ ಪ್ರತ್ಯಯಗಳನ್ನು ಅಂತ್ಯದಲ್ಲಿ ಉಳ್ಳ ಶಬ್ದರೂಪಗಳೇ ತದ್ಧಿತಾಂತ (ತದ್ಧಿತ + ಅಂತ)ಗಳು.\par

\kbvsechead{ಈ ಉದಾಹರಣೆಗಳನ್ನು ಗಮನಿಸಿ:}

\vspace{2.5pt plus 3pt minus 0.5pt}\noindent\hspace*{1.2em}\kbvnum{(1)}\hspace{0.55em}ಲಂಚ + (ಅನ್ನು) + ಕೋರ = ಲಂಚಕೋರ {[}ಲಂಚವನ್ನು ಕೋರುವವನು ಲಂಚಕೋರ{]}\par

\vspace{2.5pt plus 3pt minus 0.5pt}\noindent\hspace*{1.2em}\kbvnum{(2)}\hspace{0.55em}ಚಿತ್ರ + (ಅನ್ನು) + ಕಾರ = ಚಿತ್ರಕಾರ {[}ಚಿತ್ರವನ್ನು ಬರೆಯುವವನು ಚಿತ್ರಕಾರ{]}\par

\vspace{2.5pt plus 3pt minus 0.5pt}\noindent\hspace*{1.2em}\kbvnum{(3)}\hspace{0.55em}ಗಾಣ + (ಅನ್ನು) + ಇಗ = ಗಾಣಿಗ {[}ಗಾಣವನ್ನು ಆಡಿಸುವವನು ಗಾಣಿಗ{]}\par

\vspace{2.5pt plus 3pt minus 0.5pt}\noindent\hspace*{1.2em}\kbvnum{(4)}\hspace{0.55em}ಮಡಿ + (ಅನ್ನು) + ವಂತ = ಮಡಿವಂತ {[}ಮಡಿಯನ್ನು ಆಚರಿಸುವವನು ಮಡಿವಂತ{]}\par

\vspace{2.5pt plus 3pt minus 0.5pt}\noindent\hspace*{1.2em}\kbvnum{(5)}\hspace{0.55em}ಕೌರ + (ಅನ್ನು) + ಇಕ = ಕೌರಿಕ {[}ಕೌರವನ್ನು ಮಾಡುವವನು ಕೌರಿಕ{]}\par

\indent ಮೇಲೆ ತಿಳಿಸಿದ ಉದಾಹರಣೆಗಳಲ್ಲಿ ಕೋರ, ಕಾರ, ಇಗ, ವಂತ ಇಳಿತ್ತು. ಈ ಪ್ರತ್ಯಯಗಳು,\par

\begin{kbvrule}
\noindent\kbvbullet\hspace{0.5em}ತದ್ಧಿತ ಪ್ರತ್ಯಯಗಳು ನಾಮಪದಗಳ ಮೇಲೆ ಬಂದು ಹೊಸ ಶಬ್ದಗಳಾಗುವುದು. ಪ್ರತ್ಯಯಗಳಿಗೆ ತದ್ಧಿತಾಂತಗಳೆಂದು ಹೆಸರು. ತದ್ಧಿತಾಂತಗಳನ್ನು (1) ತದ್ಧಿತಾಂತ ನಾಮ, (2) ತದ್ಧಿತಾಂತ ಭಾವನಾಮ, (3) ತದ್ಧಿತಾಂತಾವ್ಯಯ ಎಂದು ಮೂರು ಬಗೆಯಾಗಿ ವಿಂಗಡಿಸಬಹುದು.
\end{kbvrule}

\vspace{4pt plus 4pt minus 1pt}\noindent{\kbvsecfont\fontsize{11pt}{14.5pt}\selectfont\color{kbvaccentdark} (1)\hspace{0.45em}ತದ್ಧಿತಾಂತನಾಮ:}\hspace{0.5em}ನಾಮಪದಗಳ ಮೇಲೆ ವಿವಿಧ ಅರ್ಥಗಳಲ್ಲಿ ಗಾರ, ಕಾರ, ಇಗ, ವಂತ ಮೊದಲಾದ ತದ್ಧಿತ ಪ್ರತ್ಯಯಗಳು ಸೇರಿ ಆಗುವ ಹೊಸ ಶಬ್ದಗಳಿಗೆ ತದ್ಧಿತಾಂತ ನಾಮಗಳೆಂದು ಹೆಸರು.\par
}
\clearpage

% === tex/0079.tex (Printed Page 77) ===
\renewcommand{\kbvchaptername}{ತದ್ಧಿತಾಂತಗಳು}
\setcounter{page}{77}
\thispagestyle{kbvpage}
\kbvtypesetpage{\noindent\hspace*{1.2em}\kbvlabel{ಉದಾ}\hspace{0.6em}ಜೂಜುಗಾರ, ಕಲಾಕಾರ, ಕನ್ನಡಿಗ, ಹಣವಂತ, ಮಾಂತ್ರಿಕ, ಮಡಿವಾಳ, ಚಮ್ಮಾರ, ಹಾವಾಡಿಗ ಇತ್ಯಾದಿ ; ಇವು ಪುಲ್ಲಿಂಗ ತದ್ಧಿತಾಂತ ನಾಮಗಳು.\par

\vspace{2.5pt plus 3pt minus 0.5pt}\noindent\hspace*{1.2em}\kbvnum{೧.}\hspace{0.55em}ಸ್ತ್ರೀಲಿಂಗದಲ್ಲಿ ಗಾತಿ, ವಂತೆ, ಗಿತ್ತಿ, ಇತಿ, ವತಿ, ಎ-ಇತ್ಯಾದಿ ತದ್ಧಿತ ಪ್ರತ್ಯಯಗಳು ನಾಮಪದಗಳ ಮೇಲೆ ಬಂದು ಸ್ತ್ರೀಲಿಂಗ ತದ್ಧಿತಾಂತ ನಾಮಗಳಾಗುವುವು. ಉದಾಹರಣೆಗೆ ಮೋಜುಗಾತಿ, ದಯಾವಂತೆ, ಹೂವಾಡಿಗಿತ್ತಿ, ಕನ್ನಡಿತಿ, ಗುಣವತಿ, ಜಾಣೆ ಇತ್ಯಾದಿ.\par

\kbvsechead{(೨)\hspace{0.45em}ತದ್ಧಿತಾಂತ ಭಾವನಾಮ :}

\indent ಷಷ್ಠೀ ವಿಭಕ್ತಿಗಳಾದ ನಾಮಪದಗಳ ಮೇಲೆ ಭಾವಾರ್ಥದಲ್ಲಿ ತನ, ಇಕೆ, ಮೆ, ಪು, ಉ ಇತ್ಯಾದಿ ತದ್ಧಿತ ಪ್ರತ್ಯಯಗಳು ಸೇರಿ ಆಗುವ ಹೊಸ ಶಬ್ದಗಳಿಗೆ ತದ್ಧಿತಾಂತ ಭಾವನಾಮಗಳೆಂದು ಹೆಸರು. ಉದಾ : ದಡ್ಡ ತನ.\par

\begin{kbvnote}
\noindent{\kbvsecfont ವಿವರಣೆ :}\hspace{0.5em}“ದಡ್ಡ” ಎಂಬ ನಾಮಪದದ ಷಷ್ಠೀ ವಿಭಕ್ತಿರೂಪ “ದಡ್ಡನ” ; ದಡ್ಡನ ಭಾವವೇ ದಡ್ಡ ತನ. ಇಲ್ಲಿ “ದಡ್ಡನ” ಎಂಬ ನಾಮಪದದ ಮುಂದೆ ಭಾವಾರ್ಥದಲ್ಲಿ “ತನ” ಎಂಬ ಪ್ರತ್ಯಯ ಸೇರಿ “ದಡ್ಡ ತನ” ಆಗಿರು ವುದು ; ಇದು ತದ್ಧಿತಾಂತ ಭಾವನಾಮ.
\end{kbvnote}

\indent ಇದೇ ರೀತಿಯಾಗಿ ಗೌಡಿಕೆ, ಜಾಣ್ಮೆ, ತಂಪು, ಕುರುಡು ಇವು\par

\kbvsechead{ತದ್ಧಿತಾಂತ ಭಾವನಾಮಗಳು :}

\kbvsechead{ವಿಶ್ಲೇಷಣೆ :}

\indent - ಗೌಡನ ಭಾವ ಗೌಡಿಕೆ (ಇಲ್ಲಿ “ಇಕೆ” ಪ್ರತ್ಯಯ ಸೇರಿದೆ).\par

\indent - ಜಾಣನ ಭಾವ ಜಾಣ್ಮೆ (ಇಲ್ಲಿ “ಮೆ” ಪ್ರತ್ಯಯ ಸೇರಿದೆ).\par

\indent - ತಣ್ಣನೆಯದರ ಭಾವ ತಂಪು (ಇಲ್ಲಿ “ಪು” ಪ್ರತ್ಯಯ ಸೇರಿದೆ).\par

\indent - ಕುರುಡನ ಭಾವ ಕುರುಡು (ಇಲ್ಲಿ “ಉ” ಪ್ರತ್ಯಯ ಸೇರಿದೆ).\par

\vspace{2.5pt plus 3pt minus 0.5pt}\noindent\hspace*{1.2em}\kbvnum{(೩)}\hspace{0.55em}ತದ್ಧಿತಾಂತಾವ್ಯಯ\par

\begin{kbvderiv}
\vspace{4pt plus 3pt minus 1pt}
\begin{center}\begin{adjustbox}{max width=\textwidth}
{\renewcommand{\arraystretch}{1.22}\setlength{\tabcolsep}{4pt}
\begin{tabular}{l c l}
ಶಂಕರನ + ಅಂತೆ & {\color{kbvgrey}$=$} & \kbvresult{ಶಂಕರನಂತೆ}{\color{kbvgrey} (ಪು.ಎ.)} \\
ಗೌರಿಯ + ಅಂತೆ & {\color{kbvgrey}$=$} & \kbvresult{ಗೌರಿಯಂತೆ}{\color{kbvgrey} (ಸ್ತ್ರೀ.ಎ.)} \\
ಲತೆಯ + ಅಂತೆ & {\color{kbvgrey}$=$} & \kbvresult{ಲತೆಯಂತೆ}{\color{kbvgrey} (ನ.ಎ.)} \\
ಮಕ್ಕಳ + ಅಂತೆ & {\color{kbvgrey}$=$} & \kbvresult{ಮಕ್ಕಳಂತೆ}{\color{kbvgrey} (ನ.ಬ.)} \\
\end{tabular}}
\end{adjustbox}\end{center}
\vspace{4pt plus 3pt minus 1pt}
\end{kbvderiv}
}
\clearpage

% === tex/0080.tex (Printed Page 78) ===
\renewcommand{\kbvchaptername}{ತದ್ಧಿತಾಂತಗಳು}
\setcounter{page}{78}
\thispagestyle{kbvpage}
\kbvtypesetpage{\indent ಈ ಉದಾಹರಣೆಗಳಲ್ಲಿ ನಾಮಪದಗಳ ಮುಂದೆ “ಅಂತೆ” ಎಂಬ ಪ್ರತ್ಯಯ ಸೇರಿ ಹೊಸ ಶಬ್ದಗಳಾಗಿರುವುದನ್ನು ಕಾಣುತ್ತೇವೆ ; ಈ ಪ್ರತಿ ಯೊಂದು ಶಬ್ದದಲ್ಲೂ ನಾಮಪದದಲ್ಲಿರುವ ವಿಭಕ್ತಿ ಪ್ರತ್ಯಯವು ಹಾಗೆಯೇ ಉಳಿದುಕೊಂಡಿರುವುದು ಗಮನಾರ್ಹ ಅಂಶ.\par

\indent ಹೀಗೆ ನಾಮಪದಗಳ ಮುಂದೆ ಅಂತೆ, ತನಕ, ವರೆಗೆ ಇಂತ, ಆಗಿ ವೋಲ್‌, ಮಟ್ಟಿಗೆ, ಓಸ್ಕರ, ಸಲುವಾಗಿ ಇತ್ಯಾದಿ ಪ್ರತ್ಯಯಗಳು ಸೇರಿ ಆಗುವ ಹೊಸ ಶಬ್ದಗಳಿಗೆ ತದ್ಧಿತಾಂತಾವ್ಯಯಗಳೆಂದು ಹೆಸರು.\par

\vspace{3pt plus 2pt}\noindent\hspace*{1.2em}\kbvlabel{ಉದಾ}\par\vspace{1pt}

\vspace{5pt plus 4pt minus 1pt}
\noindent{\renewcommand{\arraystretch}{1.24}\setlength{\tabcolsep}{4pt}\rowcolors{1}{kbvrow}{white}
\begin{tabularx}{\textwidth}{X X X}
ದೊರೆಯತನಕ, & ತಿಂಗಳವರೆಗೆ, & ಹಣಕ್ಕೋಸ್ಕರ, \\
ರಾಣಿಯತನಕ, & ಕೃಷ್ಣನಿಗಿಂತ, & ನನ್ನಮಟ್ಟಿಗೆ, \\
ಬಾಗಿಲತನಕ, & ಶಿಲೆಗಿಂತ, & ಅವಳಸಲುವಾಗಿ, \\
ಜನಗಳವರೆಗೆ, & ಪುಸ್ತಕಗಳಿಗಿಂತ, & ರಾಜನವೋಲ್‌ \\
ದಿಗಂತದವರೆಗೆ & ಅವರಿಗಾಗಿ & ಸಿಂಹದವೋಲ್‌ \\
\end{tabularx}}
\vspace{5pt plus 4pt minus 1pt}

\kbvexercise{ಅಭ್ಯಾಸ ಪತ್ರಿಕೆ ೧೫}

\vspace{2.5pt plus 3pt minus 0.5pt}\noindent\hspace*{1.2em}\kbvnum{(1)}\hspace{0.55em}ತದ್ಧಿತ ಪ್ರತ್ಯಯ ಎಂದರೇನು? ಒಂದು ಉದಾಹರಣೆ ಸಹಿತ ತಿಳಿಸಿ.\par

\vspace{2.5pt plus 3pt minus 0.5pt}\noindent\hspace*{1.2em}\kbvnum{(2)}\hspace{0.55em}ತದ್ಧಿತಾಂತಗಳೆಂದರೇನು? ಎರಡು ಉದಾಹರಣೆಗಳನ್ನು ಕೊಡಿ.\par

\vspace{2.5pt plus 3pt minus 0.5pt}\noindent\hspace*{1.2em}\kbvnum{(3)}\hspace{0.55em}ತದ್ಧಿತಾಂತಗಳ ಮೂರು ಪ್ರಕಾರಗಳನ್ನು ತಿಳಿಸಿ.\par

\vspace{2.5pt plus 3pt minus 0.5pt}\noindent\hspace*{1.2em}\kbvnum{(4)}\hspace{0.55em}ತದ್ಧಿತಾಂತನಾಮ ಎಂದರೇನು? ಉದಾಹರಣೆ ಕೊಡಿ.\par

\vspace{2.5pt plus 3pt minus 0.5pt}\noindent\hspace*{1.2em}\kbvnum{(5)}\hspace{0.55em}ತದ್ಧಿತಾಂತ ಭಾವನಾಮ ಎಂದರೇನು? ಉದಾಹರಣೆ ಸಹಿತ ವಿವರಿಸಿ.\par

\vspace{2.5pt plus 3pt minus 0.5pt}\noindent\hspace*{1.2em}\kbvnum{(6)}\hspace{0.55em}ಕೃದಂತ ಭಾವನಾಮಗಳಿಗೂ ತದ್ಧಿತಾಂತ ಭಾವನಾಮಗಳಿಗೂ ಇರುವ ವ್ಯತ್ಯಾಸವೇನು?\par

\vspace{2.5pt plus 3pt minus 0.5pt}\noindent\hspace*{1.2em}\kbvnum{(7)}\hspace{0.55em}ಬ್ರಾಜ್ಮಟಕೆ, ನಾಚಿಕೆ-ಈ ಪದಗಳಲ್ಲಿ ಕಂಡುಬರುವ ವ್ಯಾಕರಣ ವಿಶೇಷವೇನು?\par
}
\clearpage

% === tex/0081.tex (Printed Page 79) ===
\renewcommand{\kbvchaptername}{ಅವ್ಯಯಗಳು}
\setcounter{page}{79}
\thispagestyle{kbvpage}
\kbvtypesetpage{\vspace{2.5pt plus 3pt minus 0.5pt}\noindent\hspace*{1.2em}\kbvnum{(8)}\hspace{0.55em}ತದ್ದಿತಾಂತಾವ್ಯಯಗಳೆಂದರೇನು? ಉದಾಹರಣೆ ಸಹಿತ ವಿವರಿಸಿ.\par

\vspace{2.5pt plus 3pt minus 0.5pt}\noindent\hspace*{1.2em}\kbvnum{(9)}\hspace{0.55em}ಈ ಅರ್ಥ ಬರುವ ಒಂಟಿಪದ ಯಾವುದು ?\par

\vspace{2.5pt plus 3pt minus 0.5pt}\noindent\hspace*{1.2em}\kbvnum{(ಅ)}\hspace{0.55em}ಬಳೆಗಳನ್ನು ಮಾರುವವನು (ಬ) ಸುಂದರ ರೂಪ ಉಳ್ಳವಳು\par

\vspace{2.5pt plus 3pt minus 0.5pt}\noindent\hspace*{1.2em}\kbvnum{(ಕ)}\hspace{0.55em}ಸಿರಿ ಉಳ್ಳವನ ಭಾವ (ಡ) ಗೊಲ್ಲರ ಹೆಂಗಸು\par

\vspace{2.5pt plus 3pt minus 0.5pt}\noindent\hspace*{1.2em}\kbvnum{(10)}\hspace{0.55em}ಈ ಪದಗಳಲ್ಲಿರುವ ತದ್ದಿತ ಪ್ರತ್ಯಯಗಳನ್ನು ತಿಳಿಸಿ : ಹಾವಾಡಿಗ, ಒಕ್ಕಲಗಿತ್ತಿ, ಜಾಣತನ, ರಾಮನಂತೆ ತಾಯಿಗಿಂತ\par

\kbvchapter{ಪಾಠ ೧೬}{ಅವ್ಯಯಗಳು}

\kbvsechead{ಈ ವಾಕ್ಯಗಳನ್ನು ಓದಿ ನೋಡೋಣ :}

\noindent\hspace*{2.2em}ಗಿಡವು ಚೆನ್ನಾಗಿ ಬೆಳೆಯುತ್ತಿದೆ.\par

\indent ಜ್ಯೋತಿಯು ಚೆನ್ನಾಗಿ ಹಾಡುತ್ತಾಳೆ. ಅಕ್ಕಾ ! ಚೆನ್ನಾಗಿ ಆಲೋಚಿಸು, ನದಿಯಲ್ಲಿ ಚೆನ್ನಾಗಿ ಈಜಬಹುದು.\par

\noindent\hspace*{2.2em}ಹುಡುಗರು ಚೆನ್ನಾಗಿ ಓಡುತ್ತಿದ್ದಾರೆ.\par

\noindent\hspace*{2.2em}ಅವರಿಗೆ ಚೆನ್ನಾಗಿ ಹಣ ಸಿಕ್ಕಿತು.\par

\noindent\hspace*{2.2em}ಮಕ್ಕಳು ಚೆನ್ನಾಗಿ ಆಟವಾಡಿದರು.\par

\noindent\hspace*{2.2em}ಪಕ್ಷಿಗಳು ಚೆನ್ನಾಗಿ ಹಾರುತ್ತಿವೆ.\par

\noindent\hspace*{2.2em}ಮರಗಳಲ್ಲಿ ಚೆನ್ನಾಗಿ ಹೂಬಿಟ್ಟಿದೆ.\par

\indent ಮೇಲಿನ ವಾಕ್ಯಗಳಲ್ಲಿ ಪ್ರತಿಯೊಂದರಲ್ಲೂ ಚೆನ್ನಾಗಿ ಎಂಬ ಶಬ್ದ ಬಂದಿದೆ ; ನಾಮಪದ ಕ್ರಿಯಾಪದಗಳಲ್ಲಿ ಬದಲಾವಣೆಯಾಗಿದ್ದರೂ “ಚೆನ್ನಾಗಿ” ಎಂಬ ಶಬ್ದ ಇದ್ದಂತೆಯೇ ಇರುವುದು. ಹೀಗೆ ಲಿಂಗ ವಚನ ವಿಭಕ್ತಿಗಳಲ್ಲಿ ಬದಲಾವಣೆಯುಂಟಾದರೂ ಸ್ವತಃ ಯಾವುದೇ ವ್ಯತ್ಯಾಸ ಹೊಂದದ ಶಬ್ದವು ಅವ್ಯಯವೆನಿಸುವುದು.\par

\indent ಸುಲಭವಾಗಿ, ತಟ್ಟನೆ, ಭರ್ರನೆ, ಅಯ್ಯೋ !, ಬೇಕು, ಮತ್ತು ಅಥವಾ, ಅಂತೆ, ಬರಲು, ಛೀ !, ಧಗಧಗ—ಇವು ಅವ್ಯಯಗಳಿಗೆ ಕೆಲವು ಉದಾಹರಣೆಗಳು.\par
}
\clearpage

% === tex/0082.tex (Printed Page 80) ===
\renewcommand{\kbvchaptername}{ಅವ್ಯಯಗಳು}
\setcounter{page}{80}
\thispagestyle{kbvpage}
\kbvtypesetpage{\indent ಅವ್ಯಯಗಳಲ್ಲಿ ಎಂಟು ವಿಭೇದಗಳುಂಟು ; ಅವುಗಳನ್ನು ವಿವರವಾಗಿ ತಿಳಿಯೋಣ.\par

\vspace{4pt plus 4pt minus 1pt}\noindent{\kbvsecfont\fontsize{11pt}{14.5pt}\selectfont\color{kbvaccentdark} (೧)\hspace{0.45em}ಸಾಮಾನ್ಯಾವ್ಯಯ :}\hspace{0.5em}ಮೆಲ್ಲಗೆ, ಬೇಗನೆ ವೇಗವಾಗಿ, ಕೂಡಲೆ ಸುಮ್ಮನೆ, ಚೆನ್ನಾಗಿ ತಟ್ಟನೆ, ಎತ್ತರವಾಗಿ, ಇಂತು ಇತ್ಯಾದಿ ಪದಗಳು ಕ್ರಿಯೆ ನಡೆಯುವ ಯನ್ನು ಹೇಳುವ ಅವ್ಯಯಗಳು. ಈ ಕ್ರಿಯಾವಿಶೇಷಗಳಿಗೆ ಸಾಮಾನ್ಯಾವ್ಯಯಗಳೆಂದು ಹೆಸರು.\par

\vspace{4pt plus 4pt minus 1pt}\noindent{\kbvsecfont\fontsize{11pt}{14.5pt}\selectfont\color{kbvaccentdark} (೨)\hspace{0.45em}ಅನುಕರಣಾವ್ಯಯ :}\hspace{0.5em}ಚಟಚಟ, ಪಟಪಟ, ತಟತಟ, ದಡದಡ, ಕರಕರ, ಸರಸರ, ದಬದಬ, ಚುರುಚುರು, ಧಗಧಗ, ರೊಯ್ಯನೆ, ಭರ್ರನೆ, ಗಿರ್ರನೆ ಸುಯ ಇವೇ ಮೊದಲಾದ ಪದಗಳು ಅರ್ಥರಹಿತ ಧ್ವನಿ ವಿಶೇಷಗಳು. ಕೇಳಿದ ರೀತಿಯಲ್ಲಿ ಅನುಕರಣ ಮಾಡಿ ಹೇಳುವ ಈ ಶಬ್ದಗಳಿಗೆ ಅನುಕರಣಾವ್ಯಯಗಳೆಂದು ಹೆಸರು.\par

\vspace{4pt plus 4pt minus 1pt}\noindent{\kbvsecfont\fontsize{11pt}{14.5pt}\selectfont\color{kbvaccentdark} (೩)\hspace{0.45em}ಭಾವಸೂಚಕಾವ್ಯಯ :}\hspace{0.5em}ಕೋಪ, ದುಃಖ, ಆಶ್ಚರ್ಯ, ತಿರಸ್ಕಾರ, ಸಂತೋಷ, ಮೆಚ್ಚು ಇತ್ಯಾದಿ ಭಾವಗಳನ್ನು ಪ್ರದರ್ಶಿಸುವಾಗ ಎಲೆಲಾ! ಅಯ್ಯೋ! ಅಬ್ಬಾ! ಥೂ! ಓಹೋ! ಭಲೇ!... ಇವೇ ಮೊದಲಾದ ಅರ್ಥರಹಿತವಾದ ಶಬ್ದಗಳನ್ನು ಪ್ರಕಟಿಸುತ್ತೇವೆ. ಈ ರೀತಿಯ ಶಬ್ದಗಳಿಗೆ ಭಾವಸೂಚಕಾವ್ಯಯಗಳೆಂದು ಹೆಸರು.\par

\kbvsechead{(೪)\hspace{0.45em}ಕ್ರಿಯಾರ್ಥಕಾವ್ಯಯ :}

\kbvsubtitle{ಕ್ರಿಯಾಪದ ಪ್ರಕರಣ ನೋಡಿ}

\vspace{4pt plus 4pt minus 1pt}\noindent{\kbvsecfont\fontsize{11pt}{14.5pt}\selectfont\color{kbvaccentdark} (೫)\hspace{0.45em}ಸಂಬಂಧಾರ್ಥಕಾವ್ಯಯ :}\hspace{0.5em}ಎರಡು ಪದಗಳನ್ನು ಜೋಡಿಸಲು “ಮತ್ತು, ಅಥವಾ, ಊ” ಎಂದು.\par
}
\clearpage

% === tex/0083.tex (Printed Page 81) ===
\renewcommand{\kbvchaptername}{ಅವ್ಯಯಗಳು}
\setcounter{page}{81}
\thispagestyle{kbvpage}
\kbvtypesetpage{\indent ಪದಗಳನ್ನು ಸಾಮಾನ್ಯವಾಗಿ ಬಳಸುತ್ತೇವೆ. ಇದೇ ರೀತಿಯಾಗಿ ಎರಡು ವಾಕ್ಯಗಳನ್ನು ಜೋಡಿಸಲು “ಮತ್ತು, ಆದರೆ, ಅಥವಾ, ಆದ್ದರಿಂದ, ಹಾಗಾದರೆ, ಇನ್ನೂ, ಆದುದರಿಂದ” ಇತ್ಯಾದಿ ಪದಗಳನ್ನು ಬಳಸುತ್ತೇವೆ. ಈ ಬಗೆಯ ಪದಗಳಿಗೆ ಸಂಬಂಧಾರ್ಥಕಾವ್ಯಯಗಳೆಂದು ಹೆಸರು.\par

\vspace{3pt plus 2pt}\noindent\hspace*{1.2em}\kbvlabel{ಉದಾ}\par\vspace{1pt}

\vspace{2.5pt plus 3pt minus 0.5pt}\noindent\hspace*{1.2em}\kbvnum{(೧)}\hspace{0.55em}ಹಸು ಮತ್ತು ಕರು\par

\vspace{2.5pt plus 3pt minus 0.5pt}\noindent\hspace*{1.2em}\kbvnum{(೨)}\hspace{0.55em}ಹಸುವೂ-ಕರುವೂ\par

\vspace{2.5pt plus 3pt minus 0.5pt}\noindent\hspace*{1.2em}\kbvnum{(೩)}\hspace{0.55em}ಅಕ್ಕನೂ ತಂಗಿಯೂ\par

\vspace{2.5pt plus 3pt minus 0.5pt}\noindent\hspace*{1.2em}\kbvnum{(೪)}\hspace{0.55em}ಬೆಲ್ಲ ಅಥವಾ ಸಕ್ಕರೆ\par

\vspace{2.5pt plus 3pt minus 0.5pt}\noindent\hspace*{1.2em}\kbvnum{(೫)}\hspace{0.55em}ಕುದುರೆಯೂ ಬಂದಿತು ; ಗಾಡಿಯೂ ಬಂದಿತು.\par

\vspace{3pt plus 2pt}\noindent\hspace*{1.2em}\kbvlabel{ಉದಾ}\par\vspace{1pt}

\vspace{2.5pt plus 3pt minus 0.5pt}\noindent\hspace*{1.2em}\kbvnum{(೧)}\hspace{0.55em}ಅವಳು ಆಂದ್ರಾಗಿ ತಂದಳು ಮತ್ತು ಅದನ್ನು ಬೀಸಿ ವಳ ಗಡಿಯಿಂದ ತಂದಳು ಮತ್ತು ದನ್ನು ಬಿಹಟ್ಟು ಮಾಡಿದಳು.\par

\vspace{2.5pt plus 3pt minus 0.5pt}\noindent\hspace*{1.2em}\kbvnum{(೨)}\hspace{0.55em}ರೈಲು ನಿಲ್ದಾಣಕ್ಕೆ ಬನ್ನಿ ಅಥವಾ ನೇರವಾಗಿ ಕಲ್ಯಾಣ ಮಂಟಪಕ್ಕೆ ಬನ್ನಿ.\par

\vspace{2.5pt plus 3pt minus 0.5pt}\noindent\hspace*{1.2em}\kbvnum{(೩)}\hspace{0.55em}ಮಕ್ಕಳನ್ನು ಮುದ್ದಿಸಬೇಕು ಆದರೆ ಅತಿಯಾಗಿ ಮುದ್ದಿಸುವುದು ಸರ್ವಥಾ ತಪ್ಪು.\par

\vspace{2.5pt plus 3pt minus 0.5pt}\noindent\hspace*{1.2em}\kbvnum{(೪)}\hspace{0.55em}ರೈಲು ಬಹಳ ತಡವಾಗಿ ಬಂದಿತು ಆದುದರಿಂದ ನಾನು ಪ್ರಯಾಣವನ್ನು ಮುಂದಕ್ಕೆ ಹಾಕಿದೆ.\par

\vspace{2.5pt plus 3pt minus 0.5pt}\noindent\hspace*{1.2em}\kbvnum{(೫)}\hspace{0.55em}ನೀನು ಪ್ರಯಾಣದ ಖರ್ಚು ಕೊಡಬೇಕು ಹಾಗಾದರೆ ನಾನು ನಿನ್ನೊಂದಿಗೆ ಬರುತ್ತೇನೆ.\par

\vspace{2.5pt plus 3pt minus 0.5pt}\noindent\hspace*{1.2em}\kbvnum{(೬)}\hspace{0.55em}ನೀನು ತುಂಬಾ ಸಹಾಯ ಮಾಡಿದ್ದೀಯೆ ಇನ್ನೂ ಹೆಚ್ಚಿನ ಸಹಾಯ ನಿನ್ನಿಂದ ಆಗಬೇಕಾಗಿದೆ.\par

\kbvsechead{(೬)\hspace{0.45em}ಕೃದಂತಾವ್ಯಯ :}

\kbvsubtitle{ಕೃದಂತ ಪ್ರಕರಣ ನೋಡಿ}

\kbvsechead{(೭)\hspace{0.45em}ತದ್ದಿತಾಂತಾವ್ಯಯ :}

\kbvsubtitle{ತದ್ದಿತಾಂತ ಪ್ರಕರಣ ನೋಡಿ}
}
\clearpage

% === tex/0084.tex (Printed Page 82) ===
\renewcommand{\kbvchaptername}{ಅವ್ಯಯಗಳು}
\setcounter{page}{82}
\thispagestyle{kbvpage}
\kbvtypesetpage{\kbvsechead{(8)\hspace{0.45em}ಅವಧಾರಣಾರ್ಥಕಾವ್ಯಯ :}

\indent ನಿಶ್ಚಯಾರ್ಥದಲ್ಲಿ ಬರುವ ಅವ್ಯಯಗಳಿಗೆ ಅವಧಾರಣಾರ್ಥಕಾವ್ಯಯಗಳೆಂದು ಹೆಸರು.\par

\vspace{3pt plus 2pt}\noindent\hspace*{1.2em}\kbvlabel{ಉದಾ}\par\vspace{1pt}

\begin{kbvderiv}
\vspace{4pt plus 3pt minus 1pt}
\begin{center}\begin{adjustbox}{max width=\textwidth}
{\renewcommand{\arraystretch}{1.22}\setlength{\tabcolsep}{4pt}
\begin{tabular}{l c l}
ಅದೇ ಶಿವಗಂಗೆ (ಅದು + ಏ & {\color{kbvgrey}$=$} & \kbvresult{ಅದೇ)} \\
ನಾನೇ ಕಣ್ಣಾರೆ ನೋಡಿದೆ. ...(ನಾನು + ಏ & {\color{kbvgrey}$=$} & \kbvresult{ನಾನೇ)} \\
ಶಿವನೇ ನನ್ನ ರಕ್ಷಕ ...(ಶಿವನು + ಏ & {\color{kbvgrey}$=$} & \kbvresult{ಶಿವನೇ)} \\
ಅವಳೇ ಭಾಗ್ಯಶಾಲಿ (ಅವಳು + ಏ & {\color{kbvgrey}$=$} & \kbvresult{ಅವಳೇ)} \\
\end{tabular}}
\end{adjustbox}\end{center}
\vspace{4pt plus 3pt minus 1pt}
\end{kbvderiv}

\indent ಇಲ್ಲಿ ಅದೇ, ನಾನೇ, ಶಿವನೇ, ಅವಳೇ ಎಂಬ ಪದಗಳಲ್ಲಿ "ಏ" ಕಾರಣವಾಗಿ, ಅವಧಾರಣಾರ್ಥಕಾವ್ಯಯ. ಈ ಪದಗಳ ಪ್ರಯೋಗದಲ್ಲಿ ನಿಶ್ಚಯ ಆಗುವುದು. ನಿರ್ಣಯ ಕಂಡುಬರುವುದು.\par

\kbvexercise{ಅಭ್ಯಾಸ ಪತ್ರಿಕೆ ೧೬}

\vspace{2.5pt plus 3pt minus 0.5pt}\noindent\hspace*{1.2em}\kbvnum{1.}\hspace{0.55em}ಅವ್ಯಯ ಎಂದರೇನು ? ಉದಾಹರಣೆಗಳ ಸಹಿತ ತಿಳಿಸಿ.\par

\vspace{2.5pt plus 3pt minus 0.5pt}\noindent\hspace*{1.2em}\kbvnum{2.}\hspace{0.55em}“ಹುಡುಗರು ಸುಮ್ಮನೆ ಬಂದರು” — ಈ ವಾಕ್ಯದಲ್ಲಿ ವಚನಗಳನ್ನು ಬದಲಾಯಿಸಿ ಬರೆಯಿರಿ.\par

\vspace{2.5pt plus 3pt minus 0.5pt}\noindent\hspace*{1.2em}\kbvnum{3.}\hspace{0.55em}ಅವ್ಯಯದಲ್ಲಿ ಎಷ್ಟು ಪ್ರಕಾರಗಳಿವೆ ? ಅವು ಯಾವುವು ?\par

\vspace{2.5pt plus 3pt minus 0.5pt}\noindent\hspace*{1.2em}\kbvnum{4.}\hspace{0.55em}ಇಲ್ಲಿ ಹಲವು ವಾಕ್ಯಗಳನ್ನು ಕೊಟ್ಟಿದೆ. ಪ್ರತಿಯೊಂದ ರಲ್ಲೂ ಯಾವ ಪ್ರಕಾರದ ಅವ್ಯಯವಿರುವುದನ್ನು ತಿಳಿದು ಹೇಳಿ.\par

\vspace{2.5pt plus 3pt minus 0.5pt}\noindent\hspace*{1.2em}\kbvnum{(1)}\hspace{0.55em}ದ್ರಾಕ್ಷಿಗೊಂಚಲು ಸುಲಭವಾಗಿ ಕೈಗೆಟುಕುತ್ತದೆ.\par

\vspace{2.5pt plus 3pt minus 0.5pt}\noindent\hspace*{1.2em}\kbvnum{(2)}\hspace{0.55em}ಗಾಡಿಯು ತಟ್ಟನೆ ನಿಂತಿತು.\par

\vspace{2.5pt plus 3pt minus 0.5pt}\noindent\hspace*{1.2em}\kbvnum{(3)}\hspace{0.55em}ಗಾಳಿಯು ಭರ್ರನೆ ಬೀಸಿತು.\par

\vspace{2.5pt plus 3pt minus 0.5pt}\noindent\hspace*{1.2em}\kbvnum{(4)}\hspace{0.55em}ಅಯ್ಯೋ! ಮಗು ರೈಲಿಗೆ ಸಿಕ್ಕಿತು.\par

\vspace{2.5pt plus 3pt minus 0.5pt}\noindent\hspace*{1.2em}\kbvnum{(5)}\hspace{0.55em}ಅವನಿಂದ ಹಣ ಬರಬೇಕು.\par

\vspace{2.5pt plus 3pt minus 0.5pt}\noindent\hspace*{1.2em}\kbvnum{(6)}\hspace{0.55em}ರಾಮ ಮತ್ತು ಕೃಷ್ಣ ಅವತಾರ ಪುರುಷರು.\par

\vspace{2.5pt plus 3pt minus 0.5pt}\noindent\hspace*{1.2em}\kbvnum{(7)}\hspace{0.55em}ನನಗೆ ಬೇಕಾಗಿರುವುದು ಕಾಫಿ ಅಥವಾ ಟೀ.\par

\vspace{2.5pt plus 3pt minus 0.5pt}\noindent\hspace*{1.2em}\kbvnum{(8)}\hspace{0.55em}ನೂಲಿನಂತೆ ಸೀರೆ ತಾಯಿಯಂತೆ ಮಗಳು.\par
}
\clearpage

% === tex/0085.tex (Printed Page 83) ===
\renewcommand{\kbvchaptername}{ದ್ವಿರುಕ್ತಿ}
\setcounter{page}{83}
\thispagestyle{kbvpage}
\kbvtypesetpage{\vspace{2.5pt plus 3pt minus 0.5pt}\noindent\hspace*{1.2em}\kbvnum{(೯)}\hspace{0.55em}ಬೆಳಕು ಬರಲು ಕತ್ತಲೆ ಹೋಯಿತು.\par

\vspace{2.5pt plus 3pt minus 0.5pt}\noindent\hspace*{1.2em}\kbvnum{(೧೦)}\hspace{0.55em}ಅವಳು ಬೆಂಕಿಯ ಜ್ವಾಲೆಯಲ್ಲಿ ಧಗಧಗ ಉರಿದು ಹೋದಳು.\par

\vspace{2.5pt plus 3pt minus 0.5pt}\noindent\hspace*{1.2em}\kbvnum{(೧೧)}\hspace{0.55em}ಕಾದ ಎಣ್ಣೆಯಲ್ಲಿ ಸಾಸಿವೆ ಚಟಚಟ ಎನ್ನುವುದು.\par

\vspace{2.5pt plus 3pt minus 0.5pt}\noindent\hspace*{1.2em}\kbvnum{(೧೨)}\hspace{0.55em}ರಾಮನೂ ಸೀತೆಯೂ ಲಕ್ಷ್ಮಣನೂ ಕಾಡಿಗೆ ಹೋದರು.\par

\vspace{2.5pt plus 3pt minus 0.5pt}\noindent\hspace*{1.2em}\kbvnum{(5)}\hspace{0.55em}ಎರಡು ವಾಕ್ಯಗಳನ್ನು ಜೋಡಿಸಲು “ಆದ್ದರಿಂದ” ಎಂಬ ಅವ್ಯಯವನ್ನು ಹೇಗೆ ಬಳಸಬಹುದೆಂಬುದಕ್ಕೆ ಒಂದು ಉದಾಹರಣೆ ಕೊಡಿ.\par

\vspace{2.5pt plus 3pt minus 0.5pt}\noindent\hspace*{1.2em}\kbvnum{(6)}\hspace{0.55em}ಅವಧಾರಣಾರ್ಥಕಾವ್ಯಯ ಎಂದರೇನು? ಎರಡು ಉದಾಹರಣೆ ಗಳನ್ನು ಕೊಡಿ.\par

\kbvchapter{ಪಾಠ ೧೭}{ದ್ವಿರುಕ್ತಿ}

\indent ದ್ವಿರುಕ್ತಿ ಎಂದರೆ ಎರಡು ಬಾರಿ ಹೇಳುವುದು ಎಂದರ್ಥ. ಸಂದರ್ಭಾನು ಸಾರವಾಗಿ ಸನ್ನಿವೇಶಕ್ಕೆ ಒತ್ತುಕೊಡುವ ಉದ್ದೇಶದಿಂದ ಒಂದೇ ಶಬ್ದವನ್ನು ಎಡಬಿಡದೆ ಎರಡು ಬಾರಿ ಹೇಳುವುದು ರೂಢಿಯಲ್ಲಿದೆ. ಇದನ್ನು ದ್ವಿರುಕ್ತಿ ಎನ್ನುವರು.\par

\vspace{3pt plus 2pt}\noindent\hspace*{1.2em}\kbvlabel{ಉದಾ}\par\vspace{1pt}

\noindent\hspace*{2.2em}ರೈಲು ಹತ್ತಲು ಅವನು ಓಡಿಓಡಿ ಬಂದನು.\par

\noindent\hspace*{2.2em}ಹೌದು ! ಹೌದು ! ಅವನೇ ಸರ ಕದ್ದವನು.\par

\noindent\hspace*{2.2em}ಬೇಡಬೇಡ, ನನಗೆ ನಿನ್ನ ಸಹವಾಸವೇ ಬೇಡ.\par

\noindent\hspace*{2.2em}ತಿಂದೂತಿಂದೂ ಸಾಕಾಯಿತು.\par

\noindent\hspace*{2.2em}ನೀನು ನನ್ನನ್ನು ಆಗಾಗ ಬಂದು ಕಾಣುತ್ತಿರು.\par

\noindent\hspace*{2.2em}ಊರೂರು ಸುತ್ತಿ ನನಗೆ ಆಯಾಸವಾಗಿದೆ.\par

\indent ಅಬ್ಬಬ್ಬಾ ! ಎಷ್ಟು ದೊಡ್ಡ ಹಾವು. ಇಲ್ಲಿ ಬನ್ನಿ ಇಲ್ಲಿ ಬನ್ನಿ, ನನ್ನ ಮಾತು ಕೇಳಿ. ಈ ಮಕ್ಕಳಿಗೆ ವಿವೇಕ ಹೇಳಿಹೇಳಿ ಸಾಕಾಯಿತು.\par

\noindent\hspace*{2.2em}ಅತ್ತೂ, ಅತ್ತೂ ಕಣ್ಣು ಕೆಂಪಗಾಯಿತು.\par

\indent ಈ ವಾಕ್ಯಗಳಲ್ಲಿ ದಪ್ಪಕ್ಷರಗಳಲ್ಲಿರುವ ಪದಗಳು ದ್ವಿರುಕ್ತಿಗಳು.\par
}
\clearpage

% === tex/0086.tex (Printed Page 84) ===
\renewcommand{\kbvchaptername}{ತತ್ಸಮ-ತದ್ಭವ}
\setcounter{page}{84}
\thispagestyle{kbvpage}
\kbvtypesetpage{\kbvexercise{ಅಭ್ಯಾಸ ಪತ್ರಿಕೆ ೧೭}

\vspace{2.5pt plus 3pt minus 0.5pt}\noindent\hspace*{1.2em}\kbvnum{(1)}\hspace{0.55em}ದ್ವಿರುಕ್ತಿ ಎಂದರೇನು ? ಉದಾಹರಣೆಗಳ ಸಹಿತ ತಿಳಿಸಿ,\par

\vspace{2.5pt plus 3pt minus 0.5pt}\noindent\hspace*{1.2em}\kbvnum{(2)}\hspace{0.55em}ಕೆಳಗೆ ಕೊಟ್ಟಿರುವ ದ್ವಿರುಕ್ತಿಗಳನ್ನು ನಿಮ್ಮ ಸ್ವಂತ ಮಾತುಗಳು\par

\indent (ವಾಕ್ಯಗಳಲ್ಲಿ) ಬಳಸಿರಿ :\par

\indent ಬೇಗಬೇಗ ಈಗೀಗ ಹೆಚ್ಚುಹೆಚ್ಚು ನಡೆನಡೆ\par

\indent ನಿಲುನಿಲು ಮನೆಮನೆ  ಕಮ್ಮಿಕಮ್ಮಿ ಭಲೇಭಲೇ\par

\vspace{2.5pt plus 3pt minus 0.5pt}\noindent\hspace*{1.2em}\kbvnum{(3)}\hspace{0.55em}ಭಾವಸೂಚಕಾವ್ಯಯಗಳನ್ನು ದ್ವಿರುಕ್ತಿಗಳಾಗಿ ಬಳಸಬಹುದೇ?\par

\vspace{2.5pt plus 3pt minus 0.5pt}\noindent\hspace*{1.2em}\kbvnum{(4)}\hspace{0.55em}“ಠೂ ! ಠೂ ! ನೀನೊಬ್ಬ ಅಯೋಗ್ಯ”\par

\indent —ಇಲ್ಲಿ ಠೂ ! ಠೂ ! ಎಂಬುದು ದ್ವಿರುಕ್ತಿಯೇ ? ಅಥವಾ ಅವ್ಯಯವೇ ?\par

\vspace{2.5pt plus 3pt minus 0.5pt}\noindent\hspace*{1.2em}\kbvnum{(5)}\hspace{0.55em}“ಸತ್ಯವೇ ತಾಯಿ ತಂದೆ; ಸತ್ಯವೇ ಬಂಧು ಬಳಗ”—ಇಲ್ಲಿ ಸತ್ಯವೇ ಎಂಬ ಶಬ್ದವು ಪುನರಾವರ್ತನೆಯಾಗಿದೆ. ಇದು ದ್ವಿರುಕ್ತಿಯೇ?\par

\kbvchapter{ಪಾಠ ೧೮}{ತತ್ಸಮ-ತದ್ಭವ}

\indent ಕನ್ನಡ ಭಾಷೆಯಲ್ಲಿ ಪದ ಸಂಪತ್ತು ಬರುಬರುತ್ತಾ ಹೆಚ್ಚುತ್ತಿದೆ. ಅನ್ಯ ಭಾಷಾಪದಗಳ ಕನ್ನಡೀಕರಣ ಈ ಬೆಳವಣಿಗೆಗೆ ಬಹು ಮುಖ್ಯ ಕಾರಣ, ಆಂಗ್ಲ ಭಾಷೆಯ ಪದಗಳು, ಫಾರಸಿ ಭಾಷೆಯ ಪದಗಳು, ಸಂಸ್ಕೃತ ಪದಗಳು ಮತ್ತು ಇನ್ನಿತರ ಭಾರತೀಯ ಭಾಷೆಗಳಿಂದ ಬಂದ ಪದಗಳು ಕನ್ನಡದೊಂದಿಗೆ ಸೇರಿಕೊಂಡು ಕನ್ನಡ ಪದಗಳೇ ಆಗಿಬಿಟ್ಟಿವೆ. ಅವುಗಳಲ್ಲಿ ಸಂಸ್ಕೃತ ಪದಗಳ ಮಾತ್ರ ಹೆಚ್ಚು ಸಂಖ್ಯೆಯಲ್ಲಿವೆ,\par

\indent ಅಪ್ಪಟ ಸಂಸ್ಕೃತ ಪದಗಳನ್ನೂ ಒಂದಿಷ್ಟು ಮಾರ್ಪಾಡು ಹೊಂದಿರುವ ಸಂಸ್ಕೃತ ಪದಗಳನ್ನೂ ಕನ್ನಡದಲ್ಲಿ ಕಾಣಬಹುದು. ಯಾವುದೊಂದು ಬದ ಲಾವಣೆಯೂ ಆಗದೇ ಕನ್ನಡದಲ್ಲಿ ಬಳಕೆಗೆ ಬಂದಿರುವ ಸಂಸ್ಕೃತ ಪದಗಳಿಗೆ “ತತ್ಸಮ”ಗಳೆಂದು ಹೆಸರು; ಸಾಕಷ್ಟು ಮಟ್ಟಿಗೆ ಬದಲಾವಣೆ ಹೊಂದಿ ಕನ್ನಡ ಪದಗಳ ಗುಂಪಿಗೆ ಸೇರಿರುವ ಸಂಸ್ಕೃತ ಪದಗಳಿಗೆ “ತದ್ಭವ”ಗಳೆಂದು ಹೆಸರು.\par

\indent ಉದಾಹರಣೆಗೆ ಸ್ತ್ರೀ-ಅಪ್ಪಟ ಸಂಸ್ಕೃತ ಪದ. ಇದನ್ನು ಕನ್ನಡದಲ್ಲಿ ನೇರವಾಗಿ ಬಳಸುತ್ತೇವೆ; ಇದೊಂದು ತತ್ಸಮ ಶಬ್ದ. ಕರ್ತರಿ-ಅಪ್ಪಟ ಸಂಸ್ಕೃತ\kbvcontinues
}
\clearpage

% === tex/0087.tex (Printed Page 85) ===
\renewcommand{\kbvchaptername}{ಲೇಖನ ಚಿಹ್ನೆಗಳು}
\setcounter{page}{85}
\thispagestyle{kbvpage}
\kbvtypesetpage{\kbvcontinued ಪದ. ಇದು “ಕತ್ತರಿ”ಯಾಗಿ ಮಾರ್ಪಾಡು ಹೊಂದಿ ಕನ್ನಡದಲ್ಲಿ ಬಳಕೆಗೆ ಬಂದಿದೆ ; ಆದ್ದರಿಂದ ಇದನ್ನು ಕರ್ತರಿ ಶಬ್ದದ ತದ್ಭವ ರೂಪ ಎನ್ನುತ್ತೇವೆ. ಹೆಚ್ಚಿನ ಉದಾಹರಣೆಗಳನ್ನು ನೋಡೋಣ.\par

\vspace{5pt plus 4pt minus 1pt}
\begin{center}
{\renewcommand{\arraystretch}{1.22}\setlength{\tabcolsep}{5.5pt}
\begin{tabularx}{\textwidth}{X X | X X | X X}
{\kbvsecfont ತತ್ಸಮ} & {\kbvsecfont ತದ್ಭವ} & {\kbvsecfont ತತ್ಸಮ} & {\kbvsecfont ತದ್ಭವ} & {\kbvsecfont ತತ್ಸಮ} & {\kbvsecfont ತದ್ಭವ} \\[3pt]
ಆಶಾ & ಆಸೆ & ಗಂಗಾ & ಗಂಗೆ & ರಾಜನ್‌ & ರಾಜ \\
ಗ್ರಂಥಿ & ಗಂಟು & ಯೋಗಿ & ಜೋಗಿ & ಕುಠಾರ & ಕೊಡಲಿ \\
ಚಂದ್ರ & ಚಂದಿರ & ಶೂನ್ಯ & ಸೊನ್ನೆ & ಶರ್ಕರ & ಸಕ್ಕರೆ \\
ಧೃತಿ & ದಠಿ & ಸಂಸ್ಥಾ & ಸಂಸ್ಥೆ & ಕಾಮಿನೀ & ಕಾಮಿನಿ \\
ಬ್ರಹ್ಮ & ಬೊಮ್ಮ & ಮುಕ್ತಿ & ಮುಕುತಿ & ತೃತೀಯಾ & ತದಿಗೆ \\
ಮಾಲಾ & ಮಾಲೆ & ಘಂಟಾ & ಗಂಟೆ & ಯಶಸ್ & ಯಶಸ್ಸು \\
ಯಮ & ಜವ & ವಸತಿ & ಬಸದಿ & ರಾಕ್ಷಸ & ರಕ್ಕಸ \\
ವೀಥಿ & ಬೀದಿ & ಚತುರ & ಚದುರ & ವಿಸ್ತಾರ & ಬಿತ್ತರ \\
\end{tabularx}}
\end{center}
\vspace{5pt plus 4pt minus 1pt}

\kbvexercise{ಅಭ್ಯಾಸ ಪತ್ರಿಕೆ ೧೮}

\vspace{2.5pt plus 3pt minus 0.5pt}\noindent\hspace*{1.2em}\kbvnum{(1)}\hspace{0.55em}ತತ್ಸಮ, ತದ್ಭವ—ಈ ಪದಗಳ ಅರ್ಥವನ್ನು ವಿವರಿಸಿ.\par

\vspace{2.5pt plus 3pt minus 0.5pt}\noindent\hspace*{1.2em}\kbvnum{(2)}\hspace{0.55em}ಕೆಳಗೆ ಕಾಣಿಸಿರುವ ಪದಗಳ ತತ್ಸಮ ರೂಪಗಳನ್ನು ತಿಳಿಸಿ : ಕನ್ನಿಕೆ, ಸಿಂಗರ, ಪಯಣ, ನೇಹ, ಸಿಂಗ, ಲಕ್ಷಿ, ಸೀತೆ, ಲತೆ.\par

\vspace{2.5pt plus 3pt minus 0.5pt}\noindent\hspace*{1.2em}\kbvnum{(3)}\hspace{0.55em}ಕೆಳಗೆ ಕಾಣಿಸಿರುವ ಪದಗಳ ತದ್ಭವ ರೂಪಗಳನ್ನು ತಿಳಿಸಿ : ವಿದ್ವಾಂಸ, ಚತುರ್ಥೀ, ಕುಮಾರೀ, ಶಿರಸ್, ಸುಧಾ, ವಸುಂಧರಾ.\par

\vspace{2.5pt plus 3pt minus 0.5pt}\noindent\hspace*{1.2em}\kbvnum{(4)}\hspace{0.55em}ಕನ್ನಡದಲ್ಲಿ ಬಳಕೆಯಲ್ಲಿರುವ ಹತ್ತು ಅನ್ಯದೇಶೀಯ ಪದಗಳನ್ನು ಹೆಸರಿಸಿ. ಉದಾ : ದವಾಖಾನೆ, ಪಾದ್ರಿ.... ಇತ್ಯಾದಿ.\par

\kbvchapter{ಪಾಠ ೧೯}{ಲೇಖನ ಚಿಹ್ನೆಗಳು}

\vspace{5pt plus 4pt minus 1pt}
\noindent{\renewcommand{\arraystretch}{1.24}\setlength{\tabcolsep}{4pt}
\begin{tabularx}{\textwidth}{X X X X}
{\kbvsecfont\color{kbvaccentdark}ಪೂರ್ಣ ವಿರಾಮ ಚಿಹ್ನೆ} & {\kbvsecfont\color{kbvaccentdark}{[}.{]}} & {\kbvsecfont\color{kbvaccentdark}ಭಾವಸೂಚಕ ಚಿಹ್ನೆ} & {\kbvsecfont\color{kbvaccentdark}{[}!{]}} \\[1pt]
\arrayrulecolor{kbvline}\specialrule{0.5pt}{0pt}{2pt}\arrayrulecolor{black}
ಅರ್ಧ ವಿರಾಮ ಚಿಹ್ನೆ & {[};{]} & ಉದ್ಧರಣ ಚಿಹ್ನೆ & {[}“ ”{]} \\
\end{tabularx}}
\vspace{5pt plus 4pt minus 1pt}
}
\clearpage

% === tex/0088.tex (Printed Page 86) ===
\renewcommand{\kbvchaptername}{ಲೇಖನ ಚಿಹ್ನೆಗಳು}
\setcounter{page}{86}
\thispagestyle{kbvpage}
\kbvtypesetpage{\indent ಅಲ್ಪ ವಿರಾಮ ಚಿಹ್ನೆ {[},{]} ವಿವರಣಾತ್ಮಕ ಚಿಹ್ನೆ {[}:{]} ಪ್ರಶ್ನಾರ್ಥಕ ಚಿಹ್ನೆ {[}?{]} ಆವರಣ ಚಿಹ್ನೆ ( )\par

\indent ಇವು ಲೇಖನ ಚಿಹ್ನೆಗಳು. ಬರವಣಿಗೆಯಲ್ಲಿ ಈ ಚಿಹ್ನೆಗಳು ಎಷ್ಟು ಪ್ರಾಮುಖ್ಯವೆಂಬುದನ್ನು ತೋರಿಸುವ ಒಂದು ದೃಷ್ಟಾಂತ ಹೀಗಿದೆ—\par

\indent “{\kbvlatin Hang} {\kbvlatin him} {\kbvlatin not}, {\kbvlatin spare} {\kbvlatin him}” “ನೇಣು ಹಾಕಬೇಡಿ ಅವನನ್ನು ಉಳಿಸಿ” ಎಂಬರ್ಥ ಕೊಡುವ ವಾಕ್ಯದಲ್ಲಿ ಅಲ್ಪವಿರಾಮ ಚಿಹ್ನೆಯನ್ನು ಎಲ್ಲಿ ಬಳಸಿದೆ ಎಂಬುದನ್ನು ಗಮನಿಸಿ. ಈ ಚಿಹ್ನೆಯು ಕೈ ತಪ್ಪಿನಿಂದ ಸ್ಥಾನ ಪಲ್ಲಟವಾದಲ್ಲಿ ಎಂತಹ ಪ್ರಮಾದವಾಗುವು ದೆಂಬುದನ್ನು ನೋಡೋಣ.\par

\indent “{\kbvlatin Hang} {\kbvlatin him}, {\kbvlatin not} {\kbvlatin spare} {\kbvlatin him}” “ನೇಣು ಹಾಕಿ ಅವನನ್ನು ಉಳಿಸಬೇಡಿ” ಎನ್ನುವ ತೀರ್ಮಾನಕ್ಕೆ ಕಾರಣ ಅಜಾಗರೂಕತೆಯಿಂದ ಬಳಸಿದ ಅಲ್ಪವಿರಾಮ ಚಿಹ್ನೆ. ದೃಷ್ಟಾಂತವೊಂದರಿಂದಲೇ ಬರವಣಿಗೆಯಲ್ಲಿ ಲೇಖನ ಚಿಹ್ನೆಗಳು ಎಷ್ಟು ಮುಖ್ಯ ವೆಂಬುದು ತಿಳಿದುಬರುತ್ತದೆ, ಆದುದರಿಂದ “ಅವುಗಳ ಬಳಕೆಯ ಬಗೆ ಉತ್ತಮ ತಿಳುವಳಿಕೆ ಇರಬೇಕಾದ್ದು” ಅತ್ಯಾವಶ್ಯಕ.\par

\vspace{4pt plus 4pt minus 1pt}\noindent{\kbvsecfont\fontsize{11pt}{14.5pt}\selectfont\color{kbvaccentdark} (೧)\hspace{0.45em}ಪೂರ್ಣ ವಿರಾಮ ಚಿಹ್ನೆ:}\hspace{0.5em}ವಾಕ್ಯದ ಅಂತ್ಯದಲ್ಲಿ ಪೂರ್ಣ ವಿರಾಮ ಚಿಹ್ನೆಯನ್ನು ಇಡಬೇಕು. ಸಾಮಾನ್ಯವಾಗಿ ಕ್ರಿಯಾಪದದ ನಂತರ ಈ ಚಿಹ್ನೆ ಬರುವುದು.\par

\noindent\hspace*{1.2em}\kbvlabel{ಉದಾ}\hspace{0.6em}ಶ್ರೀರಂಗಪಟ್ಟಣವು ಕಾವೇರಿ ನದಿಯ ತೀರದಲ್ಲಿದೆ.\par

\vspace{4pt plus 4pt minus 1pt}\noindent{\kbvsecfont\fontsize{11pt}{14.5pt}\selectfont\color{kbvaccentdark} (೨)\hspace{0.45em}ಅರ್ಧ ವಿರಾಮ ಚಿಹ್ನೆ:}\hspace{0.5em}ಒಂದು ವಾಕ್ಯವು ಮತ್ತೊಂದು ವಾಕ್ಯಕ್ಕೆ ಪೂರಕವಾಗಿರುವುದುಂಟು, ಅಂತಹ ಸಂದರ್ಭದಲ್ಲಿ ಆ ಎರಡು ವಾಕ್ಯಗಳ ನಡುವೆ ಅರ್ಧ ವಿರಾಮ ಚಿಹ್ನೆ ಬಳಸಬೇಕು.\par

\begin{kbvderiv}
\vspace{4pt plus 3pt minus 1pt}
\noindent{\renewcommand{\arraystretch}{1.24}\setlength{\tabcolsep}{5pt}
\begin{tabularx}{\textwidth}{X X}
ಉದಾ: ಅನ್ನ ಹಳಸಿತ್ತು & ಆದುದರಿಂದ ನಾನು ಊಟ ಮಾಡಲಿಲ್ಲ, \\
ನಾನು ಔಷಧಿ ತೆಗೆದುಕೊಂಡೆ & ಆದರೂ ರೋಗ ಗುಣವಾಗಲಿಲ್ಲ. \\
\end{tabularx}}
\vspace{4pt plus 3pt minus 1pt}
\end{kbvderiv}
}
\clearpage

% === tex/0089.tex (Printed Page 87) ===
\renewcommand{\kbvchaptername}{ಲೇಖನ ಚಿಹ್ನೆಗಳು}
\setcounter{page}{87}
\thispagestyle{kbvpage}
\kbvtypesetpage{\vspace{4pt plus 4pt minus 1pt}\noindent{\kbvsecfont\fontsize{11pt}{14.5pt}\selectfont\color{kbvaccentdark} (೩)\hspace{0.45em}ಅಲ್ಪ ವಿರಾಮ ಚಿಹ್ನೆ:}\hspace{0.5em}({\kbvlatin i}) ಸಂಬೋಧನಾರ್ಥಕ ಶಬ್ದದ ಮುಂದೆ, ({\kbvlatin ii}) ಕರ್ತೃ, ಕರ್ಮ, ಕ್ರಿಯಾಪದಗಳಿಗೆ ಒಂದಕ್ಕಿಂತ ಹೆಚ್ಚು ವಿಶೇಷಣಗಳಿರುವಾಗ, ({\kbvlatin iii}) ಒಂದೇ ಕ್ರಿಯೆಗೆ ಒಂದಕ್ಕಿಂತ ಹೆಚ್ಚು ಕರ್ತೃ ಪದಗಳಿರುವಾಗ, ({\kbvlatin iv}) ಒಂದೇ ಕ್ರಿಯೆಗೆ ಒಂದಕ್ಕಿಂತ ಹೆಚ್ಚು ಕರ್ಮಪದ ಗಳಿರುವಾಗ ಅಲ್ಪ ವಿರಾಮ ಚಿಹ್ನೆಯನ್ನು ಕೆಳಗೆ ಕೊಟ್ಟಿರುವ ಉದಾಹರಣೆ ಗಳಲ್ಲಿ ತೋರಿಸಿರುವಂತೆ ಬಳಸಬೇಕು.\par

\vspace{3pt plus 2pt}\noindent\hspace*{1.2em}\kbvlabel{ಉದಾ}\par\vspace{1pt}

\vspace{2.5pt plus 3pt minus 0.5pt}\noindent\hspace*{1.2em}\kbvnum{(1)}\hspace{0.55em}ಪದ್ಮಲತಾ, ಬೀಗದ ಕೈ ಜೋಪಾನವಾಗಿಡು.\par

\vspace{2.5pt plus 3pt minus 0.5pt}\noindent\hspace*{1.2em}\kbvnum{(2)}\hspace{0.55em}ಕನ್ಯೆಯು ರೂಪವತಿಯೂ, ಗುಣವತಿಯೂ, ವಿದ್ಯಾವತಿಯೂ ಆಗಿದ್ದಾಳೆ.\par

\vspace{2.5pt plus 3pt minus 0.5pt}\noindent\hspace*{1.2em}\kbvnum{(3)}\hspace{0.55em}ನಮ್ಮ ಅಂಗಡಿಯಲ್ಲಿ ಶುದ್ಧವಾದ, ಹಸನಾದ, ಮಂಡ್ಯದ ಬೆಣ್ಣೆ ಮಾರಾಟಕ್ಕಿದೆ.\par

\vspace{2.5pt plus 3pt minus 0.5pt}\noindent\hspace*{1.2em}\kbvnum{(4)}\hspace{0.55em}ಚಪಾತಿಯು ಬಿಸಿಯಾಗಿಯೂ, ಮೃದುವಾಗಿಯೂ, ರುಚಿ ಯಾಗಿಯೂ ಇರಬೇಕು.\par

\vspace{2.5pt plus 3pt minus 0.5pt}\noindent\hspace*{1.2em}\kbvnum{(5)}\hspace{0.55em}ಧರ್ಮರಾಯ, ಭೀಮ, ಅರ್ಜುನ, ನಕುಲ, ಸಹದೇವ ಇವರೈ ವರೂ ದ್ರೌಪದಿಯನ್ನು ವಿವಾಹವಾದರು.\par

\vspace{2.5pt plus 3pt minus 0.5pt}\noindent\hspace*{1.2em}\kbvnum{(6)}\hspace{0.55em}ನಮ್ಮ ಕೃಷಿ ಕ್ಷೇತ್ರದಲ್ಲಿ ಕಬ್ಬು, ತೆಂಗು, ಭತ್ತ, ರಾಗಿ, ತೊಗರಿ ಬೆಳೆಯುತ್ತೇವೆ.\par

\vspace{4pt plus 4pt minus 1pt}\noindent{\kbvsecfont\fontsize{11pt}{14.5pt}\selectfont\color{kbvaccentdark} (೪)\hspace{0.45em}ಪ್ರಶ್ನಾರ್ಥಕ ಚಿಹ್ನೆ:}\hspace{0.5em}ಪ್ರಶ್ನಾರ್ಥಕ ಶಬ್ದದ ಮುಂದೆ, ಪ್ರಶ್ನಾರ್ಥಕ ವಾಕ್ಯದ ಮುಂದೆ ಈ ಚಿಹ್ನೆಯನ್ನು ಬಳಸಬೇಕು.\par

\noindent\hspace*{1.2em}\kbvlabel{ಉದಾ}\hspace{0.6em}ರೈಲು ಈ ದಿನ ತಡವಾಗಿ ಬಂದಿತು ; ಏಕೆ?\par

\noindent\hspace*{2.2em}ನೀನು ಈ ದಿನ ಎಲ್ಲಿಗೆ ಹೋಗುತ್ತೀಯೆ ?\par

\vspace{4pt plus 4pt minus 1pt}\noindent{\kbvsecfont\fontsize{11pt}{14.5pt}\selectfont\color{kbvaccentdark} (೫)\hspace{0.45em}ಭಾವಸೂಚಕ ಚಿಹ್ನೆ:}\hspace{0.5em}ಭಾವಸೂಚಕ ಶಬ್ದದ ಮುಂದೆ, ಭಾವಸೂಚಕ ವಾಕ್ಯದ ಮುಂದೆ ಈ ಚಿಹ್ನೆಯನ್ನು ಬಳಸಬೇಕು.\par
}
\clearpage

% === tex/0090.tex (Printed Page 88) ===
\renewcommand{\kbvchaptername}{ಲೇಖನ ಚಿಹ್ನೆಗಳು}
\setcounter{page}{88}
\thispagestyle{kbvpage}
\kbvtypesetpage{\vspace{3pt plus 2pt}\noindent\hspace*{1.2em}\kbvlabel{ಉದಾ}\par\vspace{1pt}

\indent ಭಲೇ ! ಈಜುವ ಪಂದ್ಯದಲ್ಲಿ ಗೆದ್ದವರಲ್ಲಿ ನೀನೇ ಮೊದಲಿಗ !\par

\noindent\hspace*{2.2em}ಅಲ್ಲಿ ನೋಡು ! ವಿಮಾನ ನೆಲಕ್ಕೆ ಅಪ್ಪಳಿಸಿತು !\par

\noindent\hspace*{2.2em}ಈ ಸೀರೆ ಎಷ್ಟು ಅಂದವಾಗಿದೆ !\par

\vspace{4pt plus 4pt minus 1pt}\noindent{\kbvsecfont\fontsize{11pt}{14.5pt}\selectfont\color{kbvaccentdark} (೬)\hspace{0.45em}ಉದ್ಧರಣ ಚಿಹ್ನೆ :}\hspace{0.5em}ನೇರ ಸಂಭಾಷಣೆಯಲ್ಲಿ ಬರುವ ಮಾತು ಗಳನ್ನು ಉದ್ಧರಿಸಿ ಬರೆಯುವಾಗ, ಮತ್ತೊಬ್ಬರು ಹೇಳಿದುದನ್ನು ಅವರು ಹೇಳಿದಂತೆಯೇ ಬರೆಯುವಾಗ, ಅನ್ಯಭಾಷಾಪದಗಳನ್ನು ಬರವಣಿಗೆಯಲ್ಲಿ ತೋರಿಸಿದಾಗ, ಹೆಚ್ಚು ಮಹತ್ವವುಳ್ಳ ಹಾಗೂ ಗಮನ ಸೆಳೆಯುವ ಶಬ್ದ ಗಳನ್ನು ಬರವಣಿಗೆಯಲ್ಲಿ ತಂದಾಗ ಈ ಚಿಹ್ನೆಯನ್ನು ಬಳಸಬೇಕು.\par

\vspace{3pt plus 2pt}\noindent\hspace*{1.2em}\kbvlabel{ಉದಾ}\par\vspace{1pt}

\vspace{2.5pt plus 3pt minus 0.5pt}\noindent\hspace*{1.2em}\kbvnum{(1)}\hspace{0.55em}ಗಾಂಧೀಜಿಯವರು ಗುಂಡಿನೇಟಿಗೆ ಬಲಿಯಾದಾಗ “ಅತಿಯಾದ ಒಳ್ಳೆಯತನವಿರುವುದು ಅಪಾಯಕಾರಿ” ಎಂದು ಜಾರ್ಜ್‌ ಬರ್ನಾರ್ಡ್‌ ಷಾ ಅಭಿಪ್ರಾಯಪಟ್ಟರು.\par

\vspace{2.5pt plus 3pt minus 0.5pt}\noindent\hspace*{1.2em}\kbvnum{(2)}\hspace{0.55em}ಶ್ರೀ ಎಸ್‌. ಎಲ್‌. ಭೈರಪ್ಪನವರ “ಪರ್ವ” ಪ್ರತಿಯೊಬ್ಬ ಕನ್ನಡಿಗನೂ ಓದಬೇಕಾದ ಪುಸ್ತಕ.\par

\vspace{2.5pt plus 3pt minus 0.5pt}\noindent\hspace*{1.2em}\kbvnum{(3)}\hspace{0.55em}ಉಷ್ಣತೆಯನ್ನು ಅಳೆಯುವ ಸಾಧನ “ಥರ್ಮಾಮೀಟರ.”\par

\vspace{2.5pt plus 3pt minus 0.5pt}\noindent\hspace*{1.2em}\kbvnum{(4)}\hspace{0.55em}“ತಾಳಿದವನು ಬಾಳಿಯಾನು” ಎಂಬ ನಾಣ್ಣುಡಿಯನ್ನು ಅರ್ಥ ಮಾಡಿಕೊಂಡರೆ ಜೀವನದಲ್ಲಿ ಸುಖವುಂಟು.\par

\vspace{4pt plus 4pt minus 1pt}\noindent{\kbvsecfont\fontsize{11pt}{14.5pt}\selectfont\color{kbvaccentdark} (೭)\hspace{0.45em}ವಿವರಣಾತ್ಮಕ ಚಿಹ್ನೆ:}\hspace{0.5em}ತಲೆಬರಹದ ಮುಂದೆ ಈ ಚಿಹ್ನೆಯನ್ನು ಹಾಕುವುದುಂಟು; ಆ ತಲೆಬರಹಕ್ಕೆ ಸಂಬಂಧಿಸಿದ ವಿವರಣೆ ಮುಂದೆ ಕಂಡೆಂದು ಇರುವುದು ಎಂಬ ಅಭಿಪ್ರಾಯವನ್ನು ಈ ಚಿಹ್ನೆ ಸೂಚಿಸುವುದು...\par

\noindent\hspace*{1.2em}\kbvlabel{ಉದಾ}\hspace{0.6em}—\par

\indent ವ್ಯಾಕರಣ ಶಾಸ್ತ್ರ : ಶಬ್ದಗಳ ಶುದ್ಧರೂಪವನ್ನು, ಶಬ್ದಗಳಿಗಿರುವ ಸಂಬಂಧವನ್ನು, ಭಾಷೆಯ ಬೆಳವಣಿಗೆಯ ಸ್ವರೂಪವನ್ನು ತಿಳಿಸುವ ಶಾಸ್ತ್ರ;\par

\noindent\hspace*{2.2em}ಗ್ರಂಥವೇ ವ್ಯಾಕರಣ. ಈ ಶಾಸ್ತ್ರದ...\par
}
\clearpage

% === tex/0091.tex (Printed Page 89) ===
\renewcommand{\kbvchaptername}{ಲೇಖನ ಚಿಹ್ನೆಗಳು}
\setcounter{page}{89}
\thispagestyle{kbvpage}
\kbvtypesetpage{\vspace{4pt plus 4pt minus 1pt}\noindent{\kbvsecfont\fontsize{11pt}{14.5pt}\selectfont\color{kbvaccentdark} (೮)\hspace{0.45em}ಆವರಣ ಚಿಹ್ನೆ :}\hspace{0.5em}ಯಾವುದಾದರೊಂದು ವಾಕ್ಯವನ್ನಾಗಲೀ, ವಾಕ್ಯಾಂಶವನ್ನಾಗಲೀ, ಶಬ್ದವನ್ನಾಗಲೀ ಪ್ರತ್ಯೇಕವಾಗಿ ತಿಳಿಸಬೇಕಾದ ಸಂದರ್ಭವೊದಗಿದಾಗ ಈ ಚಿಹ್ನೆಯನ್ನು ಬಳಸುವುದುಂಟು.\par

\vspace{3pt plus 2pt}\noindent\hspace*{1.2em}\kbvlabel{ಉದಾ}\par\vspace{1pt}

\vspace{2.5pt plus 3pt minus 0.5pt}\noindent\hspace*{1.2em}\kbvnum{(1)}\hspace{0.55em}ಸಿಹಿಮೂತ್ರ ರೋಗ (ಮಧುಮೇಹ) ಪೀಡಿತರು ಅನ್ನವನ್ನು ಆದಷ್ಟು ಕಡಿಮೆ ಸೇವಿಸಬೇಕು,\par

\vspace{2.5pt plus 3pt minus 0.5pt}\noindent\hspace*{1.2em}\kbvnum{(2)}\hspace{0.55em}ಕಂಬಳ (ಕೋಣಗಳ ಓಟದ ಸ್ಪರ್ಧೆ) ಎಂಬ ಶಬ್ದ ಕಿವಿಗೆ ಬಿದ್ದೊಡನೆ ದಕ್ಷಿಣ ಕನ್ನಡ ಜಿಲ್ಲೆ ನಮ್ಮ ನೆನಪಿಗೆ ಬರುತ್ತದೆ.\par

\kbvexercise{ಅಭ್ಯಾಸ ಪತ್ರಿಕೆ ೧೯}

\indent ಕೆಳಗೆ ಕೊಟ್ಟಿರುವ ವಾಕ್ಯಗಳಲ್ಲಿ ಅಗತ್ಯವಿದ್ದೆಡೆ ಸೂಕ್ತವಾದ ಲೇಖನ\par

\kbvsechead{ಚಿಹ್ನೆಗಳನ್ನು ಬಳಸಿರಿ :}

\vspace{2.5pt plus 3pt minus 0.5pt}\noindent\hspace*{1.2em}\kbvnum{(1)}\hspace{0.55em}ರಾಮಯ್ಯನವರೇ ರಸ್ತೆಯ ಅಂಚಿಗೆ ಬನ್ನಿ\par

\vspace{2.5pt plus 3pt minus 0.5pt}\noindent\hspace*{1.2em}\kbvnum{(2)}\hspace{0.55em}ಅಯ್ಯೋ ನಾಯಿಮರಿ ಲಾರಿಯ ಚಕ್ರಕ್ಕೆ ಸಿಕ್ಕಿ ಅಪ್ಪಚ್ಚಿ ಯಾಯಿತು.\par

\vspace{2.5pt plus 3pt minus 0.5pt}\noindent\hspace*{1.2em}\kbvnum{(3)}\hspace{0.55em}ಬೃಂದಾವನ್‌ ಎಕ್ಸ್‌ಪ್ರೆಸ್‌ ರೈಲುಗಾಡಿಯು ಬೆಂಗಳೂರಿನಿಂದ ಎಷ್ಟು ಗಂಟೆಗೆ ಹೊರಡುತ್ತದೆ\par

\vspace{2.5pt plus 3pt minus 0.5pt}\noindent\hspace*{1.2em}\kbvnum{(4)}\hspace{0.55em}ಸೀತಾಯಣ ಶ್ರೀ ಪೋಲಂಕಿ ರಾಮಮೂರ್ತಿಯವರು ಬರೆದ ಉತ್ತಮ ವೈಚಾರಿಕ ಗ್ರಂಥ\par

\vspace{2.5pt plus 3pt minus 0.5pt}\noindent\hspace*{1.2em}\kbvnum{(5)}\hspace{0.55em}ಡೆಕ್ಕನ್‌ ಹೆರಾಲ್ಡ್‌ ಪತ್ರಿಕೆಗೆ ನೀಡಿದ ಸಂದರ್ಶನದಲ್ಲಿ ರಾಜ್ಯದ ಮುಖ್ಯ ಮಂತ್ರಿಯವರು ಹಿಂಸಾಚಾರ ಹಾಗೂ ಭ್ರಷ್ಟಾಚಾರ ನಿರ್ನಾಮ ವಾಗದ ಹೊರತು ಜನರು ನೆಮ್ಮದಿಯಿಂದಿರಲು ಸಾಧ್ಯವಿಲ್ಲ ಎಂದು ಹೇಳಿದರು,\par

\vspace{2.5pt plus 3pt minus 0.5pt}\noindent\hspace*{1.2em}\kbvnum{(6)}\hspace{0.55em}ಆದಾಯ ತೆರಿಗೆ ತೆಗೆದು ಹಾಕಬೇಕು ಜೀವನಾವಶ್ಯಕ ವಸ್ತು ಗಳನ್ನು ಬಿಟ್ಟು ಉಳಿದೆಲ್ಲ ವಸ್ತುಗಳ ಮೇಲೂ ಮಾರಾಟ ತೆರಿಗೆಯನ್ನು ಹೆಚ್ಚಿನ ದರದಲ್ಲಿ ವಿಧಿಸಬೇಕು ಆಗ ಜನರ ನೈತಿಕ ಮಟ್ಟ ಸುಧಾರಿಸುವುದು ಜನಜೀವನ ಉತ್ತಮಗೊಳ್ಳುವುದು ದೇಶದಲ್ಲಿ ಆರ್ಥಿಕ ಸುಭದ್ರತೆಯುಂಟಾಗುವುದು,\par
}
\clearpage

% === tex/0092.tex (Printed Page 90) ===
\renewcommand{\kbvchaptername}{ಅಲಂಕಾರಗಳು}
\setcounter{page}{90}
\thispagestyle{kbvpage}
\kbvtypesetpage{\kbvchapter{ಪಾಠ ೨೦}{ಅಲಂಕಾರಗಳು}

\indent ಬರವಣಿಗೆಯು ಅರ್ಥ ಮಾಧುರ್ಯದಿಂದ ಮತ್ತು ಶಬ್ದ ಮಾಧುರ್ಯದಿಂದ ಕೂಡಿದ್ದರೆ ಓದಲು ಕೇಳಲು ಹಿತವಾಗಿರುವುದು. ಮನಸೆಳೆಯುವ ಶಬ್ದಾರ್ಥ ಗಳಿಂದಲೂ, ಕಿವಿಗಿಂಪಾದ ಶಬ್ದಗಳಿಂದಲೂ ಬರವಣಿಗೆಯು ಕೂಡಿದ್ದರೆ ಓದ ಬೇಕೆಂಬ ಇಚ್ಛೆ ಪ್ರಬಲಿಸುವುದು. ಹೀಗೆ, ಅಲಂಕಾರಮಯವಾಗಿರುವ ಬರ ವಣಿಗೆಯು ಚಿತ್ತಾಕರ್ಷಣೀಯವಾಗಿದ್ದು ರಸದೌತಣಕ್ಕೆ ಮಿಗಿಲಾಗಿರುವುದು,\par

\indent ಶಬ್ದಾರ್ಥದಲ್ಲಿ ಚಮತ್ಕಾರವಿದ್ದು ಆ ಮೂಲಕವಾಗಿ ಬರಹದಲ್ಲಿ ಸೊಬಗು ಮೂಡಿಬಂದರೆ ಅದು ಅರ್ಥಾಲಂಕಾರದಿಂದ ಕೂಡಿದ ಬರವಣಿಗೆಯೆನಿಸುವುದು.\par

\indent ಶಬ್ದಗಳ ಪ್ರಯೋಗದಲ್ಲಿ ಜಾಣತನವಿದ್ದು ಆ ಮೂಲಕವಾಗಿ ಬರಹ ದಲ್ಲಿ ಸೊಬಗು ಮೂಡಿಬಂದರೆ, ಅದು “ಶಬ್ದಾಲಂಕಾರ”ದಿಂದ ಕೂಡಿದ ಬರವಣಿಗೆ ಎನಿಸುವುದು.\par

\vspace{7pt plus 6pt minus 1pt}\begin{center}{\kbvsecfont\fontsize{13pt}{16.5pt}\selectfont\color{kbvaccentdark} 1. ಅರ್ಥಾಲಂಕಾರ}\end{center}\vspace{3pt plus 3pt minus 1pt}

\indent ಅರ್ಥಾಲಂಕಾರದಲ್ಲಿ ಹಲವು ಪ್ರಕಾರಗಳುಂಟು. ಅವುಗಳು ಯಾವು ವೆಂಬುದನ್ನು ನೋಡೋಣ.\par

\kbvsechead{(೧)\hspace{0.45em}ಉಪಮಾಲಂಕಾರ :}

\indent “ಇಂದ್ರಾಣಿಯ ಮುಖವು ತಾವರೆಯ ಹೂವಿನಂತೆ ಕಂಗೊಳಿಸುತ್ತಿತ್ತು.”\par

\noindent\hfill{\fontsize{9.5pt}{12.5pt}\selectfont\color{kbvgrey}- ಈ ವಾಕ್ಯದಲ್ಲಿ ಇಂದ್ರಾಣಿಯ ಮುಖವನ್ನು ತಾವರೆಯ ಹೂವಿಗೆ ಹೋಲಿಸಿ}\par\vspace{3pt plus 2pt}

\indent ರುವುದನ್ನು ಕಾಣುತ್ತೇವೆ. ಇಂತಹ ಹೋಲಿಕೆಗಳಲ್ಲಿ ಯಾವುದಕ್ಕೆ ಹೋಲಿಕೆ\par

\indent ಯನ್ನು ಕೊಡಲಾಗುವುದೋ ಅದನ್ನು “ಉಪಮೇಯ” ಎಂದು, ಯಾವು\par

\indent ದನ್ನು ಹೋಲಿಸಲಾಗುವುದೋ ಅದನ್ನು “ಉಪಮಾನ” ಎಂದು ಕರೆಯು ವುದು ವಾಡಿಕೆಯಲ್ಲಿದೆ. ಈ ಉದಾಹರಣೆಯಲ್ಲಿ “ಇಂದ್ರಾಣಿಯ ಮುಖ” ಉಪಮೇಯವಾಗಿರುವುದು, “ತಾವರೆಯ ಹೂವು” ಉಪಮಾನವಾಗಿರುವುದು. ಈ ಬಗೆಯ ಅರ್ಥಾಲಂಕಾರವನ್ನು ಉಪಮಾಲಂಕಾರವೆನ್ನುವರು.\par
}
\clearpage

% === tex/0093.tex (Printed Page 91) ===
\renewcommand{\kbvchaptername}{ಅಲಂಕಾರಗಳು}
\setcounter{page}{91}
\thispagestyle{kbvpage}
\kbvtypesetpage{\noindent\hspace*{1.2em}\kbvlabel{ಉದಾ}\hspace{0.6em}ಪುಷ್ಪ ಸಮುದಾಯದಲ್ಲಿ ಗುಲಾಬಿಯು ಹೇಗೋ ಕವಿ ಸಮುದಾಯದಲ್ಲಿ ಲಕ್ಷ್ಮೀಶನು ಹಾಗೆ.\par

\kbvsechead{(೨)\hspace{0.45em}ರೂಪಕಾಲಂಕಾರ :}

\indent “ಸೀತಮ್ಮನವರ ಸೊಸೆ ಮಹಾಲಕ್ಷ್ಮಿಯೇ ಸರಿ.”—ಈ ವಾಕ್ಯದಲ್ಲಿ, ಸೀತಮ್ಮನವರ ಸೊಸೆಯೇ ಮಹಾಲಕ್ಷ್ಮಿ ಎಂಬ ಅರ್ಥ ಬರುತ್ತದೆ. ಅಂದರೆ ಸೀತಮ್ಮನವರ ಸೊಸೆಗೂ ಮಹಾಲಕ್ಷ್ಮಿಗೂ ಯಾವುದೇ ಭೇದವಿರುವುದಾಗಿ ತಿಳಿದುಬರುವುದಿಲ್ಲ... ಉಪಮಾನ ಉಪಮೇಯಗಳೆರಡೂ ಒಂದೇ ಎಂಬ ಅಭೇದರೂಪಕ ಇಲ್ಲಿ ಕಂಡುಬರುತ್ತದೆ. ಈ ಬಗೆಯ ಅರ್ಥಾಲಂಕಾರವನ್ನು ರೂಪಕಾಲಂಕಾರವೆನ್ನುವರು.\par

\noindent\hspace*{1.2em}\kbvlabel{ಉದಾ}\hspace{0.6em}ಗಾನವಿಶಾರದೆ ಶ್ರೀಮತಿ ಸುಬ್ಬಲಕ್ಷ್ಮಿಯು ಸಾಕ್ಷಾತ್‌ ಸರಸ್ವತಿ.\par

\kbvsechead{(೩)\hspace{0.45em}ಉತ್ಪ್ರೇಕ್ಷಾಲಂಕಾರ :}

\indent “ಹೈನುಗಾರಿಕೆಗೆ ಪ್ರೋತ್ಸಾಹ ಸಿಕ್ಕಿದ ಕಾರಣ ಕರ್ನಾಟಕದ ರಾಜ್ಯದಲ್ಲಿ ಹಾಲಿನ ಹೊಳೆ ಹರಿಯುತ್ತಿದೆ.”—ಈ ವಾಕ್ಯದಲ್ಲಿ ಉತ್ಪ್ರೇಕ್ಷೆ ಎದ್ದು ಕಾಣುತ್ತಿದೆ. ರಾಜ್ಯದಲ್ಲಿ ಹಾಲಿನ ಉತ್ಪಾದನೆ ಗಣನೀಯವಾಗಿ ಹೆಚ್ಚಿದೆಯೆಂಬುದನ್ನು ತಿಳಿಸುವ ಉದ್ದೇಶವಿಟ್ಟುಕೊಂಡು ವಸ್ತುಸ್ಥಿತಿಯನ್ನು ಅತಿಶಯವಾಗಿ ಬಣ್ಣಿಸುವುದು ಇಲ್ಲಿನ ವೈಶಿಷ್ಟ್ಯ... ಈ ಬಗೆಯ ಅರ್ಥಾಲಂಕಾರವನ್ನು ಉತ್ಪ್ರೇಕ್ಷಾಲಂಕಾರವೆನ್ನುವರು.\par

\noindent\hspace*{1.2em}\kbvlabel{ಉದಾ}\hspace{0.6em}ಭೀಮನು ಶಿಶುವಾಗಿದ್ದಾಗ ತಾಯ ತೊಡೆಯಿಂದ ಕೆಳಗುರುಳಿ ಬಿದ್ದುದರಿಂದ ಆ ಅಡಿಗಲ್ಲೇ ಚೂರ್ಣವಾಯಿತಂತೆ !\par

\kbvsechead{(೪)\hspace{0.45em}ದೃಷ್ಟಾಂತಾಲಂಕಾರ :}

\indent “ಮನೆಗೊಂದು ಮುದುಕಿ, ಒಲೆಗೊಂದು ಕೊರಡು.”—ಈ ಗಾದೆ ಮಾತಿನಲ್ಲಿ ಪ್ರಧಾನವಾದ ಅಂಶವು ಪೂರ್ವಾರ್ಧದಲ್ಲಿ ಅಡಕವಾಗಿದೆ. “ಮನೆ ಗೊಂದು ಮುದುಕಿ” ಅವಶ್ಯಕವಾಗಿ ಇರಬೇಕು. ಆ ಮುದುಕಿಯು ತನ್ನ ಅನುಭವಗಳನ್ನು ಮುಂದಿನ ಪೀಳಿಗೆಗೆ ಧಾರೆ ಎರೆಯುವ ಮೂಲಕ ಆ ಮನೆಯ\kbvcontinues
}
\clearpage

% === tex/0094.tex (Printed Page 92) ===
\renewcommand{\kbvchaptername}{ಅಲಂಕಾರಗಳು}
\setcounter{page}{92}
\thispagestyle{kbvpage}
\kbvtypesetpage{\kbvcontinued ಸಂಸ್ಕೃತಿಯನ್ನು ಉಳಿಸಿ ಬೆಳೆಸುತ್ತಾಳೆ. ಈ ಅಂಶಕ್ಕೆ ಒತ್ತುಕೊಡುವ ಉದ್ದೇಶದಿಂದ ಗಾದೆಯ ಉತ್ತರಾರ್ಧದಲ್ಲಿ “ಒಲೆಗೊಂದು ಕೊರಡು” ಎಂಬ ಮಾತನ್ನು ಸೇರಿಸಿದೆ. ಒಲೆಯಲ್ಲಿ ಸದಾಕಾಲ ಬೆಂಕಿ ಇದ್ದೇ ತೀರಬೇಕು ; ಇದೊಂದು ದೃಷ್ಟಾಂತ. ಹೀಗೆ ಸ್ಪಷ್ಟೀಕರಣ ನೀಡಬೇಕಾದ ವಿಷಯಕ್ಕೆ ಮತ್ತೊಂದು ವಿಷಯದ ದೃಷ್ಟಾಂತ ಪ್ರತಿಬಿಂಬದಂತೆ ಕಾಣುವಂತಿದ್ದರೆ ಅರ್ಥಸಾದೃಶ್ಯ ಸಮಂಜಸವಾಗಿರುವುದು.\par

\indent ಈ ರೀತಿಯಾಗಿ ಎರಡು ವಾಕ್ಯಗಳು ಅಥವಾ ವಾಕ್ಯಾಂಶಗಳು ಅರ್ಥಸಾದೃಶ್ಯದಿಂದ ಬಿಂಬ ಪ್ರತಿಬಿಂಬದಂತೆ ಕಂಡುಬಂದರೆ, ಅಂತಹ ಅಲಂಕಾರವೆನಿಸುವುದು.\par

\noindent\hspace*{1.2em}\kbvlabel{ಉದಾ}\hspace{0.6em}ಊಟ ಬಲ್ಲವನಿಗೆ ರೋಗವಿಲ್ಲ; ಮಾತು ಬಲ್ಲವನಿಗೆ ಜಗಳವಿಲ್ಲ.\par

\kbvsechead{(೫)\hspace{0.45em}ಅರ್ಥಾಂತರನ್ಯಾಸಾಲಂಕಾರ :}

\indent “ಮಹಾಭಾರತದ ಯುದ್ಧದಲ್ಲಿ “ಸತ್ಯವೇ ದೇವರು; ಸತ್ಯಪಾಂಡವರಿಗೆ ಜಯ ಲಭಿಸಿತು; ಮಾರ್ಗವನ್ನೇ ಅನುಸರಿಸಿ ನಡೆದ ಸತ್ಯಧರ್ಮಗಳನ್ನು ನಂಬಿದವರಿಗೆ ಗಾಂಧೀಜಿಯವರು ಮಹಾತ್ಮರೆನಿಸಿ ಜಯ ಕಟ್ಟಿಟ್ಟಿದವರು.” ಇಲ್ಲಿ ಕೊಂಡು ಲೋಕಪ್ರಸಿದ್ಧರಾದರು” ಎರಡು ವಾಕ್ಯಗಳಿವೆ. ಮೊದಲನೆಯ ವಾಕ್ಯದಲ್ಲಿ ಕಂಡುಬರುವ ವಿಷಯ ಕಂಡುಬರುವ ಲೋಕೋಕ್ತಿಯನ್ನು ಎರಡನೆಯ ವಾಕ್ಯದಲ್ಲಿ ಕಂಡುಬರುವ ಲೋಕೋಕ್ತಿಯು `ಬರುವ ದೃಷ್ಟಾಂತವು ಸಮರ್ಥಿಸುವುದು, ವುದು.\par

\indent ಈ ಪ್ರಕಾರವಾಗಿ, ಒಂದು ವಿಶೇಷ ವಾಕ್ಯವನ್ನು ಸರ್ವಕಾಲಕ್ಕೂ ಅನ್ವಯಿಸುವ ಹೇಳಿಕೆಯಿಂದಾಗಲೀ ಅಥವಾ ಸರ್ವಕಾಲಕ್ಕೂ ಅನ್ವಯಿಸುವ ಹೇಳಿಕೆಯನ್ನು ಒಂದು ವಿಶೇಷ ವಾಕ್ಯದಿಂದಾಗಲೀ ಸಮರ್ಥಿಸುವ ತೆರನಾದ ಅರ್ಥಾಲಂಕಾರವನ್ನು ಅರ್ಥಾಂತರನ್ಯಾಸಾಲಂಕಾರವೆನ್ನುವರು.\par
}
\clearpage

% === tex/0095.tex (Printed Page 93) ===
\renewcommand{\kbvchaptername}{ಅಲಂಕಾರಗಳು}
\setcounter{page}{93}
\thispagestyle{kbvpage}
\kbvtypesetpage{\kbvsechead{(೬)\hspace{0.45em}ಶ್ಲೇಷಾಲಂಕಾರ :}

\indent “ಆ ಉದ್ಯಾನವು ಅಮರಾವತಿಯಂತೆ ಸುಮನಭರಿತವಾಗಿತ್ತು.”—ಈ ವಾಕ್ಯದಲ್ಲಿ ಬಳಸಿರುವ “ಸುಮನ” ಎಂಬ ಶಬ್ದಕ್ಕೆ ಎರಡು ಅರ್ಥಗಳುಂಟು ; ವರ್ಣಿಸಬೇಕಾದ ಉದ್ಯಾನದ ಪರವಾಗಿ ಹೇಳುವುದಾದರೆ, ಆ ಉದ್ಯಾ ನವು ಹೂವುಗಳಿಂದ ತುಂಬಿತ್ತು ಎಂದಾಗುವುದು ; ವರ್ಣಿಸುವ ವಿಷಯವಲ್ಲದ ಅಮರಾವತಿಯ ಪರವಾಗಿ ಹೇಳುವುದಾದರೆ, ಅಮರಾವತಿಯು ದೇವತೆಗಳಿಂದ ತುಂಬಿತ್ತು ಎಂದಾಗುವುದು. ಇಲ್ಲಿ ಶ್ಲೇಷೆಯಿರುವುದನ್ನು ಕಾಣಬಹುದು, {[}ಅಮರಾವತಿಯು ದೇವತೆಗಳಿಂದ ತುಂಬಿರುವಂತೆ, ಆ ಉದ್ಯಾನವು ಹೂವುಗಳಿಂದ ತುಂಬಿತ್ತು ಎನ್ನುವುದು ಒಟ್ಟಾರೆ ಅಭಿಪ್ರಾಯ{]}. ಹೀಗೆ, ಒಂದೇ ಶಬ್ದವು ವರ್ಣಿಸಬೇಕಾದ ವಿಷಯದ ಪರವಾಗಿ ಒಂದು ಅರ್ಥವನ್ನೂ ವರ್ಣನೆಯ ಅಗತ್ಯವಿರದ ವಿಷಯದ ಪರವಾಗಿ ಮತ್ತೊಂದು ಅರ್ಥವನ್ನು ಕೊಡುವುದಾದರೆ, ಅಂತಹ ಅರ್ಥಾಲಂಕಾರವು ಶ್ಲೇಷಾಲಂಕಾರ ವೆನಿಸುವುದು.\par

\vspace{7pt plus 6pt minus 1pt}\begin{center}{\kbvsecfont\fontsize{13pt}{16.5pt}\selectfont\color{kbvaccentdark} 2. ಶಬ್ದಾಲಂಕಾರ}\end{center}\vspace{3pt plus 3pt minus 1pt}

\indent ಶಬ್ದಾಲಂಕಾರದಲ್ಲಿ ಹಲವು ಪ್ರಕಾರಗಳುಂಟು, ಅವುಗಳು ಯಾವು ವೆಂಬುದನ್ನು ನೋಡೋಣ.\par

\kbvsechead{(೧)\hspace{0.45em}ವೃತ್ತಾನುಪ್ರಾಸ :}

\indent “ಕೊಳುಗುಳದ ಕೋಳ್ಗುದಿಯ ಕೋಲಾಹಲದ ಕೆಸರಿನಲಿ”-ಇದನ್ನು ಕರ್ಣಾಟಕ ಭಾರತ ಕಥಾಮಂಜರಿಯಿಂದ ಆರಿಸಿ ತೆಗೆದಿದೆ. ಇಲ್ಲಿ ಒಂದೇ ಅಕ್ಷರ (ಕ-ಕಾರ)ವು ನಾಲ್ಕು ಬಾರಿ ಪುನರುಕ್ತವಾಗಿ ವಾಕ್ಯದ ಸೊಬಗನ್ನು ಹೆಚ್ಚಿಸಿದೆ, ಈ ರೀತಿಯಾಗಿ, ಒಂದು ಅಥವಾ ಎರಡು ವ್ಯಂಜನಾಕ್ಷರಗಳು ಹಲವು ಬಾರಿ ಪುನರುಕ್ತವಾಗಿದ್ದರೆ ಅದನ್ನು ವೃತ್ತಾನುಪ್ರಾಸವೆನ್ನುವರು,\par
}
\clearpage

% === tex/0096.tex (Printed Page 94) ===
\renewcommand{\kbvchaptername}{ಅಲಂಕಾರಗಳು}
\setcounter{page}{94}
\thispagestyle{kbvpage}
\kbvtypesetpage{\noindent\hspace*{1.2em}\kbvlabel{ಉದಾ}\hspace{0.6em}ನಿರವಧಿಕ ನಿರ್ಮಾಯ ನಿಸ್ಪೃಹ ನಿರುಪಮಿತ ನಿರ್ದ್ವಂದ್ವ ನಿರ್ಗುಣ\par

\begin{kbvverse}
\noindent ನಿರವಯವ ನಿರ್ಲೇಪ ನಿರಾಗ್ರಹ ನಿರಾಧಾರ ।
\end{kbvverse}

\indent ನಿರುಪಮ ನಿರಾಮಯ ನಿರಂತರ ನಿರವಶೇಷ ನಿರಂಗ ನಿರ್ಮಲ\par

\begin{kbvverse}
\noindent ನಿರತಿಶಯ ನಿಷ್ಕಳ ಮಹೇಶ್ವರ ಕರುಣಿಸುವುದೆಂದ ॥
\end{kbvverse}

\indent (ಕರ್ಣಾಟ ಭಾರತ ಕಥಾಮಂಜರಿ : ಅರಣ್ಯಪರ್ವ) ಇದು ವೃತ್ತಾನುಪ್ರಾಸಕ್ಕೆ ಸೊಗಸಾದ ಉದಾಹರಣೆ. ಇಲ್ಲಿ ನಕಾರವು ಹಲವಾರು ಬಾರಿ ಪುನರುಕ್ತವಾಗಿರುವುದರಿಂದ ಪದ್ಯದ ಮಾಧುರ್ಯ ಹೆಚ್ಚಿದೆ.\par

\kbvsechead{(೨)\hspace{0.45em}ಛೇಕಾನುಪ್ರಾಸ:}

\begin{kbvverse}
\noindent ರುಣ ನಿರೋಧಿ ವಿಭಾಡ ರುಣ ರುಣ । \\
\noindent ರುಣು ಜಗತ್ರಯ ಜಟ್ಟಿ ರುಣ ರುಣ । \\
\noindent ರುಣು ರುಣೆಂಬಬ್ಬರಣೆ ಮಸಗಿದುದೆರಡು ಬಾಹೆಯಲಿ ॥
\end{kbvverse}

\indent ಹೀಗೆ, ಎರಡೆರಡು ವ್ಯಂಜನಗಳಿಂದ ಕೂಡಿದ ಶಬ್ದಗಳ ಪುನರಾವರ್ತನೆ ಯಿಂದ ಪದ್ಯದ ಮಾಧುರ್ಯ ಹಾಗು ಸೌಂದರ್ಯ ಹೆಚ್ಚಿದಲ್ಲಿ ಅದು ಛೇಕಾನು ಪ್ರಾಸವೆನಿಸುವುದು.\par

\vspace{3pt plus 2pt}\noindent\hspace*{1.2em}\kbvlabel{ಉದಾ}\par\vspace{1pt}

\begin{kbvverse}
\noindent ರಾಮನಂ ಭುವನಾಭಿ ರಾಮನಂ ಗುಣರತ್ನ । \\
\noindent ಧಾಮನಂ ರಘುಕುಲೋದ್ಧಾಮನಂ ರೂಪಜಿತ । \\
\noindent ಕಾಮನಂ ಸತ್ಕೀರ್ತಿ ಕಾಮನಂ ಶರಣಜನವಾರ್ಧಿಯಂ ಮಿಗೆ ಪೆರ್ಚಿನ । \\
\noindent ಸೋಮನಂ ಸೌಭಾಗ್ಯ ಸೋಮನಂ ಕುವಲಯ । \\
\noindent ಶ್ಯಾಮನಂ ನಿಜತನು ಶ್ಯಾಮನಂ ಘನಪುಣ್ಯ । \\
\noindent ನಾಮನಂ ಸಂತತಂ ನಾ ಮನಂದಣಿಯೆ ರಮಿಸದೆ ಬಾಳೆನೆಂದಳು ।
\end{kbvverse}

\noindent\hfill{\fontsize{9.5pt}{12.5pt}\selectfont\color{kbvgrey}--ಕನ್ನಡ ಜೈಮಿನಿ ಭಾರತ : ಸೀತಾ ಪರಿತ್ಯಾಗ ವೃತ್ತಾಂತ}\par\vspace{3pt plus 2pt}

\indent ಇಲ್ಲಿ "ಮನಂ" ಎಂಬುದು ಹಲವು ಬಾರಿ ಪುನರುಕ್ತವಾಗಿರುವ ದೆಸೆಯಿಂದ ಪದ್ಯದ ಮಾಧುರ್ಯ ಹಾಗೂ ಸೌಂದರ್ಯ ಹೆಚ್ಚಿದೆ ; ಇದು ಛೇಕಾನು ಪ್ರಾಸಕ್ಕೆ ಉತ್ತಮ ಉದಾಹರಣೆ.\par
}
\clearpage

% === tex/0097.tex (Printed Page 95) ===
\renewcommand{\kbvchaptername}{ಅಲಂಕಾರಗಳು}
\setcounter{page}{95}
\thispagestyle{kbvpage}
\kbvtypesetpage{\kbvsechead{(3)\hspace{0.45em}ಯಮಕಾಲಂಕಾರ :}

\indent ಒಂದು ಪದ್ಯದಲ್ಲಿ ಮೂರು ಅಥವಾ ಹೆಚ್ಚು ಅಕ್ಷರಗಳಿಂದ ಕೂಡಿದ ಪದವು ನಿಯತವಾಗಿ ಬಂದಿದ್ದರೆ, ಅದು ಯಮಕಾಲಂಕಾರವೆನಿಸುವುದು.\par

\noindent\hspace*{1.2em}\kbvlabel{ಉದಾ}\hspace{0.6em}ಯತಿ ಕೆಡುಗು ದುಸ್ಸಂಗದಲಿ ಭೂ ಪತಿ ಕೆಡುಗು ದುರ್ಮಂತ್ರಿಯಲಿ ನರ ಸುತ ಕೆಡುಗು ಲಾಲನೆಗಳಲಿ ಕೃಷಿ ಕೆಡುಗುಪೇಕ್ಷೆಯಲಿ ಮತಿ ಕೆಡುಗು ಮಧುಪಾನದಲಿ ಸ ದ್ಗತಿ ಕೆಡುಗು ದುಶ್ಚರಿತದಲಿ ನಿಜ\par

\begin{kbvverse}
\noindent ಸತಿ ಕೆಡುಗು ದುರ್ವ್ಯಸನದಲಿ ಭೂಪಾಲ ಕೇಳೆಂದ ॥
\end{kbvverse}

\noindent\hfill{\fontsize{9.5pt}{12.5pt}\selectfont\color{kbvgrey}—ಕರ್ಣಾಟ ಭಾರತ ಕಥಾ ಮಂಜರಿ : ಉದ್ಯೋಗ ಪರ್ವ}\par\vspace{3pt plus 2pt}

\indent ಇಲ್ಲಿ “ಕೆಡುಗು” ಎಂಬ ಶಬ್ದ ಪದ್ಯದ ಪ್ರತಿಸಾಲಿನಲ್ಲಿ ಎರಡನೆಯ ಪದವಾಗಿ ಪುನರುಕ್ತವಾಗಿದ್ದು, ಆ ಪದ್ಯದ ಸೊಬಗನ್ನು ಹೆಚ್ಚಿಸಿದೆ. “ಕೆಡುಗು” ಪದದಲ್ಲಿ ಮೂರು ಅಕ್ಷರಗಳಿರುವ ಕಾರಣ ಇದು ಯಮಕಾಲಂಕಾರ ವೆನಿಸುವುದು.\par

\kbvexercise{ಅಭ್ಯಾಸ ಪತ್ರಿಕೆ ೨೦}

\vspace{2.5pt plus 3pt minus 0.5pt}\noindent\hspace*{1.2em}\kbvnum{(1)}\hspace{0.55em}ಭಾಷಾಲಂಕಾರದ ಎರಡು ಶಾಖೆಗಳು ಯಾವುವು?\par

\vspace{2.5pt plus 3pt minus 0.5pt}\noindent\hspace*{1.2em}\kbvnum{(2)}\hspace{0.55em}ಅರ್ಥಾಲಂಕಾರ ಎಂದರೇನು? ಅತ್ಯಂತ ಸಮಂಜಸವಾದ ಉದಾಹರಣೆ ಸಹಿತ ವಿವರಿಸಿ.\par

\vspace{2.5pt plus 3pt minus 0.5pt}\noindent\hspace*{1.2em}\kbvnum{(3)}\hspace{0.55em}ಶಬ್ದಾಲಂಕಾರ ಎಂದರೇನು? ಅತ್ಯಂತ ಸಮಂಜಸವಾದ ಉದಾಹರಣೆ ಸಹಿತ ವಿವರಿಸಿ.\par

\vspace{2.5pt plus 3pt minus 0.5pt}\noindent\hspace*{1.2em}\kbvnum{(4)}\hspace{0.55em}ಅರ್ಥಾಲಂಕಾರದ ಯಾವುದಾದರೂ ಎರಡು ಪ್ರಕಾರಗಳನ್ನು ಹೆಸರಿಸಿ.\par

\vspace{2.5pt plus 3pt minus 0.5pt}\noindent\hspace*{1.2em}\kbvnum{(5)}\hspace{0.55em}ಶಬ್ದಾಲಂಕಾರದ ಯಾವುದಾದರೂ ಎರಡು ಪ್ರಕಾರಗಳನ್ನು ಹೆಸರಿಸಿ.\par
}
\clearpage

% === tex/0098.tex (Printed Page 96) ===
\renewcommand{\kbvchaptername}{ಅಲಂಕಾರಗಳು}
\setcounter{page}{96}
\thispagestyle{kbvpage}
\kbvtypesetpage{\vspace{2.5pt plus 3pt minus 0.5pt}\noindent\hspace*{1.2em}\kbvnum{(6)}\hspace{0.55em}“ನಕ್ಷತ್ರ ಸಮುದಾಯದಲ್ಲಿ ಧ್ರುವ ನಕ್ಷತ್ರವು ಹೇಗೋ, ಹೊಸ ಗನ್ನಡಭಾಷೆಯ ಕವಿ ಸಮುದಾಯದಲ್ಲಿ ದ. ರಾ. ಬೇಂದ್ರೆಯವರು ಹಾಗೆ.”\par

\noindent\hfill{\fontsize{9.5pt}{12.5pt}\selectfont\color{kbvgrey}— ಈ ವಾಕ್ಯದಲ್ಲಿರುವ ಉಪಮಾನ, ಉಪಮೇಯಗಳನ್ನು ತಿಳಿಸಿರಿ.}\par\vspace{3pt plus 2pt}

\vspace{2.5pt plus 3pt minus 0.5pt}\noindent\hspace*{1.2em}\kbvnum{(7)}\hspace{0.55em}ಇವುಗಳ ಸೂತ್ರ ಬರೆದು, ಉದಾಹರಣೆ ಸಹಿತ ವಿವರಿಸಿ.\par

\vspace{2.5pt plus 3pt minus 0.5pt}\noindent\hspace*{1.2em}\kbvnum{(ಅ)}\hspace{0.55em}ಉತ್ಪ್ರೇಕ್ಷಾಲಂಕಾರ, (ಬ) ಶ್ಲೇಷಾಲಂಕಾರ,\par

\vspace{2.5pt plus 3pt minus 0.5pt}\noindent\hspace*{1.2em}\kbvnum{(ಕ)}\hspace{0.55em}ವೃತ್ತಾನುಪ್ರಾಸ, (ಡ) ಛೇಕಾನುಪ್ರಾಸ.\par

\vspace{2.5pt plus 3pt minus 0.5pt}\noindent\hspace*{1.2em}\kbvnum{(8)}\hspace{0.55em}ಇವುಗಳಲ್ಲಿ ಕಂಡುಬರುವ ಭಾಷಾಲಂಕಾರವನ್ನು ತಿಳಿಸಿರಿ.\par

\vspace{2.5pt plus 3pt minus 0.5pt}\noindent\hspace*{1.2em}\kbvnum{(೧)}\hspace{0.55em}ತಾಯಿಯೇ ದೇವರು.\par

\vspace{2.5pt plus 3pt minus 0.5pt}\noindent\hspace*{1.2em}\kbvnum{(೨)}\hspace{0.55em}ವಿಶ್ವಕನ್ನಡ ಸಮ್ಮೇಳನದ ಅವಧಿಯಲ್ಲಿ ಮೈಸೂರು ನಗರವು ಇಂದ್ರಲೋಕದಂತೆ ಕಂಗೊಳಿಸುತ್ತಿತ್ತು.\par

\vspace{2.5pt plus 3pt minus 0.5pt}\noindent\hspace*{1.2em}\kbvnum{(೩)}\hspace{0.55em}ಜಾಲಿಯ ಮರವು ನೆರಳಲ್ಲ-ಹೆಣ್ಣೀಗೆ\par

\begin{kbvverse}
\noindent ತಾಯೀಯ ಮನೆಯು ಸ್ಥಿರವಲ್ಲ ॥
\end{kbvverse}

\vspace{2.5pt plus 3pt minus 0.5pt}\noindent\hspace*{1.2em}\kbvnum{(೪)}\hspace{0.55em}ಬಡತನದಲ್ಲಿ ಹುಟ್ಟಿ ಬೆಳೆದ ಮೋಕ್ಷಗುಂಡಂ ವಿಶ್ವೇಶ್ವರಯ್ಯ ನವರು ತಮ್ಮ ಸ್ವಂತ ಪರಿಶ್ರಮದಿಂದ ವಿದ್ಯಾರ್ಜನೆ ಹಾಗೂ ಜ್ಞಾನಾರ್ಜನೆ ಮಾಡಿ ಮೈಸೂರಿನ ದಿವಾನರಾದರು ಮತ್ತು ಪ್ರತಿಭಾವಂತ ಇಂಜನಿಯರ್‌ ಎನಿಸಿಕೊಂಡು ಲೋಕ ವಿಖ್ಯಾತರಾದರು ಕೈಕೆಸರಾದರೆ ಬಾಯಿ ಮೊಸರು ಎಂಬ ಮಾತು ಸುಳ್ಳಲ್ಲ.\par

\vspace{2.5pt plus 3pt minus 0.5pt}\noindent\hspace*{1.2em}\kbvnum{(೫)}\hspace{0.55em}ಹಾರುತಿದೇ ನೋಡಿರಲ್ಲಿ ನಮ್ಮ ಬಾವುಟ\par

\begin{kbvverse}
\noindent ಮೂರು ಬಣ್ಣದಿಂದ ನಲಿವ ನಮ್ಮ ಬಾವುಟ ।
\end{kbvverse}

\indent ಗಗನವೇರಿ ಭರದಿ ಕುಣಿವ ದಿವ್ಯ ಬಾವುಟ\par

\begin{kbvverse}
\noindent ಗಾಂಧಿ ತಾತ ಗಳಿಸಿಕೊಟ್ಟ ಭವ್ಯ ಬಾವುಟ ॥
\end{kbvverse}

\vspace{2.5pt plus 3pt minus 0.5pt}\noindent\hspace*{1.2em}\kbvnum{(೬)}\hspace{0.55em}ಮಹಾಪುರುಷನು ವಿಪತ್ತಿಗೆ ಒಳಗಾದಾಗ ಹಿಂದೆಂದೂ ಕಾಣದ ಉದಾತ್ತ ಗುಣವನ್ನು ಪ್ರಕಟಿಸುತ್ತಾರೆ; ಧೂಪವು ಬೆಂಕಿಯ ನಡುವೆ ಇದ್ದಾಗ ಲೋಕೋತ್ತರವಾದ ಪರಿಮಳವನ್ನು ಹರಡುತ್ತದೆ. -ರಸಗಂಗಾಧರ.\par

\vspace{2.5pt plus 3pt minus 0.5pt}\noindent\hspace*{1.2em}\kbvnum{(೭)}\hspace{0.55em}ಭೀಮಸೇನನು ಬಾಲ್ಯದಲ್ಲಿ ನದಿ ತೀರದಲ್ಲಿ ಆಟವಾಡುತ್ತಿದ್ದಾಗ ಆನೆಗಳನ್ನು ಕೈಯಿಂದ ಎತ್ತಿ ನೀರಿಗೆ ಎಸೆಯುತ್ತಿದ್ದನಂತೆ!\par

\kbvsubtitle{0}
}
\clearpage

% === tex/0099.tex (Printed Page 97) ===
\renewcommand{\kbvchaptername}{ಛಂದಸ್ಸು}
\setcounter{page}{97}
\thispagestyle{kbvpage}
\kbvtypesetpage{\kbvchapter{ಪಾಠ ೨೧}{ಛಂದಸ್ಸು}

\kbvsechead{ಪದ್ಯಗಳನ್ನು ರಚಿಸುವುದು ಎಲ್ಲರಿಗೂ ಸಾಧ್ಯವಾಗಲಾರದು :}

\indent ಇದೊಂದು ಹುಟ್ಟು ಕಲೆ, ಈ ಕಲೆ ಕರಗತವಾಗಬೇಕಾದರೆ ಕೆಲವೊಂದು ನಿಯಮಗಳನ್ನು ತಿಳಿದು ಪಾಲಿಸಬೇಕು, ಹೀಗೆ ಪದ್ಯರಚನಾಕೌಶಲ್ಯವನ್ನು ತಿಳಿಸುವ ಶಾಸ್ತ್ರವೇ ಛಂದಶಾಸ್ತ್ರ, ಛಂದಶಾಸ್ತ್ರವನ್ನು ಸಾಕಷ್ಟು ಮಟ್ಟಿಗೆ ಅರ್ಥಮಾಡಿಕೊಳ್ಳದಿದ್ದರೆ ಪದ್ಯಗಳಲ್ಲಿ ಅಡಗಿರುವ ಸೊಬಗನ್ನು ಕಾಣ ಲಾಗುವುದಿಲ್ಲ.\par

\indent “ಛಂದೋಂಬುಧಿ” ಎಂಬ ಹಳಗನ್ನಡ ಗ್ರಂಥದಲ್ಲಿ ಪದ್ಯಗಳ ರಚನೆಗೆ ಸಂಬಂಧಪಟ್ಟ ನಿಯಮಗಳನ್ನು ಮತ್ತು ಪದ್ಯಗಳ ಲಕ್ಷಣಗಳನ್ನು ವಿಸ್ತಾರವಾಗಿ ಹೇಳಲಾಗಿದೆ. ಇದು ಬಹುಮುಖ್ಯವಾದ ಆಧಾರ ಗ್ರಂಥ; ಈ ಗ್ರಂಥದ ಕರ್ತೃ ಕ್ರಿ.ಶ. ಹತ್ತನೆಯ ಶತಮಾನದಲ್ಲಿದ್ದ “ಕವಿರಾಜಹಂಸ” ಬಿರುದಾಂಕಿತ “ನಾಗವರ್ಮ”.\par

\indent ಪದ್ಯಗಳಲ್ಲಿ ನಿರ್ದಿಷ್ಟ ಸಂಖ್ಯೆಯ ಸಾಲುಗಳಿರುವುವು. ಈ ಸಾಲುಗಳಿಗೆ "ಪಾದ”ಗಳು ಅಥವಾ “ಚರಣ”ಗಳು ಎಂದು ಹೆಸರು, ಕನ್ನಡ ಸಾಹಿತ್ಯ ಕ್ಷೇತ್ರದಲ್ಲಿ ಎರಡು, ಮೂರು, ನಾಲ್ಕು ಮತ್ತು ಆರು ಪಾದಗಳ ಪದ್ಯಗಳನ್ನು ಬಹುಮಟ್ಟಿಗೆ ಕಾಣಬಹುದು; ಇವುಗಳನ್ನು ಕ್ರಮವಾಗಿ ದ್ವಿಪದಿ, ತ್ರಿಪದಿ, ಚೌಪದಿ ಮತ್ತು ಷಟ್ಪದಿ ಎನ್ನುವರು.\par

\indent ಯಾವುದೇ ಪದ್ಯದ ಪ್ರತಿಯೊಂದು ಪಾದವೂ ಯತಿ, ಗಣ, ಪ್ರಾಸ ನಿಯಮಗಳಿಗೆ ಅನುಸಾರವಾಗಿರುವುದು.\par

\vspace{4pt plus 4pt minus 1pt}\noindent{\kbvsecfont\fontsize{11pt}{14.5pt}\selectfont\color{kbvaccentdark} (೧)\hspace{0.45em}ಯತಿ:}\hspace{0.5em}ಪದ್ಯವನ್ನು ಓದುವಾಗ ಅವಶ್ಯಕವಾಗಿ ನಿಲ್ಲಿಸುವ ಹಂತವೇ “ಯತಿ” ; ಗೊತ್ತಾದ ಹಂತದಲ್ಲಿ ಓದು ನಿಲ್ಲಿಸಿದಾಗ ಅರ್ಥ ಕೆಡುವಂತಿರಬಾರದು. ಆದರೆ ಕನ್ನಡದಲ್ಲಿ ಕಾವ್ಯರಚನೆ ಮಾಡುವಾಗ ಯತಿನಿಯಮವನ್ನು ಕಟ್ಟುನಿಟ್ಟಾಗಿ ಪಾಲಿಸುವ ಸಂಪ್ರದಾಯ ನಮ್ಮಲ್ಲಿ ಬೆಳೆದುಬಂದಿಲ್ಲ.\par
}
\clearpage

% === tex/0100.tex (Printed Page 98) ===
\renewcommand{\kbvchaptername}{ಛಂದಸ್ಸು}
\setcounter{page}{98}
\thispagestyle{kbvpage}
\kbvtypesetpage{\vspace{4pt plus 4pt minus 1pt}\noindent{\kbvsecfont\fontsize{11pt}{14.5pt}\selectfont\color{kbvaccentdark} (೨)\hspace{0.45em}ಪ್ರಾಸ:}\hspace{0.5em}ಯಾವುದೇ ಪದ್ಯದ ಪ್ರತಿಯೊಂದು ಪಾದದ ಮೊದಲನೆಯ ಮತ್ತು ಎರಡನೆಯ ಸ್ವರಗಳ ನಡುವೆ. ಒಂದೇ ಜಾತಿಯ ವ್ಯಂಜನಾಕ್ಷರ ಬರುವುದನ್ನು “ಪ್ರಾಸ” ಎನ್ನುತ್ತೇವೆ. ಮೊದಲನೆಯ ಪಾದದ ಪ್ರಾಸಾಕ್ಷರದ ಹಿಂದಿನ ಅಕ್ಷರ ಹ್ರಸ್ವಸ್ವರವಾಗಿರಲೀ, ದೀರ್ಘಸ್ವರವಾಗಿರಲೀ, ಅನುಸ್ವಾರವಾಗಿರಲೀ ವಿಸರ್ಗವಾಗಿರಲೀ ಆಗಿದ್ದ ಪಕ್ಷದಲ್ಲಿ ಉಳಿದೆಲ್ಲ ಪಾದ ಗಳಲ್ಲಿಯೂ ಅದೇ ಲಕ್ಷಣ ಕಂಡುಬರುವುದು. ಇದು ಪ್ರಾಸ ನಿಯಮ.\par

\vspace{3pt plus 2pt}\noindent\hspace*{1.2em}\kbvlabel{ಉದಾ}\par\vspace{1pt}

\indent ಮಡದಿ ನಿಜ ನಿಳಯವ ಬಿಟ್ಟಡಿ\par

\kbvsechead{ಗಡಿಗೆ ಪರಗೃಹದೊಳಗೆ ಬಾಯನು :}

\begin{kbvverse}
\noindent ಬಡಿದು ಮನೆಮುಕೆನಾರ್ತೆಯೆನ್ನದೆ ಬೀದಿಗಲಹವನು ।
\end{kbvverse}

\indent ಒಡರಿಚುವ ಸತಿಯೊಬ್ಬನುಂಬೆಂ\par

\begin{kbvverse}
\noindent ದೆಡಹಿ ಕಾಣದ ದಿಟ್ಟೆ ಹತ್ತನು ।
\end{kbvverse}

\indent ಹಡೆದೊಡೆಯು ವರ್ಜಿಸುವುದುತ್ತಮ ಪುರುಷರುಗಳೆಂದ\par

\noindent\hfill{\fontsize{9.5pt}{12.5pt}\selectfont\color{kbvgrey}—ಕರ್ಣಾಟ ಭಾರತ ಕಥಾಮಂಜರಿ: ಉದ್ಯೋಗಪರ್ವ}\par\vspace{3pt plus 2pt}

\indent ಮಡ = ಮ್ + ಅ + ಡ್ + ಅ; ಒಡ = ಒ + ಡ್ + ಅ\par

\indent ಗಡಿ = ಗ್ + ಅ + ಡ್ + ಇ; ದೆಡ = ದ್ + ಎ + ಡ್ + ಅ\par

\indent ಬಡಿ = ಬ್ + ಅ + ಡ್ + ಇ; ಹಡೆ = ಹ್ + ಅ + ಡ್ + ಎ\par

\begin{kbvnote}
\noindent{\kbvsecfont ವಿ.ಸೂ. :}\hspace{0.5em}“ರಗಳೆ”ಗಳಲ್ಲಿ ಪಾದಾಂತ್ಯ (ಸಾಲಿನ ಕೊನೆ) ದಲ್ಲೂ ಪ್ರಾಸ ನಿಯಮ ಕಂಡುಬರುವುದು ; ಇದೊಂದು ವಿಶಿಷ್ಟ ಲಕ್ಷಣವೇ ಸರಿ.
\end{kbvnote}

\vspace{4pt plus 4pt minus 1pt}\noindent{\kbvsecfont\fontsize{11pt}{14.5pt}\selectfont\color{kbvaccentdark} (೩)\hspace{0.45em}ಗಣ:}\hspace{0.5em}ಪದ್ಯದ ಪ್ರತಿಪಾದದಲ್ಲಿರುವ ಅಕ್ಷರಗಳ ಗುಂಪಿಗೆ “ಗಣ” ಎಂದು ಹೆಸರು. ಗಣಗಳಲ್ಲಿ ಪ್ರಧಾನವಾಗಿ ಮಾತ್ರಾ ಗಣಗಳು ಮತ್ತು ಅಕ್ಷರ ಗಣಗಳು ಎಂದು ಎರಡು ವಿಧ.\par
}
\clearpage

% === tex/0101.tex (Printed Page 99) ===
\renewcommand{\kbvchaptername}{ಛಂದಸ್ಸು}
\setcounter{page}{99}
\thispagestyle{kbvpage}
\kbvtypesetpage{\begin{kbvrule}
\noindent\kbvbullet\hspace{0.5em}ಕೆಲವು ಪದ್ಯಗಳಲ್ಲಿ ಮೂರು, ನಾಲ್ಕು ಮತ್ತು ಐದು ಮಾತ್ರೆಗಳ ಗುಂಪುಗಳಿರುವುವು. ಇವುಗಳಿಗೆ ಮಾತ್ರಾಗಣಗಳೆಂದು ಹೆಸರು. ಕಂದ, ಷಟ್ಪದಿ, ರಗಳೆ-ಇವು ಮಾತ್ರಾಗಣಗಳಿಂದ ಕೂಡಿದ ಪದ್ಯಗಳು.\par
\noindent\kbvbullet\hspace{0.5em}ಕೆಲವು ಪದ್ಯಗಳಲ್ಲಿ ಮೂರು ಅಕ್ಷರಗಳ ಗುಂಪುಗಳಿರುವುವು ; ಅವುಗಳನ್ನು ಅಕ್ಷರಗಣಗಳೆಂದು ಕರೆಯುತ್ತೇವೆ. ಅಕ್ಷರಗಣಗಳಲ್ಲಿ ಮಾತ್ರೆಗಳ ಸಂಖ್ಯೆ ಹೆಚ್ಚು ಕಡಿಮೆಯಾಗಬಹುದು; ಇಲ್ಲಿ ಅಕ್ಷರಗಳ ಸಂಖ್ಯೆ ಗಣನೀಯವೇ ಹೊರತು ಮಾತ್ರೆಗಳಲ್ಲ.
\end{kbvrule}

\indent ಲಘು : ಅಕ್ಷರಗಳನ್ನು ಉಚ್ಚರಿಸಲು ಬೇಕಾಗುವ ಕಾಲದ ಆಧಾರದ ಮೇಲೆ ಅವುಗಳನ್ನು ಒಂದು, ಎರಡು, ಮೂರು ಮಾತ್ರಾಕಾಲದ ಅಕ್ಷರ ಗಳೆಂದು ವಿಂಗಡಿಸಬಹುದು. “ಅ” ಎಂಬ ಹ್ರಸ್ವ ಸ್ವರವನ್ನು ಉಚ್ಚರಿಸುವುದಕ್ಕೆ ಬೇಕಾಗುವ ಕಾಲವು ಒಂದು ಮಾತ್ರಾಕಾಲವಾಗಿರುವುದು. ಹೀಗೆ ಒಂದು ಮಾತ್ರಾಕಾಲದಲ್ಲಿ ಉಚ್ಚರಿಸಬಹುದಾದ ಅಕ್ಷರಗಳಿಗೆ “ಲಘು”ಗಳೆಂದು ಹೆಸರು. ಹ್ರಸ್ವ ಸ್ವರಗಳು ಮತ್ತು ಹ್ರಸ್ವ ಸ್ವರಗಳಿಂದ ಕೂಡಿದ ಗುಣಿತಾಕ್ಷರ ಗಳು ಲಘುಗಳಾಗಿರುವುವು. ಉದಾಹರಣೆಗೆ ಅ, ಇ, ಉ, ಋ, ಎ, ಒ; ಕ, ಕಿ, ಕು, ಕೃ, ಕೆ, ಕೊ,...ಲಘುಗಳು.\par

\indent ಪದ್ಯದ ಸಾಲಿನಲ್ಲಿ ಬರುವ ಅಕ್ಷರಗಳ ಮೇಲೆ “$\mbox{\laghu}$” ಚಿಹ್ನೆ ಹಾಕಿ ಲಘು ಗಳನ್ನು ಗುರ್ತಿಸುವುದು ವಾಡಿಕೆ.\par

\noindent\hspace*{1.2em}\kbvlabel{ಉದಾ}\hspace{0.6em}$\mbox{\laghu}\mbox{\laghu}\mbox{\laghu}\mbox{\laghu}$ $\mbox{\laghu}\mbox{\laghu}\mbox{\laghu}\mbox{\laghu}$ $\mbox{\laghu}\mbox{\laghu}$\par

\indent ಅವನಿಪತಿಯ ಕಮಲನಯನಕೆ\par

\indent ಗುರು : “ಆ” ಎಂಬ ದೀರ್ಘಸ್ವರವನ್ನು ಉಚ್ಚರಿಸುವುದಕ್ಕೆ ಬೇಕಾಗುವ ಕಾಲವು ಎರಡು ಮಾತ್ರಾಕಾಲವಾಗಿರುವುದು. ಹೀಗೆ ಎರಡು ಮಾತ್ರಾಕಾಲದಲ್ಲಿ ಉಚ್ಚರಿಸಬಹುದಾದ ಅಕ್ಷರಗಳಿಗೆ “ಗುರು”ಗಳೆಂದು ಹೆಸರು. ದೀರ್ಘಸ್ವರಗಳು, ದೀರ್ಘಸ್ವರಗಳಿಂದ ಕೂಡಿದ ಗುಣಿತಾಕ್ಷರಗಳು, ಅನುಸ್ವಾರ ಹಾಗೂ ವಿಸರ್ಗಗಳಿಂದ ಕೂಡಿದ ಅಕ್ಷರಗಳು, ಒತ್ತಕ್ಷರಗಳ ಹಿಂದಿನ ಅಕ್ಷರಗಳು\kbvcontinues
}
\clearpage

% === tex/0102.tex (Printed Page 100) ===
\renewcommand{\kbvchaptername}{ಛಂದಸ್ಸು}
\setcounter{page}{100}
\thispagestyle{kbvpage}
\kbvtypesetpage{\kbvcontinued ನ್ಯಾ, ವ್ಯಂಜನಾಕ್ಷರಗಳ ಹಿಂದಿನ ಅಕ್ಷರಗಳು, ಷಟ್ಪದಿಗಳಲ್ಲಿನ ಮೂರು ಮತ್ತು ಆರನೆಯ ಸಾಲುಗಳ ಅಂತ್ಯದಲ್ಲಿರುವ ಅಕ್ಷರಗಳು ‘ಗುರು’ಗಳೆನಿಸುವುವು. ಉದಾಹರಣೆಗೆ ಆ, ಈ, ಊ, ೠ, ಏ, ಐ, ಓ, ಔ; ಕಾ, ಕೀ, ಕೂ, ಕೄ, ಕೈ, ಕೇ, ಕೋ, ಕೌ, ಕಂ, ಕಃ, —ಗುರುಗಳು.\par

\indent ಪದ್ಯದ ಸಾಲಿನಲ್ಲಿ ಬರುವ ಅಕ್ಷರಗಳ ಮೇಲೆ ‘\guru{}’ ಚಿಹ್ನೆ ಹಾಕಿ ಗುರುಗಳನ್ನು ಗುರ್ತಿಸುವುದು ವಾಡಿಕೆ.\par

\vspace{3pt plus 2pt}\noindent\hspace*{1.2em}\kbvlabel{ಉದಾ}\par\vspace{1pt}

\vspace{2.5pt plus 3pt minus 0.5pt}\noindent\hspace*{1.2em}\kbvnum{(1)}\hspace{0.55em}ಓಷಧ\par

\vspace{2.5pt plus 3pt minus 0.5pt}\noindent\hspace*{1.2em}\kbvnum{(2)}\hspace{0.55em}ಕಾಗದ\par

\vspace{2.5pt plus 3pt minus 0.5pt}\noindent\hspace*{1.2em}\kbvnum{(3)}\hspace{0.55em}ಮಂಗಳ\par

\vspace{2.5pt plus 3pt minus 0.5pt}\noindent\hspace*{1.2em}\kbvnum{(4)}\hspace{0.55em}ದುಃಖ\par

\vspace{2pt plus 2pt}\noindent\hspace*{1.5em}(5) ಉತ್ಸವ\hfill (6) ಈಶ್ವರ\hfill (7) ಕೊರಳ್\hspace*{1.5em}\par\vspace{2pt plus 2pt}

\vspace{2.5pt plus 3pt minus 0.5pt}\noindent\hspace*{1.2em}\kbvnum{(1)}\hspace{0.55em}\guru{}\laghu{}\laghu{} (ಇಲ್ಲಿ ‘ಓ’ ಎಂಬುದು ದೀರ್ಘಸ್ವರ)\par

\vspace{2.5pt plus 3pt minus 0.5pt}\noindent\hspace*{1.2em}\kbvnum{(2)}\hspace{0.55em}\guru{}\laghu{}\laghu{} (ಇಲ್ಲಿ ‘ಕಾ’ ಎಂಬುದು ದೀರ್ಘಸ್ವರದಿಂದ ಕೂಡಿದ ಗುಣಿತಾಕ್ಷರ)\par

\vspace{2.5pt plus 3pt minus 0.5pt}\noindent\hspace*{1.2em}\kbvnum{(3)}\hspace{0.55em}\guru{}\laghu{}\laghu{} (ಇಲ್ಲಿ ‘ಮಂ’ ಎಂಬುದು ಅನುಸ್ವಾರ)\par

\vspace{2.5pt plus 3pt minus 0.5pt}\noindent\hspace*{1.2em}\kbvnum{(4)}\hspace{0.55em}\guru{} (ಇಲ್ಲಿ ‘ದುಃ’ ಎಂಬುದು ವಿಸರ್ಗ)\par

\vspace{2.5pt plus 3pt minus 0.5pt}\noindent\hspace*{1.2em}\kbvnum{(5)}\hspace{0.55em}\guru{}\laghu{}\laghu{} (ಇಲ್ಲಿ ‘ಉ’ ಎಂಬುದು ಒತ್ತಕ್ಷರದ ಹಿಂದಿನ ಅಕ್ಷರ)\par

\vspace{2.5pt plus 3pt minus 0.5pt}\noindent\hspace*{1.2em}\kbvnum{(6)}\hspace{0.55em}\guru{}\laghu{}\laghu{} (ಇಲ್ಲಿ ‘ಈ’ ಎಂಬುದು ಒತ್ತಕ್ಷರದ ಹಿಂದಿನ ಅಕ್ಷರ ಹಾಗೂ ದೀರ್ಘಸ್ವರ ಆಗಿರುವುದು)\par

\vspace{2.5pt plus 3pt minus 0.5pt}\noindent\hspace*{1.2em}\kbvnum{(7)}\hspace{0.55em}\laghu{}\guru{} (ಇಲ್ಲಿ ‘ರ’ ಎಂಬುದು ‘ಳ್’ ವ್ಯಂಜನಾಕ್ಷರದ ಹಿಂದಿನ ಅಕ್ಷರ)\par

\indent ಸೂಚನೆ: ವ್ಯಂಜನಾಕ್ಷರಗಳನ್ನು (ಉದಾ: ಕ್, ಟ್, ನ್, ಭ್, ಣ್ ಇತ್ಯಾದಿ) ಲಘುಗಳೆಂದಾಗಲೀ, ಗುರುಗಳೆಂದಾಗಲೀ ಪರಿಗಣಿಸಲಾಗದು ಆದುದರಿಂದ ಅವುಗಳಿಗೆ ಯಾವ ಚಿಹ್ನೆಯನ್ನೂ ಹಾಕಬಾರದು.\par
}
\clearpage

% === tex/0103.tex (Printed Page 101) ===
\renewcommand{\kbvchaptername}{ಛಂದಸ್ಸು}
\setcounter{page}{101}
\thispagestyle{kbvpage}
\kbvtypesetpage{\indent ಪ್ರಸ್ತಾರ ಹಾಕುವುದು : ಯಾವುದೇ ಪದ್ಯದ ಪ್ರತಿಯೊಂದು ಸಾಲಿನಲ್ಲಿರುವ ಅಕ್ಷರಗಳಿಗೆ ಗುರು-ಲಘುಗಳನ್ನು ಗುರುತು ಮಾಡುವ ಕ್ರಮವನ್ನು ಪ್ರಸ್ತಾರ ಹಾಕುವುದು ಎನ್ನುತ್ತೇವೆ. ಪ್ರಸ್ತಾರ ಹಾಕಿದ ನಂತರ ಗುರುವಿಗೆ ಎರಡು ಮಾತ್ರೆ ಮತ್ತು ಲಘುವಿಗೆ ಒಂದು ಮಾತ್ರೆ-ಈ ಪ್ರಕಾರವಾಗಿ ಎಣಿಕೆ ಮಾಡಿ ಮೂರು ಮಾತ್ರೆ, ನಾಲ್ಕು ಮಾತ್ರೆ ಅಥವಾ ಐದು ಮಾತ್ರೆಯ ಗುಂಪುಗಳನ್ನು ವಿಂಗಡಿಸಬೇಕು; ಈ ಗುಂಪುಗಳಿಗೆ “ಮಾತ್ರಾಗಣ” ಗಳೆಂದು ಹೆಸರು.\par

\kbvsechead{ಮಾತ್ರಾಗಣದ ಪದ್ಯಗಳ ಲಕ್ಷಣ :}

\vspace{4pt plus 4pt minus 1pt}\noindent{\kbvsecfont\fontsize{11pt}{14.5pt}\selectfont\color{kbvaccentdark} (೧)\hspace{0.45em}ಕಂದ:}\hspace{0.5em}ಕಂದ ಪದ್ಯದಲ್ಲಿ ನಾಲ್ಕು ಪಾದ (ಸಾಲು) ಗಳಿರುವುವು. ಇವುಗಳ ಪೈಕಿ ಒಂದು ಮತ್ತು ಮೂರನೆಯ ಪಾದಗಳಲ್ಲಿ ನಾಲ್ಕು ಮಾತ್ರೆಯ ಮೂರು ಗಣಗಳಿರುವುವು; ಎರಡು ಮತ್ತು ನಾಲ್ಕನೆಯ ಪಾದಗಳಲ್ಲಿ ನಾಲ್ಕು ಮಾತ್ರೆಯ ಐದು ಗಣಗಳಿರುವುವು.\par

\vspace{3pt plus 2pt}\noindent\hspace*{1.2em}\kbvlabel{ಉದಾ}\par\vspace{1pt}

\indent \laghu{}\laghu{}\laghu{} \laghu{} \laghu{}\laghu{}\laghu{} \laghu{} \laghu{}\laghu{}\guru{}\par

\indent 4 \ensuremath{\times} 3....ಭವದನು ಜನರುಣ ಜಲಮಂ \laghu{}\laghu{}\guru{} \laghu{}\laghu{}\guru{} \laghu{}\guru{}\laghu{} \laghu{}\laghu{} \laghu{}\laghu{}\guru{}\par

\begin{kbvverse}
\noindent 4 \ensuremath{\times} 5....ಸವಿನೋ ಡಿದನಿಂ ತು ನಿನ್ನ ಬಲಜಲ ನಿಧಿಯಂ ।
\end{kbvverse}

\indent \laghu{}\laghu{}\guru{} \laghu{}\laghu{}\guru{} \laghu{} \laghu{}\guru{}\par

\indent 4 \ensuremath{\times} 3....ಸವಿನೋ ಡಿದೆನೀ ಕೊ ಳನಂ \laghu{}\laghu{}\guru{} \laghu{}\guru{}\laghu{}\laghu{} ೧ \guru{}\laghu{} \laghu{}\laghu{}\guru{} \guru{} \guru{}\par

\begin{kbvverse}
\noindent 4 \ensuremath{\times} 5....ತವೆಪೀ ರ್ದುಬಳಿ ಕೈ ನಿನ್ನ ಸನಿಯಂ ನೋಳ್ಪಂ ॥
\end{kbvverse}

\indent ವಿಶ್ಲೇಷಣೆ : (1) ಈ ಪದ್ಯದಲ್ಲಿ ನಾಲ್ಕು ಪಾದಗಳಿವೆ. (2) ಒಂದು ಮತ್ತು ಮೂರನೆಯ ಪಾದಗಳಲ್ಲಿ ನಾಲ್ಕು ಮಾತ್ರೆಯ ಮೂರು ಗಣಗಳಿವೆ. (3) ಎರಡು ಮತ್ತು ನಾಲ್ಕನೆಯ ಪಾದಗಳಲ್ಲಿ ನಾಲ್ಕು ಮಾತ್ರೆಯ ಐದು ಗಣಗಳಿವೆ. (4) ಪೂರ್ವಾರ್ಧ ಮತ್ತು ಉತ್ತರಾರ್ಧಗಳಲ್ಲಿ ಕ್ರಮವಾಗಿ ತಲಾ ಎಂಟು ಮಾತ್ರಾಗಣಗಳಿವೆ. (5) ಪೂರ್ವಾರ್ಧ ಹಾಗೂ ಉತ್ತರಾರ್ಧದ ಮಾತ್ರಾಗಣಗಳ ವಿಷಮ ಸ್ಥಾನಗಳಲ್ಲಿ ಮಧ್ಯಗುರುವುಳ್ಳ “\laghu{}\guru{}\laghu{}” ಈ.\par
}
\clearpage

% === tex/0104.tex (Printed Page 102) ===
\renewcommand{\kbvchaptername}{ಛಂದಸ್ಸು}
\setcounter{page}{102}
\thispagestyle{kbvpage}
\kbvtypesetpage{\indent ರೀತಿಯ ಗಣವು ಬಂದಿರುವುದಿಲ್ಲ. (6) ಪೂರ್ವಾರ್ಧ ಹಾಗೂ ಉತ್ತರಾರ್ಧದಲ್ಲಿ ಆರನೆಯ ಮಾತ್ರಾಗಣವು ಮಧ್ಯಗುರುವುಳ್ಳ ಗಣವಾಗಿರುವುದು; ಇದು “ಲಲ” ಈ ರೀತಿಯಾಗಿಯೂ ಇರಬಹುದು, (7) ಪೂರ್ವಾರ್ಧದ ಅಂತ್ಯದಲ್ಲಿ ಮತ್ತು ಉತ್ತರಾರ್ಧದ ಅಂತ್ಯದಲ್ಲಿ ಗುರು ಬರುವ “ರ್ರ” ಅಥವಾ “\guru{}\guru{}” ರೀತಿಯ ಮಾತ್ರಾಗಣವಿರುವುದು.\par

\vspace{4pt plus 4pt minus 1pt}\noindent{\kbvsecfont\fontsize{11pt}{14.5pt}\selectfont\color{kbvaccentdark} (೨)\hspace{0.45em}ಷಟ್ಪದಿ:}\hspace{0.5em}ಇವು ಆರು ಸಾಲಿನ ಪದ್ಯಗಳು. ಷಟ್ಪದಿ ಪದ್ಯಗಳಲ್ಲಿ ಶರ ಷಟ್ಪದಿ, ಕುಸುಮ ಷಟ್ಪದಿ, ಭೋಗ ಷಟ್ಪದಿ, ಭಾಮಿನೀ ಷಟ್ಪದಿ, ಪರಿವರ್ಧಿನೀ ಷಟ್ಪದಿ, ವಾರ್ಧಕ ಷಟ್ಪದಿ ಎಂಬುದಾಗಿ ಆರು ಪ್ರಕಾರಗಳಿವೆ, ಇವುಗಳೆಲ್ಲವೂ ವಿಶಿಷ್ಟ ಲಕ್ಷಣಗಳಿಂದ ಕೂಡಿರುವುವು.\par

\indent ಷಟ್ಪದಿಗಳ ಲಕ್ಷಣಗಳು\par

\begin{kbvrule}
\noindent\kbvbullet\hspace{0.5em}ಕೆಳಗೆ ಕಂಡ ಭಾಮಿನೀ ಷಟ್ಪದಿಯ ಪದ್ಯಕ್ಕೆ ಪ್ರಸ್ತಾರ ಹಾಕಿ ಮಾತ್ರಾಗಣಗಳನ್ನು ವಿಂಗಡಿಸಿರುವುದನ್ನು ಗಮನಿಸಿ :
\end{kbvrule}

\indent \laghu{}\laghu{} \laghu{}\laghu{} \laghu{}\laghu{} \laghu{} \laghu{}\laghu{} \guru{}\laghu{}\laghu{} ಮಡದಿ ನಿಜ ನಿಲಯವನು ಬಿಟ್ಟಡಿ ಅ\par

\indent \laghu{}\laghu{} \laghu{}\laghu{} \laghu{}\laghu{}\laghu{}\laghu{}\laghu{} \guru{} \laghu{}\laghu{} ಗಡಿಗೆ ಪರಗೃಹದೊಳಗೆ ಬಾಯನು\par

\indent \laghu{}\laghu{}\laghu{}\laghu{} \laghu{}\laghu{}\laghu{}\laghu{}\guru{}\laghu{}\guru{}\laghu{}\laghu{}\guru{}\laghu{}\laghu{}\laghu{} \laghu{}\laghu{}\guru{}\par

\begin{kbvverse}
\noindent ಬಡಿದು ಮನೆಮನೆವಾರ್ತೆಯೆನ್ನದೆ ಬೀದಿಗಲಹವನು ।
\end{kbvverse}

\indent \laghu{}\laghu{}\laghu{}\laghu{} \laghu{}\laghu{}\laghu{}\laghu{}\guru{}\laghu{}\guru{}\guru{} ಒಡರಿಚುವ ಪತಿಯೊಬ್ಬನುಂಬೆಂ\par

\indent \laghu{}\laghu{}\laghu{} \guru{}\laghu{}\laghu{} \guru{}\laghu{} \guru{}\laghu{}\laghu{} ಬೆಡಹಿ ಕಾಣದ ದಿಟ್ಟೆ ಹತ್ತನು\par

\indent \laghu{}\laghu{}\laghu{} \laghu{}\laghu{} \guru{}\laghu{} \laghu{}\laghu{}\guru{}\laghu{}\laghu{} \laghu{}\laghu{}\laghu{}\laghu{}\guru{} \guru{}\par

\begin{kbvverse}
\noindent ಹಡೆದೊಡೆಯು ವರ್ಜಿಸುವುದುತ್ತಮ ಪುರುಷರುಗಳೆಂದ ॥
\end{kbvverse}

\indent (ಕರ್ಣಾಟ ಭಾರತ ಕಥಾಮಂಜರಿ)\par
}
\clearpage

% === tex/0105.tex (Printed Page 103) ===
\renewcommand{\kbvchaptername}{ಛಂದಸ್ಸು}
\setcounter{page}{103}
\thispagestyle{kbvpage}
\kbvtypesetpage{% === File 0105.md (Printed Page 103: ಷಟ್ಪದಿಗಳ ಲಕ್ಷಣಗಳು, full-width upright table) ===
\par\vspace{6pt plus 4pt}\noindent\begin{minipage}{\textwidth}% keep title and table together
\begin{center}{\kbvchapfont\fontsize{14pt}{18pt}\selectfont ಷಟ್ಪದಿಗಳ ಲಕ್ಷಣಗಳು}\end{center}
\vspace{4pt plus 3pt minus 1pt}
\noindent{\fontsize{8.8pt}{11pt}\selectfont\renewcommand{\arraystretch}{1.05}\setlength{\tabcolsep}{3pt}\arrayrulecolor{kbvline}%
\setlength{\kbvcolw}{\dimexpr(\textwidth-17mm-21mm-6\tabcolsep)/2\relax}%
\begin{tabular}{@{} >{\raggedright\arraybackslash}p{17mm} >{\raggedright\arraybackslash}p{\kbvcolw} >{\raggedright\arraybackslash}p{\kbvcolw} >{\raggedright\arraybackslash}p{21mm} @{}}
      \specialrule{0.5pt}{0pt}{3pt}
      {\kbvsecfont ಷಟ್ಪದಿ} & {\kbvsecfont ೧, ೨, ೪, ೫ ನೆಯ ಪಾದಗಳು} & {\kbvsecfont ೩ ಮತ್ತು ೬ ನೆಯ ಪಾದಗಳು} & {\kbvsecfont ವಿಶೇಷ ವಿಚಾರ} \\
      \specialrule{0.5pt}{3pt}{3pt}
      {\kbvsecfont ೧.}\,ಶರ ಷಟ್ಪದಿ & ಪ್ರತಿ ಪಾದದಲ್ಲಿ ನಾಲ್ಕು ಮಾತ್ರೆಯ ಎರಡು ಗಣಗಳಿರುವುವು. & ಪ್ರತಿ ಪಾದದಲ್ಲಿ 4 ಮಾತ್ರೆಯ 3 ಗಣಗಳು ಮತ್ತು ಮೇಲೊಂದು ಗುರು ಇರುವುದು. & “\laghu\,\guru\,\laghu” ಈ ರೀತಿಯ ಗಣ ಬರುವುದಿಲ್ಲ. \\
      \specialrule{0.3pt}{2.5pt}{2.5pt}
      {\kbvsecfont ೨.}\,ಕುಸುಮ ಷಟ್ಪದಿ & ಪ್ರತಿ ಪಾದದಲ್ಲಿ ಐದು ಮಾತ್ರೆಯ ಎರಡು ಗಣಗಳಿರುವುವು. & ಪ್ರತಿ ಪಾದದಲ್ಲಿ 5 ಮಾತ್ರೆಯ 3 ಗಣಗಳು ಮತ್ತು ಮೇಲೊಂದು ಗುರು ಇರುವುದು. & “\laghu\,\guru\,\laghu\laghu” ಮತ್ತು “\laghu\,\guru\,\guru” ರೀತಿಯ ಗಣವಿರುವುದಿಲ್ಲ. \\
      \specialrule{0.3pt}{2.5pt}{2.5pt}
      {\kbvsecfont ೩.}\,ಭೋಗ ಷಟ್ಪದಿ & ಪ್ರತಿ ಪಾದದಲ್ಲಿ ಮೂರು ಮಾತ್ರೆಯ ನಾಲ್ಕು ಗಣಗಳಿರುವುವು. & ಪ್ರತಿ ಪಾದದಲ್ಲಿ 3 ಮಾತ್ರೆಯ 6 ಗಣಗಳಿರುವುವು ಮತ್ತು ಮೇಲೊಂದು ಗುರುವಿರುವುದು. & “\laghu\,\guru” ಈ ರೀತಿಯ ಗಣಗಳಿರುವುದಿಲ್ಲ. \\
      \specialrule{0.3pt}{2.5pt}{2.5pt}
      {\kbvsecfont ೪.}\,ಭಾಮಿನೀ ಷಟ್ಪದಿ & ಪ್ರತಿ ಪಾದದಲ್ಲಿ 3 ಮತ್ತು 4 ಮಾತ್ರೆಯ 4 ಗಣಗಳು ಪರ್ಯಾಯವಾಗಿ ಬರುವುವು. & ಪ್ರತಿ ಪಾದದಲ್ಲಿ 3 ಮತ್ತು 4 ಮಾತ್ರೆಯ 6 ಗಣಗಳು ಪರ್ಯಾಯವಾಗಿ ಬರುವುವು ಮತ್ತು ಮೇಲೊಂದು ಗುರುವಿರುವುದು. & “\laghu\,\guru\,\laghu” ಈ ರೀತಿಯ ಗಣವಿರುವುದಿಲ್ಲ. \\
      \specialrule{0.3pt}{2.5pt}{2.5pt}
      {\kbvsecfont ೫.}\,ಪರಿವರ್ಧಿನೀ ಷಟ್ಪದಿ & ಪ್ರತಿ ಪಾದದಲ್ಲಿ 4 ಮಾತ್ರೆಯ 4 ಗಣಗಳಿರುವುವು. & ಪ್ರತಿ ಪಾದದಲ್ಲಿ 4 ಮಾತ್ರೆಯ 6 ಗಣಗಳಿರುವುವು ಮತ್ತು ಮೇಲೊಂದು ಗುರುವಿರುವುದು. & “\laghu\,\guru\,\laghu” ಈ ರೀತಿಯ ಗಣವಿರುವುದಿಲ್ಲ. \\
      \specialrule{0.3pt}{2.5pt}{2.5pt}
      {\kbvsecfont ೬.}\,ವಾರ್ಧಕ ಷಟ್ಪದಿ & ಪ್ರತಿ ಪಾದದಲ್ಲಿ 5 ಮಾತ್ರೆಯ 4 ಗಣಗಳಿರುವುವು. & ಪ್ರತಿ ಪಾದದಲ್ಲಿ 5 ಮಾತ್ರೆಯ 6 ಗಣಗಳು ಮತ್ತು ಮೇಲೊಂದು ಗುರು ಇರುವುದು. & {\color{kbvgrey}\ldots\ldots} \\
      \specialrule{0.5pt}{3pt}{0pt}
\end{tabular}\arrayrulecolor{black}}
\end{minipage}\par\vspace{6pt plus 4pt}
}
\clearpage

% === tex/0106.tex (Printed Page 104) ===
\renewcommand{\kbvchaptername}{ಛಂದಸ್ಸು}
\setcounter{page}{104}
\thispagestyle{kbvpage}
\kbvtypesetpage{\indent ವಿ. ಸೂ. :ಷಟ್ಪದಿಯ ಮೂರು ಮತ್ತು ಆರನೆಯ ಸಾಲುಗಳಲ್ಲಿನ ಕೊನೆಯ ಅಕ್ಷರ “ಲಘು” ಆಗಿದ್ದರೂ, ಅದನ್ನು ಗುರುವೆಂತಲೇ ತಿಳಿಯ ಬೇಕು.\par

\vspace{4pt plus 4pt minus 1pt}\noindent{\kbvsecfont\fontsize{11pt}{14.5pt}\selectfont\color{kbvaccentdark} (4)\hspace{0.45em}ರಗಳೆ :}\hspace{0.5em}ಇದು ಮಾತ್ರಾಗಣಕ್ಕೆ ಸೇರಿದ ಪದ್ಯ ಜಾತಿ, “ರಘಟಾ” ಎಂಬ ಸಂಸ್ಕೃತ ಶಬ್ದದ ತದ್ಭವ ರೂಪವೇ “ರಗಳೆ”. ರಗಳೆಗಳಲ್ಲಿ ಉತ್ಸಾಹ ರಗಳೆ, ಮಂದಾನಿಲ ರಗಳೆ, ಲಲಿತ ರಗಳೆ ಎಂಬ ಮೂರು ಪ್ರಕಾರ ಗಳು ಮುಖ್ಯವಾದವುಗಳು. ರಗಳೆಗಳಲ್ಲಿ ಇಂತಿಷ್ಟು ಪಾದಗಳಿರಬೇಕೆಂಬ ನಿಯಮವಿಲ್ಲ; ಆದರೆ ಮಾತ್ರಾಗಣಗಳು ಪ್ರತಿಯೊಂದು ಪಾದದಲ್ಲೂ ಒಂದೇ ಸಮನಾಗಿರುವುದು; ಮತ್ತು ಪ್ರತಿ ಎರಡು ಪಾದಗಳಲ್ಲಿ ಪ್ರಾಸನಿಯಮವಿರುವುದು. ಕೆಲವೊಂದು ಸಂದರ್ಭಗಳಲ್ಲಿ ಪಾದದ ಆದ್ಯಂತಗಳಲ್ಲಿ ಪ್ರಾಸವಿರುವುದೂ ಉಂಟು.\par

\indent ರಗಳೆಗಳ ಲಕ್ಷಣಗಳು\par

\vspace{4pt plus 4pt minus 1pt}\noindent{\kbvsecfont\fontsize{11pt}{14.5pt}\selectfont\color{kbvaccentdark} (1)\hspace{0.45em}ಉತ್ಸಾಹ ರಗಳೆ:}\hspace{0.5em}ಈ ಜಾತಿಯ ಪದ್ಯಗಳಲ್ಲಿ ಪ್ರತಿಯೊಂದು ( ೧೧ ೬) ಪಾದದಲ್ಲೂ ಮೂರು ಮಾತ್ರೆಯ ನಾಲ್ಕು ಮಾತ್ರಾಗಣಗಳಿರುವುವು.\par

\vspace{4pt plus 4pt minus 1pt}\noindent{\kbvsecfont\fontsize{11pt}{14.5pt}\selectfont\color{kbvaccentdark} (2)\hspace{0.45em}ಮಂದಾನಿಲ ರಗಳೆ:}\hspace{0.5em}ಈ ಜಾತಿಯ ಪದ್ಯಗಳಲ್ಲಿ, ಪ್ರತಿ ಯೊಂದು ಪಾದದಲ್ಲೂ ನಾಲ್ಕು ಮಾತ್ರೆಯ ನಾಲ್ಕು ಮಾತ್ರಾಗಣಗಳಿರು ವುವು; ಅಥವಾ ಪ್ರತಿಯೊಂದು ಪಾದದಲ್ಲೂ ಮೂರು ಮತ್ತು ಐದು\par

\indent ಮಾತ್ರೆಯ ಮಾತ್ರಾಗಣಗಳು (3 + 5 + 3 + 5) ತೆರನಾಗಿ ಸಂಭವಿಸುವುವು.\par

\vspace{4pt plus 4pt minus 1pt}\noindent{\kbvsecfont\fontsize{11pt}{14.5pt}\selectfont\color{kbvaccentdark} (3)\hspace{0.45em}ಲಲಿತ ರಗಳೆ:}\hspace{0.5em}ಈ ಜಾತಿಯ ಪದ್ಯಗಳಲ್ಲಿ ಪ್ರತಿಯೊಂದು ಪಾದದಲ್ಲೂ ಐದು ಮಾತ್ರೆಯ ನಾಲ್ಕು ಮಾತ್ರಾಗಣಗಳಿರುವುವು.\par

\indent ಸೂಚನೆ: ರಗಳೆಗಳಲ್ಲಿ ಪ್ರತಿ ಎರಡು ಪಾದಗಳಲ್ಲಿ ಆದಿ ಪ್ರಾಸವಿರು ವುದು; ಕೆಲವೊಂದು ಸಂದರ್ಭಗಳಲ್ಲಿ ಆದಿಪ್ರಾಸದ ಜೊತೆಗೆ ಅಂತ್ಯ ಪ್ರಾಸವೂ ಇರುವುದುಂಟು. ಅಪರೂಪವಾಗಿ ಅಂತ್ಯಪ್ರಾಸ ಮಾತ್ರವಿರುವ ಸಂದರ್ಭಗಳೂ ಉಂಟು.\par
}
\clearpage

% === tex/0107.tex (Printed Page 105) ===
\renewcommand{\kbvchaptername}{ಛಂದಸ್ಸು}
\setcounter{page}{105}
\thispagestyle{kbvpage}
\kbvtypesetpage{\kbvsechead{ಅಕ್ಷರ ಗಣಗಳು :}

\indent ಅಕ್ಷರ ಗಣಗಳಲ್ಲಿ ಅಕ್ಷರಗಳ ಸಂಖ್ಯೆ ಗಣನೀಯವೇ ಹೊರತು. ಮೂರು ಅಕ್ಷರಗಳ ಗುಂಪು ಒಂದು ಗಣವೆನಿಸುವುದು. ಆದುದರಿಂದ ಪದ್ಯದ ಪ್ರತಿ ಸಾಲಿನಲ್ಲಿರುವ ಅಕ್ಷರಗಳನ್ನು ಮೂರು ಮೂರು ಅಕ್ಷರಗಳಂತೆ ಲೆಕ್ಕಮಾಡಿ ಗಣಗಳನ್ನು ವಿಂಗಡಿಸಲಾಗುವುದು. ಈ ಬಗೆಯ ವಿಂಗಡನೆಯು “ವೃತ್ತ ಜಾತಿಯ ಛಂದಸ್ಸಿಗೆ” ಸೀಮಿತವಾಗಿರುವುದು. ವೃತ್ತ ಜಾತಿಯ ಛಂದಸ್ಸಿನಲ್ಲಿ (೧) ಉತ್ಪಲಮಾಲಾ ವೃತ್ತ, (೨) ಚಂಪಕಮಾಲಾ ವೃತ್ತ, (೩) ಶಾರ್ದೂಲ ವಿಕ್ರೀಡಿತ ವೃತ್ತ, (೪) ಮತ್ತೇಭ ವಿಕ್ರೀಡಿತ ವೃತ್ತ, (೫) ಸ್ರಗ್ಧರಾ ವೃತ್ತ, (೬) ಮಹಾಸ್ರಗ್ಧರಾ ವೃತ್ತ ಎಂಬ ಆರು ಜಾತಿಯ ವೃತ್ತಗಳು ಮುಖ್ಯವಾದವುಗಳು.\par

\indent ಅಕ್ಷರ-ಗಣಗಳಲ್ಲಿ ಎಂಟು ವಿಧ: ಯಗಣ, ಮಗಣ, ತಗಣ, ರಗಣ, ಜಗಣ, ಭಗಣ, ನಗಣ ಮತ್ತು ಸಗಣ-ಇವೇ ಆ ಎಂಟು ಗಣಗಳು. “ಯಮಾತಾರಾಜಭಾನಸಲಗಂ” ಎಂಬ ಅಕ್ಷರ ಸೂತ್ರದ ಸಹಾಯದಿಂದ ಮೇಲೆ ತಿಳಿಸಿದ ಗಣಗಳ ಗುರು-ಲಘು ವಿನ್ಯಾಸವನ್ನು ತಿಳಿಯಬಹುದು. ಸೂತ್ರದಲ್ಲಿನ ಯಕಾರದಿಂದ ಪ್ರಾರಂಭಿಸಿ ಪ್ರತಿ ಮೂರು ಕ್ರಮಗತ ಅಕ್ಷರಗಳಿಗೆ ಪ್ರಸ್ತಾರ ಹಾಕಿದರೆ ಯಗಣದಿಂದ ಸಗಣದವರೆಗಿನ ಅಕ್ಷರ ಗಣಗಳ ವಿನ್ಯಾಸ ದೊರಕುತ್ತದೆ.\par

\kbvsechead{ಗಣಗಳ ವಿನ್ಯಾಸ}

\vspace{3pt plus 2pt}\noindent\hspace*{1.2em}\kbvlabel{ಉದಾ}\par\vspace{1pt}

\vspace{5pt plus 4pt minus 1pt}
\noindent{\renewcommand{\arraystretch}{1.24}\setlength{\tabcolsep}{4pt}\rowcolors{1}{kbvrow}{white}
\begin{tabularx}{\textwidth}{X X X X X}
ಗಣ & ಸೂತ್ರ & ವಿನ್ಯಾಸ & ಲಕ್ಷಣ & ಉದಾಹರಣೆ \\
೧. ಯಗಣ & ಯಮಾತಾ & $\mbox{\laghu}$ \guru{} \guru{} & (ಆದಿಲಘು) & : ಯಥಾರ್ಥಂ \\
೨. ಮಗಣ & ಮಾತಾರಾ & \guru{} \guru{} \guru{} & (ಸರ್ವಗುರು) & : ಮಸ್ತಿಷ್ಕಂ \\
೩. ತಗಣ & ತಾರಾಜ & \guru{} \guru{} $\mbox{\laghu}$ & (ಅಂತ್ಯಲಘು) & : ತಣ್ಣೀರು \\
೪. ರಗಣ & ರಾಜಭಾ & \guru{} $\mbox{\laghu}$ \guru{} & (ಮಧ್ಯಲಘು) & : ರಕ್ಕಸೀ \\
೫. ಜಗಣ & ಜಭಾನ & $\mbox{\laghu}$ \guru{} $\mbox{\laghu}$ & (ಮಧ್ಯಗುರು) & : ಜಟಾಯು \\
೬. ಭಗಣ & ಭಾನಸ & \guru{} $\mbox{\laghu}$ $\mbox{\laghu}$ & (ಆದಿಗುರು) & : ಭವ್ಯತೆ \\
೭. ನಗಣ & ನಸಲ & $\mbox{\laghu}$ $\mbox{\laghu}$ $\mbox{\laghu}$ & (ಸರ್ವಲಘು) & : ನಕುಲ \\
೮. ಸಗಣ & ಸಲಗಂ & $\mbox{\laghu}$ $\mbox{\laghu}$ \guru{} & (ಅಂತ್ಯಗುರು) & : ಸರ್ವಥಾ \\
\end{tabularx}}
\vspace{5pt plus 4pt minus 1pt}
}
\clearpage

% === tex/0108.tex (Printed Page 106) ===
\renewcommand{\kbvchaptername}{ಛಂದಸ್ಸು}
\setcounter{page}{106}
\thispagestyle{kbvpage}
\kbvtypesetpage{\kbvsechead{ವೃತ್ತ ಜಾತಿಯ ಛಂದಸ್ಸಿನ ಲಕ್ಷಣ :}

\vspace{4pt plus 4pt minus 1pt}\noindent{\kbvsecfont\fontsize{11pt}{14.5pt}\selectfont\color{kbvaccentdark} (೧)\hspace{0.45em}ಉತ್ಪಲಮಾಲಾ ವೃತ್ತ :}\hspace{0.5em}ಇದು ನಾಲ್ಕು ಪಾದಗಳ ವೃತ್ತ, ಪ್ರತಿಯೊಂದು ಪಾದದಲ್ಲಿ ಆರು ಗಣಗಳು, ಒಂದು ಲಘು ಮತ್ತು ಒಂದು ಗುರು ಸೇರಿ ಒಟ್ಟು ಇಪ್ಪತ್ತು ಅಕ್ಷರಗಳಿರುವುವು.\par

\begin{kbvnote}
\noindent{\kbvsecfont ವಿನ್ಯಾಸ :}\hspace{0.5em}(ಭ--ರ..-ನ..-ಭ..-ಭ..-ರ..-೧ಲ--೧ಗು)
\end{kbvnote}

\vspace{4pt plus 4pt minus 1pt}\noindent{\kbvsecfont\fontsize{11pt}{14.5pt}\selectfont\color{kbvaccentdark} (೨)\hspace{0.45em}ಚಂಪಕಮಾಲಾ ವೃತ್ತ :}\hspace{0.5em}ಇದು ನಾಲ್ಕು ಪಾದಗಳ ವೃತ್ತ, ಪ್ರತಿಯೊಂದು ಪಾದದಲ್ಲಿ ಏಳು ಗಣಗಳಿದ್ದು ಒಟ್ಟು ಇಪ್ಪತ್ತೊಂದು ಅಕ್ಷರ ಗಳಿರುವುವು.\par

\begin{kbvnote}
\noindent{\kbvsecfont ವಿನ್ಯಾಸ :}\hspace{0.5em}(ನ..-ಜ..-ಭ..-ಜ..-ಜ..-ಜ..-ರ..)
\end{kbvnote}

\vspace{4pt plus 4pt minus 1pt}\noindent{\kbvsecfont\fontsize{11pt}{14.5pt}\selectfont\color{kbvaccentdark} (೩)\hspace{0.45em}ಶಾರ್ದೂಲ ವಿಕ್ರೀಡಿತ ವೃತ್ತ :}\hspace{0.5em}ಇದು ನಾಲ್ಕು ಪಾದಗಳ ವೃತ್ತ. ಪ್ರತಿಯೊಂದು ಪಾದದಲ್ಲಿ ಆರು ಗಣಗಳು ಮತ್ತು ಒಂದು ಗುರು ಸೇರಿ ಒಟ್ಟು ಹತ್ತೊಂಬತ್ತು ಅಕ್ಷರಗಳಿರುವುವು.\par

\begin{kbvnote}
\noindent{\kbvsecfont ವಿನ್ಯಾಸ :}\hspace{0.5em}(ಮ--ಸ..-ಜ..-ಸ..-ತ..-ತ..-೧ಗು)
\end{kbvnote}

\vspace{4pt plus 4pt minus 1pt}\noindent{\kbvsecfont\fontsize{11pt}{14.5pt}\selectfont\color{kbvaccentdark} (೪)\hspace{0.45em}ಮತ್ತೇಭ ವಿಕ್ರೀಡಿತ ವೃತ್ತ :}\hspace{0.5em}ಇದು ನಾಲ್ಕು ಪಾದಗಳ ವೃತ್ತ. ಪ್ರತಿಯೊಂದು ಪಾದದಲ್ಲಿ ಆರು ಗಣಗಳು, ಒಂದು ಲಘು ಮತ್ತು ಒಂದು ಗುರು ಸೇರಿ ಒಟ್ಟು ಇಪ್ಪತ್ತು ಅಕ್ಷರಗಳಿರುವುವು.\par

\begin{kbvnote}
\noindent{\kbvsecfont ವಿನ್ಯಾಸ :}\hspace{0.5em}(ಸ..-ಭ..-ರ..-ನ..-ಮ..-ಯ..-೧ಲ--೧ಗು)
\end{kbvnote}

\vspace{4pt plus 4pt minus 1pt}\noindent{\kbvsecfont\fontsize{11pt}{14.5pt}\selectfont\color{kbvaccentdark} (೫)\hspace{0.45em}ಸ್ರಗ್ಧರಾ ವೃತ್ತ :}\hspace{0.5em}ಇದು ನಾಲ್ಕು ಪಾದಗಳ ವೃತ್ತ. ಪ್ರತಿ ಯೊಂದು ಪಾದದಲ್ಲಿ ಏಳು ಗಣಗಳಿದ್ದು ಒಟ್ಟು ಇಪ್ಪತ್ತೊಂದು ಅಕ್ಷರಗಳಿರು ವುವು.\par

\begin{kbvnote}
\noindent{\kbvsecfont ವಿನ್ಯಾಸ :}\hspace{0.5em}(ಮ..-ರ..-ಭ..-ನ..-ಯ..-ಯ..-ಯ..)
\end{kbvnote}
}
\clearpage

% === tex/0109.tex (Printed Page 107) ===
\renewcommand{\kbvchaptername}{ಛಂದಸ್ಸು}
\setcounter{page}{107}
\thispagestyle{kbvpage}
\kbvtypesetpage{\vspace{4pt plus 4pt minus 1pt}\noindent{\kbvsecfont\fontsize{11pt}{14.5pt}\selectfont\color{kbvaccentdark} (೬)\hspace{0.45em}ಮಹಾಸ್ರಗ್ಧರಾ ವೃತ್ತ :}\hspace{0.5em}ಇದು ನಾಲ್ಕು ಪಾದಗಳ ವೃತ್ತ. ಪ್ರತಿಯೊಂದು ಪಾದದಲ್ಲಿ ಏಳು ಗಣಗಳು ಮತ್ತು ಒಂದು ಗುರು ಸೇರಿ ಒಟ್ಟು ಇಪ್ಪತ್ತೆರಡು ಅಕ್ಷರಗಳಿರುವುವು.\par

\begin{kbvnote}
\noindent{\kbvsecfont ವಿನ್ಯಾಸ :}\hspace{0.5em}(ಸ + ತ + ತ + ನ + ಸ + ರ + ರ + ೧ಗು)
\end{kbvnote}

\vspace{6pt plus 5pt minus 1pt}\begin{center}{\kbvsecfont\fontsize{12.5pt}{16pt}\selectfont\color{kbvaccentdark} ಅಕ್ಷರ ಗಣಗಳ ವಿನ್ಯಾಸ ತೋರಿಸುವ ಪದಕ}\end{center}\vspace{3pt plus 3pt minus 1pt}

\vspace{5pt plus 4pt minus 1pt}
\noindent\begin{adjustbox}{max width=\textwidth}
{\renewcommand{\arraystretch}{1.30}\setlength{\tabcolsep}{4.5pt}
\begin{tabular}{l c c c c c c c c c c}
\toprule
\textbf{ವೃತ್ತ} & \textbf{ಗ1} & \textbf{ಗ2} & \textbf{ಗ3} & \textbf{ಗ4} & \textbf{ಗ5} & \textbf{ಗ6} & \textbf{ಗ7} & \textbf{ಲ} & \textbf{ಗು} & \textbf{ಒಟ್ಟು ಅಕ್ಷರಗಳು} \\
\midrule
ಉ.ಮಾ. & ಭ & ರ & ನ & ಭ & ಭ & ರ & .... & ೧ & ೧ & 20 \\
\hline
ಚಂ.ಮಾ. & ನ & ಜ & ಭ & ಜ & ಜ & ಜ & ರ & .... & .... & 21 \\
\hline
ಶಾ.ವಿ. & ಮ & ಸ & ಜ & ಸ & ತ & ತ & .... & .... & ೧ & 19 \\
\hline
ಮ.ವಿ. & ಸ & ಭ & ರ & ನ & ಮ & ಯ & .... & ೧ & ೧ & 20 \\
\hline
ಸ್ರಗ್ಧರಾ & ಮ & ರ & ಭ & ನ & ಯ & ಯ & ಯ & .... & .... & 21 \\
\hline
ಮ.ಸ್ರ. & ಸ & ತ & ತ & ನ & ಸ & ರ & ರ & .... & ೧ & 22 \\
\bottomrule
\end{tabular}}
\end{adjustbox}
\vspace{6pt plus 4pt minus 1pt}

\kbvexercise{ಅಭ್ಯಾಸ ಪತ್ರಿಕೆ : ೨೧}

\vspace{2.5pt plus 3pt minus 0.5pt}\noindent\hspace*{1.2em}\kbvnum{1.}\hspace{0.55em}ಛಂದಶ್ಶಾಸ್ತ್ರ ಎಂದರೇನು ?\par

\vspace{2.5pt plus 3pt minus 0.5pt}\noindent\hspace*{1.2em}\kbvnum{2.}\hspace{0.55em}“ಛಂದೋಂಬುಧಿ” ಎಂಬ ಗ್ರಂಥದ ಕರ್ತೃ ಯಾರು ?\par

\vspace{2.5pt plus 3pt minus 0.5pt}\noindent\hspace*{1.2em}\kbvnum{3.}\hspace{0.55em}“ತ್ರಿಪದಿ” ಎಂದರೇನು ?\par

\kbvsechead{4. ಇವುಗಳನ್ನು ಕುರಿತು ಟಿಪ್ಪಣಿ ಬರೆಯಿರಿ :}

\vspace{2.5pt plus 3pt minus 0.5pt}\noindent\hspace*{1.2em}\kbvnum{(೧)}\hspace{0.55em}ಯತಿ, (೨) ಗಣ, (೩) ಪ್ರಾಸ\par

\vspace{2.5pt plus 3pt minus 0.5pt}\noindent\hspace*{1.2em}\kbvnum{5.}\hspace{0.55em}ಇವುಗಳಲ್ಲಿ ಪ್ರತಿಯೊಂದಕ್ಕೆ ನಾಲ್ಕು ಉದಾಹರಣೆಗಳನ್ನು ಕೊಡಿ :\par

\vspace{2.5pt plus 3pt minus 0.5pt}\noindent\hspace*{1.2em}\kbvnum{(೧)}\hspace{0.55em}ಒಂದು ಮಾತ್ರಾಕಾಲದ ಅಕ್ಷರಗಳು\par

\vspace{2.5pt plus 3pt minus 0.5pt}\noindent\hspace*{1.2em}\kbvnum{(೨)}\hspace{0.55em}ಎರಡು ಮಾತ್ರಾಕಾಲದ ಅಕ್ಷರಗಳು\par
}
\clearpage

% === tex/0110.tex (Printed Page 108) ===
\renewcommand{\kbvchaptername}{ಛಂದಸ್ಸು}
\setcounter{page}{108}
\thispagestyle{kbvpage}
\kbvtypesetpage{\kbvsechead{6. ಈ ಪದಗಳ ಅರ್ಥ ವಿವರಣೆ ಕೊಡಿ :}

\vspace{2.5pt plus 3pt minus 0.5pt}\noindent\hspace*{1.2em}\kbvnum{(೧)}\hspace{0.55em}ಲಘು, (೨) ಗುರು\par

\vspace{2.5pt plus 3pt minus 0.5pt}\noindent\hspace*{1.2em}\kbvnum{7.}\hspace{0.55em}ಪ್ರಸ್ತಾರ ಹಾಕುವುದು ಎಂದರೇನು ?\par

\vspace{2.5pt plus 3pt minus 0.5pt}\noindent\hspace*{1.2em}\kbvnum{8.}\hspace{0.55em}ಈ ಕಂದ ಪದ್ಯಕ್ಕೆ ಪ್ರಸ್ತಾರ ಹಾಕಿ ಮಾತ್ರಾಗಣಗಳನ್ನು ವಿಂಗಡಿಸಿ.\par

\begin{kbvverse}
\noindent ಗಿರಿಜೆ ಶಶಿವದನೆ ಕೊಳನಂ । \\
\noindent ಭರದಿಂ ಪೊಗೆ ನೆಯ್ದಿಲೆಯ್ದಿದುವು ವಿಕಸನಮಂ । \\
\noindent ವಿರಹವನಾಂತವು ಕೋಳಂ । \\
\noindent ಸರಸಿರುಹಂ ಮುಗಿದುವರರೆ ಪಗಲೊಳ ಚಿತ್ರಂ ॥
\end{kbvverse}

\noindent\hfill{\fontsize{9.5pt}{12.5pt}\selectfont\color{kbvgrey}—ಹರೀಶ್ವರನ ಗಿರಿಜಾ ಕಲ್ಯಾಣ}\par\vspace{3pt plus 2pt}

\vspace{2.5pt plus 3pt minus 0.5pt}\noindent\hspace*{1.2em}\kbvnum{9.}\hspace{0.55em}ಮೇಲಿನ ಪದ್ಯವನ್ನು ಆಧಾರವಾಗಿಟ್ಟುಕೊಂಡು ಕಂದ ಪದ್ಯದ ಲಕ್ಷಣಗಳನ್ನು ತಿಳಿಸಿ.\par

\vspace{2.5pt plus 3pt minus 0.5pt}\noindent\hspace*{1.2em}\kbvnum{10.}\hspace{0.55em}ಷಟ್ಪದಿ ಎಂದರೇನು ? ಷಟ್ಪದಿಯಲ್ಲಿ ಎಷ್ಟು ಪ್ರಕಾರಗಳಿವೆ ? ಮತ್ತು ಅವು ಯಾವುವು ?\par

\vspace{2.5pt plus 3pt minus 0.5pt}\noindent\hspace*{1.2em}\kbvnum{11.}\hspace{0.55em}ನಿಮ್ಮ ಪಠ್ಯಪುಸ್ತಕದಲ್ಲಿರುವ ಷಟ್ಪದಿ ಜಾತಿಯ ಯಾವುದಾದ ರೊಂದು ಪದ್ಯಕ್ಕೆ ಪ್ರಸ್ತಾರ ಹಾಕಿ.\par

\vspace{2.5pt plus 3pt minus 0.5pt}\noindent\hspace*{1.2em}\kbvnum{12.}\hspace{0.55em}“ರಗಳೆ” ಎಂದರೇನು ? ಮುಖ್ಯವೆನಿಸಿದ ಮೂರು ಪ್ರಕಾರದ ರಗಳೆಗಳನ್ನು ಹೆಸರಿಸಿ.\par

\vspace{2.5pt plus 3pt minus 0.5pt}\noindent\hspace*{1.2em}\kbvnum{13.}\hspace{0.55em}ಲಲಿತ ರಗಳೆಯ ಲಕ್ಷಣಗಳನ್ನು ತಿಳಿಸಿ.\par

\vspace{2.5pt plus 3pt minus 0.5pt}\noindent\hspace*{1.2em}\kbvnum{14.}\hspace{0.55em}ವಾರ್ಧಕ ಷಟ್ಪದಿಯ ಲಕ್ಷಣಗಳನ್ನು ತಿಳಿಸಿರಿ.\par

\vspace{2.5pt plus 3pt minus 0.5pt}\noindent\hspace*{1.2em}\kbvnum{15.}\hspace{0.55em}ಯಾವ ಪ್ರಕಾರದ ಷಟ್ಪದಿ ಪದ್ಯಗಳಲ್ಲಿ ಮಧ್ಯೆ ಗುರುವುಳ್ಳ ಗಣ ಕಂಡುಬರುವುದಿಲ್ಲ ?\par

\vspace{2.5pt plus 3pt minus 0.5pt}\noindent\hspace*{1.2em}\kbvnum{16.}\hspace{0.55em}ಅಕ್ಷರ ಗಣ ಎಂದರೇನು ? ಅಕ್ಷರಗಣಕ್ಕೂ ಮಾತ್ರಾಗಣಕ್ಕೂ ಇರುವ ವ್ಯತ್ಯಾಸವೇನು ?\par

\vspace{2.5pt plus 3pt minus 0.5pt}\noindent\hspace*{1.2em}\kbvnum{17.}\hspace{0.55em}ವೃತ್ತಜಾತಿಯ ಛಂದಸ್ಸಿಗೆ ಸೇರಿದ ಆರು ವಿಧವಾದ ವೃತ್ತ ಗಳನ್ನು ಹೆಸರಿಸಿ.\par
}
\clearpage

% === tex/0111.tex (Printed Page 109) ===
\renewcommand{\kbvchaptername}{ಛಂದಸ್ಸು}
\setcounter{page}{109}
\thispagestyle{kbvpage}
\kbvtypesetpage{\vspace{2.5pt plus 3pt minus 0.5pt}\noindent\hspace*{1.2em}\kbvnum{18.}\hspace{0.55em}ಎಂಟು ವಿಧವಾದ ಅಕ್ಷರ ಗಣಗಳನ್ನು ಮತ್ತು ಅವುಗಳ ವಿನ್ಯಾಸವನ್ನು ಬರೆದು ತಿಳಿಸಿ\par

\vspace{2.5pt plus 3pt minus 0.5pt}\noindent\hspace*{1.2em}\kbvnum{19.}\hspace{0.55em}ಕೆಳಗೆ ಕೊಟ್ಟಿರುವ ಪದಗಳಿಗೆ ಪ್ರಸ್ತಾರ ಹಾಕಿ ಅವುಗಳು ಯಾವ ಅಕ್ಷರ ಗಣಕ್ಕೆ ಸೇರಿದವುಗಳೆಂಬುದನ್ನು ತಿಳಿಸಿರಿ.\par

\kbvsechead{ಕೆಂಜೇಳು ಸಾಧನೆ, ವಿನಾಶ ಮಕರ, ರೇಣುಕಾ, ಜನತಾ :}

\indent ಮಹಂತೇಶ್‌ ಶ್ರೀಲಂಕಾ, ಮಾಣಿಕ್ಯಂ\par

\vspace{2.5pt plus 3pt minus 0.5pt}\noindent\hspace*{1.2em}\kbvnum{20.}\hspace{0.55em}ಉತ್ಪಲಮಾಲಾ ವೃತ್ತ ಜಾತಿಯ ಪದ್ಯದ ಲಕ್ಷಣಗಳನ್ನು ತಿಳಿಸಿ.\par

\vspace{2.5pt plus 3pt minus 0.5pt}\noindent\hspace*{1.2em}\kbvnum{21.}\hspace{0.55em}ನಿಂದಿಸುವರ್‌ ದುರ್ಜನರ್‌ ಎಂ ।\par

\begin{kbvverse}
\noindent ಬೊಂದು ಭಯಂಜಿತ್ತು ಸುಕವಿ ರಚಿಸನೆ ಕೃತಿಯಂ ॥ \\
\noindent ಮಂದೇಹರ ಭಯದಿಂದರ । \\
\noindent ವಿಂದಸಖಂ ನಿಜ ಮಯೂಖಮಂ ಪ್ರಸರಿಸದೇ ? ॥
\end{kbvverse}

\noindent\hfill{\fontsize{9.5pt}{12.5pt}\selectfont\color{kbvgrey}— ರಾಜಶೇಖರ ವಿಲಾಸ}\par\vspace{3pt plus 2pt}

\vspace{2.5pt plus 3pt minus 0.5pt}\noindent\hspace*{1.2em}\kbvnum{೧)}\hspace{0.55em}ಈ ಪದ್ಯದ ತಾತ್ಪರ್ಯವೇನು ?\par

\vspace{2.5pt plus 3pt minus 0.5pt}\noindent\hspace*{1.2em}\kbvnum{೨)}\hspace{0.55em}ಈ ಪದ್ಯದಲ್ಲಿ ಯಾವ ಅಲಂಕಾರವಿರುವುದು ಎಂಬುದನ್ನು... ತಿಳಿಸಿರಿ,\par

\vspace{2.5pt plus 3pt minus 0.5pt}\noindent\hspace*{1.2em}\kbvnum{೩)}\hspace{0.55em}ಈ ಪದ್ಯಕ್ಕೆ ಪ್ರಸ್ತಾರ ಹಾಕಿ, ಮಾತ್ರಾ ಗಣಗಳನ್ನು ವಿಂಗಡಿಸಿರಿ,\par

\vspace{2.5pt plus 3pt minus 0.5pt}\noindent\hspace*{1.2em}\kbvnum{೪)}\hspace{0.55em}ಇದು “ಕಂದ” ಪದ್ಯ ಎಂಬುದನ್ನು ಹೇಗೆ ಸಮರ್ಥಿಸುವಿರಿ?\par
}
\clearpage

% === tex/0112.tex (Printed Page 110) ===
\renewcommand{\kbvchaptername}{ದಿನಬಳಕೆಯ ಗಾದೆಗಳು}
\setcounter{page}{110}
\thispagestyle{kbvpage}
\kbvtypesetpage{\kbvchapter{ಪಾಠ ೨೨}{ದಿನಬಳಕೆಯ ಗಾದೆಗಳು}

\noindent\hspace*{0.6em}\kbvbullet\hspace{0.55em}ಅಕ್ಕಿ ಮೇಲೆ ಆಸೆ ನೆಂಟರ ಮೇಲೆ ಇಷ್ಟ.\par

\noindent\hspace*{0.6em}\kbvbullet\hspace{0.55em}ಅತ್ತೆಗೊಂದು ಕಾಲ ಸೊಸೆಗೊಂದು ಕಾಲ.\par

\noindent\hspace*{0.6em}\kbvbullet\hspace{0.55em}ಅತಿ ಆಸೆ ಗತಿ ಕೆಡಿಸಿತು.\par

\noindent\hspace*{0.6em}\kbvbullet\hspace{0.55em}ಅಲ್ಪರ ಸಂಗ ಅಭಿಮಾನ ಭಂಗ.\par

\noindent\hspace*{0.6em}\kbvbullet\hspace{0.55em}ಅಗ್ಗದ ಶೆಟ್ಟಿ ಮುಗ್ಗಿದ ಜೋಳ.\par

\noindent\hspace*{0.6em}\kbvbullet\hspace{0.55em}ಅಜ್ಜಿ ನೂತದ್ದೆಲ್ಲಾ ಅಜ್ಜನ ಉಡಿದಾರಕ್ಕೆ.\par

\noindent\hspace*{0.6em}\kbvbullet\hspace{0.55em}ಅನಿಷ್ಟಕ್ಕೆ ಶನೇಶ್ವರ ಗುರಿ,\par

\noindent\hspace*{0.6em}\kbvbullet\hspace{0.55em}ಅಹಂಕಾರಕ್ಕೆ ಉದಾಸೀನವೇ ಮದ್ದು.\par

\noindent\hspace*{0.6em}\kbvbullet\hspace{0.55em}ಅಟ್ಟ ಹತ್ತದವನು ಬೆಟ್ಟ ಹತ್ತಿಯಾನೆ?\par

\noindent\hspace*{0.6em}\kbvbullet\hspace{0.55em}ಅಲ್ಪ ವಿದ್ಯಾ ಮಹಾಗರ್ವಿ.\par

\noindent\hspace*{0.6em}\kbvbullet\hspace{0.55em}ಅಲ್ಪನಿಗೆ ಐಶ್ವರ್ಯ ಬಂದರೆ ನಡುರಾತ್ರೀಲಿ ಕೊಡೆ ಹಿಡಿಸಿಕೊಂಡ.\par

\noindent\hspace*{0.6em}\kbvbullet\hspace{0.55em}ಅಂಗೈ ಹುಣ್ಣಿಗೆ ಕನ್ನಡಿಯೇಕೆ?\par

\noindent\hspace*{0.6em}\kbvbullet\hspace{0.55em}ಆರು ದೋಸೆ ಕೊಟ್ಟರೆ ಅತ್ತೆ ಕಡೆ ಮೂರು ದೋಸೆ ಕೊಟ್ಟರೆ ಸೊಸೆ ಕಡೆ.\par

\noindent\hspace*{0.6em}\kbvbullet\hspace{0.55em}ಆಕಳು ಕಪ್ಪಾದರೆ ಹಾಲು ಕಪ್ಪೇ?\par

\noindent\hspace*{0.6em}\kbvbullet\hspace{0.55em}ಆಳಾಗಬಲ್ಲವನು ಅರಸಾಗಬಲ್ಲ.\par

\noindent\hspace*{0.6em}\kbvbullet\hspace{0.55em}ಆಪತ್ತಿಗೆ ಆದವನೇ ನೆಂಟ.\par

\noindent\hspace*{0.6em}\kbvbullet\hspace{0.55em}ಇಬ್ಬರ ಜಗಳದಲ್ಲಿ ಮೂರನೆಯವನಿಗೆ ಲಾಭ.\par

\noindent\hspace*{0.6em}\kbvbullet\hspace{0.55em}ಇರುಳು ಕಂಡ ಬಾವೀಲಿ ಹಗಲು ಬಿದ್ದ.\par

\noindent\hspace*{0.6em}\kbvbullet\hspace{0.55em}ಇಲಿಗೆ ಪ್ರಾಣ ಸಂಕಟ, ಬೆಕ್ಕಿಗೆ ಚೆಲ್ಲಾಟ.\par

\noindent\hspace*{0.6em}\kbvbullet\hspace{0.55em}ಉರಿಯೋ ಹುಣ್ಣಿನ ಮೇಲೆ ಬರೆ ಹಾಕಿದ ಹಾಗೆ.\par

\noindent\hspace*{0.6em}\kbvbullet\hspace{0.55em}ಉಗುರಿನಲ್ಲಿ ಹೋಗೋದಕ್ಕೆ ಕೊಡಲಿ ಪೆಟ್ಟು.\par

\noindent\hspace*{0.6em}\kbvbullet\hspace{0.55em}ಉಪ್ಪು ತಿಂದವನು ನೀರು ಕುಡೀತಾನೆ.\par

\noindent\hspace*{0.6em}\kbvbullet\hspace{0.55em}ಉಪ್ಪುಂಡ ಮನೆಗೆ ಎರಡು ಬಗೆಯಬಾರದು.\par

\noindent\hspace*{0.6em}\kbvbullet\hspace{0.55em}ಊರಿದ್ದ ಕಡೆ ಹಗರಿ,\par

\noindent\hspace*{0.6em}\kbvbullet\hspace{0.55em}ಊಟ ಬಲ್ಲವನಿಗೆ ರೋಗವಿಲ್ಲ, ಮಾತು ಬಲ್ಲವನಿಗೆ ಜಗಳವಿಲ್ಲ.\par
}
\clearpage

% === tex/0113.tex (Printed Page 111) ===
\renewcommand{\kbvchaptername}{ದಿನಬಳಕೆಯ ಗಾದೆಗಳು}
\setcounter{page}{111}
\thispagestyle{kbvpage}
\kbvtypesetpage{\noindent\hspace*{0.6em}\kbvbullet\hspace{0.55em}ಊರು ಹಾಳಾದ ಮೇಲೆ ಕೋಟೆ ಬಾಗಿಲು ಹಾಕಿದಂತೆ,\par

\noindent\hspace*{0.6em}\kbvbullet\hspace{0.55em}ಎಲ್ಲರ ಮನೆ ದೋಸೆ ತೂತು ಅಂದ್ರೆ, ನಮ್ಮ ಮನೆ ಕಾವಲಿಯೇ ತೂತು.\par

\noindent\hspace*{0.6em}\kbvbullet\hspace{0.55em}ಎತ್ತು ಏರಿಗೆ ಎಳೆದರೆ, ಕೋಣ ನೀರಿಗೆ ಎಳೀತು.\par

\noindent\hspace*{0.6em}\kbvbullet\hspace{0.55em}ಎತ್ತಿಗೆ ಜ್ವರ ಬಂದರೆ ಎಮ್ಮೆಗೆ ಬರೆ ಹಾಕಿದರು.\par

\noindent\hspace*{0.6em}\kbvbullet\hspace{0.55em}ಕಚ್ಚುವ ನಾಯಿ ಬೊಗಳೋಲ್ಲ, ಬೊಗಳುವ ನಾಯಿ ಕಚ್ಚೊಲ್ಲ.\par

\noindent\hspace*{0.6em}\kbvbullet\hspace{0.55em}ಕಜ್ಜಿ ಆದವರಿಗೆ ಲಜ್ಜೆ ಇಲ್ಲ.\par

\noindent\hspace*{0.6em}\kbvbullet\hspace{0.55em}ಕಳ್ಳನ ಹೆಂಡ್ತಿ ಎಂದಿದ್ದರೂ ಮುಂಡೆ,\par

\noindent\hspace*{0.6em}\kbvbullet\hspace{0.55em}ಕಟುಕನ ಮುಂದೆ ಬಸವನ ಪುರಾಣ\par

\noindent\hspace*{0.6em}\kbvbullet\hspace{0.55em}ಕಾರ್ಯವಾಸಿ ಕತ್ತೆ ಕಾಲು ಕಟ್ಟು.\par

\noindent\hspace*{0.6em}\kbvbullet\hspace{0.55em}ಕುಣಿಯಲಾರದ ಸೂಳೆ ನೆಲ ಡೊಂಕು ಅಂದಳು,\par

\noindent\hspace*{0.6em}\kbvbullet\hspace{0.55em}ಕೂತು ಉಣ್ಣುವವರಿಗೆ ಕುಡಿಕೆ ಹಣ ಸಾಲದು.\par

\noindent\hspace*{0.6em}\kbvbullet\hspace{0.55em}ಕೂಸು ಬಡವಾಗಬಾರದು, ಕಾಸು ಖರ್ಚಾಗಬಾರದು.\par

\noindent\hspace*{0.6em}\kbvbullet\hspace{0.55em}ಕೂಸು ಹುಟ್ಟೋಕೆ ಮುಂಚೆ ಕುಲಾವಿ ಹೊಲಿಸಿದರು.\par

\noindent\hspace*{0.6em}\kbvbullet\hspace{0.55em}ಕೈಲಾಗದವ ಮೈ ಪರಚಿಕೊಂಡ.\par

\noindent\hspace*{0.6em}\kbvbullet\hspace{0.55em}ಕೈ ಕೆಸರಾದರೆ ಬಾಯಿ ಮೊಸರು.\par

\noindent\hspace*{0.6em}\kbvbullet\hspace{0.55em}ಕೈಗೆ ಬಂದ ತುತ್ತು ಬಾಯಿಗೆ ಬರಲಿಲ್ಲ.\par

\noindent\hspace*{0.6em}\kbvbullet\hspace{0.55em}ಕೊಟ್ಟ ಹೆಣ್ಣು ಕುಲಕ್ಕೆ ಹೊರಗು.\par

\noindent\hspace*{0.6em}\kbvbullet\hspace{0.55em}ಕೋಳೀನ ಕೇಳಿ ಮೆಣಸು ಅರೆಯೋದು.\par

\noindent\hspace*{0.6em}\kbvbullet\hspace{0.55em}ಕೋಣನ ಮುಂದೆ ಕಿನ್ನರಿ ಬಾರಿಸಿದರು.\par

\noindent\hspace*{0.6em}\kbvbullet\hspace{0.55em}ಕೋತಿ ಕೈಗೆ ಮಾಣಿಕ್ಯ ಕೊಟ್ಟ ಹಾಗೆ.\par

\noindent\hspace*{0.6em}\kbvbullet\hspace{0.55em}ಕಂತೆಗೆ ತಕ್ಕಂತೆ ಬೊಂತೆ.\par

\noindent\hspace*{0.6em}\kbvbullet\hspace{0.55em}ಕಂಕುಳಲ್ಲಿ ದೊಣ್ಣೆ, ಕೈಯಲ್ಲಿ ಶರಣಾರ್ಥಿ.\par

\noindent\hspace*{0.6em}\kbvbullet\hspace{0.55em}ಕಂಕುಳಲ್ಲಿ ಮಗು ಇಟ್ಟುಕೊಂಡು ಊರೆಲ್ಲಾ ಹುಡುಕಿದರು.\par

\noindent\hspace*{0.6em}\kbvbullet\hspace{0.55em}ಕುಂಬಳಕಾಯಿ ಕಳ್ಳ ಅಂದರೆ ಹೆಗಲು ಮುಟ್ಟಿ ನೋಡಿಕೊಂಡ.\par

\noindent\hspace*{0.6em}\kbvbullet\hspace{0.55em}ಕುಂಡಿ ಮೇಲೆ ಹೊಡೆದರೆ ದವಡೆ ಹಲ್ಲು ಉದುರಿತು.\par

\noindent\hspace*{0.6em}\kbvbullet\hspace{0.55em}ಕುಂಬಾರನಿಗೆ ವರ್ಷ, ದೊಣ್ಣೆಗೆ ಒಂದು ನಿಮಿಷ.\par
}
\clearpage

% === tex/0114.tex (Printed Page 112) ===
\renewcommand{\kbvchaptername}{ದಿನಬಳಕೆಯ ಗಾದೆಗಳು}
\setcounter{page}{112}
\thispagestyle{kbvpage}
\kbvtypesetpage{\noindent\hspace*{0.6em}\kbvbullet\hspace{0.55em}ಕುಂಟೆತ್ತು ಹೊಲಕ್ಕೆ ಊನ, ತನ್ನವರು ತನಗೇ ಊನ.\par

\noindent\hspace*{0.6em}\kbvbullet\hspace{0.55em}ಕೆಂಬೂತ ನೋಡಿ ಕಾಗೆ ಪುಕ್ಕ ತರಕೊಂಡ ಹಾಗೆ.\par

\noindent\hspace*{0.6em}\kbvbullet\hspace{0.55em}ಗಾಣಗಿತ್ತಿ ಅಯ್ಯೋ ಅಂದ್ರೆ ನೆತ್ತಿ ತಂಪಾದೀತೆ?\par

\noindent\hspace*{0.6em}\kbvbullet\hspace{0.55em}ಗಾಳಿ ಬಂದಾಗ ತೂರಿಕೋ.\par

\noindent\hspace*{0.6em}\kbvbullet\hspace{0.55em}ಗುಬ್ಬಿ ಮೇಲೆ ಬ್ರಹ್ಮಾಸ್ತ್ರ.\par

\noindent\hspace*{0.6em}\kbvbullet\hspace{0.55em}ಗಂಜಿ ಕುಡಿಯೋನಿಗೆ ಗಡ್ಡ ನೀವೋನೊಬ್ಬ.\par

\noindent\hspace*{0.6em}\kbvbullet\hspace{0.55em}ಚೋರ ಗುರು ಚಂಡಾಲ ಶಿಷ್ಯ.\par

\noindent\hspace*{0.6em}\kbvbullet\hspace{0.55em}ಚಿಂತೆಯಿಲ್ಲದವನಿಗೆ ಸಂತೇಲಿ ನಿದ್ದೆ.\par

\noindent\hspace*{0.6em}\kbvbullet\hspace{0.55em}ಜಟ್ಟಿ ಜಾರಿಬಿದ್ದರೂ ಮೀಸೆ ಮಣ್ಣಾಗಲಿಲ್ಲ.\par

\noindent\hspace*{0.6em}\kbvbullet\hspace{0.55em}ಜಾತಿಗೆ ಜಾತಿ ವೈರಿ, ನೀರಿಗೆ ಪಾಚಿ ವೈರಿ.\par

\noindent\hspace*{0.6em}\kbvbullet\hspace{0.55em}ತಕ್ಕಡೀಲಿ ಕಪ್ಪೆ ತೂಗಿದ ಹಾಗೆ.\par

\noindent\hspace*{0.6em}\kbvbullet\hspace{0.55em}ತಾಳಿದವನು ಬಾಳಿಯಾನು.\par

\noindent\hspace*{0.6em}\kbvbullet\hspace{0.55em}ತೋಳ ಹಳ್ಳಕ್ಕೆ ಬಿದ್ದರೆ ಆಳಿಗೊಂದು ಕಲ್ಲು.\par

\noindent\hspace*{0.6em}\kbvbullet\hspace{0.55em}ತುಂಬಿದ ಕೊಡ ತುಳುಕುವುದಿಲ್ಲ.\par

\noindent\hspace*{0.6em}\kbvbullet\hspace{0.55em}ದಡ್ಡನಿಗೆ ದೊಣ್ಣೆ ಪೆಟ್ಟು, ಜಾಣನಿಗೆ ಮಾತಿನ ಪೆಟ್ಟು.\par

\noindent\hspace*{0.6em}\kbvbullet\hspace{0.55em}ದಿನ ಸಾಯೋರಿಗೆ ಅಳೋರ್ಯಾರು?\par

\noindent\hspace*{0.6em}\kbvbullet\hspace{0.55em}ದುಡ್ಡಿದ್ದವನೆ ದೊಡ್ಡಪ್ಪ.\par

\noindent\hspace*{0.6em}\kbvbullet\hspace{0.55em}ದುಷ್ಟರ ಕಂಡರೆ ದೂರ ಇರು.\par

\noindent\hspace*{0.6em}\kbvbullet\hspace{0.55em}ದೂರದ ಬೆಟ್ಟ ನುಣ್ಣಗೆ.\par

\noindent\hspace*{0.6em}\kbvbullet\hspace{0.55em}ದೇವರು ವರ ಕೊಟ್ಟರೂ ಪೂಜಾರಿ ಕೊಡ.\par

\noindent\hspace*{0.6em}\kbvbullet\hspace{0.55em}ಧರ್ಮಕ್ಕೆ ದಟ್ಟಿಕೊಟ್ಟರೆ ಹಿತ್ತಲಿಗೆ ಹೋಗಿ ಮೊಳ ಹಾಕಿದನಂತೆ.\par

\noindent\hspace*{0.6em}\kbvbullet\hspace{0.55em}ನಾಯಿ ಮೊಲೇಲಿ ಖಂಡುಗ ಹಾಲು.\par

\noindent\hspace*{0.6em}\kbvbullet\hspace{0.55em}ನಾಯಿ ಬಾಲ ಡೊಂಕು.\par

\noindent\hspace*{0.6em}\kbvbullet\hspace{0.55em}ನೆಕ್ಕುವ ನಾಯಿಗೆ ಲಿಂಗವೇನು? ಪೀಠವೇನು?\par

\noindent\hspace*{0.6em}\kbvbullet\hspace{0.55em}ಪಾಪಿ ಸಮುದ್ರಕ್ಕೆ ಹೋದರೂ ಮೊಣಕಾಲುದ್ದ ನೀರೇ.\par
}
\clearpage

% === tex/0115.tex (Blank Page) ===
\thispagestyle{empty}
\mbox{}
\clearpage

% === tex/0116.tex ===
% === File 0116.md: Back Cover Page ===
\thispagestyle{empty}
\vspace*{2mm}
\begin{center}
\begin{tikzpicture}
  % Main content node inside the framed box
  \node[
    inner xsep=6.5mm,
    inner ysep=8mm,
    text width=90mm,
    align=justify
  ] (box) {%
    \begin{center}
      {\Large\textbf{ಗಮನಿಸಿ :-}}
    \end{center}
    \vspace{4mm}
    \setlength{\parindent}{1.8em}%
    \setlength{\parskip}{4.5mm}%
    \linespread{1.28}\selectfont
    ಪ್ರಾಥಮಿಕ ಹಾಗೂ ಮಾಧ್ಯಮಿಕ ಶಾಲೆಯ ವಿದ್ಯಾರ್ಥಿಗಳಿಗೂ, ಪ್ರೌಢ ಶಾಲೆಯ ವಿದ್ಯಾರ್ಥಿಗಳಿಗೂ, ಕನ್ನಡವನ್ನು ತೃತೀಯ ಭಾಷೆಯಾಗಿ ತೆಗೆದುಕೊಂಡಿರುವ ವಿದ್ಯಾರ್ಥಿಗಳಿಗೂ, ಖಾಸಗಿಯಾಗಿ ಕನ್ನಡ ಕಲಿಯಲು ಬಯಸುವ ಕನ್ನಡೇತರರಿಗೂ, ವಯಸ್ಕರ ಶಾಲೆ ಸೇರಿ ಓದುಬರಹ ಕಲಿಯಬಯಸುವ ಹಿರಿಯರಿಗೂ ಉಪಯುಕ್ತವಾಗುವ ರೀತಿಯಲ್ಲಿ ಈ ವ್ಯಾಕರಣ ಶಾಸ್ತ್ರವನ್ನು ರಚಿಸಲಾಗಿದೆ.\par
    ಪ್ರಾರಂಭಿಕ ಹಂತದಲ್ಲಿ ಕಲಿಯುವವರ ಮನಸ್ಸಿಗೆ ತೊಡಕು ಉಂಟಾಗಬಾರದು. ಈ ಉದ್ದೇಶದಿಂದ ಆಯಾಯ ಪ್ರಕರಣಗಳಿಗೆ ಸಂಬಂಧಪಟ್ಟ ಪಠ್ಯ ವಿಷಯಗಳನ್ನು ತಿಳಿಯಾಗಿ ವಿವರಿಸಲಾಗಿದೆ ಮತ್ತು ಅಧಿಕ ಸಂಖ್ಯೆಯಲ್ಲಿ ಅತ್ಯಂತ ಸಮಂಜಸವಾದ ಉದಾಹರಣೆಗಳನ್ನು ಕೊಡಲಾಗಿದೆ. ಆದುದರಿಂದ ಈ ಪುಸ್ತಕವು \textbf{5, 6, 7, 8, 9} ಮತ್ತು \textbf{10} ನೆಯ ತರಗತಿಯ ವಿದ್ಯಾರ್ಥಿಗಳಿಗೆ ಹೆಚ್ಚು ಉಪಯುಕ್ತವೆನ್ನಬಹುದು.\par
    ವಿದ್ಯಾರ್ಥಿಗಳು ಆಯಾಯ ದರ್ಜೆಗೆ ಸೀಮಿತವಾದ ಪಠ್ಯಕ್ರಮಕ್ಕೆ ಅನುಸಾರವಾಗಿ ವಿಷಯವನ್ನು ಆರಿಸಿಕೊಂಡು ಅಭ್ಯಾಸ ಮಾಡಬೇಕಾದ್ದು ಅತಿಮುಖ್ಯ. ಅಭ್ಯಾಸ ಕಷ್ಟವೆನಿಸಿದಾಗ ಶಿಕ್ಷಕರ ನೆರವು ಪಡೆಯಲು ಹಿಂತೆಗೆಯಬಾರದು. ಆಸಕ್ತಿಯಿಂದ ಕಲಿಕೆ ಮುಂದುವರಿಸಿದಲ್ಲಿ ಕಷ್ಟ ಎನಿಸಿದ್ದು ಸುಲಭ ಸಾಧ್ಯವಾಗುವುದು.\par
    \vspace{2mm}
  };
  % Draw letterpress-style border rules with slight corner gaps like full/0116.webp
  \draw[line width=0.9pt] ([xshift=1.2mm]box.north west) -- ([xshift=-1.2mm]box.north east);
  \draw[line width=0.9pt] ([xshift=1.2mm]box.south west) -- ([xshift=-1.2mm]box.south east);
  \draw[line width=0.9pt] ([yshift=-1.2mm]box.north west) -- ([yshift=1.2mm]box.south west);
  \draw[line width=0.9pt] ([yshift=-1.2mm]box.north east) -- ([yshift=1.2mm]box.south east);
  % Bottom title and price aligned with box edges
  \node[anchor=north west, inner sep=0pt] at ([yshift=-3.5mm]box.south west) {\large\textbf{ಕನ್ನಡ ಬಾಲ ವ್ಯಾಕರಣ}};
  \node[anchor=north east, inner sep=0pt] at ([yshift=-3.5mm]box.south east) {\large\textbf{ರೂ. 10-00}};
\end{tikzpicture}
\end{center}
\clearpage


\end{document}
